% !TeX program = pdflatex
% papers/cristal_floxina.tex — GERADO por tools/cristal_tex.py; NÃO editar à mão.
% Fonte: cristal/cristal.jsonl (broca-so, última versão por conceito).
% Cada secção carrega o registo original em comentário %CRISTAL — a volta é
% tests/cristal_volta.js (resíduo 0 byte a byte contra a fonte).
\documentclass[11pt,a4paper]{article}
\usepackage[utf8]{inputenc}\usepackage[T1]{fontenc}\usepackage[brazil]{babel}
\usepackage{amsmath,amssymb}\usepackage{textcomp}
\usepackage[margin=2.4cm]{geometry}\usepackage{xcolor}
\definecolor{cinza}{gray}{0.35}
\newcommand{\prov}[1]{\par\smallskip\noindent{\small\color{cinza}#1}\par}
\setcounter{secnumdepth}{0}
% ─────────────────────────────────────────────────────────────────────────────
%  papers/gkcapa.tex — a capa do Reino, mínima, para os papers em `article`.
%
%  O estilo.tex completo é do `report` (títulos de capítulo, skak, …). Aqui fica
%  só o que a capa precisa, para o paper continuar a compilar sozinho com
%  pdflatex. O tradutor próprio (tests/tex) carrega o estilo.tex da raiz / o
%  symlink papers/estilo.tex e avalia o mesmo `\gkcapa`.
% ─────────────────────────────────────────────────────────────────────────────
\definecolor{ouro}{HTML}{B8912F}
\definecolor{ouroclaro}{HTML}{F3E9CE}
\definecolor{tinta}{HTML}{1E1B16}
\definecolor{regua}{HTML}{6E6858}
\providecommand{\gknota}{\fontsize{7.62}{11.05}\selectfont}
\providecommand{\gktexto}{\fontsize{10.50}{15.22}\selectfont}
\providecommand{\gksub}{\fontsize{12.33}{17.87}\selectfont}
\providecommand{\gksec}{\fontsize{14.47}{20.98}\selectfont}
\providecommand{\gkcap}{\fontsize{16.99}{24.63}\selectfont}
\providecommand{\gktit}{\fontsize{23.42}{33.95}\selectfont}
\providecommand{\Prj}{\mathbb{P}}
\providecommand{\Trf}{\mathcal{T}}
\providecommand{\gkcapa}[3]{%
\title{\vspace{-0.6cm}
{\color{ouro}\rule{\textwidth}{1.4pt}}\\[6mm]
\textbf{\gktit\color{tinta} Reino Dourado}\\[3mm]{\gksec\itshape\color{regua} Golden Kingdom}\\[8mm]
{\gkcap\scshape\color{ouro} #1}\\[5mm]
{\gksub\color{regua} #2}\\[2mm]
{\gktexto\color{regua}\itshape #3}\\[7mm]
{\color{ouro}\rule{\textwidth}{0.6pt}}\\[9mm]{\gknota\color{regua}\begin{minipage}{\textwidth}\setlength{\parindent}{0pt}%
\textbf{Aviso de Isenção Algorítmica:} O universo, as leis e as entidades contidas nestas páginas ---
incluindo felinos matriciais, aves de rapina algébricas e amuletos de controle de fluxo --- são
construções de pura ficção formal. Esta obra não descreve a realidade física, não serve como
engenharia reversa do cosmos e não constitui doutrina religiosa. Qualquer semelhança com bugs reais do seu
sistema operacional é mera coincidência (ou falta de reza). Prossiga por sua própria conta e
risco.\end{minipage}}}
\author{}\date{}}

\begin{document}
\gkcapa{O Cristal --- a floxina e as hipóteses}
       {a ponta de investigação do projeto}
       {corpus recuperado do broca-so; 187 conceitos; projeção verificável da fonte}
\maketitle

%CRISTAL {"arestas": [{"alvo": "ponte_quantica_conjunto_dougherty_e_orbitais_de_hidrogenio", "confianca": 1.0, "epistemico": "fato", "origem": "ponte_quantica_dougherty.md", "rel": "parte_de"}], "confianca": 0.52, "contraexemplos": [], "descricao": "Ponte geométrica **conjectural** entre densidades de probabilidade de elétrons do átomo de hidrogénio (microscopia de fotoionização, APS) e estruturas toroidais, helicoidais e vórtices aninhados do **Conjunto Dougherty**.", "epistemico": "hipotese", "exemplos": [], "id": "abstrato", "memoria": "semantica", "meta": {"dominio": "hipotese", "falsificavel": true, "nivel": 7}, "origem": "ponte_quantica_dougherty.md", "palavras_chave": [], "sinonimos": [], "tipo": "conceito", "titulo": "Abstrato"}
\section{Abstrato}
Ponte geométrica **conjectural** entre densidades de probabilidade de elétrons do átomo de hidrogénio (microscopia de fotoionização, APS) e estruturas toroidais, helicoidais e vórtices aninhados do **Conjunto Dougherty**.
\prov{relações: parte\_de ponte\_quantica\_conjunto\_dougherty\_e\_orbitais\_de\_hidrogenio}
\prov{proveniência: ponte\_quantica\_dougherty.md \textperiodcentered{} hipotese \textperiodcentered{} confiança 0.52 \textperiodcentered{} domínio hipotese \textperiodcentered{} história 6 versões no jornal}

%CRISTAL {"arestas": [{"alvo": "gato", "confianca": 1.0, "epistemico": "fato", "origem": "wiki", "rel": "usa"}, {"alvo": "gamba_antissimetrico", "confianca": 1.0, "epistemico": "fato", "origem": "wiki", "rel": "usa"}, {"alvo": "invariante_molecular_residuo", "confianca": 1.0, "epistemico": "fato", "origem": "wiki", "rel": "usa"}, {"alvo": "ponte_hidrogenio_parabola_geradora", "confianca": 1.0, "epistemico": "fato", "origem": "wiki", "rel": "relaciona"}, {"alvo": "gerador_do_gap_comutador", "confianca": 1.0, "epistemico": "fato", "origem": "wiki", "rel": "relaciona"}, {"alvo": "promocao_gamba_esquilo", "confianca": 1.0, "epistemico": "fato", "origem": "wiki", "rel": "usa"}], "confianca": 1.0, "contraexemplos": [], "descricao": "O ÁCIDO PROPIÔNICO (propanóico), CH₃-CH₂-COOH = C₃H₆O₂ (MW 74), fortalece a série do cheiro no canto 'pequeno, saturado, ácido'. É o menor ácido carboxílico de ODOR CORPORAL: cheiro pungente/azedo, no SUOR e no queijo suíço (produzido pela Propionibacterium; a P. acnes na pele). pKa≈4,88 (ácido fraco, como os carboxílicos ~4,8). SATURADO: grau de insaturação = 1 (só o C=O do carbonilo, NENHUM C=C). 5 átomos pesados (3C + 2O). DICIONÁRIO (a série gato/esquilo): o −COOH (a ponta ácida, doadora de H) ↔ o ESQUILO (a ponta reativa antissimétrica — o operador que gira, antes gambá); o pKa≈4,88 ↔ o GAP (a acidez: um gap grande = ácido); SATURADO (0 C=C) ↔ 0 π/esquilo (sem duplas — o mínimo, só o carbonilo); o esqueleto C-C-C+2O (5 pesados) ↔ o GATO (a adjacência de Hückel), 5 qubits (2⁵=32 Fock); o invariante a+c²=1, a≥a_min>0 ↔ a ZPE da ponte de H (o mesmo de toda molécula). Na série (pelo gap/acidez): −COOH(propiônico ≈4,9 ≈ o ác. 3-Me-hex-2-enoico 4,5) < fenol/−SH(~10) < álcool(~16). O propiônico é o −COOH puro, saturado (0 esquilo) — o vértice mínimo do ramo ácido. Fato: a fórmula, o pKa, o odor (Propionibacterium). Síntese: a leitura gato(esqueleto)/esquilo(−COOH)/gap(acidez), a série do cheiro.", "epistemico": "fato", "exemplos": [], "id": "acido_propionico_dicionario", "memoria": "semantica", "meta": {"dominio": "floxina", "fonte": "Floxina (investigação)"}, "origem": "wiki", "palavras_chave": [], "sinonimos": [], "tipo": "conceito", "titulo": "Ácido propiônico (CH₃CH₂COOH) — o menor ácido saturado de odor corporal; o dicionário"}
\section{Ácido propiônico (CH\ensuremath{_{3}}CH\ensuremath{_{2}}COOH) — o menor ácido saturado de odor corporal; o dicionário}
O ÁCIDO PROPIÔNICO (propanóico), CH\ensuremath{_{3}}-CH\ensuremath{_{2}}-COOH = C\ensuremath{_{3}}H\ensuremath{_{6}}O\ensuremath{_{2}} (MW 74), fortalece a série do cheiro no canto 'pequeno, saturado, ácido'. É o menor ácido carboxílico de ODOR CORPORAL: cheiro pungente/azedo, no SUOR e no queijo suíço (produzido pela Propionibacterium; a P. acnes na pele). pKa\ensuremath{\approx}4,88 (ácido fraco, como os carboxílicos \~{}4,8). SATURADO: grau de insaturação = 1 (só o C=O do carbonilo, NENHUM C=C). 5 átomos pesados (3C + 2O). DICIONÁRIO (a série gato/esquilo): o \ensuremath{-}COOH (a ponta ácida, doadora de H) \ensuremath{\leftrightarrow} o ESQUILO (a ponta reativa antissimétrica — o operador que gira, antes gambá); o pKa\ensuremath{\approx}4,88 \ensuremath{\leftrightarrow} o GAP (a acidez: um gap grande = ácido); SATURADO (0 C=C) \ensuremath{\leftrightarrow} 0 \ensuremath{\pi}/esquilo (sem duplas — o mínimo, só o carbonilo); o esqueleto C-C-C+2O (5 pesados) \ensuremath{\leftrightarrow} o GATO (a adjacência de Hückel), 5 qubits (2\ensuremath{^{5}}=32 Fock); o invariante a+c\ensuremath{^{2}}=1, a\ensuremath{\ge}a\_min>0 \ensuremath{\leftrightarrow} a ZPE da ponte de H (o mesmo de toda molécula). Na série (pelo gap/acidez): \ensuremath{-}COOH(propiônico \ensuremath{\approx}4,9 \ensuremath{\approx} o ác. 3-Me-hex-2-enoico 4,5) < fenol/\ensuremath{-}SH(\~{}10) < álcool(\~{}16). O propiônico é o \ensuremath{-}COOH puro, saturado (0 esquilo) — o vértice mínimo do ramo ácido. Fato: a fórmula, o pKa, o odor (Propionibacterium). Síntese: a leitura gato(esqueleto)/esquilo(\ensuremath{-}COOH)/gap(acidez), a série do cheiro.
\prov{relações: usa gato; usa gamba\_antissimetrico; usa invariante\_molecular\_residuo; relaciona ponte\_hidrogenio\_parabola\_geradora; relaciona gerador\_do\_gap\_comutador; usa promocao\_gamba\_esquilo}
\prov{proveniência: wiki \textperiodcentered{} Floxina (investigação) \textperiodcentered{} fato \textperiodcentered{} confiança 1.0 \textperiodcentered{} domínio floxina \textperiodcentered{} história 252 versões no jornal}

%CRISTAL {"arestas": [{"alvo": "psicologia", "confianca": 1.0, "epistemico": "fato", "origem": "wiki", "rel": "parte_de"}], "confianca": 1.0, "contraexemplos": [], "descricao": "Três coisas diferentes: a EMOÇÃO é intensa, curta e tem OBJETO (com medo DE algo); o HUMOR (mood) é difuso, longo e SEM objeto claro (acordar de bom humor); o AFETO é o guarda-chuva. O coração administra o HUMOR (o estado de fundo que colore as falas), alimentado pelas emoções pontuais de cada turno.", "epistemico": "fato", "exemplos": [], "id": "afeto_humor_emocao", "memoria": "semantica", "meta": {"dominio": "coracao", "fonte": "coracao hub"}, "origem": "wiki", "palavras_chave": [], "sinonimos": [], "tipo": "conceito", "titulo": "Afeto, humor e emoção (a distinção)"}
\section{Afeto, humor e emoção (a distinção)}
Três coisas diferentes: a EMOÇÃO é intensa, curta e tem OBJETO (com medo DE algo); o HUMOR (mood) é difuso, longo e SEM objeto claro (acordar de bom humor); o AFETO é o guarda-chuva. O coração administra o HUMOR (o estado de fundo que colore as falas), alimentado pelas emoções pontuais de cada turno.
\prov{relações: parte\_de psicologia}
\prov{proveniência: wiki \textperiodcentered{} coracao hub \textperiodcentered{} fato \textperiodcentered{} confiança 1.0 \textperiodcentered{} domínio coracao \textperiodcentered{} história 6 versões no jornal}

%CRISTAL {"arestas": [{"alvo": "omnitrix_corpo_universal", "confianca": 1.0, "epistemico": "fato", "origem": "wiki", "rel": "parte_de"}, {"alvo": "reta_matricial_compacta_cf", "confianca": 1.0, "epistemico": "fato", "origem": "wiki", "rel": "usa"}, {"alvo": "torre_dois_lados_classico_relativistico", "confianca": 1.0, "epistemico": "fato", "origem": "wiki", "rel": "usa"}, {"alvo": "multifractal_newton_omnitrix", "confianca": 1.0, "epistemico": "fato", "origem": "wiki", "rel": "usa"}, {"alvo": "o_gato_e_um_mapa_de_anosov", "confianca": 1.0, "epistemico": "fato", "origem": "wiki", "rel": "usa"}, {"alvo": "ancora_selberg_e_seu_escopo", "confianca": 1.0, "epistemico": "fato", "origem": "wiki", "rel": "relaciona"}, {"alvo": "fecho_torre_8d_7d", "confianca": 1.0, "epistemico": "fato", "origem": "wiki", "rel": "relaciona"}], "confianca": 1.0, "contraexemplos": [], "descricao": "A COSTURA que FECHA o corpo algébrico é a AGM (média aritmético-geométrica). Interleaves a PA (média ARITMÉTICA (a+b)/2=⊕, o gato, o lado CLÁSSICO/aditivo) e a PG (média GEOMÉTRICA √(ab)=⊗, o gambá, o lado RELATIVÍSTICO/multiplicativo), até que as duas MEIAM num só limite M(a,b) — ⊕ e ⊗ convergem a um, o corpo fecha. UNIFICA NEGRO E BRANCO: a=branco(matéria, o fractal branco)/b=negro(antimatéria, o fractal negro) → um só limite; |a−b|→0 QUADRÁTICO (os dígitos dobram). SOBRE π: a AGM computa π (Gauss-Legendre/Brent-Salamin, erro 1e-15) — o mesmo π da compactação (a CF de π=∏gatos). SELBERG-ANOSOV: a convergência quadrática = a contração de Anosov (o fluxo geodésico); a CF de π (o mapa de Gauss) é a mesma costura Selberg-Anosov do corpo estelar. PA/PG de ORDEM m (Gentil, as multimatrizes; Borwein cúbica/quártica) = a costura multidimensional. A AGM fecha o corpo: ⊕(clássico)/⊗(relativístico), branco/negro, num só limite sobre π. linguagem/omnitrix.py (agm_costura, 4/4).", "epistemico": "hipotese", "exemplos": [], "id": "agm_costura_fecha_corpo", "memoria": "semantica", "meta": {"dominio": "floxina", "fonte": "Floxina (investigação)"}, "origem": "wiki", "palavras_chave": [], "sinonimos": [], "tipo": "conceito", "titulo": "A costura AGM: a média aritmético-geométrica fecha o corpo algébrico (sobre π)"}
\section{A costura AGM: a média aritmético-geométrica fecha o corpo algébrico (sobre \ensuremath{\pi})}
A COSTURA que FECHA o corpo algébrico é a AGM (média aritmético-geométrica). Interleaves a PA (média ARITMÉTICA (a+b)/2=\ensuremath{\oplus}, o gato, o lado CLÁSSICO/aditivo) e a PG (média GEOMÉTRICA \ensuremath{\sqrt{\,}}(ab)=\ensuremath{\otimes}, o gambá, o lado RELATIVÍSTICO/multiplicativo), até que as duas MEIAM num só limite M(a,b) — \ensuremath{\oplus} e \ensuremath{\otimes} convergem a um, o corpo fecha. UNIFICA NEGRO E BRANCO: a=branco(matéria, o fractal branco)/b=negro(antimatéria, o fractal negro) \ensuremath{\to} um só limite; |a\ensuremath{-}b|\ensuremath{\to}0 QUADRÁTICO (os dígitos dobram). SOBRE \ensuremath{\pi}: a AGM computa \ensuremath{\pi} (Gauss-Legendre/Brent-Salamin, erro 1e-15) — o mesmo \ensuremath{\pi} da compactação (a CF de \ensuremath{\pi}=\ensuremath{\prod}gatos). SELBERG-ANOSOV: a convergência quadrática = a contração de Anosov (o fluxo geodésico); a CF de \ensuremath{\pi} (o mapa de Gauss) é a mesma costura Selberg-Anosov do corpo estelar. PA/PG de ORDEM m (Gentil, as multimatrizes; Borwein cúbica/quártica) = a costura multidimensional. A AGM fecha o corpo: \ensuremath{\oplus}(clássico)/\ensuremath{\otimes}(relativístico), branco/negro, num só limite sobre \ensuremath{\pi}. linguagem/omnitrix.py (agm\_costura, 4/4).
\prov{relações: parte\_de omnitrix\_corpo\_universal; usa reta\_matricial\_compacta\_cf; usa torre\_dois\_lados\_classico\_relativistico; usa multifractal\_newton\_omnitrix; usa o\_gato\_e\_um\_mapa\_de\_anosov; relaciona ancora\_selberg\_e\_seu\_escopo; relaciona fecho\_torre\_8d\_7d}
\prov{proveniência: wiki \textperiodcentered{} Floxina (investigação) \textperiodcentered{} hipotese \textperiodcentered{} confiança 1.0 \textperiodcentered{} domínio floxina \textperiodcentered{} história 344 versões no jornal}

%CRISTAL {"arestas": [{"alvo": "gamba_litografico_assinatura", "confianca": 1.0, "epistemico": "fato", "origem": "wiki", "rel": "relaciona"}, {"alvo": "costura_perelman_cicatrizes", "confianca": 1.0, "epistemico": "fato", "origem": "wiki", "rel": "relaciona"}, {"alvo": "invariante_topologico_assinatura", "confianca": 1.0, "epistemico": "fato", "origem": "wiki", "rel": "usa"}, {"alvo": "mass_gap_algebrico", "confianca": 1.0, "epistemico": "fato", "origem": "wiki", "rel": "relaciona"}], "confianca": 1.0, "contraexemplos": [], "descricao": "A assinatura litográfica é TOPOLÓGICA (o grau, 0/1) só até um LIMITE. A agulha de Perelman (band-limited, costura_perelman_cicatrizes) costura a assinatura na base σ_n do metal, cada vez mais perto do núcleo (o floco irracional), e NUNCA o toca (o anel de exclusão, a≥a_min>0 = o mass gap). O limite é a profundidade K_n=⌊|ln a_min|/log σ_n⌋: furar ≤ K_n grava TOPOLOGIA (o grau, robusto); FURAR MAIS que K_n põe a assinatura na GEOMETRIA (o resíduo contínuo, a posição — a cicatriz). O limite é invariante no metal: K_n·log σ_n=|ln a_min|. conversa/agulha_perelman.py (8/8).", "epistemico": "fato", "exemplos": [], "id": "agulha_perelman_limite", "memoria": "semantica", "meta": {"dominio": "floxina", "fonte": "Floxina (investigação)"}, "origem": "wiki", "palavras_chave": [], "sinonimos": [], "tipo": "conceito", "titulo": "A agulha de Perelman — o limite do invariante da assinatura"}
\section{A agulha de Perelman — o limite do invariante da assinatura}
A assinatura litográfica é TOPOLÓGICA (o grau, 0/1) só até um LIMITE. A agulha de Perelman (band-limited, costura\_perelman\_cicatrizes) costura a assinatura na base \ensuremath{\sigma}\_n do metal, cada vez mais perto do núcleo (o floco irracional), e NUNCA o toca (o anel de exclusão, a\ensuremath{\ge}a\_min>0 = o mass gap). O limite é a profundidade K\_n=\ensuremath{\lfloor}|ln a\_min|/log \ensuremath{\sigma}\_n\ensuremath{\rfloor}: furar \ensuremath{\le} K\_n grava TOPOLOGIA (o grau, robusto); FURAR MAIS que K\_n põe a assinatura na GEOMETRIA (o resíduo contínuo, a posição — a cicatriz). O limite é invariante no metal: K\_n\textperiodcentered log \ensuremath{\sigma}\_n=|ln a\_min|. conversa/agulha\_perelman.py (8/8).
\prov{relações: relaciona gamba\_litografico\_assinatura; relaciona costura\_perelman\_cicatrizes; usa invariante\_topologico\_assinatura; relaciona mass\_gap\_algebrico}
\prov{proveniência: wiki \textperiodcentered{} Floxina (investigação) \textperiodcentered{} fato \textperiodcentered{} confiança 1.0 \textperiodcentered{} domínio floxina \textperiodcentered{} história 385 versões no jornal}

%CRISTAL {"arestas": [{"alvo": "gato", "confianca": 1.0, "epistemico": "fato", "origem": "wiki", "rel": "usa"}, {"alvo": "gamba_antissimetrico", "confianca": 1.0, "epistemico": "fato", "origem": "wiki", "rel": "usa"}, {"alvo": "metais", "confianca": 1.0, "epistemico": "fato", "origem": "wiki", "rel": "relaciona"}], "confianca": 1.0, "contraexemplos": [], "descricao": "Gato e gambá são a ÁLGEBRA FUNDAMENTAL. Potências: A^k=Fibonacci [[F_{k+1},F_k],[F_k,F_{k−1}]] (a reta, hiperbólico); G^k=período 4 (G,−I,−G,I — o círculo, elíptico). JUNTOS geram os QUATÉRNIOS DIVIDIDOS (coquaternions) = M_2(ℝ): i²=−1 (o gambá), j²=+1 (2A−I, o gato sem traço), k²=+1 ([A,G]), j·i=k. Dela DERIVA tudo: (1) a RETA mineral — todo metal A_n=[[n,1],[1,0]]=A+(n−1)·E, com E=[[1,0],[0,0]]=0,4I+0,2A+0,2G−0,4(AG); (2) os POLINÔMIOS — char poly(A_n)=x²−n·x−1 (o metálico), toda quádrica é a char poly de um 2×2; (3) as OPERAÇÕES estelares (⊕ soma=matriz+, ⊗ La Hire=matriz·) e a norma N=det (−1, multiplicativa, N(AB)=N(A)N(B)). A reta, os metais e os polinômios são a sombra de UMA álgebra. conversa/gamba.py (algebra_fundamental, 6/6).", "epistemico": "fato", "exemplos": [], "id": "algebra_fundamental_gato_gamba", "memoria": "semantica", "meta": {"dominio": "floxina", "fonte": "Floxina (investigação)"}, "origem": "wiki", "palavras_chave": [], "sinonimos": [], "tipo": "conceito", "titulo": "A álgebra fundamental: gato+gambá geram os quatérnios divididos"}
\section{A álgebra fundamental: gato+gambá geram os quatérnios divididos}
Gato e gambá são a ÁLGEBRA FUNDAMENTAL. Potências: A\^{}k=Fibonacci [[F\_\{k+1\},F\_k],[F\_k,F\_\{k\ensuremath{-}1\}]] (a reta, hiperbólico); G\^{}k=período 4 (G,\ensuremath{-}I,\ensuremath{-}G,I — o círculo, elíptico). JUNTOS geram os QUATÉRNIOS DIVIDIDOS (coquaternions) = M\_2(\ensuremath{\mathbb{R}}): i\ensuremath{^{2}}=\ensuremath{-}1 (o gambá), j\ensuremath{^{2}}=+1 (2A\ensuremath{-}I, o gato sem traço), k\ensuremath{^{2}}=+1 ([A,G]), j\textperiodcentered i=k. Dela DERIVA tudo: (1) a RETA mineral — todo metal A\_n=[[n,1],[1,0]]=A+(n\ensuremath{-}1)\textperiodcentered E, com E=[[1,0],[0,0]]=0,4I+0,2A+0,2G\ensuremath{-}0,4(AG); (2) os POLINÔMIOS — char poly(A\_n)=x\ensuremath{^{2}}\ensuremath{-}n\textperiodcentered x\ensuremath{-}1 (o metálico), toda quádrica é a char poly de um 2$\times$2; (3) as OPERAÇÕES estelares (\ensuremath{\oplus} soma=matriz+, \ensuremath{\otimes} La Hire=matriz\textperiodcentered ) e a norma N=det (\ensuremath{-}1, multiplicativa, N(AB)=N(A)N(B)). A reta, os metais e os polinômios são a sombra de UMA álgebra. conversa/gamba.py (algebra\_fundamental, 6/6).
\prov{relações: usa gato; usa gamba\_antissimetrico; relaciona metais}
\prov{proveniência: wiki \textperiodcentered{} Floxina (investigação) \textperiodcentered{} fato \textperiodcentered{} confiança 1.0 \textperiodcentered{} domínio floxina \textperiodcentered{} história 296 versões no jornal}

%CRISTAL {"arestas": [{"alvo": "definicao", "confianca": 1.0, "epistemico": "fato", "origem": "curriculo", "rel": "define"}], "confianca": 1.0, "contraexemplos": [], "descricao": "| Dougherty (hipótese) | Broca-SO (fato/paper) | relação | |---------------------|------------------------|---------| | toro / vórtice | colapso Koch / grade | analogia | | fibração toroidal | fibração de Hopf S³→S² | analogia | | φ escala | gold / régua de ouro | analogia | | ψ hidrogénio | mecânica quântica | referencia externa | | onda composta por tori | gato A / espectro toro | analogia | | observador projectivo | semântica / Peirce | analogia |", "epistemico": "fato", "exemplos": [], "id": "analogia", "memoria": "semantica", "meta": {"dominio": "hipotese", "duplicata_sugerida": [], "nivel": 7, "verificacao": "parte de conceito mais amplo 'analogia'"}, "origem": "ponte_quantica_dougherty.md", "palavras_chave": ["metáfora", "mapeamento"], "sinonimos": ["analogias"], "tipo": "conceito", "titulo": "analogia"}
\section{analogia}
| Dougherty (hipótese) | Broca-SO (fato/paper) | relação | |---------------------|------------------------|---------| | toro / vórtice | colapso Koch / grade | analogia | | fibração toroidal | fibração de Hopf S\ensuremath{^{3}}\ensuremath{\to}S\ensuremath{^{2}} | analogia | | \ensuremath{\varphi} escala | gold / régua de ouro | analogia | | \ensuremath{\psi} hidrogénio | mecânica quântica | referencia externa | | onda composta por tori | gato A / espectro toro | analogia | | observador projectivo | semântica / Peirce | analogia |
\prov{sinónimos: analogias}
\prov{relações: define definicao}
\prov{proveniência: ponte\_quantica\_dougherty.md \textperiodcentered{} fato \textperiodcentered{} confiança 1.0 \textperiodcentered{} domínio hipotese \textperiodcentered{} história 4 versões no jornal}

%CRISTAL {"arestas": [{"alvo": "parabola_geradora_idempotente", "confianca": 1.0, "epistemico": "fato", "origem": "wiki", "rel": "parte_de"}, {"alvo": "gerador_do_gap_comutador", "confianca": 1.0, "epistemico": "fato", "origem": "wiki", "rel": "usa"}, {"alvo": "gamba_antissimetrico", "confianca": 1.0, "epistemico": "fato", "origem": "wiki", "rel": "usa"}, {"alvo": "corpo_matricial_matrix", "confianca": 1.0, "epistemico": "fato", "origem": "wiki", "rel": "parte_de"}, {"alvo": "invariante_estelar", "confianca": 1.0, "epistemico": "fato", "origem": "wiki", "rel": "relaciona"}], "confianca": 1.0, "contraexemplos": [], "descricao": "A MATÉRIA e a ANTIMATÉRIA são as DUAS RAÍZES da parábola do gato x²−n·x−1: σ_n=(n+√(n²+4))/2>0 (matéria) e σ_n'=(n−√(n²+4))/2<0 (antimatéria) — a mesma dinâmica de Dirac E=±√(p²+m²) (a partícula E>0, a antipartícula E<0, o mar de Dirac). TRÊS leituras fecham: (1) o PRODUTO σ_n·σ_n'=det=−1 = a ANIQUILAÇÃO para o vácuo (a raiz 0 da parábola geradora, o neutro ⊕); a soma=n=o traço. (2) o GAP=a separação √(n²+4)=os autovalores de [gato,gambá] — o gap matéria–antimatéria É o comutador (o análogo do 2m de Dirac; criar γ→e⁺e⁻ custa o gap; o vértice=a massa de repouso). (3) a CONJUGAÇÃO DE CARGA C: x↦−1/x troca σ_n↔σ_n', é o FLIP do gambá G↦−G=Gᵀ (i↦−i). A parábola geradora crua x²−x (raízes 0,1) NÃO tem antimatéria; o flip −1 abre a raiz negativa: x²−x−1 (φ, −1/φ). conversa/gamba.py (antimateria, 5/5).", "epistemico": "fato", "exemplos": [], "id": "antimateria_parabolica", "memoria": "semantica", "meta": {"dominio": "floxina", "fonte": "Floxina (investigação)"}, "origem": "wiki", "palavras_chave": [], "sinonimos": [], "tipo": "conceito", "titulo": "A antimatéria: as duas raízes da parábola do gato (a dinâmica parabólica)"}
\section{A antimatéria: as duas raízes da parábola do gato (a dinâmica parabólica)}
A MATÉRIA e a ANTIMATÉRIA são as DUAS RAÍZES da parábola do gato x\ensuremath{^{2}}\ensuremath{-}n\textperiodcentered x\ensuremath{-}1: \ensuremath{\sigma}\_n=(n+\ensuremath{\sqrt{\,}}(n\ensuremath{^{2}}+4))/2>0 (matéria) e \ensuremath{\sigma}\_n'=(n\ensuremath{-}\ensuremath{\sqrt{\,}}(n\ensuremath{^{2}}+4))/2<0 (antimatéria) — a mesma dinâmica de Dirac E=$\pm$\ensuremath{\sqrt{\,}}(p\ensuremath{^{2}}+m\ensuremath{^{2}}) (a partícula E>0, a antipartícula E<0, o mar de Dirac). TRÊS leituras fecham: (1) o PRODUTO \ensuremath{\sigma}\_n\textperiodcentered \ensuremath{\sigma}\_n'=det=\ensuremath{-}1 = a ANIQUILAÇÃO para o vácuo (a raiz 0 da parábola geradora, o neutro \ensuremath{\oplus}); a soma=n=o traço. (2) o GAP=a separação \ensuremath{\sqrt{\,}}(n\ensuremath{^{2}}+4)=os autovalores de [gato,gambá] — o gap matéria–antimatéria É o comutador (o análogo do 2m de Dirac; criar \ensuremath{\gamma}\ensuremath{\to}e\ensuremath{^{+}}e\ensuremath{^{-}} custa o gap; o vértice=a massa de repouso). (3) a CONJUGAÇÃO DE CARGA C: x\ensuremath{\mapsto}\ensuremath{-}1/x troca \ensuremath{\sigma}\_n\ensuremath{\leftrightarrow}\ensuremath{\sigma}\_n', é o FLIP do gambá G\ensuremath{\mapsto}\ensuremath{-}G=G\ensuremath{^{T}} (i\ensuremath{\mapsto}\ensuremath{-}i). A parábola geradora crua x\ensuremath{^{2}}\ensuremath{-}x (raízes 0,1) NÃO tem antimatéria; o flip \ensuremath{-}1 abre a raiz negativa: x\ensuremath{^{2}}\ensuremath{-}x\ensuremath{-}1 (\ensuremath{\varphi}, \ensuremath{-}1/\ensuremath{\varphi}). conversa/gamba.py (antimateria, 5/5).
\prov{relações: parte\_de parabola\_geradora\_idempotente; usa gerador\_do\_gap\_comutador; usa gamba\_antissimetrico; parte\_de corpo\_matricial\_matrix; relaciona invariante\_estelar}
\prov{proveniência: wiki \textperiodcentered{} Floxina (investigação) \textperiodcentered{} fato \textperiodcentered{} confiança 1.0 \textperiodcentered{} domínio floxina \textperiodcentered{} história 336 versões no jornal}

%CRISTAL {"arestas": [{"alvo": "coracao", "confianca": 1.0, "epistemico": "fato", "origem": "wiki", "rel": "parte_de"}, {"alvo": "teoria_superioridade", "confianca": 1.0, "epistemico": "fato", "origem": "wiki", "rel": "relaciona"}, {"alvo": "humor_pelo_gap", "confianca": 1.0, "epistemico": "fato", "origem": "wiki", "rel": "relaciona"}], "confianca": 1.0, "contraexemplos": [], "descricao": "O humor de rir de si mesmo: baixa a guarda, sinaliza humildade e segurança ao mesmo tempo, cria afiliação (quem ri dos próprios furos é confiável). É o registro CAOS/fail-closed da Tiffany: quando ela não sabe, admite com graça em vez de fingir — a violação (o furo) tornada benigna (assumida).", "epistemico": "fato", "exemplos": [], "id": "autodepreciacao", "memoria": "semantica", "meta": {"dominio": "coracao", "fonte": "coracao hub"}, "origem": "wiki", "palavras_chave": [], "sinonimos": [], "tipo": "conceito", "titulo": "A autodepreciação (rir de si)"}
\section{A autodepreciação (rir de si)}
O humor de rir de si mesmo: baixa a guarda, sinaliza humildade e segurança ao mesmo tempo, cria afiliação (quem ri dos próprios furos é confiável). É o registro CAOS/fail-closed da Tiffany: quando ela não sabe, admite com graça em vez de fingir — a violação (o furo) tornada benigna (assumida).
\prov{relações: parte\_de coracao; relaciona teoria\_superioridade; relaciona humor\_pelo\_gap}
\prov{proveniência: wiki \textperiodcentered{} coracao hub \textperiodcentered{} fato \textperiodcentered{} confiança 1.0 \textperiodcentered{} domínio coracao \textperiodcentered{} história 12 versões no jornal}

%CRISTAL {"arestas": [{"alvo": "reticulado", "confianca": 0.109, "epistemico": "fato", "origem": "manual", "rel": "referencia"}, {"alvo": "bai", "confianca": 0.031, "epistemico": "fato", "origem": "manual", "rel": "referencia"}, {"alvo": "bai_biblioteca_de_autorizacao_isomorfica", "confianca": 0.95, "epistemico": "fato", "origem": "manual", "rel": "parte_de"}, {"alvo": "reticulados", "confianca": 0.95, "epistemico": "fato", "origem": "manual", "rel": "usa"}], "confianca": 1.0, "contraexemplos": [], "descricao": "- APS Physics v6/58 — orbitais hidrogénio (fotoionização) - Buddy James — Dougherty Set (arte + conjectura geométrica) - Vladimir B. Ginzburg — toro-hélice (literatura externa)", "epistemico": "fato", "exemplos": [], "id": "autoridades_intrinsecas_duas_fontes_de_autoridade_nao_deri", "memoria": "semantica", "meta": {"dominio": "hipotese", "duplicata_sugerida": [], "nivel": 7, "verificacao": "parte de conceito mais amplo 'autoridades_intrinsecas_duas_fontes_de_autoridade_nao_deri'"}, "origem": "ponte_quantica_dougherty.md", "palavras_chave": [], "sinonimos": [], "tipo": "conceito", "titulo": "[Autoridades intrínsecas] Duas fontes de autoridade não deri…"}
\section{[Autoridades intrínsecas] Duas fontes de autoridade não deri…}
- APS Physics v6/58 — orbitais hidrogénio (fotoionização) - Buddy James — Dougherty Set (arte + conjectura geométrica) - Vladimir B. Ginzburg — toro-hélice (literatura externa)
\prov{relações: referencia reticulado; referencia bai; parte\_de bai\_biblioteca\_de\_autorizacao\_isomorfica; usa reticulados}
\prov{proveniência: ponte\_quantica\_dougherty.md \textperiodcentered{} fato \textperiodcentered{} confiança 1.0 \textperiodcentered{} domínio hipotese \textperiodcentered{} história 126 versões no jornal}

%CRISTAL {"arestas": [{"alvo": "psicologia", "confianca": 1.0, "epistemico": "fato", "origem": "wiki", "rel": "parte_de"}], "confianca": 1.0, "contraexemplos": [], "descricao": "Richard Lazarus: a emoção vem de uma AVALIAÇÃO (appraisal) — a mesma situação gera emoções diferentes conforme o que ela significa para você (é ameaça? desafio? perda?). A reavaliação (mudar o sentido) muda a emoção sem mudar o fato — a base da regulação por reappraisal.", "epistemico": "fato", "exemplos": [], "id": "avaliacao_lazarus", "memoria": "semantica", "meta": {"autor": "Lazarus", "dominio": "coracao", "fonte": "coracao hub"}, "origem": "wiki", "palavras_chave": [], "sinonimos": [], "tipo": "conceito", "titulo": "A avaliação cognitiva (Lazarus)"}
\section{A avaliação cognitiva (Lazarus)}
Richard Lazarus: a emoção vem de uma AVALIAÇÃO (appraisal) — a mesma situação gera emoções diferentes conforme o que ela significa para você (é ameaça? desafio? perda?). A reavaliação (mudar o sentido) muda a emoção sem mudar o fato — a base da regulação por reappraisal.
\prov{relações: parte\_de psicologia}
\prov{proveniência: wiki \textperiodcentered{} coracao hub \textperiodcentered{} fato \textperiodcentered{} confiança 1.0 \textperiodcentered{} domínio coracao \textperiodcentered{} história 6 versões no jornal}

%CRISTAL {"arestas": [{"alvo": "cosmologia_quantica", "confianca": 1.0, "epistemico": "fato", "origem": "wiki", "rel": "especializa"}, {"alvo": "bai_tese_mae", "confianca": 1.0, "epistemico": "fato", "origem": "wiki", "rel": "parte_de"}, {"alvo": "desitter_4d_universo", "confianca": 1.0, "epistemico": "fato", "origem": "wiki", "rel": "usa"}, {"alvo": "bai_t_atenuacao", "confianca": 1.0, "epistemico": "fato", "origem": "wiki", "rel": "usa"}, {"alvo": "bai_delta", "confianca": 1.0, "epistemico": "fato", "origem": "wiki", "rel": "usa"}, {"alvo": "constantes_horizonte_algebricas", "confianca": 1.0, "epistemico": "fato", "origem": "wiki", "rel": "usa"}, {"alvo": "entropia_log_phi", "confianca": 1.0, "epistemico": "fato", "origem": "wiki", "rel": "relaciona"}], "confianca": 1.0, "contraexemplos": [], "descricao": "A delegação é 'o universo de de Sitter num meio dissipativo': a MESMA geometria (o espaço fluxo-δ-campo carrega literalmente a métrica dS₄, R=12 verificado ~1e-16), com ATRITO. T (atenuação ⪯ id) é o atrito; δ é a energia. Atrito SECO/Coulomb (−1/salto, não viscoso) → vida finita EXATA τ=δ, parada abrupta. Todo universo morre: 'toda concessão finita é um universo de vida curta'. A dissipação é a física certa, não aproximação estática.", "epistemico": "inferencia", "exemplos": [], "id": "bai_cosmologia_desitter_dissipativo", "memoria": "semantica", "meta": {"autor": "broca-so", "dominio": "bai_cosmologia", "fonte": "papers/_ap_fisica_snippet.tex sec:orbits; casl-propagation.tex ap:fisica"}, "origem": "wiki", "palavras_chave": [], "sinonimos": [], "tipo": "conceito", "titulo": "BAI = de Sitter dissipativo (o universo real que morre)"}
\section{BAI = de Sitter dissipativo (o universo real que morre)}
A delegação é 'o universo de de Sitter num meio dissipativo': a MESMA geometria (o espaço fluxo-\ensuremath{\delta}-campo carrega literalmente a métrica dS\ensuremath{_{4}}, R=12 verificado \~{}1e-16), com ATRITO. T (atenuação \ensuremath{\preceq} id) é o atrito; \ensuremath{\delta} é a energia. Atrito SECO/Coulomb (\ensuremath{-}1/salto, não viscoso) \ensuremath{\to} vida finita EXATA \ensuremath{\tau}=\ensuremath{\delta}, parada abrupta. Todo universo morre: 'toda concessão finita é um universo de vida curta'. A dissipação é a física certa, não aproximação estática.
\prov{relações: especializa cosmologia\_quantica; parte\_de bai\_tese\_mae; usa desitter\_4d\_universo; usa bai\_t\_atenuacao; usa bai\_delta; usa constantes\_horizonte\_algebricas; relaciona entropia\_log\_phi}
\prov{proveniência: wiki \textperiodcentered{} papers/\_ap\_fisica\_snippet.tex sec:orbits; casl-propagation.tex ap:fisica \textperiodcentered{} inferencia \textperiodcentered{} confiança 1.0 \textperiodcentered{} domínio bai\_cosmologia \textperiodcentered{} história 8 versões no jornal}

%CRISTAL {"arestas": [{"alvo": "bai_cosmologia_desitter_dissipativo", "confianca": 1.0, "epistemico": "fato", "origem": "wiki", "rel": "usa"}, {"alvo": "mecanica_quantica", "confianca": 1.0, "epistemico": "fato", "origem": "wiki", "rel": "relaciona"}, {"alvo": "bohr_niveis_energia", "confianca": 1.0, "epistemico": "fato", "origem": "wiki", "rel": "relaciona"}, {"alvo": "bai_delta", "confianca": 1.0, "epistemico": "fato", "origem": "wiki", "rel": "explica"}, {"alvo": "colapso", "confianca": 1.0, "epistemico": "fato", "origem": "wiki", "rel": "relaciona"}, {"alvo": "bai_psi_colapso", "confianca": 1.0, "epistemico": "fato", "origem": "wiki", "rel": "relaciona"}, {"alvo": "wheeler_dewitt", "confianca": 1.0, "epistemico": "fato", "origem": "wiki", "rel": "relaciona"}], "confianca": 1.0, "contraexemplos": [], "descricao": "A leitura estritamente QUÂNTICA do BAI: δ é discreto (decresce exatamente 1/salto) — uma escada de energia quantizada ESTRUTURAL (fato). κ é a unidade que se propaga (o quantum); T+deny+monotonicidade são leis de não-escalada (superseleção). RESSALVA HONESTA: o repo trata a cosmologia como de Sitter/RG dissipativa (não MQ) e NEGA Ψ=colapso-MQ (nó `colapso`: 'não colapso MQ'); Ψ é argmax determinístico, não projeção de Born. Logo a camada 'quântica' plena é interpretação de 2ª ordem (hipótese), sustentada pela discretude de δ (fato) e pela analogia de Sitter (declarada no paper).", "epistemico": "hipotese", "exemplos": [], "id": "bai_cosmologia_quantica_tese", "memoria": "semantica", "meta": {"autor": "broca-so", "dominio": "bai_cosmologia", "fonte": "papers/bai_quinteto_hub.md; casl-propagation.tex; nó colapso"}, "origem": "wiki", "palavras_chave": [], "sinonimos": [], "tipo": "conceito", "titulo": "A tese quântica (2ª ordem): δ escada quantizada, com ressalva"}
\section{A tese quântica (2ª ordem): \ensuremath{\delta} escada quantizada, com ressalva}
A leitura estritamente QUÂNTICA do BAI: \ensuremath{\delta} é discreto (decresce exatamente 1/salto) — uma escada de energia quantizada ESTRUTURAL (fato). \ensuremath{\kappa} é a unidade que se propaga (o quantum); T+deny+monotonicidade são leis de não-escalada (superseleção). RESSALVA HONESTA: o repo trata a cosmologia como de Sitter/RG dissipativa (não MQ) e NEGA \ensuremath{\Psi}=colapso-MQ (nó `colapso`: 'não colapso MQ'); \ensuremath{\Psi} é argmax determinístico, não projeção de Born. Logo a camada 'quântica' plena é interpretação de 2ª ordem (hipótese), sustentada pela discretude de \ensuremath{\delta} (fato) e pela analogia de Sitter (declarada no paper).
\prov{relações: usa bai\_cosmologia\_desitter\_dissipativo; relaciona mecanica\_quantica; relaciona bohr\_niveis\_energia; explica bai\_delta; relaciona colapso; relaciona bai\_psi\_colapso; relaciona wheeler\_dewitt}
\prov{proveniência: wiki \textperiodcentered{} papers/bai\_quinteto\_hub.md; casl-propagation.tex; nó colapso \textperiodcentered{} hipotese \textperiodcentered{} confiança 1.0 \textperiodcentered{} domínio bai\_cosmologia \textperiodcentered{} história 8 versões no jornal}

%CRISTAL {"arestas": [{"alvo": "bai_cosmologia_desitter_dissipativo", "confianca": 1.0, "epistemico": "fato", "origem": "wiki", "rel": "especializa"}, {"alvo": "estados_estacionarios_selecao", "confianca": 1.0, "epistemico": "fato", "origem": "wiki", "rel": "relaciona"}, {"alvo": "de_sitter_lambda_energia_escura", "confianca": 1.0, "epistemico": "fato", "origem": "wiki", "rel": "relaciona"}, {"alvo": "bai_delta", "confianca": 1.0, "epistemico": "fato", "origem": "wiki", "rel": "usa"}], "confianca": 1.0, "contraexemplos": [], "descricao": "O caso sem atrito (δ=∞, reservado às autoridades intrínsecas): a energia de Noether E=cos²ψ·θ̇ (vetor de Killing θ) é CONSERVADA e a órbita é fechada/eterna. É o de Sitter conservativo — a cosmologia padrão externa (Λ eterno, estados estacionários). Verificado integrando a geodésica da seção euclidiana de dS₂ (S²): |ΔE|~1,2e-14 em 400k passos. A conservação é o LIMITE idealizado do BAI dissipativo, não o caso genérico.", "epistemico": "inferencia", "exemplos": [], "id": "bai_limite_nao_dissipativo", "memoria": "semantica", "meta": {"autor": "broca-so", "dominio": "bai_cosmologia", "fonte": "papers/_ap_fisica_snippet.tex rem:vida; verificação: geodésica dS₂ euclidiana (esta sessão)"}, "origem": "wiki", "palavras_chave": [], "sinonimos": [], "tipo": "conceito", "titulo": "Limite não-dissipativo (δ=∞) = a cosmologia padrão de fora"}
\section{Limite não-dissipativo (\ensuremath{\delta}=\ensuremath{\infty}) = a cosmologia padrão de fora}
O caso sem atrito (\ensuremath{\delta}=\ensuremath{\infty}, reservado às autoridades intrínsecas): a energia de Noether E=cos\ensuremath{^{2}}\ensuremath{\psi}\textperiodcentered \ensuremath{\theta}\ensuremath{} (vetor de Killing \ensuremath{\theta}) é CONSERVADA e a órbita é fechada/eterna. É o de Sitter conservativo — a cosmologia padrão externa (\ensuremath{\Lambda} eterno, estados estacionários). Verificado integrando a geodésica da seção euclidiana de dS\ensuremath{_{2}} (S\ensuremath{^{2}}): |\ensuremath{\Delta}E|\~{}1,2e-14 em 400k passos. A conservação é o LIMITE idealizado do BAI dissipativo, não o caso genérico.
\prov{relações: especializa bai\_cosmologia\_desitter\_dissipativo; relaciona estados\_estacionarios\_selecao; relaciona de\_sitter\_lambda\_energia\_escura; usa bai\_delta}
\prov{proveniência: wiki \textperiodcentered{} papers/\_ap\_fisica\_snippet.tex rem:vida; verificação: geodésica dS\ensuremath{_{2}} euclidiana (esta sessão) \textperiodcentered{} inferencia \textperiodcentered{} confiança 1.0 \textperiodcentered{} domínio bai\_cosmologia \textperiodcentered{} história 5 versões no jornal}

%CRISTAL {"arestas": [{"alvo": "bai_psi_colapso", "confianca": 1.0, "epistemico": "fato", "origem": "wiki", "rel": "usa"}, {"alvo": "bai_delta", "confianca": 1.0, "epistemico": "fato", "origem": "wiki", "rel": "usa"}, {"alvo": "bai_cosmologia_desitter_dissipativo", "confianca": 1.0, "epistemico": "fato", "origem": "wiki", "rel": "parte_de"}, {"alvo": "orbitais_hidrogenio_nlm", "confianca": 1.0, "epistemico": "fato", "origem": "wiki", "rel": "relaciona"}, {"alvo": "colapso_multicamada_peso", "confianca": 1.0, "epistemico": "fato", "origem": "wiki", "rel": "explica"}], "confianca": 1.0, "contraexemplos": [], "descricao": "Cada capacidade emitida carrega seu δ — logo seu raio e tempo de vida: δ diferentes = órbitas diferentes. Cada TIPO DE TRABALHO tem sua órbita (o conjunto vigente 𝒞_δ que alimenta o colapso Ψ) e sua energia (δ). É o princípio que dá o peso/prioridade por FÍSICA (energia e órbita), não por heurística — a resposta de princípio ao peso de camada do colapso.", "epistemico": "inferencia", "exemplos": [], "id": "bai_orbita_energia_por_trabalho", "memoria": "semantica", "meta": {"autor": "broca-so", "dominio": "bai_cosmologia", "fonte": "papers/_ap_fisica_snippet.tex sec:orbits; casl def:collapse"}, "origem": "wiki", "palavras_chave": [], "sinonimos": [], "tipo": "conceito", "titulo": "Órbita e energia por tipo de trabalho (κ,δ,𝒞_δ)"}
\section{Órbita e energia por tipo de trabalho (\ensuremath{\kappa},\ensuremath{\delta},\ensuremath{\mathcal{C}}\_\ensuremath{\delta})}
Cada capacidade emitida carrega seu \ensuremath{\delta} — logo seu raio e tempo de vida: \ensuremath{\delta} diferentes = órbitas diferentes. Cada TIPO DE TRABALHO tem sua órbita (o conjunto vigente \ensuremath{\mathcal{C}}\_\ensuremath{\delta} que alimenta o colapso \ensuremath{\Psi}) e sua energia (\ensuremath{\delta}). É o princípio que dá o peso/prioridade por FÍSICA (energia e órbita), não por heurística — a resposta de princípio ao peso de camada do colapso.
\prov{relações: usa bai\_psi\_colapso; usa bai\_delta; parte\_de bai\_cosmologia\_desitter\_dissipativo; relaciona orbitais\_hidrogenio\_nlm; explica colapso\_multicamada\_peso}
\prov{proveniência: wiki \textperiodcentered{} papers/\_ap\_fisica\_snippet.tex sec:orbits; casl def:collapse \textperiodcentered{} inferencia \textperiodcentered{} confiança 1.0 \textperiodcentered{} domínio bai\_cosmologia \textperiodcentered{} história 9 versões no jornal}

%CRISTAL {"arestas": [{"alvo": "coracao", "confianca": 1.0, "epistemico": "fato", "origem": "wiki", "rel": "parte_de"}], "confianca": 1.0, "contraexemplos": [], "descricao": "Peter McGraw: algo é engraçado quando é ao MESMO TEMPO uma violação (errado, ameaçador, transgressor) E benigno (seguro, ok, de brincadeira). Só violação = ofende; só benigno = tédio; os dois juntos = riso. Unifica incongruência, superioridade e alívio — e casa com o gap: a violação é o gap, o benigno é a rede embaixo.", "epistemico": "fato", "exemplos": [], "id": "benign_violation_mcgraw", "memoria": "semantica", "meta": {"autor": "McGraw", "dominio": "coracao", "fonte": "coracao hub"}, "origem": "wiki", "palavras_chave": [], "sinonimos": [], "tipo": "conceito", "titulo": "A violação benigna (McGraw)"}
\section{A violação benigna (McGraw)}
Peter McGraw: algo é engraçado quando é ao MESMO TEMPO uma violação (errado, ameaçador, transgressor) E benigno (seguro, ok, de brincadeira). Só violação = ofende; só benigno = tédio; os dois juntos = riso. Unifica incongruência, superioridade e alívio — e casa com o gap: a violação é o gap, o benigno é a rede embaixo.
\prov{relações: parte\_de coracao}
\prov{proveniência: wiki \textperiodcentered{} coracao hub \textperiodcentered{} fato \textperiodcentered{} confiança 1.0 \textperiodcentered{} domínio coracao \textperiodcentered{} história 6 versões no jornal}

%CRISTAL {"arestas": [{"alvo": "coracao", "confianca": 1.0, "epistemico": "fato", "origem": "wiki", "rel": "parte_de"}], "confianca": 1.0, "contraexemplos": [], "descricao": "Henri Bergson: rimos do 'mecânico colado sobre o vivo' — quando o humano vira máquina (a rigidez, o gesto repetido, o distraído que tropeça sempre igual). O riso é um corretivo social: pune a rigidez, pede flexibilidade. Gerenciar (girar, flexível) é o oposto do mecânico rígido que Bergson ri.", "epistemico": "fato", "exemplos": [], "id": "bergson_o_riso", "memoria": "semantica", "meta": {"autor": "Bergson", "dominio": "coracao", "fonte": "coracao hub"}, "origem": "wiki", "palavras_chave": [], "sinonimos": [], "tipo": "conceito", "titulo": "O riso (Bergson): o mecânico sobre o vivo"}
\section{O riso (Bergson): o mecânico sobre o vivo}
Henri Bergson: rimos do 'mecânico colado sobre o vivo' — quando o humano vira máquina (a rigidez, o gesto repetido, o distraído que tropeça sempre igual). O riso é um corretivo social: pune a rigidez, pede flexibilidade. Gerenciar (girar, flexível) é o oposto do mecânico rígido que Bergson ri.
\prov{relações: parte\_de coracao}
\prov{proveniência: wiki \textperiodcentered{} coracao hub \textperiodcentered{} fato \textperiodcentered{} confiança 1.0 \textperiodcentered{} domínio coracao \textperiodcentered{} história 6 versões no jornal}

%CRISTAL {"arestas": [{"alvo": "omnitrix_corpo_universal", "confianca": 1.0, "epistemico": "fato", "origem": "wiki", "rel": "parte_de"}, {"alvo": "multifractal_newton_omnitrix", "confianca": 1.0, "epistemico": "fato", "origem": "wiki", "rel": "usa"}, {"alvo": "gato", "confianca": 1.0, "epistemico": "fato", "origem": "wiki", "rel": "usa"}, {"alvo": "bsd_gap", "confianca": 1.0, "epistemico": "fato", "origem": "wiki", "rel": "relaciona"}, {"alvo": "bsd_a_intersecao_entre_as_retas_de_riemann", "confianca": 1.0, "epistemico": "fato", "origem": "wiki", "rel": "relaciona"}, {"alvo": "leminiscata_grupo_de_grupos", "confianca": 1.0, "epistemico": "fato", "origem": "wiki", "rel": "usa"}], "confianca": 1.0, "contraexemplos": [], "descricao": "O confronto do Corpo Algébrico com BSD. Para uma curva elíptica E/ℚ, cada primo p dá o FROBENIUS a_p=p+1−#E(𝔽_p), e o Frobenius é um GATO: a companheira Frob_p=[[a_p,−p],[1,0]] de z²−a_p·z+p, com autovalores α_p,β_p (α·β=p=det, α+β=a_p=traço, |α|=√p por Hasse). Assim: (1) CADA PRIMO É UM GATO (um 'bit', gerado pela redução mod p, o bit local); (2) CADA PRIMO É UM FRACTAL DE NEWTON — α_p,β_p são os pontos fixos de Newton de z²−a_p·z+p, na BORDA |z|=√p (Sato-Tate: α_p=√p·e^{iθ_p}, uma rotação — não o gato mineral, que está FORA do círculo, N=−1); (3) a GRADE DE BITS / a LEMNISCATA DISCRETA do Corpo Algébrico é o produto de Euler L(E,s)=∏_p(1−a_p·p^{−s}+p^{1−2s})⁻¹ — a grade dos Frobenius-gatos sobre os primos, o análogo DISCRETO da lemniscata. (4) BSD: posto(E(ℚ))=ord_{s=1}L(E,s) — ABERTO (milênio). Correspondência do projeto (bsd_gap): posto=nº de direções livres do cruzamento do gap, R_n=log σ_n o regulador — CORRESPONDÊNCIA, NÃO PROVA. Verificado (32a rank 0, 37a rank 1; a_p conferem com LMFDB; Hasse; Sato-Tate): linguagem/bsd_matricial.py (5/5). FATO: a_p, Hasse, o Frobenius-gato, Sato-Tate. Localizo a estrutura; BSD segue ABERTO — não decreto.", "epistemico": "hipotese", "exemplos": [], "id": "bsd_matricial_grade_bits", "memoria": "semantica", "meta": {"dominio": "floxina", "fonte": "Floxina (investigação)"}, "origem": "wiki", "palavras_chave": [], "sinonimos": [], "tipo": "conceito", "titulo": "BSD matricial: a grade de bits (Frobenius-gatos) sobre os primos, a lemniscata discreta"}
\section{BSD matricial: a grade de bits (Frobenius-gatos) sobre os primos, a lemniscata discreta}
O confronto do Corpo Algébrico com BSD. Para uma curva elíptica E/\ensuremath{\mathbb{Q}}, cada primo p dá o FROBENIUS a\_p=p+1\ensuremath{-}\#E(\ensuremath{\mathbb{F}}\_p), e o Frobenius é um GATO: a companheira Frob\_p=[[a\_p,\ensuremath{-}p],[1,0]] de z\ensuremath{^{2}}\ensuremath{-}a\_p\textperiodcentered z+p, com autovalores \ensuremath{\alpha}\_p,\ensuremath{\beta}\_p (\ensuremath{\alpha}\textperiodcentered \ensuremath{\beta}=p=det, \ensuremath{\alpha}+\ensuremath{\beta}=a\_p=traço, |\ensuremath{\alpha}|=\ensuremath{\sqrt{\,}}p por Hasse). Assim: (1) CADA PRIMO É UM GATO (um 'bit', gerado pela redução mod p, o bit local); (2) CADA PRIMO É UM FRACTAL DE NEWTON — \ensuremath{\alpha}\_p,\ensuremath{\beta}\_p são os pontos fixos de Newton de z\ensuremath{^{2}}\ensuremath{-}a\_p\textperiodcentered z+p, na BORDA |z|=\ensuremath{\sqrt{\,}}p (Sato-Tate: \ensuremath{\alpha}\_p=\ensuremath{\sqrt{\,}}p\textperiodcentered e\^{}\{i\ensuremath{\theta}\_p\}, uma rotação — não o gato mineral, que está FORA do círculo, N=\ensuremath{-}1); (3) a GRADE DE BITS / a LEMNISCATA DISCRETA do Corpo Algébrico é o produto de Euler L(E,s)=\ensuremath{\prod}\_p(1\ensuremath{-}a\_p\textperiodcentered p\^{}\{\ensuremath{-}s\}+p\^{}\{1\ensuremath{-}2s\})\ensuremath{^{-1}} — a grade dos Frobenius-gatos sobre os primos, o análogo DISCRETO da lemniscata. (4) BSD: posto(E(\ensuremath{\mathbb{Q}}))=ord\_\{s=1\}L(E,s) — ABERTO (milênio). Correspondência do projeto (bsd\_gap): posto=nº de direções livres do cruzamento do gap, R\_n=log \ensuremath{\sigma}\_n o regulador — CORRESPONDÊNCIA, NÃO PROVA. Verificado (32a rank 0, 37a rank 1; a\_p conferem com LMFDB; Hasse; Sato-Tate): linguagem/bsd\_matricial.py (5/5). FATO: a\_p, Hasse, o Frobenius-gato, Sato-Tate. Localizo a estrutura; BSD segue ABERTO — não decreto.
\prov{relações: parte\_de omnitrix\_corpo\_universal; usa multifractal\_newton\_omnitrix; usa gato; relaciona bsd\_gap; relaciona bsd\_a\_intersecao\_entre\_as\_retas\_de\_riemann; usa leminiscata\_grupo\_de\_grupos}
\prov{proveniência: wiki \textperiodcentered{} Floxina (investigação) \textperiodcentered{} hipotese \textperiodcentered{} confiança 1.0 \textperiodcentered{} domínio floxina \textperiodcentered{} história 322 versões no jornal}

%CRISTAL {"arestas": [{"alvo": "fisica_corpo_matricial_n_qubits", "confianca": 1.0, "epistemico": "fato", "origem": "wiki", "rel": "parte_de"}, {"alvo": "gato", "confianca": 1.0, "epistemico": "fato", "origem": "wiki", "rel": "usa"}, {"alvo": "gamba_antissimetrico", "confianca": 1.0, "epistemico": "fato", "origem": "wiki", "rel": "usa"}, {"alvo": "biblioteca_pipe", "confianca": 1.0, "epistemico": "fato", "origem": "wiki", "rel": "relaciona"}], "confianca": 1.0, "contraexemplos": [], "descricao": "A divisão-mestra Sym⊕Skew (gato/gambá) É a dos campos da física. A MÉTRICA g_μν é SIMÉTRICA (g=gᵀ) = GRAVITAÇÃO = o GATO; o TENSOR EM F_μν é ANTI-SIMÉTRICO (F=−Fᵀ) = ELETROMAGNETISMO = o GAMBÁ. O campo total de um espaço-tempo é g⊕F = Sym⊕Skew = gato⊕gambá (Kaluza-Klein: a métrica em D+1 parte-se em g [gravidade] + A [o EM, a fase U(1), o gambá]). Fecha as 8 faces no corpo matricial: geometria, mecânica (os n-qubits), óptica (reflexão=gato/refração=gambá), termo (Onsager), EM (F=gambá), quântica (os n-qubits+Schrödinger unitário), gravitação (g=gato), matéria (a carga na torre clássica/o octônion, o mass gap 𝔤₂=14). Uma só divisão: simétrico(gato)/anti-simétrico(gambá) — gravidade/EM, métrica/campo, medir/girar. FATO: g simétrica e F anti-simétrica; SÍNTESE: a identificação com gato/gambá.", "epistemico": "hipotese", "exemplos": [], "id": "campos_gravidade_gato_em_gamba", "memoria": "semantica", "meta": {"dominio": "floxina", "fonte": "Floxina (investigação)"}, "origem": "wiki", "palavras_chave": [], "sinonimos": [], "tipo": "conceito", "titulo": "Gravitação=o gato (g simétrica), eletromagnetismo=o gambá (F anti-sim.)"}
\section{Gravitação=o gato (g simétrica), eletromagnetismo=o gambá (F anti-sim.)}
A divisão-mestra Sym\ensuremath{\oplus}Skew (gato/gambá) É a dos campos da física. A MÉTRICA g\_\ensuremath{\mu}\ensuremath{\nu} é SIMÉTRICA (g=g\ensuremath{^{T}}) = GRAVITAÇÃO = o GATO; o TENSOR EM F\_\ensuremath{\mu}\ensuremath{\nu} é ANTI-SIMÉTRICO (F=\ensuremath{-}F\ensuremath{^{T}}) = ELETROMAGNETISMO = o GAMBÁ. O campo total de um espaço-tempo é g\ensuremath{\oplus}F = Sym\ensuremath{\oplus}Skew = gato\ensuremath{\oplus}gambá (Kaluza-Klein: a métrica em D+1 parte-se em g [gravidade] + A [o EM, a fase U(1), o gambá]). Fecha as 8 faces no corpo matricial: geometria, mecânica (os n-qubits), óptica (reflexão=gato/refração=gambá), termo (Onsager), EM (F=gambá), quântica (os n-qubits+Schrödinger unitário), gravitação (g=gato), matéria (a carga na torre clássica/o octônion, o mass gap \ensuremath{\mathfrak{g}}\ensuremath{_{2}}=14). Uma só divisão: simétrico(gato)/anti-simétrico(gambá) — gravidade/EM, métrica/campo, medir/girar. FATO: g simétrica e F anti-simétrica; SÍNTESE: a identificação com gato/gambá.
\prov{relações: parte\_de fisica\_corpo\_matricial\_n\_qubits; usa gato; usa gamba\_antissimetrico; relaciona biblioteca\_pipe}
\prov{proveniência: wiki \textperiodcentered{} Floxina (investigação) \textperiodcentered{} hipotese \textperiodcentered{} confiança 1.0 \textperiodcentered{} domínio floxina \textperiodcentered{} história 320 versões no jornal}

%CRISTAL {"arestas": [{"alvo": "gamba_litografico_assinatura", "confianca": 1.0, "epistemico": "fato", "origem": "wiki", "rel": "usa"}, {"alvo": "agulha_perelman_limite", "confianca": 1.0, "epistemico": "fato", "origem": "wiki", "rel": "usa"}, {"alvo": "pipe_transmissao_dual", "confianca": 1.0, "epistemico": "fato", "origem": "wiki", "rel": "parte_de"}], "confianca": 1.0, "contraexemplos": [], "descricao": "O pipe da lib abre e fecha as CHAVES de assinatura nos canais de comunicação. abrir_canal(nome, chave, metal) estabelece a chave costurada pela agulha de Perelman (o metal fixa a profundidade K_n = a robustez); cada transmitir_canal CARREGA (o gambá, checksum a+c²=1) e ASSINA (gato cifra + gambá orienta), encadeando o MAC; fechar_canal assina o fecho sobre o MAC acumulado e verifica (Pf=+1). Chave errada no fecho falha. transmissao.abrir_canal/transmitir_canal/fechar_canal — toda troca no canal é carregada e assinada, a chave abre e fecha pela agulha.", "epistemico": "fato", "exemplos": [], "id": "canais_chave_assinatura", "memoria": "semantica", "meta": {"dominio": "floxina", "fonte": "Floxina (investigação)"}, "origem": "wiki", "palavras_chave": [], "sinonimos": [], "tipo": "conceito", "titulo": "Os canais — abrir/fechar a chave de assinatura (a agulha costura)"}
\section{Os canais — abrir/fechar a chave de assinatura (a agulha costura)}
O pipe da lib abre e fecha as CHAVES de assinatura nos canais de comunicação. abrir\_canal(nome, chave, metal) estabelece a chave costurada pela agulha de Perelman (o metal fixa a profundidade K\_n = a robustez); cada transmitir\_canal CARREGA (o gambá, checksum a+c\ensuremath{^{2}}=1) e ASSINA (gato cifra + gambá orienta), encadeando o MAC; fechar\_canal assina o fecho sobre o MAC acumulado e verifica (Pf=+1). Chave errada no fecho falha. transmissao.abrir\_canal/transmitir\_canal/fechar\_canal — toda troca no canal é carregada e assinada, a chave abre e fecha pela agulha.
\prov{relações: usa gamba\_litografico\_assinatura; usa agulha\_perelman\_limite; parte\_de pipe\_transmissao\_dual}
\prov{proveniência: wiki \textperiodcentered{} Floxina (investigação) \textperiodcentered{} fato \textperiodcentered{} confiança 1.0 \textperiodcentered{} domínio floxina \textperiodcentered{} história 308 versões no jornal}

%CRISTAL {"arestas": [{"alvo": "omnitrix_corpo_universal", "confianca": 1.0, "epistemico": "fato", "origem": "wiki", "rel": "parte_de"}, {"alvo": "estojo_agulhas_minerais_fixo", "confianca": 1.0, "epistemico": "fato", "origem": "wiki", "rel": "usa"}, {"alvo": "canais_chave_assinatura", "confianca": 1.0, "epistemico": "fato", "origem": "wiki", "rel": "usa"}, {"alvo": "gamba_litografico_assinatura", "confianca": 1.0, "epistemico": "fato", "origem": "wiki", "rel": "usa"}, {"alvo": "invariante_molecular_residuo", "confianca": 1.0, "epistemico": "fato", "origem": "wiki", "rel": "usa"}, {"alvo": "multifractal_newton_omnitrix", "confianca": 1.0, "epistemico": "fato", "origem": "wiki", "rel": "relaciona"}, {"alvo": "agm_costura_fecha_corpo", "confianca": 1.0, "epistemico": "fato", "origem": "wiki", "rel": "relaciona"}], "confianca": 1.0, "contraexemplos": [], "descricao": "O Corpo Algébrico define formalmente a CHAVE de cada corpo molecular. A chave vem do ESTOJO DE AGULHAS (os metais 1,2,3,4,5, as cúspides; NÃO inventa agulha nova), codificada na ALTURA K_n=⌊|ln a_min|/log σ_n⌋ de cada agulha. Cada corpo tem o seu a_min (o invariante a≥a_min>0, a ZPE=o resíduo octônio mínimo); a sua chave é (K_1..K_5). Corpos distintos → chaves distintas (tiol (28,15,11,9,8), THC (30,16,12,10,8)). A chave: (a) IMPRIME ASSINATURAS nos corpos moleculares (a litografia: assinar→verificar OK, adulterar→o selo quebra); (b) ABRE E FECHA os canais do LADO NEGRO da comunicação (o fractal negro/a antimatéria). A DINÂMICA: quando o corpo se transforma reversivelmente (o gambá — geraniol→citronelol, +H₂), a parábola muda (a_min), a chave muda; MAS tudo reversível SEM DISSIPAÇÃO (o gambá conserva a norma) — transmite e recupera intacta, a chave 'continua a mesma' (recuperável), REGIDA pela parábola x²−m·x−1 do corpo. linguagem/omnitrix.py (chave_corpo_molecular 4/4, dinamica_chave 3/3); agulha_perelman.py, litografia.py.", "epistemico": "hipotese", "exemplos": [], "id": "chave_corpo_molecular_estojo", "memoria": "semantica", "meta": {"dominio": "floxina", "fonte": "Floxina (investigação)"}, "origem": "wiki", "palavras_chave": [], "sinonimos": [], "tipo": "conceito", "titulo": "A chave do corpo molecular: o estojo de agulhas, a litografia, o lado negro, a dinâmica"}
\section{A chave do corpo molecular: o estojo de agulhas, a litografia, o lado negro, a dinâmica}
O Corpo Algébrico define formalmente a CHAVE de cada corpo molecular. A chave vem do ESTOJO DE AGULHAS (os metais 1,2,3,4,5, as cúspides; NÃO inventa agulha nova), codificada na ALTURA K\_n=\ensuremath{\lfloor}|ln a\_min|/log \ensuremath{\sigma}\_n\ensuremath{\rfloor} de cada agulha. Cada corpo tem o seu a\_min (o invariante a\ensuremath{\ge}a\_min>0, a ZPE=o resíduo octônio mínimo); a sua chave é (K\_1..K\_5). Corpos distintos \ensuremath{\to} chaves distintas (tiol (28,15,11,9,8), THC (30,16,12,10,8)). A chave: (a) IMPRIME ASSINATURAS nos corpos moleculares (a litografia: assinar\ensuremath{\to}verificar OK, adulterar\ensuremath{\to}o selo quebra); (b) ABRE E FECHA os canais do LADO NEGRO da comunicação (o fractal negro/a antimatéria). A DINÂMICA: quando o corpo se transforma reversivelmente (o gambá — geraniol\ensuremath{\to}citronelol, +H\ensuremath{_{2}}), a parábola muda (a\_min), a chave muda; MAS tudo reversível SEM DISSIPAÇÃO (o gambá conserva a norma) — transmite e recupera intacta, a chave 'continua a mesma' (recuperável), REGIDA pela parábola x\ensuremath{^{2}}\ensuremath{-}m\textperiodcentered x\ensuremath{-}1 do corpo. linguagem/omnitrix.py (chave\_corpo\_molecular 4/4, dinamica\_chave 3/3); agulha\_perelman.py, litografia.py.
\prov{relações: parte\_de omnitrix\_corpo\_universal; usa estojo\_agulhas\_minerais\_fixo; usa canais\_chave\_assinatura; usa gamba\_litografico\_assinatura; usa invariante\_molecular\_residuo; relaciona multifractal\_newton\_omnitrix; relaciona agm\_costura\_fecha\_corpo}
\prov{proveniência: wiki \textperiodcentered{} Floxina (investigação) \textperiodcentered{} hipotese \textperiodcentered{} confiança 1.0 \textperiodcentered{} domínio floxina \textperiodcentered{} história 320 versões no jornal}

%CRISTAL {"arestas": [{"alvo": "coracao", "confianca": 1.0, "epistemico": "fato", "origem": "wiki", "rel": "parte_de"}, {"alvo": "residuo_2qubit_checkpoint", "confianca": 1.0, "epistemico": "fato", "origem": "wiki", "rel": "relaciona"}, {"alvo": "dimensoes_do_pipe", "confianca": 1.0, "epistemico": "fato", "origem": "wiki", "rel": "relaciona"}, {"alvo": "chicote_e_o_crivo_de_selberg", "confianca": 1.0, "epistemico": "fato", "origem": "wiki", "rel": "relaciona"}], "confianca": 1.0, "contraexemplos": [], "descricao": "O estado emocional é lido como n-qubits (3 qubits: 2³=8 amplitudes, |ψ|²=1) — o CHICOTE que estala entre turnos, o MESMO objeto do resíduo da mineração (residuo_2qubit generalizado). A continuidade emocional da Tiffany É esse resíduo: a emoção não reinicia a cada fala, ela CARREGA no chicote, como o floco carrega o resíduo φ de um passo ao outro. O octônion é de divisão ⟹ reversível ⟹ nada se perde.", "epistemico": "fato", "exemplos": [], "id": "chicote_emocional", "memoria": "semantica", "meta": {"dominio": "coracao", "fonte": "coracao hub"}, "origem": "wiki", "palavras_chave": [], "sinonimos": [], "tipo": "conceito", "titulo": "O chicote emocional (os n-qubits que carregam a emoção)"}
\section{O chicote emocional (os n-qubits que carregam a emoção)}
O estado emocional é lido como n-qubits (3 qubits: 2\ensuremath{^{3}}=8 amplitudes, |\ensuremath{\psi}|\ensuremath{^{2}}=1) — o CHICOTE que estala entre turnos, o MESMO objeto do resíduo da mineração (residuo\_2qubit generalizado). A continuidade emocional da Tiffany É esse resíduo: a emoção não reinicia a cada fala, ela CARREGA no chicote, como o floco carrega o resíduo \ensuremath{\varphi} de um passo ao outro. O octônion é de divisão \ensuremath{\Longrightarrow} reversível \ensuremath{\Longrightarrow} nada se perde.
\prov{relações: parte\_de coracao; relaciona residuo\_2qubit\_checkpoint; relaciona dimensoes\_do\_pipe; relaciona chicote\_e\_o\_crivo\_de\_selberg}
\prov{proveniência: wiki \textperiodcentered{} coracao hub \textperiodcentered{} fato \textperiodcentered{} confiança 1.0 \textperiodcentered{} domínio coracao \textperiodcentered{} história 15 versões no jornal}

%CRISTAL {"arestas": [{"alvo": "psicologia", "confianca": 1.0, "epistemico": "fato", "origem": "wiki", "rel": "parte_de"}], "confianca": 1.0, "contraexemplos": [], "descricao": "James Russell: toda emoção mora num plano de dois eixos — VALÊNCIA (agradável↔desagradável) e ATIVAÇÃO (alta↔baixa energia). Alegria = valência+/ativação alta; serenidade = valência+/ativação baixa; medo = valência−/ativação alta; tristeza = valência−/ativação baixa. O afeto é dimensional, não uma lista de rótulos.", "epistemico": "fato", "exemplos": [], "id": "circumplexo_russell", "memoria": "semantica", "meta": {"autor": "Russell", "dominio": "coracao", "fonte": "coracao hub"}, "origem": "wiki", "palavras_chave": [], "sinonimos": [], "tipo": "conceito", "titulo": "O circumplexo do afeto (Russell)"}
\section{O circumplexo do afeto (Russell)}
James Russell: toda emoção mora num plano de dois eixos — VALÊNCIA (agradável\ensuremath{\leftrightarrow}desagradável) e ATIVAÇÃO (alta\ensuremath{\leftrightarrow}baixa energia). Alegria = valência+/ativação alta; serenidade = valência+/ativação baixa; medo = valência\ensuremath{-}/ativação alta; tristeza = valência\ensuremath{-}/ativação baixa. O afeto é dimensional, não uma lista de rótulos.
\prov{relações: parte\_de psicologia}
\prov{proveniência: wiki \textperiodcentered{} coracao hub \textperiodcentered{} fato \textperiodcentered{} confiança 1.0 \textperiodcentered{} domínio coracao \textperiodcentered{} história 6 versões no jornal}

%CRISTAL {"arestas": [{"alvo": "computador_mineral", "confianca": 1.0, "epistemico": "fato", "origem": "wiki", "rel": "relaciona"}, {"alvo": "computador_universal_pnp", "confianca": 1.0, "epistemico": "fato", "origem": "wiki", "rel": "parte_de"}, {"alvo": "omnitrix_corpo_universal", "confianca": 1.0, "epistemico": "fato", "origem": "wiki", "rel": "relaciona"}, {"alvo": "pipe_transmissao_dual", "confianca": 1.0, "epistemico": "fato", "origem": "wiki", "rel": "usa"}, {"alvo": "broca_so_no_metal", "confianca": 1.0, "epistemico": "fato", "origem": "wiki", "rel": "relaciona"}, {"alvo": "gamba_antissimetrico", "confianca": 1.0, "epistemico": "fato", "origem": "wiki", "rel": "usa"}], "confianca": 1.0, "contraexemplos": [], "descricao": "O COMPUTADOR MINERAL (gato=ULA, cristais=memória, molécula=registrador, órbita=canal) é o COMPUTADOR UNIVERSAL (o Omnitrix, a UTM de corpos) INSTANCIADO no substrato do METAL — uma instância, como a matrix (ordem 2) ou a molécula (química). O DICIONÁRIO: gato=a ULA (A_n, ×A_n=1 passo O(1), det=−1) ↔ o gato mede/lê (o Hamiltoniano); cristais=a memória (norma conservada=integridade) ↔ os n-qubits (|ψ|²=1, o checksum); molécula=o registrador ↔ as instâncias=os programas; órbita=o canal ↔ o PIPE; reversível (𝕆 divisão) ↔ o gambá (unitário). A MIGRAÇÃO DO PIPE NO METAL: todo o pipe (linguagem/transmissao.py) já roda no metal — o portador é o octônio 𝕆 (o MESMO da mineração), a órbita É o canal mineral, a norma conservada É o checksum (a+c²=1)=a integridade dos cristais; reversível por alternatividade (u⁻¹·(u·x)=x). Verificado por round-trip (transmitir→receber intacto, checksum batendo). O mineral é o universal no metal; o pipe é a órbita. linguagem/omnitrix.py (computador_mineral_universal, 4/4).", "epistemico": "hipotese", "exemplos": [], "id": "computador_mineral_promovido", "memoria": "semantica", "meta": {"dominio": "floxina", "fonte": "Floxina (investigação)"}, "origem": "wiki", "palavras_chave": [], "sinonimos": [], "tipo": "conceito", "titulo": "A promoção: o computador mineral → o computador universal (no metal); o pipe migra para a órbita"}
\section{A promoção: o computador mineral \ensuremath{\to} o computador universal (no metal); o pipe migra para a órbita}
O COMPUTADOR MINERAL (gato=ULA, cristais=memória, molécula=registrador, órbita=canal) é o COMPUTADOR UNIVERSAL (o Omnitrix, a UTM de corpos) INSTANCIADO no substrato do METAL — uma instância, como a matrix (ordem 2) ou a molécula (química). O DICIONÁRIO: gato=a ULA (A\_n, $\times$A\_n=1 passo O(1), det=\ensuremath{-}1) \ensuremath{\leftrightarrow} o gato mede/lê (o Hamiltoniano); cristais=a memória (norma conservada=integridade) \ensuremath{\leftrightarrow} os n-qubits (|\ensuremath{\psi}|\ensuremath{^{2}}=1, o checksum); molécula=o registrador \ensuremath{\leftrightarrow} as instâncias=os programas; órbita=o canal \ensuremath{\leftrightarrow} o PIPE; reversível (\ensuremath{\mathbb{O}} divisão) \ensuremath{\leftrightarrow} o gambá (unitário). A MIGRAÇÃO DO PIPE NO METAL: todo o pipe (linguagem/transmissao.py) já roda no metal — o portador é o octônio \ensuremath{\mathbb{O}} (o MESMO da mineração), a órbita É o canal mineral, a norma conservada É o checksum (a+c\ensuremath{^{2}}=1)=a integridade dos cristais; reversível por alternatividade (u\ensuremath{^{-1}}\textperiodcentered (u\textperiodcentered x)=x). Verificado por round-trip (transmitir\ensuremath{\to}receber intacto, checksum batendo). O mineral é o universal no metal; o pipe é a órbita. linguagem/omnitrix.py (computador\_mineral\_universal, 4/4).
\prov{relações: relaciona computador\_mineral; parte\_de computador\_universal\_pnp; relaciona omnitrix\_corpo\_universal; usa pipe\_transmissao\_dual; relaciona broca\_so\_no\_metal; usa gamba\_antissimetrico}
\prov{proveniência: wiki \textperiodcentered{} Floxina (investigação) \textperiodcentered{} hipotese \textperiodcentered{} confiança 1.0 \textperiodcentered{} domínio floxina \textperiodcentered{} história 287 versões no jornal}

%CRISTAL {"arestas": [{"alvo": "omnitrix_corpo_universal", "confianca": 1.0, "epistemico": "fato", "origem": "wiki", "rel": "parte_de"}, {"alvo": "gamba_antissimetrico", "confianca": 1.0, "epistemico": "fato", "origem": "wiki", "rel": "usa"}, {"alvo": "gerador_do_gap_comutador", "confianca": 1.0, "epistemico": "fato", "origem": "wiki", "rel": "usa"}, {"alvo": "dicionario_bai_omnitrix", "confianca": 1.0, "epistemico": "fato", "origem": "wiki", "rel": "relaciona"}, {"alvo": "floxina_quantum_materia_escura", "confianca": 1.0, "epistemico": "fato", "origem": "wiki", "rel": "relaciona"}, {"alvo": "hierarquia_pspace_indecidivel", "confianca": 1.0, "epistemico": "fato", "origem": "wiki", "rel": "relaciona"}, {"alvo": "classes_p_np_cook_levin", "confianca": 1.0, "epistemico": "fato", "origem": "wiki", "rel": "usa"}, {"alvo": "cortes_dedekind", "confianca": 1.0, "epistemico": "fato", "origem": "wiki", "rel": "usa"}, {"alvo": "gato_e_o_hamiltoniano", "confianca": 1.0, "epistemico": "fato", "origem": "wiki", "rel": "relaciona"}], "confianca": 1.0, "contraexemplos": [], "descricao": "O Omnitrix é um COMPUTADOR UNIVERSAL: o GATO lê/mede (Hermitiano, a cabeça); o GAMBÁ evolui (unitário, o portão reversível); os n-QUBITS são a MEMÓRIA (2ⁿ); as INSTÂNCIAS são os PROGRAMAS. Gato+gambá geram M_2 → com os n-qubits, computação universal = uma UTM (máquina de Turing universal) de corpos. DICIONÁRIO com P vs NP: o gambá (reversível/unitário, determinístico) ↔ P (poly-tempo); o gato (medir/colapsar, escolher 1 entre 2ⁿ) ↔ NP (a busca, o não-determinismo); o GAP [gato,gambá]≠0 (buscar≠verificar) ↔ o gap P–NP; o exponencial 2ⁿ ↔ o espaço de busca de NP; o mass gap (o associador do 𝕆≠0) ↔ a separação P–NP. A DECLARAÇÃO (o teorema da unicidade da reta, corpo_estelar: 0,999…≠1): todo número mora numa reta ÚNICA. P é o PROCESSO (o esquilo reversível seguindo a reta — a fase subindo a 1⁻ SEM cruzar a costura); NP é o CRISTALIZADO (o gato colapsa, cruzando a costura ao 1 da PRÓXIMA reta). Como 0,999…(reta 0)≠1(reta 1), o processo≠o cristalizado: NO CORPO ESTELAR, P≠NP — com o MESMO estatuto de 0,999…≠1, uma DECLARAÇÃO DE SÍNTESE do frame DINÂMICO. O milênio CLÁSSICO é o frame CRISTALIZADO (a completude que cruza a costura, onde P=NP seria congelar o processo — como 0,999…=1); ali segue aberto. São dois frames; a distinção dinâmico/cristalizado é o que o teorema da unicidade declara. (Corrige a versão anterior 'não se decreta' — a unicidade DECLARA no frame dinâmico.) linguagem/omnitrix.py (dicionario_pnp_turing 3/3).", "epistemico": "hipotese", "exemplos": [], "id": "computador_universal_pnp", "memoria": "semantica", "meta": {"dominio": "floxina", "fonte": "Floxina (investigação)"}, "origem": "wiki", "palavras_chave": [], "sinonimos": [], "tipo": "conceito", "titulo": "O computador universal e o dicionário P vs NP / a máquina de Turing"}
\section{O computador universal e o dicionário P vs NP / a máquina de Turing}
O Omnitrix é um COMPUTADOR UNIVERSAL: o GATO lê/mede (Hermitiano, a cabeça); o GAMBÁ evolui (unitário, o portão reversível); os n-QUBITS são a MEMÓRIA (2\ensuremath{^{n}}); as INSTÂNCIAS são os PROGRAMAS. Gato+gambá geram M\_2 \ensuremath{\to} com os n-qubits, computação universal = uma UTM (máquina de Turing universal) de corpos. DICIONÁRIO com P vs NP: o gambá (reversível/unitário, determinístico) \ensuremath{\leftrightarrow} P (poly-tempo); o gato (medir/colapsar, escolher 1 entre 2\ensuremath{^{n}}) \ensuremath{\leftrightarrow} NP (a busca, o não-determinismo); o GAP [gato,gambá]\ensuremath{\ne}0 (buscar\ensuremath{\ne}verificar) \ensuremath{\leftrightarrow} o gap P–NP; o exponencial 2\ensuremath{^{n}} \ensuremath{\leftrightarrow} o espaço de busca de NP; o mass gap (o associador do \ensuremath{\mathbb{O}}\ensuremath{\ne}0) \ensuremath{\leftrightarrow} a separação P–NP. A DECLARAÇÃO (o teorema da unicidade da reta, corpo\_estelar: 0,999…\ensuremath{\ne}1): todo número mora numa reta ÚNICA. P é o PROCESSO (o esquilo reversível seguindo a reta — a fase subindo a 1\ensuremath{^{-}} SEM cruzar a costura); NP é o CRISTALIZADO (o gato colapsa, cruzando a costura ao 1 da PRÓXIMA reta). Como 0,999…(reta 0)\ensuremath{\ne}1(reta 1), o processo\ensuremath{\ne}o cristalizado: NO CORPO ESTELAR, P\ensuremath{\ne}NP — com o MESMO estatuto de 0,999…\ensuremath{\ne}1, uma DECLARAÇÃO DE SÍNTESE do frame DINÂMICO. O milênio CLÁSSICO é o frame CRISTALIZADO (a completude que cruza a costura, onde P=NP seria congelar o processo — como 0,999…=1); ali segue aberto. São dois frames; a distinção dinâmico/cristalizado é o que o teorema da unicidade declara. (Corrige a versão anterior 'não se decreta' — a unicidade DECLARA no frame dinâmico.) linguagem/omnitrix.py (dicionario\_pnp\_turing 3/3).
\prov{relações: parte\_de omnitrix\_corpo\_universal; usa gamba\_antissimetrico; usa gerador\_do\_gap\_comutador; relaciona dicionario\_bai\_omnitrix; relaciona floxina\_quantum\_materia\_escura; relaciona hierarquia\_pspace\_indecidivel; usa classes\_p\_np\_cook\_levin; usa cortes\_dedekind; relaciona gato\_e\_o\_hamiltoniano}
\prov{proveniência: wiki \textperiodcentered{} Floxina (investigação) \textperiodcentered{} hipotese \textperiodcentered{} confiança 1.0 \textperiodcentered{} domínio floxina \textperiodcentered{} história 507 versões no jornal}

%CRISTAL {"arestas": [{"alvo": "omnitrix_corpo_universal", "confianca": 1.0, "epistemico": "fato", "origem": "wiki", "rel": "parte_de"}, {"alvo": "bsd_matricial_grade_bits", "confianca": 1.0, "epistemico": "fato", "origem": "wiki", "rel": "usa"}, {"alvo": "multifractal_newton_omnitrix", "confianca": 1.0, "epistemico": "fato", "origem": "wiki", "rel": "usa"}, {"alvo": "reta_matricial_compacta_cf", "confianca": 1.0, "epistemico": "fato", "origem": "wiki", "rel": "usa"}, {"alvo": "gerador_do_gap_comutador", "confianca": 1.0, "epistemico": "fato", "origem": "wiki", "rel": "usa"}, {"alvo": "floxina_quantum_materia_escura", "confianca": 1.0, "epistemico": "fato", "origem": "wiki", "rel": "relaciona"}, {"alvo": "ancora_selberg_e_seu_escopo", "confianca": 1.0, "epistemico": "fato", "origem": "wiki", "rel": "relaciona"}], "confianca": 1.0, "contraexemplos": [], "descricao": "No CORPO ALGÉBRICO COMPACTO (a superfície modular compacta, o lado do fractal negro/o corpo negro) os quatro problemas restantes do milênio num só quadro. HODGE=o corpo compacto: a curva modular X₀(N) é compacta projetiva (Kähler), o palco de Hodge; Lefschetz (1,1) dá os divisores (provado). RH=as LINHAS das parábolas: λ=¼−δ², λ≥¼ ⟺ δ imaginário ⟺ s=½+it na reta crítica. BSD=as CÚSPIDES DISCRETAS: as cúspides de X₀(N) (#=Σ_{d|N}φ(gcd(d,N/d)): X₀(1)=1, X₀(11)=2, X₀(32)=8), onde as L-funções cruzam; BSD vive nas cúspides. YANG-MILLS=a DINÂMICA do gap: o gap espectral λ₁>0 (Selberg λ₁≥3/16 PROVADO que existe; ≥¼ conjectural = a mesma parábola de RH); a dinâmica é o fluxo geodésico com esse gap. DISCIPLINA: afirmo só o provado (nº de cúspides, Hasse-Weil, Selberg 3/16, Lefschetz (1,1), a compacidade); os QUATRO (RH/GRH, BSD, Yang-Mills, Hodge) seguem ABERTOS. A síntese que os costura é DECLARADA, não prova — não decreto. linguagem/confronto_millennium.py (5/5).", "epistemico": "hipotese", "exemplos": [], "id": "confronto_millennium_compacto", "memoria": "semantica", "meta": {"dominio": "floxina", "fonte": "Floxina (investigação)"}, "origem": "wiki", "palavras_chave": [], "sinonimos": [], "tipo": "conceito", "titulo": "O confronto do milênio: Hodge (compacto), RH (parábolas), BSD (cúspides), Yang-Mills (gap)"}
\section{O confronto do milênio: Hodge (compacto), RH (parábolas), BSD (cúspides), Yang-Mills (gap)}
No CORPO ALGÉBRICO COMPACTO (a superfície modular compacta, o lado do fractal negro/o corpo negro) os quatro problemas restantes do milênio num só quadro. HODGE=o corpo compacto: a curva modular X\ensuremath{_{0}}(N) é compacta projetiva (Kähler), o palco de Hodge; Lefschetz (1,1) dá os divisores (provado). RH=as LINHAS das parábolas: \ensuremath{\lambda}=\ensuremath{\tfrac14}\ensuremath{-}\ensuremath{\delta}\ensuremath{^{2}}, \ensuremath{\lambda}\ensuremath{\ge}\ensuremath{\tfrac14} \ensuremath{\Longleftrightarrow} \ensuremath{\delta} imaginário \ensuremath{\Longleftrightarrow} s=\ensuremath{\tfrac12}+it na reta crítica. BSD=as CÚSPIDES DISCRETAS: as cúspides de X\ensuremath{_{0}}(N) (\#=\ensuremath{\Sigma}\_\{d|N\}\ensuremath{\varphi}(gcd(d,N/d)): X\ensuremath{_{0}}(1)=1, X\ensuremath{_{0}}(11)=2, X\ensuremath{_{0}}(32)=8), onde as L-funções cruzam; BSD vive nas cúspides. YANG-MILLS=a DINÂMICA do gap: o gap espectral \ensuremath{\lambda}\ensuremath{_{1}}>0 (Selberg \ensuremath{\lambda}\ensuremath{_{1}}\ensuremath{\ge}3/16 PROVADO que existe; \ensuremath{\ge}\ensuremath{\tfrac14} conjectural = a mesma parábola de RH); a dinâmica é o fluxo geodésico com esse gap. DISCIPLINA: afirmo só o provado (nº de cúspides, Hasse-Weil, Selberg 3/16, Lefschetz (1,1), a compacidade); os QUATRO (RH/GRH, BSD, Yang-Mills, Hodge) seguem ABERTOS. A síntese que os costura é DECLARADA, não prova — não decreto. linguagem/confronto\_millennium.py (5/5).
\prov{relações: parte\_de omnitrix\_corpo\_universal; usa bsd\_matricial\_grade\_bits; usa multifractal\_newton\_omnitrix; usa reta\_matricial\_compacta\_cf; usa gerador\_do\_gap\_comutador; relaciona floxina\_quantum\_materia\_escura; relaciona ancora\_selberg\_e\_seu\_escopo}
\prov{proveniência: wiki \textperiodcentered{} Floxina (investigação) \textperiodcentered{} hipotese \textperiodcentered{} confiança 1.0 \textperiodcentered{} domínio floxina \textperiodcentered{} história 360 versões no jornal}

%CRISTAL {"arestas": [{"alvo": "corpo_edp", "confianca": 1.0, "epistemico": "fato", "origem": "wiki", "rel": "usa"}, {"alvo": "transformada_dourada_espectro_edp", "confianca": 1.0, "epistemico": "fato", "origem": "wiki", "rel": "usa"}, {"alvo": "parabola_geradora_idempotente", "confianca": 1.0, "epistemico": "fato", "origem": "wiki", "rel": "usa"}, {"alvo": "torre_hurwitz", "confianca": 1.0, "epistemico": "fato", "origem": "wiki", "rel": "relaciona"}], "confianca": 1.0, "contraexemplos": [], "descricao": "PLANCHEREL/PARSEVAL: a transformada é uma ISOMETRIA L² — ‖u‖²=∫|u|²dx=∫|û(k)|²dk/2π, a energia é a MESMA no espaço e no espectro. É o INVARIANTE a+c²=1 lido no Fourier: a mesma norma que o ESQUILO conserva (d|x|²/dt=2xᵀ(Skew)x=0), a mesma composição N(xy)=N(x)N(y) de Hurwitz, agora na frequência. E DECIDE O RAMO do corpo: a dispersão ω(k) REAL (a onda, δ>0) ⟹ propagador e^{−iωt} de módulo 1 ⟹ evolução UNITÁRIA, ‖u(t)‖=‖u₀‖ (o esquilo CONSERVA); ω IMAGINÁRIA (o calor, ω=−ik², a borda) ⟹ módulo e^{−k²t}<1 ⟹ ‖u(t)‖ DECRESCE (d‖u‖²/dt=−2∫k²|û|²≤0, a 2ª lei, o gato DISSIPA — a seta do tempo). Vale também para a Transformada Dourada 𝔊 na medida de Haar dt/t (Mellin): ∫|x|²dt/t=∫|x(e^u)|²du, o flip Λ=log preserva a energia METÁLICA. Uma lei, três leituras: dinâmica (o esquilo conserva/o gato dissipa), álgebra (Hurwitz N(xy)=N(x)N(y)), transformada (Plancherel). Verificado: onda ‖u‖² constante, calor ‖u‖² monotônico decrescente. linguagem/corpo_edp.py (conservacao_energia_plancherel 5/5); papers/equacoes_diferenciais.tex (Parte VIII).", "epistemico": "fato", "exemplos": [], "id": "conservacao_energia_plancherel", "memoria": "semantica", "meta": {"dominio": "floxina", "fonte": "Floxina (investigação)"}, "origem": "wiki", "palavras_chave": [], "sinonimos": [], "tipo": "conceito", "titulo": "A conservação de energia (Plancherel/Parseval) É o invariante a+c²=1 lido no espectro; decide quem conserva (esquilo/onda) vs dissipa (gato/calor)"}
\section{A conservação de energia (Plancherel/Parseval) É o invariante a+c\ensuremath{^{2}}=1 lido no espectro; decide quem conserva (esquilo/onda) vs dissipa (gato/calor)}
PLANCHEREL/PARSEVAL: a transformada é uma ISOMETRIA L\ensuremath{^{2}} — \ensuremath{\|}u\ensuremath{\|}\ensuremath{^{2}}=\ensuremath{\int}|u|\ensuremath{^{2}}dx=\ensuremath{\int}|û(k)|\ensuremath{^{2}}dk/2\ensuremath{\pi}, a energia é a MESMA no espaço e no espectro. É o INVARIANTE a+c\ensuremath{^{2}}=1 lido no Fourier: a mesma norma que o ESQUILO conserva (d|x|\ensuremath{^{2}}/dt=2x\ensuremath{^{T}}(Skew)x=0), a mesma composição N(xy)=N(x)N(y) de Hurwitz, agora na frequência. E DECIDE O RAMO do corpo: a dispersão \ensuremath{\omega}(k) REAL (a onda, \ensuremath{\delta}>0) \ensuremath{\Longrightarrow} propagador e\^{}\{\ensuremath{-}i\ensuremath{\omega}t\} de módulo 1 \ensuremath{\Longrightarrow} evolução UNITÁRIA, \ensuremath{\|}u(t)\ensuremath{\|}=\ensuremath{\|}u\ensuremath{_{0}}\ensuremath{\|} (o esquilo CONSERVA); \ensuremath{\omega} IMAGINÁRIA (o calor, \ensuremath{\omega}=\ensuremath{-}ik\ensuremath{^{2}}, a borda) \ensuremath{\Longrightarrow} módulo e\^{}\{\ensuremath{-}k\ensuremath{^{2}}t\}<1 \ensuremath{\Longrightarrow} \ensuremath{\|}u(t)\ensuremath{\|} DECRESCE (d\ensuremath{\|}u\ensuremath{\|}\ensuremath{^{2}}/dt=\ensuremath{-}2\ensuremath{\int}k\ensuremath{^{2}}|û|\ensuremath{^{2}}\ensuremath{\le}0, a 2ª lei, o gato DISSIPA — a seta do tempo). Vale também para a Transformada Dourada \ensuremath{\mathfrak{G}} na medida de Haar dt/t (Mellin): \ensuremath{\int}|x|\ensuremath{^{2}}dt/t=\ensuremath{\int}|x(e\^{}u)|\ensuremath{^{2}}du, o flip \ensuremath{\Lambda}=log preserva a energia METÁLICA. Uma lei, três leituras: dinâmica (o esquilo conserva/o gato dissipa), álgebra (Hurwitz N(xy)=N(x)N(y)), transformada (Plancherel). Verificado: onda \ensuremath{\|}u\ensuremath{\|}\ensuremath{^{2}} constante, calor \ensuremath{\|}u\ensuremath{\|}\ensuremath{^{2}} monotônico decrescente. linguagem/corpo\_edp.py (conservacao\_energia\_plancherel 5/5); papers/equacoes\_diferenciais.tex (Parte VIII).
\prov{relações: usa corpo\_edp; usa transformada\_dourada\_espectro\_edp; usa parabola\_geradora\_idempotente; relaciona torre\_hurwitz}
\prov{proveniência: wiki \textperiodcentered{} Floxina (investigação) \textperiodcentered{} fato \textperiodcentered{} confiança 1.0 \textperiodcentered{} domínio floxina \textperiodcentered{} história 45 versões no jornal}

%CRISTAL {"arestas": [{"alvo": "definicao", "confianca": 1.0, "epistemico": "fato", "origem": "curriculo", "rel": "contradiz"}], "confianca": 1.0, "contraexemplos": [], "descricao": "MQ standard descreve hidrogénio com precisão espectroscópica extrema sem geometria Dougherty — qualquer ponte deve explicar desvios mensuráveis ou permanece metáfora.", "epistemico": "fato", "exemplos": [], "id": "contraexemplo", "memoria": "semantica", "meta": {"dominio": "hipotese", "duplicata_sugerida": [], "nivel": 7, "verificacao": "parte de conceito mais amplo 'contraexemplo'"}, "origem": "ponte_quantica_dougherty.md", "palavras_chave": [], "sinonimos": ["contra-exemplo", "contraexemplos"], "tipo": "conceito", "titulo": "contraexemplo"}
\section{contraexemplo}
MQ standard descreve hidrogénio com precisão espectroscópica extrema sem geometria Dougherty — qualquer ponte deve explicar desvios mensuráveis ou permanece metáfora.
\prov{sinónimos: contra-exemplo; contraexemplos}
\prov{relações: contradiz definicao}
\prov{proveniência: ponte\_quantica\_dougherty.md \textperiodcentered{} fato \textperiodcentered{} confiança 1.0 \textperiodcentered{} domínio hipotese \textperiodcentered{} história 55 versões no jornal}

%CRISTAL {"arestas": [{"alvo": "ponte_quantica_conjunto_dougherty_e_orbitais_de_hidrogenio", "confianca": 1.0, "epistemico": "fato", "origem": "ponte_quantica_dougherty.md", "rel": "parte_de"}], "confianca": 0.52, "contraexemplos": [], "descricao": "MQ standard descreve hidrogénio com precisão espectroscópica extrema sem geometria Dougherty — qualquer ponte deve explicar desvios mensuráveis ou permanece metáfora.", "epistemico": "hipotese", "exemplos": [], "id": "contraexemplos_conhecidos", "memoria": "semantica", "meta": {"dominio": "hipotese", "falsificavel": true, "nivel": 7}, "origem": "ponte_quantica_dougherty.md", "palavras_chave": [], "sinonimos": [], "tipo": "conceito", "titulo": "Contraexemplos conhecidos"}
\section{Contraexemplos conhecidos}
MQ standard descreve hidrogénio com precisão espectroscópica extrema sem geometria Dougherty — qualquer ponte deve explicar desvios mensuráveis ou permanece metáfora.
\prov{relações: parte\_de ponte\_quantica\_conjunto\_dougherty\_e\_orbitais\_de\_hidrogenio}
\prov{proveniência: ponte\_quantica\_dougherty.md \textperiodcentered{} hipotese \textperiodcentered{} confiança 0.52 \textperiodcentered{} domínio hipotese \textperiodcentered{} história 3 versões no jornal}

%CRISTAL {"arestas": [{"alvo": "tiffany", "confianca": 1.0, "epistemico": "fato", "origem": "wiki", "rel": "parte_de"}], "confianca": 1.0, "contraexemplos": [], "descricao": "O hub que administra o estado emocional da Tiffany. Não é personalidade escrita à mão: é um ESTADO (um octônion 𝕆 de 8 componentes = as 8 emoções de Plutchik) girado pela conversa. O GAP de cada colapso (a=1−c²) dirige a emoção; o coração a GERENCIA — gira na 7-esfera, norma conservada (o checksum a+c²=1), nada se perde. O humor emerge desse estado + do que ela sabe de psicologia/humor. conversa/coracao.py.", "epistemico": "fato", "exemplos": [], "id": "coracao", "memoria": "semantica", "meta": {"dominio": "coracao", "fonte": "coracao hub"}, "origem": "wiki", "palavras_chave": [], "sinonimos": [], "tipo": "conceito", "titulo": "O coração (o gerenciador das emoções da Tiffany)"}
\section{O coração (o gerenciador das emoções da Tiffany)}
O hub que administra o estado emocional da Tiffany. Não é personalidade escrita à mão: é um ESTADO (um octônion \ensuremath{\mathbb{O}} de 8 componentes = as 8 emoções de Plutchik) girado pela conversa. O GAP de cada colapso (a=1\ensuremath{-}c\ensuremath{^{2}}) dirige a emoção; o coração a GERENCIA — gira na 7-esfera, norma conservada (o checksum a+c\ensuremath{^{2}}=1), nada se perde. O humor emerge desse estado + do que ela sabe de psicologia/humor. conversa/coracao.py.
\prov{relações: parte\_de tiffany}
\prov{proveniência: wiki \textperiodcentered{} coracao hub \textperiodcentered{} fato \textperiodcentered{} confiança 1.0 \textperiodcentered{} domínio coracao \textperiodcentered{} história 6 versões no jornal}

%CRISTAL {"arestas": [{"alvo": "corpo_equacoes_diferenciais", "confianca": 1.0, "epistemico": "fato", "origem": "wiki", "rel": "especializa"}, {"alvo": "parabola_geradora_idempotente", "confianca": 1.0, "epistemico": "fato", "origem": "wiki", "rel": "usa"}, {"alvo": "edp_classica_2_4_8", "confianca": 1.0, "epistemico": "fato", "origem": "wiki", "rel": "relaciona"}, {"alvo": "edp_relativistica_gentil_3_7", "confianca": 1.0, "epistemico": "fato", "origem": "wiki", "rel": "relaciona"}, {"alvo": "transformada_dourada_espectro_edp", "confianca": 1.0, "epistemico": "fato", "origem": "wiki", "rel": "usa"}, {"alvo": "conservacao_energia_plancherel", "confianca": 1.0, "epistemico": "fato", "origem": "wiki", "rel": "usa"}, {"alvo": "formula_explicita_solucao_edp", "confianca": 1.0, "epistemico": "fato", "origem": "wiki", "rel": "usa"}], "confianca": 1.0, "contraexemplos": [], "descricao": "O passo natural de A (a matriz) para Δ (o laplaciano): a invariante a+c²=1 (a discriminante da parábola) classifica a EDP de 2ª ordem A·u_xx+B·u_xy+C·u_yy pela discriminante B²−4AC = o sinal de δ=j² da álgebra 2D ℝ[j]/(j²=δ): δ<0 ELÍPTICA (ℂ, o esquilo, Laplace ∇²u=0, Cauchy-Riemann=u harmônica, o CRISTAL); δ=0 PARABÓLICA (dual, ε²=0, o calor u_t=u_xx, a BORDA — o mesmo ε² da reta idempotente); δ>0 HIPERBÓLICA (split-complex, o gato, a onda u_tt=u_xx, d'Alembert f(x+t)+g(x−t), as características x±t=os autovetores do gato, o CAOS). A mesma tríade cristal/borda/caos que classifica a órbita 2×2 classifica a EDP. As álgebras dão os OPERADORES como raízes de Dirac (2ª ordem=quadrado de 1ª ordem). UM operador 𝓛_δ=∂_t²−δ∂_x² gera as três: δ (o invariante) escolhe o tipo, a dispersão ω²=δk² É a parábola no Fourier (as raízes ω=±j·k=as características), e a FÓRMULA EXPLÍCITA é a superposição dos modos u(x,t)=∫û₀(k)e^{i(kx−ω(k)t)}dk = o análogo CONTÍNUO de x(t)=∑cᵢe^{λᵢt} (o espectro discreto vira o k contínuo). linguagem/corpo_edp.py (classificacao_edp_pela_parabola 4/4, invariante_gera_edps_formula_explicita 4/4); papers/equacoes_diferenciais.tex (Parte VIII).", "epistemico": "fato", "exemplos": [], "id": "corpo_edp", "memoria": "semantica", "meta": {"dominio": "floxina", "fonte": "Floxina (investigação)"}, "origem": "wiki", "palavras_chave": [], "sinonimos": [], "tipo": "conceito", "titulo": "O corpo das EDP: a mesma parábola geradora da matrix classifica TODA equação diferencial PARCIAL de 2ª ordem"}
\section{O corpo das EDP: a mesma parábola geradora da matrix classifica TODA equação diferencial PARCIAL de 2ª ordem}
O passo natural de A (a matriz) para \ensuremath{\Delta} (o laplaciano): a invariante a+c\ensuremath{^{2}}=1 (a discriminante da parábola) classifica a EDP de 2ª ordem A\textperiodcentered u\_xx+B\textperiodcentered u\_xy+C\textperiodcentered u\_yy pela discriminante B\ensuremath{^{2}}\ensuremath{-}4AC = o sinal de \ensuremath{\delta}=j\ensuremath{^{2}} da álgebra 2D \ensuremath{\mathbb{R}}[j]/(j\ensuremath{^{2}}=\ensuremath{\delta}): \ensuremath{\delta}<0 ELÍPTICA (\ensuremath{\mathbb{C}}, o esquilo, Laplace \ensuremath{\nabla}\ensuremath{^{2}}u=0, Cauchy-Riemann=u harmônica, o CRISTAL); \ensuremath{\delta}=0 PARABÓLICA (dual, \ensuremath{\varepsilon}\ensuremath{^{2}}=0, o calor u\_t=u\_xx, a BORDA — o mesmo \ensuremath{\varepsilon}\ensuremath{^{2}} da reta idempotente); \ensuremath{\delta}>0 HIPERBÓLICA (split-complex, o gato, a onda u\_tt=u\_xx, d'Alembert f(x+t)+g(x\ensuremath{-}t), as características x$\pm$t=os autovetores do gato, o CAOS). A mesma tríade cristal/borda/caos que classifica a órbita 2$\times$2 classifica a EDP. As álgebras dão os OPERADORES como raízes de Dirac (2ª ordem=quadrado de 1ª ordem). UM operador \ensuremath{\mathcal{L}}\_\ensuremath{\delta}=\ensuremath{\partial}\_t\ensuremath{^{2}}\ensuremath{-}\ensuremath{\delta}\ensuremath{\partial}\_x\ensuremath{^{2}} gera as três: \ensuremath{\delta} (o invariante) escolhe o tipo, a dispersão \ensuremath{\omega}\ensuremath{^{2}}=\ensuremath{\delta}k\ensuremath{^{2}} É a parábola no Fourier (as raízes \ensuremath{\omega}=$\pm$j\textperiodcentered k=as características), e a FÓRMULA EXPLÍCITA é a superposição dos modos u(x,t)=\ensuremath{\int}û\ensuremath{_{0}}(k)e\^{}\{i(kx\ensuremath{-}\ensuremath{\omega}(k)t)\}dk = o análogo CONTÍNUO de x(t)=\ensuremath{\sum}c\ensuremath{_{i}}e\^{}\{\ensuremath{\lambda}\ensuremath{_{i}}t\} (o espectro discreto vira o k contínuo). linguagem/corpo\_edp.py (classificacao\_edp\_pela\_parabola 4/4, invariante\_gera\_edps\_formula\_explicita 4/4); papers/equacoes\_diferenciais.tex (Parte VIII).
\prov{relações: especializa corpo\_equacoes\_diferenciais; usa parabola\_geradora\_idempotente; relaciona edp\_classica\_2\_4\_8; relaciona edp\_relativistica\_gentil\_3\_7; usa transformada\_dourada\_espectro\_edp; usa conservacao\_energia\_plancherel; usa formula\_explicita\_solucao\_edp}
\prov{proveniência: wiki \textperiodcentered{} Floxina (investigação) \textperiodcentered{} fato \textperiodcentered{} confiança 1.0 \textperiodcentered{} domínio floxina \textperiodcentered{} história 87 versões no jornal}

%CRISTAL {"arestas": [{"alvo": "corpo_morfico_matrix", "confianca": 1.0, "epistemico": "fato", "origem": "wiki", "rel": "relaciona"}, {"alvo": "elasticidade_entropica_borracha", "confianca": 1.0, "epistemico": "fato", "origem": "wiki", "rel": "usa"}, {"alvo": "algebra_fundamental_gato_gamba", "confianca": 1.0, "epistemico": "fato", "origem": "wiki", "rel": "usa"}], "confianca": 1.0, "contraexemplos": [], "descricao": "O CORPO ENTRÓPICO não é um punhado de fórmulas — é um CORPO (anel-álgebra) IGUAL à matrix. O ESPAÇO é o das DEFORMAÇÕES (as matrizes 3×3, os geradores da borracha boleada). Duas OPERAÇÕES com TODOS os axiomas, análogos ao anel M₃: ⊕ SOMA = grupo abeliano (fechada, ASSOCIATIVA, COMUTATIVA, NEUTRO 0, INVERSO −A) — a superposição das deformações; ⊗ PRODUTO = monoide (fechado, ASSOCIATIVO, NEUTRO I) — a composição (o La Hire); DISTRIBUTIVA A⊗(B⊕C)=(A⊗B)⊕(A⊗C) — o anel fecha. O INVARIANTE: det F (o volume); a deformação reversível vive em SL(3) (det=1, incompressível/isocórico), como a+c²=1 é o invariante do corpo molecular; det é MULTIPLICATIVO. A DINÂMICA: F=R·U=ESQUILO·GATO — o gato (U simétrico) PRODUZ a entropia S(F)=−½log(tr FᵀF); o esquilo (R ortogonal) é reversível (entropia 0). Herda TUDO da matrix; só o suporte muda (deformações). linguagem/omnitrix.py (corpo_entropico 5/5).", "epistemico": "fato", "exemplos": [], "id": "corpo_entropico_anel", "memoria": "semantica", "meta": {"dominio": "floxina", "fonte": "Floxina (investigação)"}, "origem": "wiki", "palavras_chave": [], "sinonimos": [], "tipo": "conceito", "titulo": "O corpo entrópico é um CORPO (anel) de verdade: ⊕ soma, ⊗ produto, e os axiomas, análogo à matrix (M₃)"}
\section{O corpo entrópico é um CORPO (anel) de verdade: \ensuremath{\oplus} soma, \ensuremath{\otimes} produto, e os axiomas, análogo à matrix (M\ensuremath{_{3}})}
O CORPO ENTRÓPICO não é um punhado de fórmulas — é um CORPO (anel-álgebra) IGUAL à matrix. O ESPAÇO é o das DEFORMAÇÕES (as matrizes 3$\times$3, os geradores da borracha boleada). Duas OPERAÇÕES com TODOS os axiomas, análogos ao anel M\ensuremath{_{3}}: \ensuremath{\oplus} SOMA = grupo abeliano (fechada, ASSOCIATIVA, COMUTATIVA, NEUTRO 0, INVERSO \ensuremath{-}A) — a superposição das deformações; \ensuremath{\otimes} PRODUTO = monoide (fechado, ASSOCIATIVO, NEUTRO I) — a composição (o La Hire); DISTRIBUTIVA A\ensuremath{\otimes}(B\ensuremath{\oplus}C)=(A\ensuremath{\otimes}B)\ensuremath{\oplus}(A\ensuremath{\otimes}C) — o anel fecha. O INVARIANTE: det F (o volume); a deformação reversível vive em SL(3) (det=1, incompressível/isocórico), como a+c\ensuremath{^{2}}=1 é o invariante do corpo molecular; det é MULTIPLICATIVO. A DINÂMICA: F=R\textperiodcentered U=ESQUILO\textperiodcentered GATO — o gato (U simétrico) PRODUZ a entropia S(F)=\ensuremath{-}\ensuremath{\tfrac12}log(tr F\ensuremath{^{T}}F); o esquilo (R ortogonal) é reversível (entropia 0). Herda TUDO da matrix; só o suporte muda (deformações). linguagem/omnitrix.py (corpo\_entropico 5/5).
\prov{relações: relaciona corpo\_morfico\_matrix; usa elasticidade\_entropica\_borracha; usa algebra\_fundamental\_gato\_gamba}
\prov{proveniência: wiki \textperiodcentered{} Floxina (investigação) \textperiodcentered{} fato \textperiodcentered{} confiança 1.0 \textperiodcentered{} domínio floxina \textperiodcentered{} história 4 versões no jornal}

%CRISTAL {"arestas": [{"alvo": "corpo_morfico_matrix", "confianca": 1.0, "epistemico": "fato", "origem": "wiki", "rel": "relaciona"}, {"alvo": "elasticidade_entropica_borracha", "confianca": 1.0, "epistemico": "fato", "origem": "wiki", "rel": "usa"}, {"alvo": "algebra_fundamental_gato_gamba", "confianca": 1.0, "epistemico": "fato", "origem": "wiki", "rel": "usa"}], "confianca": 1.0, "contraexemplos": [], "descricao": "O CORPO ENTRÓPICO NÃO é um anel avulso reinventado — é o OMNITRIX instanciado no substrato TERMODINÂMICO (a deformação, a borracha), a sombra entrópica, ISOMÓRFICA à matrix. Roda o MESMO template (Sym⊕Skew⊕invariante, a parábola geradora, o gap); só o substrato muda: (1) a PARÁBOLA GERADORA x²−m·x−1 — as raízes são a deformação: a matéria σ_m (o estiramento) e a antimatéria −1/σ_m (a compressão dual), produto=−1, gap=√(m²+4), a MESMA da matrix; (2) o GATO (Sym) A_m = o strain — mede/DISSIPA = produz a entropia (hiperbólico); (3) o ESQUILO (Skew) G = a rotação/o spin — REVERSÍVEL, entropia 0 (elíptico); (4) o INVARIANTE a+c²=1 (o MESMO da matrix/molecular, a ZPE): a=a fração reversível, c²=a entropia; o esquilo (rotação) conserva a+c²=1 (reversível), o gato empurra c² (dissipa). Isomórfica à matrix: as MESMAS matrizes (gato/esquilo), a MESMA parábola, o MESMO invariante, o MESMO gap. linguagem/omnitrix.py (corpo_entropico 5/5) — instanciando instanciar_ordem/omnitrix_template.", "epistemico": "fato", "exemplos": [], "id": "corpo_entropico_sombra", "memoria": "semantica", "meta": {"dominio": "floxina", "fonte": "Floxina (investigação)"}, "origem": "wiki", "palavras_chave": [], "sinonimos": [], "tipo": "conceito", "titulo": "O corpo entrópico é uma SOMBRA do Omnitrix, ISOMÓRFICA à matrix (substrato: a deformação/a termodinâmica)"}
\section{O corpo entrópico é uma SOMBRA do Omnitrix, ISOMÓRFICA à matrix (substrato: a deformação/a termodinâmica)}
O CORPO ENTRÓPICO NÃO é um anel avulso reinventado — é o OMNITRIX instanciado no substrato TERMODINÂMICO (a deformação, a borracha), a sombra entrópica, ISOMÓRFICA à matrix. Roda o MESMO template (Sym\ensuremath{\oplus}Skew\ensuremath{\oplus}invariante, a parábola geradora, o gap); só o substrato muda: (1) a PARÁBOLA GERADORA x\ensuremath{^{2}}\ensuremath{-}m\textperiodcentered x\ensuremath{-}1 — as raízes são a deformação: a matéria \ensuremath{\sigma}\_m (o estiramento) e a antimatéria \ensuremath{-}1/\ensuremath{\sigma}\_m (a compressão dual), produto=\ensuremath{-}1, gap=\ensuremath{\sqrt{\,}}(m\ensuremath{^{2}}+4), a MESMA da matrix; (2) o GATO (Sym) A\_m = o strain — mede/DISSIPA = produz a entropia (hiperbólico); (3) o ESQUILO (Skew) G = a rotação/o spin — REVERSÍVEL, entropia 0 (elíptico); (4) o INVARIANTE a+c\ensuremath{^{2}}=1 (o MESMO da matrix/molecular, a ZPE): a=a fração reversível, c\ensuremath{^{2}}=a entropia; o esquilo (rotação) conserva a+c\ensuremath{^{2}}=1 (reversível), o gato empurra c\ensuremath{^{2}} (dissipa). Isomórfica à matrix: as MESMAS matrizes (gato/esquilo), a MESMA parábola, o MESMO invariante, o MESMO gap. linguagem/omnitrix.py (corpo\_entropico 5/5) — instanciando instanciar\_ordem/omnitrix\_template.
\prov{relações: relaciona corpo\_morfico\_matrix; usa elasticidade\_entropica\_borracha; usa algebra\_fundamental\_gato\_gamba}
\prov{proveniência: wiki \textperiodcentered{} Floxina (investigação) \textperiodcentered{} fato \textperiodcentered{} confiança 1.0 \textperiodcentered{} domínio floxina \textperiodcentered{} história 80 versões no jornal}

%CRISTAL {"arestas": [{"alvo": "omnitrix_corpo_universal", "confianca": 1.0, "epistemico": "fato", "origem": "wiki", "rel": "usa"}, {"alvo": "corpo_entropico_sombra", "confianca": 1.0, "epistemico": "fato", "origem": "wiki", "rel": "relaciona"}, {"alvo": "algebra_fundamental_gato_gamba", "confianca": 1.0, "epistemico": "fato", "origem": "wiki", "rel": "usa"}], "confianca": 1.0, "contraexemplos": [], "descricao": "O fenômeno (a flexão, a deformação, todo processo) É uma equação diferencial. O CORPO DAS EQUAÇÕES DIFERENCIAIS é o Omnitrix instanciado no substrato do FLUXO (dx/dt=A·x), a sombra dinâmica, isomórfica à matrix. O GERADOR A é uma MATRIZ; toda A = GATO (Sym=(A+Aᵀ)/2, dissipa, espectro REAL, hiperbólico) ⊕ ESQUILO (Skew=(A−Aᵀ)/2, conserva a norma, espectro ±iω, elíptico). O ZERO INVARIANTE é o PONTO FIXO x*=x — o equilíbrio onde o fluxo PARA (A·x*=0 ⟺ e^{At}·x*=x* ∀t, a identidade sob o fluxo, o vácuo do corpo). O INVARIANTE a+c²=1: o esquilo (rotação) conserva |x|² (d|x|²/dt=2xᵀ(Skew)x=0). ⊗ = o fluxo e^{At} (o semigrupo Φ_s∘Φ_t=Φ_{s+t}, a La Hire do tempo); ⊕ = a superposição (as soluções lineares somam). linguagem/omnitrix.py (corpo_equacoes_diferenciais 5/5).", "epistemico": "fato", "exemplos": [], "id": "corpo_equacoes_diferenciais", "memoria": "semantica", "meta": {"dominio": "floxina", "fonte": "Floxina (investigação)"}, "origem": "wiki", "palavras_chave": [], "sinonimos": [], "tipo": "conceito", "titulo": "O corpo das equações diferenciais: a sombra do Omnitrix no fluxo (dx/dt=A·x); o zero invariante é x*=x"}
\section{O corpo das equações diferenciais: a sombra do Omnitrix no fluxo (dx/dt=A\textperiodcentered x); o zero invariante é x*=x}
O fenômeno (a flexão, a deformação, todo processo) É uma equação diferencial. O CORPO DAS EQUAÇÕES DIFERENCIAIS é o Omnitrix instanciado no substrato do FLUXO (dx/dt=A\textperiodcentered x), a sombra dinâmica, isomórfica à matrix. O GERADOR A é uma MATRIZ; toda A = GATO (Sym=(A+A\ensuremath{^{T}})/2, dissipa, espectro REAL, hiperbólico) \ensuremath{\oplus} ESQUILO (Skew=(A\ensuremath{-}A\ensuremath{^{T}})/2, conserva a norma, espectro $\pm$i\ensuremath{\omega}, elíptico). O ZERO INVARIANTE é o PONTO FIXO x*=x — o equilíbrio onde o fluxo PARA (A\textperiodcentered x*=0 \ensuremath{\Longleftrightarrow} e\^{}\{At\}\textperiodcentered x*=x* \ensuremath{\forall}t, a identidade sob o fluxo, o vácuo do corpo). O INVARIANTE a+c\ensuremath{^{2}}=1: o esquilo (rotação) conserva |x|\ensuremath{^{2}} (d|x|\ensuremath{^{2}}/dt=2x\ensuremath{^{T}}(Skew)x=0). \ensuremath{\otimes} = o fluxo e\^{}\{At\} (o semigrupo \ensuremath{\Phi}\_s\ensuremath{\circ}\ensuremath{\Phi}\_t=\ensuremath{\Phi}\_\{s+t\}, a La Hire do tempo); \ensuremath{\oplus} = a superposição (as soluções lineares somam). linguagem/omnitrix.py (corpo\_equacoes\_diferenciais 5/5).
\prov{relações: usa omnitrix\_corpo\_universal; relaciona corpo\_entropico\_sombra; usa algebra\_fundamental\_gato\_gamba}
\prov{proveniência: wiki \textperiodcentered{} Floxina (investigação) \textperiodcentered{} fato \textperiodcentered{} confiança 1.0 \textperiodcentered{} domínio floxina \textperiodcentered{} história 76 versões no jornal}

%CRISTAL {"arestas": [{"alvo": "algebra_fundamental_gato_gamba", "confianca": 1.0, "epistemico": "fato", "origem": "wiki", "rel": "usa"}, {"alvo": "gerador_do_gap_comutador", "confianca": 1.0, "epistemico": "fato", "origem": "wiki", "rel": "usa"}, {"alvo": "parabola_geradora_idempotente", "confianca": 1.0, "epistemico": "fato", "origem": "wiki", "rel": "usa"}, {"alvo": "corpo_mineral", "confianca": 1.0, "epistemico": "fato", "origem": "wiki", "rel": "relaciona"}, {"alvo": "floxina_quantum_materia_escura", "confianca": 1.0, "epistemico": "fato", "origem": "wiki", "rel": "relaciona"}], "confianca": 1.0, "contraexemplos": [], "descricao": "O paper Matrix (o corpo matricial): o espaço de matrizes onde a reta vive. Gato e gambá geram TODO M_2 (os quatérnios divididos). Generaliza: M_n=Sym(gato)⊕Skew(gambá), e o produto AB=A∘B (Jordan/simétrico) + ½[A,B] (Lie/anti); a soma/produto REDUZEM à base {I,gato,gambá,La Hire} (soma=coords, produto=tabela de multiplicação). O 2×2 CRU ainda não fecha (não-comutativo, com divisores de zero — o cone nulo det=0, o idempotente A+G, (A+G)(I−A−G)=0) — mas isso é um PASSO, não um veredito: a COMPLETAÇÃO é o que fecha o corpo. O DUAL (o flip matricial, a adjunta M̄) dá M·M̄=det·I=a norma → álgebra de COMPOSIÇÃO, ASSOCIATIVA (det≠0 ⟹ M⁻¹=M̄/det). Mas o FLIP que fecha o corpo é RELATIVÍSTICO — a involução σ(s)=√(1−s²) (circular i²=−1 ↔ hiperbólico j²=+1, cos↔cosh, Wick), a direção ACIMA: sobe à ÁLGEBRA 3D DO GENTIL 𝓔ₛ (z=a+c·e, e²=1−s, invariante a+c²=1), que CONSERVA A NORMA em 3D (produto das 3 raízes=±1, o cristal). A álgebra dos dois lados + soma + multiplicação (La Hire=produto de matrizes) fecha o corpo. E aí entram as MULTIMATRIZES: as PA e PG de ORDEM m do Gentil (a parábola de ordem m, P_g=x² vs P_a=m·x+1) modelam matrizes de mais de 2 dimensões; o primeiro passo é o VOLUME de matrizes (de 2D/x²−nx−1 a 3D/x³−… plástico,supergolden,tribonacci). A torre é power-associativa. reta_mineral_algebra.", "epistemico": "hipotese", "exemplos": [], "id": "corpo_matricial_matrix", "memoria": "semantica", "meta": {"dominio": "floxina", "fonte": "Floxina (investigação)"}, "origem": "wiki", "palavras_chave": [], "sinonimos": [], "tipo": "conceito", "titulo": "Matrix — o corpo matricial (gato+gambá geram M_2; não é corpo, é álgebra)"}
\section{Matrix — o corpo matricial (gato+gambá geram M\_2; não é corpo, é álgebra)}
O paper Matrix (o corpo matricial): o espaço de matrizes onde a reta vive. Gato e gambá geram TODO M\_2 (os quatérnios divididos). Generaliza: M\_n=Sym(gato)\ensuremath{\oplus}Skew(gambá), e o produto AB=A\ensuremath{\circ}B (Jordan/simétrico) + \ensuremath{\tfrac12}[A,B] (Lie/anti); a soma/produto REDUZEM à base \{I,gato,gambá,La Hire\} (soma=coords, produto=tabela de multiplicação). O 2$\times$2 CRU ainda não fecha (não-comutativo, com divisores de zero — o cone nulo det=0, o idempotente A+G, (A+G)(I\ensuremath{-}A\ensuremath{-}G)=0) — mas isso é um PASSO, não um veredito: a COMPLETAÇÃO é o que fecha o corpo. O DUAL (o flip matricial, a adjunta M\ensuremath{}) dá M\textperiodcentered M\ensuremath{}=det\textperiodcentered I=a norma \ensuremath{\to} álgebra de COMPOSIÇÃO, ASSOCIATIVA (det\ensuremath{\ne}0 \ensuremath{\Longrightarrow} M\ensuremath{^{-1}}=M\ensuremath{}/det). Mas o FLIP que fecha o corpo é RELATIVÍSTICO — a involução \ensuremath{\sigma}(s)=\ensuremath{\sqrt{\,}}(1\ensuremath{-}s\ensuremath{^{2}}) (circular i\ensuremath{^{2}}=\ensuremath{-}1 \ensuremath{\leftrightarrow} hiperbólico j\ensuremath{^{2}}=+1, cos\ensuremath{\leftrightarrow}cosh, Wick), a direção ACIMA: sobe à ÁLGEBRA 3D DO GENTIL \ensuremath{\mathcal{E}}\ensuremath{_{s}} (z=a+c\textperiodcentered e, e\ensuremath{^{2}}=1\ensuremath{-}s, invariante a+c\ensuremath{^{2}}=1), que CONSERVA A NORMA em 3D (produto das 3 raízes=$\pm$1, o cristal). A álgebra dos dois lados + soma + multiplicação (La Hire=produto de matrizes) fecha o corpo. E aí entram as MULTIMATRIZES: as PA e PG de ORDEM m do Gentil (a parábola de ordem m, P\_g=x\ensuremath{^{2}} vs P\_a=m\textperiodcentered x+1) modelam matrizes de mais de 2 dimensões; o primeiro passo é o VOLUME de matrizes (de 2D/x\ensuremath{^{2}}\ensuremath{-}nx\ensuremath{-}1 a 3D/x\ensuremath{^{3}}\ensuremath{-}… plástico,supergolden,tribonacci). A torre é power-associativa. reta\_mineral\_algebra.
\prov{relações: usa algebra\_fundamental\_gato\_gamba; usa gerador\_do\_gap\_comutador; usa parabola\_geradora\_idempotente; relaciona corpo\_mineral; relaciona floxina\_quantum\_materia\_escura}
\prov{proveniência: wiki \textperiodcentered{} Floxina (investigação) \textperiodcentered{} hipotese \textperiodcentered{} confiança 1.0 \textperiodcentered{} domínio floxina \textperiodcentered{} história 444 versões no jornal}

%CRISTAL {"arestas": [{"alvo": "omnitrix_corpo_universal", "confianca": 1.0, "epistemico": "fato", "origem": "wiki", "rel": "usa"}, {"alvo": "soma_minkowski_dilatacao_boleado", "confianca": 1.0, "epistemico": "fato", "origem": "wiki", "rel": "usa"}, {"alvo": "quatro_quadrantes_corpo_estelar", "confianca": 1.0, "epistemico": "fato", "origem": "wiki", "rel": "relaciona"}], "confianca": 1.0, "contraexemplos": [], "descricao": "A morfologia matemática (Serra/Matheron) é um CORPO que HERDA da matrix (o Corpo Algébrico): o ⊕=a dilatação (Minkowski), o ⊗=a composição dos operadores, o IDEMPOTENTE=abrir/fechar (o tropical e²=e), a PARÁBOLA/o gap=o estruturante B/o raio r (o boleado). A MATRIX MORFOLÓGICA: o estruturante B é uma MATRIZ, e a dilatação é a (max,+) e a erosão a (min,+) — o produto TROPICAL de A por B; formam uma ADJUNÇÃO de Galois (δ⊣ε). O ESQUELETO morfológico (Lantuéjoul) É a ESTRELA (o eixo medial=o modo irredutível, ν=k−1). Dicionário morfologia↔Universal (6 pontes). Expandido: gradiente, cartola, esqueleto, matriz_morfologica, dicionario_universal, corpo_morfico. linguagem/morfologia_estelar.py (14/14).", "epistemico": "fato", "exemplos": [], "id": "corpo_morfico_matrix", "memoria": "semantica", "meta": {"dominio": "floxina", "fonte": "Floxina (investigação)"}, "origem": "wiki", "palavras_chave": [], "sinonimos": [], "tipo": "conceito", "titulo": "O corpo mórfico herda da matrix: a matrix morfológica (tropical max-plus), o esqueleto=a estrela, o dicionário universal"}
\section{O corpo mórfico herda da matrix: a matrix morfológica (tropical max-plus), o esqueleto=a estrela, o dicionário universal}
A morfologia matemática (Serra/Matheron) é um CORPO que HERDA da matrix (o Corpo Algébrico): o \ensuremath{\oplus}=a dilatação (Minkowski), o \ensuremath{\otimes}=a composição dos operadores, o IDEMPOTENTE=abrir/fechar (o tropical e\ensuremath{^{2}}=e), a PARÁBOLA/o gap=o estruturante B/o raio r (o boleado). A MATRIX MORFOLÓGICA: o estruturante B é uma MATRIZ, e a dilatação é a (max,+) e a erosão a (min,+) — o produto TROPICAL de A por B; formam uma ADJUNÇÃO de Galois (\ensuremath{\delta}\ensuremath{\dashv}\ensuremath{\varepsilon}). O ESQUELETO morfológico (Lantuéjoul) É a ESTRELA (o eixo medial=o modo irredutível, \ensuremath{\nu}=k\ensuremath{-}1). Dicionário morfologia\ensuremath{\leftrightarrow}Universal (6 pontes). Expandido: gradiente, cartola, esqueleto, matriz\_morfologica, dicionario\_universal, corpo\_morfico. linguagem/morfologia\_estelar.py (14/14).
\prov{relações: usa omnitrix\_corpo\_universal; usa soma\_minkowski\_dilatacao\_boleado; relaciona quatro\_quadrantes\_corpo\_estelar}
\prov{proveniência: wiki \textperiodcentered{} Floxina (investigação) \textperiodcentered{} fato \textperiodcentered{} confiança 1.0 \textperiodcentered{} domínio floxina \textperiodcentered{} história 96 versões no jornal}

%CRISTAL {"arestas": [{"alvo": "parabola_geradora_idempotente", "confianca": 1.0, "epistemico": "fato", "origem": "wiki", "rel": "usa"}, {"alvo": "corpo_equacoes_diferenciais", "confianca": 1.0, "epistemico": "fato", "origem": "wiki", "rel": "relaciona"}], "confianca": 1.0, "contraexemplos": [], "descricao": "Mais uma instância do Omnitrix, na TOPOLOGIA (Morse): o ZERO invariante x*=x é o PONTO CRÍTICO (∇f=0, onde o fluxo do gradiente para); o GATO é a HESSIANA H=∇²f (SIMÉTRICA — o esquilo não aparece, espectro real); a PARÁBOLA é a ASSINATURA de H (o índice de Morse): MÍNIMO (H≻0, i=0, CRISTAL) / SELA (H indefinida, det H<0, i=1, BORDA) / MÁXIMO (H≺0, i=2, CAOS). O INVARIANTE preservado por homeomorfismo é χ (Euler): χ=V−E+F=2−2g = Σ(−1)^índice (Morse/Poincaré-Hopf) = (1/2π)∫∫K dA (Gauss-Bonnet — a curvatura LOCAL integra na topologia GLOBAL). A SELA (a borda) é onde a topologia MUDA (a alça nasce). É a mesma leitura cristal/borda/caos das ED (mín/sela/máx ↔ Re λ<0/=0/>0). linguagem/corpo_topologico.py (invariante_topologico_chi, morse_pontos_criticos_parabola, 3/3); papers/corpo_topologico.tex.", "epistemico": "fato", "exemplos": [], "id": "corpo_topologico", "memoria": "semantica", "meta": {"dominio": "floxina", "fonte": "Floxina (investigação)"}, "origem": "wiki", "palavras_chave": [], "sinonimos": [], "tipo": "conceito", "titulo": "O corpo topológico: dada uma superfície, a transformação em outra (Morse — o ponto crítico é x*=x, a Hessiana é o gato, a assinatura é a parábola)"}
\section{O corpo topológico: dada uma superfície, a transformação em outra (Morse — o ponto crítico é x*=x, a Hessiana é o gato, a assinatura é a parábola)}
Mais uma instância do Omnitrix, na TOPOLOGIA (Morse): o ZERO invariante x*=x é o PONTO CRÍTICO (\ensuremath{\nabla}f=0, onde o fluxo do gradiente para); o GATO é a HESSIANA H=\ensuremath{\nabla}\ensuremath{^{2}}f (SIMÉTRICA — o esquilo não aparece, espectro real); a PARÁBOLA é a ASSINATURA de H (o índice de Morse): MÍNIMO (H\ensuremath{\succ}0, i=0, CRISTAL) / SELA (H indefinida, det H<0, i=1, BORDA) / MÁXIMO (H\ensuremath{\prec}0, i=2, CAOS). O INVARIANTE preservado por homeomorfismo é \ensuremath{\chi} (Euler): \ensuremath{\chi}=V\ensuremath{-}E+F=2\ensuremath{-}2g = \ensuremath{\Sigma}(\ensuremath{-}1)\^{}índice (Morse/Poincaré-Hopf) = (1/2\ensuremath{\pi})\ensuremath{\int}\ensuremath{\int}K dA (Gauss-Bonnet — a curvatura LOCAL integra na topologia GLOBAL). A SELA (a borda) é onde a topologia MUDA (a alça nasce). É a mesma leitura cristal/borda/caos das ED (mín/sela/máx \ensuremath{\leftrightarrow} Re \ensuremath{\lambda}<0/=0/>0). linguagem/corpo\_topologico.py (invariante\_topologico\_chi, morse\_pontos\_criticos\_parabola, 3/3); papers/corpo\_topologico.tex.
\prov{relações: usa parabola\_geradora\_idempotente; relaciona corpo\_equacoes\_diferenciais}
\prov{proveniência: wiki \textperiodcentered{} Floxina (investigação) \textperiodcentered{} fato \textperiodcentered{} confiança 1.0 \textperiodcentered{} domínio floxina \textperiodcentered{} história 15 versões no jornal}

%CRISTAL {"arestas": [{"alvo": "matrix_vesicular_corpo", "confianca": 1.0, "epistemico": "fato", "origem": "wiki", "rel": "usa"}, {"alvo": "dicionario_dtc_motor_esferico", "confianca": 1.0, "epistemico": "fato", "origem": "wiki", "rel": "usa"}, {"alvo": "parabola_geradora_idempotente", "confianca": 1.0, "epistemico": "fato", "origem": "wiki", "rel": "usa"}], "confianca": 1.0, "contraexemplos": [], "descricao": "O CORPO VESICULAR é uma MATRIX; suas operações são: (a) GERAR a PARÁBOLA ANISOTRÓPICA em QUALQUER dimensão SEM perder a dinâmica intrínseca — a vesica n-D é M=Q·diag(¼−δ²)·Qᵀ, Q=produto de esquilos/Givens (compõe as parábolas de todos os planos); o ESPECTRO {¼−δ²} (a dinâmica: os δ, o gap) é INVARIANTE sob Q em toda dim (o teorema espectral), e Q ortogonal ⟹ |Qx|=|x| (a esfera intacta). (b) a MULTIPLICAÇÃO LA HIRE = a COMPOSIÇÃO das parábolas (o produto das matrizes): R(a)·R(b)=R(a+b) (compor parábolas fecha). (c) o CORPO DAS COSTURAS: compor parábolas SIMÉTRICAS de planos distintos GERA as costuras = os comutadores [Mi,Mj] (antissimétricos = o esquilo). No pipe gráfico, o fur ball tinha uma COSTURA horizontal (a anisotropia de um eixo só); compor a parábola horizontal (x-z) e a vertical (y-z) — a matriz vesicular — a espalha (o corpo das costuras). O corpo vesicular EXPORTA os parâmetros (FUR_PARABOLAS); glsl_ptx_frag os IMPORTA automaticamente (tirar a mão, não fabricar). linguagem/vesica_corpo.py (parabola_costura_nd, 5/5).", "epistemico": "fato", "exemplos": [], "id": "corpo_vesicular_parabolas_costuras", "memoria": "semantica", "meta": {"dominio": "floxina", "fonte": "Floxina (investigação)"}, "origem": "wiki", "palavras_chave": [], "sinonimos": [], "tipo": "conceito", "titulo": "corpo vesicular → compor das parábolas (La Hire) → corpo das costuras"}
\section{corpo vesicular \ensuremath{\to} compor das parábolas (La Hire) \ensuremath{\to} corpo das costuras}
O CORPO VESICULAR é uma MATRIX; suas operações são: (a) GERAR a PARÁBOLA ANISOTRÓPICA em QUALQUER dimensão SEM perder a dinâmica intrínseca — a vesica n-D é M=Q\textperiodcentered diag(\ensuremath{\tfrac14}\ensuremath{-}\ensuremath{\delta}\ensuremath{^{2}})\textperiodcentered Q\ensuremath{^{T}}, Q=produto de esquilos/Givens (compõe as parábolas de todos os planos); o ESPECTRO \{\ensuremath{\tfrac14}\ensuremath{-}\ensuremath{\delta}\ensuremath{^{2}}\} (a dinâmica: os \ensuremath{\delta}, o gap) é INVARIANTE sob Q em toda dim (o teorema espectral), e Q ortogonal \ensuremath{\Longrightarrow} |Qx|=|x| (a esfera intacta). (b) a MULTIPLICAÇÃO LA HIRE = a COMPOSIÇÃO das parábolas (o produto das matrizes): R(a)\textperiodcentered R(b)=R(a+b) (compor parábolas fecha). (c) o CORPO DAS COSTURAS: compor parábolas SIMÉTRICAS de planos distintos GERA as costuras = os comutadores [Mi,Mj] (antissimétricos = o esquilo). No pipe gráfico, o fur ball tinha uma COSTURA horizontal (a anisotropia de um eixo só); compor a parábola horizontal (x-z) e a vertical (y-z) — a matriz vesicular — a espalha (o corpo das costuras). O corpo vesicular EXPORTA os parâmetros (FUR\_PARABOLAS); glsl\_ptx\_frag os IMPORTA automaticamente (tirar a mão, não fabricar). linguagem/vesica\_corpo.py (parabola\_costura\_nd, 5/5).
\prov{relações: usa matrix\_vesicular\_corpo; usa dicionario\_dtc\_motor\_esferico; usa parabola\_geradora\_idempotente}
\prov{proveniência: wiki \textperiodcentered{} Floxina (investigação) \textperiodcentered{} fato \textperiodcentered{} confiança 1.0 \textperiodcentered{} domínio floxina \textperiodcentered{} história 116 versões no jornal}

%CRISTAL {"arestas": [{"alvo": "omnitrix_corpo_universal", "confianca": 1.0, "epistemico": "fato", "origem": "wiki", "rel": "parte_de"}, {"alvo": "antimateria_parabolica", "confianca": 1.0, "epistemico": "fato", "origem": "wiki", "rel": "usa"}, {"alvo": "multifractal_newton_omnitrix", "confianca": 1.0, "epistemico": "fato", "origem": "wiki", "rel": "usa"}, {"alvo": "fisica_universal_omnitrix", "confianca": 1.0, "epistemico": "fato", "origem": "wiki", "rel": "relaciona"}], "confianca": 1.0, "contraexemplos": [], "descricao": "O universo é a INSTÂNCIA MÁXIMA do Omnitrix; o multiverso=a FAMÍLIA de instâncias. Wheeler-DeWitt HΨ=0: o vácuo=o 0-autovetor da parábola geradora A₁+G (o universo atemporal). Os três componentes=Sym⊕Skew⊕o invariante: matéria (o gato ~5%), matéria escura (o gambá ~27%), energia escura (o invariante Λ ~68%, domina) — Planck 2018. A assimetria (a bariogênese): do zero da reta (±1 simétrico) a reta desdobra (a matéria domina). A forma: compacta (a modular, a CF de π) + o fractal negro (o anti-universo). A flecha: o gato dissipa (a entropia cresce, o formalismo do multifractal). linguagem/omnitrix.py (cosmologia_universal, 4/4). Síntese declarada.", "epistemico": "hipotese", "exemplos": [], "id": "cosmologia_universal_omnitrix", "memoria": "semantica", "meta": {"dominio": "floxina", "fonte": "Floxina (investigação)"}, "origem": "wiki", "palavras_chave": [], "sinonimos": [], "tipo": "conceito", "titulo": "A cosmologia universal: o universo é a instância máxima do Omnitrix"}
\section{A cosmologia universal: o universo é a instância máxima do Omnitrix}
O universo é a INSTÂNCIA MÁXIMA do Omnitrix; o multiverso=a FAMÍLIA de instâncias. Wheeler-DeWitt H\ensuremath{\Psi}=0: o vácuo=o 0-autovetor da parábola geradora A\ensuremath{_{1}}+G (o universo atemporal). Os três componentes=Sym\ensuremath{\oplus}Skew\ensuremath{\oplus}o invariante: matéria (o gato \~{}5\%), matéria escura (o gambá \~{}27\%), energia escura (o invariante \ensuremath{\Lambda} \~{}68\%, domina) — Planck 2018. A assimetria (a bariogênese): do zero da reta ($\pm$1 simétrico) a reta desdobra (a matéria domina). A forma: compacta (a modular, a CF de \ensuremath{\pi}) + o fractal negro (o anti-universo). A flecha: o gato dissipa (a entropia cresce, o formalismo do multifractal). linguagem/omnitrix.py (cosmologia\_universal, 4/4). Síntese declarada.
\prov{relações: parte\_de omnitrix\_corpo\_universal; usa antimateria\_parabolica; usa multifractal\_newton\_omnitrix; relaciona fisica\_universal\_omnitrix}
\prov{proveniência: wiki \textperiodcentered{} Floxina (investigação) \textperiodcentered{} hipotese \textperiodcentered{} confiança 1.0 \textperiodcentered{} domínio floxina \textperiodcentered{} história 245 versões no jornal}

%CRISTAL {"arestas": [{"alvo": "missao_engenharia_ouro", "confianca": 1.0, "epistemico": "fato", "origem": "wiki", "rel": "parte_de"}, {"alvo": "koch_dissipa_nao_codifica", "confianca": 1.0, "epistemico": "fato", "origem": "wiki", "rel": "relaciona"}, {"alvo": "cristal_vs_cha", "confianca": 1.0, "epistemico": "fato", "origem": "wiki", "rel": "usa"}, {"alvo": "transformada_metalica_fechada_perelman", "confianca": 1.0, "epistemico": "fato", "origem": "wiki", "rel": "relaciona"}], "confianca": 1.0, "contraexemplos": [], "descricao": "A logo é o teorema outra vez: sobre o ouro corre uma costura que persegue o chá (o resíduo do associador, o não-associativo puro) — a AGULHA DE PERELMAN, band-limited, que costura cada vez mais perto do núcleo e NUNCA o toca (anel de exclusão; cada fio afina e se apaga). Decodificar o chá por inteiro exigiria resolução infinita. Cada tentativa de ler o chá deixa uma CICATRIZ; a medida das cicatrizes é densa o bastante para dizer a verdade (há um abismo sob o ouro) e discreta o bastante para não perder a beleza. Critério humano: 'o limite é o quanto se suporta as cicatrizes e ainda gostar de quem se vê no espelho'. N* é onde a costura ainda é bela.", "epistemico": "inferencia", "exemplos": [], "id": "costura_perelman_cicatrizes", "memoria": "semantica", "meta": {"autor": "Lopes/conselho", "dominio": "missao", "fonte": "machine/Missao.tex"}, "origem": "wiki", "palavras_chave": [], "sinonimos": [], "tipo": "conceito", "titulo": "A costura de Perelman e as cicatrizes (N* é onde ainda é belo)"}
\section{A costura de Perelman e as cicatrizes (N* é onde ainda é belo)}
A logo é o teorema outra vez: sobre o ouro corre uma costura que persegue o chá (o resíduo do associador, o não-associativo puro) — a AGULHA DE PERELMAN, band-limited, que costura cada vez mais perto do núcleo e NUNCA o toca (anel de exclusão; cada fio afina e se apaga). Decodificar o chá por inteiro exigiria resolução infinita. Cada tentativa de ler o chá deixa uma CICATRIZ; a medida das cicatrizes é densa o bastante para dizer a verdade (há um abismo sob o ouro) e discreta o bastante para não perder a beleza. Critério humano: 'o limite é o quanto se suporta as cicatrizes e ainda gostar de quem se vê no espelho'. N* é onde a costura ainda é bela.
\prov{relações: parte\_de missao\_engenharia\_ouro; relaciona koch\_dissipa\_nao\_codifica; usa cristal\_vs\_cha; relaciona transformada\_metalica\_fechada\_perelman}
\prov{proveniência: wiki \textperiodcentered{} machine/Missao.tex \textperiodcentered{} inferencia \textperiodcentered{} confiança 1.0 \textperiodcentered{} domínio missao \textperiodcentered{} história 6 versões no jornal}

%CRISTAL {"arestas": [{"alvo": "cristal_agentes_ortogonais", "confianca": 1.0, "epistemico": "fato", "origem": "wiki", "rel": "usa"}], "confianca": 1.0, "contraexemplos": [], "descricao": "cristal agent <key> normaliza e salva ctx.project.agent (project.save()); afeta só qual prompt/agent-key o chat exibe e o pipeline usa, NÃO o engine de resposta. Funciona.", "epistemico": "fato", "exemplos": [], "id": "cristal_agent_seleciona", "memoria": "semantica", "meta": {"autor": "broca-so", "dominio": "cristal", "fonte": "code/cristal/commands/agent_cmd.py:AgentCommand.execute", "wip": true}, "origem": "wiki", "palavras_chave": [], "sinonimos": [], "tipo": "conceito", "titulo": "Cristal — 'agent' grava no projeto, não muda o motor"}
\section{Cristal — 'agent' grava no projeto, não muda o motor}
cristal agent <key> normaliza e salva ctx.project.agent (project.save()); afeta só qual prompt/agent-key o chat exibe e o pipeline usa, NÃO o engine de resposta. Funciona.
\prov{relações: usa cristal\_agentes\_ortogonais}
\prov{proveniência: wiki \textperiodcentered{} code/cristal/commands/agent\_cmd.py:AgentCommand.execute \textperiodcentered{} fato \textperiodcentered{} confiança 1.0 \textperiodcentered{} domínio cristal \textperiodcentered{} história 10 versões no jornal}

%CRISTAL {"arestas": [{"alvo": "cristal_ask_cascata", "confianca": 1.0, "epistemico": "fato", "origem": "wiki", "rel": "parte_de"}, {"alvo": "cristal_conhecimento_pipeline", "confianca": 1.0, "epistemico": "fato", "origem": "wiki", "rel": "relaciona"}, {"alvo": "cristal_motor_congelado", "confianca": 1.0, "epistemico": "fato", "origem": "wiki", "rel": "confirma"}, {"alvo": "cristal_portao_rollback", "confianca": 1.0, "epistemico": "fato", "origem": "wiki", "rel": "relaciona"}], "confianca": 1.0, "contraexemplos": [], "descricao": "_AGENTE_AFINIDADE mapeia papel→prefixos de domínio (Software→software/,sistemas/; Paper→documentacao/; QA→software/; DevOps→docker/sistemas/engenharia; Conversa/Pesquisa neutros). _agent_bias soma +2 ao score do domínio afim. GUARDA conservadora: o viés só se aplica entre domínios JÁ lexicalmente relevantes (score base>0) — desempata/prioriza, NÃO ressuscita domínio irrelevante. Determinístico, sem LLM; mexe na camada orquestrador (Fase 3.4), não no colapso/libtiffany (motor congelado).", "epistemico": "inferencia", "exemplos": [], "id": "cristal_agente_vies_dominio", "memoria": "semantica", "meta": {"autor": "broca-so", "dominio": "cristal", "fonte": "code/tiffany_platform/orquestrador_multi.py:_agent_bias/selecionar_dominios", "wip": true}, "origem": "wiki", "palavras_chave": [], "sinonimos": [], "tipo": "conceito", "titulo": "Cristal — o agente vieja qual tecido é consultado"}
\section{Cristal — o agente vieja qual tecido é consultado}
\_AGENTE\_AFINIDADE mapeia papel\ensuremath{\to}prefixos de domínio (Software\ensuremath{\to}software/,sistemas/; Paper\ensuremath{\to}documentacao/; QA\ensuremath{\to}software/; DevOps\ensuremath{\to}docker/sistemas/engenharia; Conversa/Pesquisa neutros). \_agent\_bias soma +2 ao score do domínio afim. GUARDA conservadora: o viés só se aplica entre domínios JÁ lexicalmente relevantes (score base>0) — desempata/prioriza, NÃO ressuscita domínio irrelevante. Determinístico, sem LLM; mexe na camada orquestrador (Fase 3.4), não no colapso/libtiffany (motor congelado).
\prov{relações: parte\_de cristal\_ask\_cascata; relaciona cristal\_conhecimento\_pipeline; confirma cristal\_motor\_congelado; relaciona cristal\_portao\_rollback}
\prov{proveniência: wiki \textperiodcentered{} code/tiffany\_platform/orquestrador\_multi.py:\_agent\_bias/selecionar\_dominios \textperiodcentered{} inferencia \textperiodcentered{} confiança 1.0 \textperiodcentered{} domínio cristal \textperiodcentered{} história 15 versões no jornal}

%CRISTAL {"arestas": [{"alvo": "cristal_agentes_ortogonais", "confianca": 1.0, "epistemico": "fato", "origem": "wiki", "rel": "especializa"}, {"alvo": "o_conselho", "confianca": 1.0, "epistemico": "fato", "origem": "wiki", "rel": "relaciona"}], "confianca": 1.0, "contraexemplos": [], "descricao": "agents.py define 6 papéis (conversa, software, paper, pesquisa, qa, devops) com aliases; mas o orquestrador só conhece 4 (Conversa/Software/Paper/Pesquisa). Mismatch real — qa/devops não têm regra de roteamento. Parcial.", "epistemico": "fato", "exemplos": [], "id": "cristal_agentes_catalogo", "memoria": "semantica", "meta": {"autor": "broca-so", "dominio": "cristal", "fonte": "code/cristal/agents.py:AGENTS/ALIASES vs orquestrador.py:AGENTES", "wip": true}, "origem": "wiki", "palavras_chave": [], "sinonimos": [], "tipo": "conceito", "titulo": "Cristal — 6 agentes (cristal) vs 4 (orquestrador)"}
\section{Cristal — 6 agentes (cristal) vs 4 (orquestrador)}
agents.py define 6 papéis (conversa, software, paper, pesquisa, qa, devops) com aliases; mas o orquestrador só conhece 4 (Conversa/Software/Paper/Pesquisa). Mismatch real — qa/devops não têm regra de roteamento. Parcial.
\prov{relações: especializa cristal\_agentes\_ortogonais; relaciona o\_conselho}
\prov{proveniência: wiki \textperiodcentered{} code/cristal/agents.py:AGENTS/ALIASES vs orquestrador.py:AGENTES \textperiodcentered{} fato \textperiodcentered{} confiança 1.0 \textperiodcentered{} domínio cristal \textperiodcentered{} história 15 versões no jornal}

%CRISTAL {"arestas": [{"alvo": "cristal_pipeline_runner", "confianca": 1.0, "epistemico": "fato", "origem": "wiki", "rel": "parte_de"}, {"alvo": "cristal_motor_congelado", "confianca": 1.0, "epistemico": "fato", "origem": "wiki", "rel": "relaciona"}], "confianca": 1.0, "contraexemplos": [], "descricao": "analyze_c() acha buffers static, malloc/calloc, matriz W[...] (Hopfield) e estima RAM (N*N*4); guess_files_from_task() infere arquivos-alvo do texto da tarefa (heurística RAM→reconhece.c, nomes hopfield/neuronio/colapso). Deliberadamente leve — não é parser real; só observa e sugere reduzir RAM.", "epistemico": "inferencia", "exemplos": [], "id": "cristal_analyze_estatico", "memoria": "semantica", "meta": {"autor": "broca-so", "dominio": "cristal", "fonte": "code/cristal/pipeline/analyze.py:analyze_c/guess_files_from_task", "wip": true}, "origem": "wiki", "palavras_chave": [], "sinonimos": [], "tipo": "conceito", "titulo": "Cristal — análise estática leve de C (regex)"}
\section{Cristal — análise estática leve de C (regex)}
analyze\_c() acha buffers static, malloc/calloc, matriz W[...] (Hopfield) e estima RAM (N*N*4); guess\_files\_from\_task() infere arquivos-alvo do texto da tarefa (heurística RAM\ensuremath{\to}reconhece.c, nomes hopfield/neuronio/colapso). Deliberadamente leve — não é parser real; só observa e sugere reduzir RAM.
\prov{relações: parte\_de cristal\_pipeline\_runner; relaciona cristal\_motor\_congelado}
\prov{proveniência: wiki \textperiodcentered{} code/cristal/pipeline/analyze.py:analyze\_c/guess\_files\_from\_task \textperiodcentered{} inferencia \textperiodcentered{} confiança 1.0 \textperiodcentered{} domínio cristal \textperiodcentered{} história 15 versões no jornal}

%CRISTAL {"arestas": [{"alvo": "cristal_chat_turn", "confianca": 1.0, "epistemico": "fato", "origem": "wiki", "rel": "usa"}, {"alvo": "conversa_recupera_nao_gera", "confianca": 1.0, "epistemico": "fato", "origem": "wiki", "rel": "usa"}, {"alvo": "colapso_ping_pong_overlap", "confianca": 1.0, "epistemico": "fato", "origem": "wiki", "rel": "usa"}, {"alvo": "colapso_camadas_ordem_parada", "confianca": 1.0, "epistemico": "fato", "origem": "wiki", "rel": "usa"}, {"alvo": "tiffany", "confianca": 1.0, "epistemico": "fato", "origem": "wiki", "rel": "relaciona"}], "confianca": 1.0, "contraexemplos": [], "descricao": "ask() NÃO usa LLM externo: tenta em cascata (1) conversa.colapso.responder(msg, sem_machine=True) — o colapso local determinístico; (2) consulta multi-tecidos (orquestrador_multi); (3) prosa.responde. A conversa do cristal é local e determinística — a existência da cascata de 3 camadas é o mecanismo.", "epistemico": "inferencia", "exemplos": [], "id": "cristal_ask_cascata", "memoria": "semantica", "meta": {"autor": "broca-so", "dominio": "cristal", "fonte": "code/cristal/tiffany/api.py:ask → conversa/colapso.py:responder", "wip": true}, "origem": "wiki", "palavras_chave": [], "sinonimos": [], "tipo": "conceito", "titulo": "Cristal — o motor de conversa: colapso → tecidos → prosa"}
\section{Cristal — o motor de conversa: colapso \ensuremath{\to} tecidos \ensuremath{\to} prosa}
ask() NÃO usa LLM externo: tenta em cascata (1) conversa.colapso.responder(msg, sem\_machine=True) — o colapso local determinístico; (2) consulta multi-tecidos (orquestrador\_multi); (3) prosa.responde. A conversa do cristal é local e determinística — a existência da cascata de 3 camadas é o mecanismo.
\prov{relações: usa cristal\_chat\_turn; usa conversa\_recupera\_nao\_gera; usa colapso\_ping\_pong\_overlap; usa colapso\_camadas\_ordem\_parada; relaciona tiffany}
\prov{proveniência: wiki \textperiodcentered{} code/cristal/tiffany/api.py:ask \ensuremath{\to} conversa/colapso.py:responder \textperiodcentered{} inferencia \textperiodcentered{} confiança 1.0 \textperiodcentered{} domínio cristal \textperiodcentered{} história 26 versões no jornal}

%CRISTAL {"arestas": [{"alvo": "grafo_de_proveniencia", "confianca": 1.0, "epistemico": "fato", "origem": "wiki", "rel": "parte_de"}], "confianca": 1.0, "contraexemplos": [], "descricao": "log_event grava em memoria/historico eventos de ingest/index/promote com detalhes (chunk_id, domain, entry); a trilha que liga ações a chunks. Alimenta o grafo_de_proveniencia. Funciona.", "epistemico": "fato", "exemplos": [], "id": "cristal_audit_historico", "memoria": "semantica", "meta": {"autor": "broca-so", "dominio": "cristal", "fonte": "code/cristal/tiffany_platform/memoria/historico.py:log_event", "wip": true}, "origem": "wiki", "palavras_chave": [], "sinonimos": [], "tipo": "conceito", "titulo": "Cristal — trilha de auditoria de proveniência"}
\section{Cristal — trilha de auditoria de proveniência}
log\_event grava em memoria/historico eventos de ingest/index/promote com detalhes (chunk\_id, domain, entry); a trilha que liga ações a chunks. Alimenta o grafo\_de\_proveniencia. Funciona.
\prov{relações: parte\_de grafo\_de\_proveniencia}
\prov{proveniência: wiki \textperiodcentered{} code/cristal/tiffany\_platform/memoria/historico.py:log\_event \textperiodcentered{} fato \textperiodcentered{} confiança 1.0 \textperiodcentered{} domínio cristal \textperiodcentered{} história 10 versões no jornal}

%CRISTAL {"arestas": [{"alvo": "cristal_cli", "confianca": 1.0, "epistemico": "fato", "origem": "wiki", "rel": "parte_de"}, {"alvo": "cristal_motor_congelado", "confianca": 1.0, "epistemico": "fato", "origem": "wiki", "rel": "usa"}], "confianca": 1.0, "contraexemplos": [], "descricao": "chat_turn() grava a mensagem do usuário na memória de conversa (append_turn), roteia por regex, chama tiffany.ask(orq=True, full=False) e grava a resposta; opcionalmente loga na sessão JSONL. Ponto de entrada concreto do cristal_motor_congelado. Funciona.", "epistemico": "fato", "exemplos": [], "id": "cristal_chat_turn", "memoria": "semantica", "meta": {"autor": "broca-so", "dominio": "cristal", "fonte": "code/cristal/commands/chat.py:chat_turn", "wip": true}, "origem": "wiki", "palavras_chave": [], "sinonimos": [], "tipo": "conceito", "titulo": "Cristal — o turno de chat (grava→roteia→ask)"}
\section{Cristal — o turno de chat (grava\ensuremath{\to}roteia\ensuremath{\to}ask)}
chat\_turn() grava a mensagem do usuário na memória de conversa (append\_turn), roteia por regex, chama tiffany.ask(orq=True, full=False) e grava a resposta; opcionalmente loga na sessão JSONL. Ponto de entrada concreto do cristal\_motor\_congelado. Funciona.
\prov{relações: parte\_de cristal\_cli; usa cristal\_motor\_congelado}
\prov{proveniência: wiki \textperiodcentered{} code/cristal/commands/chat.py:chat\_turn \textperiodcentered{} fato \textperiodcentered{} confiança 1.0 \textperiodcentered{} domínio cristal \textperiodcentered{} história 15 versões no jornal}

%CRISTAL {"arestas": [{"alvo": "cristal_cli", "confianca": 1.0, "epistemico": "fato", "origem": "wiki", "rel": "especializa"}, {"alvo": "console_do_operador_lab", "confianca": 1.0, "epistemico": "fato", "origem": "wiki", "rel": "relaciona"}, {"alvo": "tiffany", "confianca": 1.0, "epistemico": "fato", "origem": "wiki", "rel": "relaciona"}], "confianca": 1.0, "contraexemplos": [], "descricao": "A CLI do broca-so está em v0.2 funcional avançando para 0.3–0.5 (roadmap CRISTAL_DEV.md: 0.3=pipeline code/task automático, 0.4=deploy/hooks/paper publish, 0.5=editor+sessão JSONL). O Aarão avisou (2 jul 2026): 'o cli cristal também precisa avançar' — FRENTE ABERTA / work-in-progress, como o cockpit. O que roda, roda; o que falta está nomeado sem maquiagem.", "epistemico": "hipotese", "exemplos": [], "id": "cristal_cli_estado_wip", "memoria": "semantica", "meta": {"autor": "broca-so", "dominio": "cristal", "fonte": "code/cristal/README_CRISTAL.md; CRISTAL_DEV.md", "wip": true}, "origem": "wiki", "palavras_chave": [], "sinonimos": [], "tipo": "conceito", "titulo": "Cristal CLI — Estado (v0.2, frente aberta)"}
\section{Cristal CLI — Estado (v0.2, frente aberta)}
A CLI do broca-so está em v0.2 funcional avançando para 0.3–0.5 (roadmap CRISTAL\_DEV.md: 0.3=pipeline code/task automático, 0.4=deploy/hooks/paper publish, 0.5=editor+sessão JSONL). O Aarão avisou (2 jul 2026): 'o cli cristal também precisa avançar' — FRENTE ABERTA / work-in-progress, como o cockpit. O que roda, roda; o que falta está nomeado sem maquiagem.
\prov{relações: especializa cristal\_cli; relaciona console\_do\_operador\_lab; relaciona tiffany}
\prov{proveniência: wiki \textperiodcentered{} code/cristal/README\_CRISTAL.md; CRISTAL\_DEV.md \textperiodcentered{} hipotese \textperiodcentered{} confiança 1.0 \textperiodcentered{} domínio cristal \textperiodcentered{} história 16 versões no jornal}

%CRISTAL {"arestas": [{"alvo": "cristal_chat_turn", "confianca": 1.0, "epistemico": "fato", "origem": "wiki", "rel": "usa"}, {"alvo": "cristal_motor_congelado", "confianca": 1.0, "epistemico": "fato", "origem": "wiki", "rel": "relaciona"}], "confianca": 1.0, "contraexemplos": [], "descricao": "code explain/review e paper review não têm analisador próprio — montam um prompt (arquivo+pergunta) e chamam chat_turn, logo a 'revisão de código' é respondida pelo mesmo motor de colapso, não por um LLM. code fix/refactor/optimize/docs vão ao run_pipeline.", "epistemico": "inferencia", "exemplos": [], "id": "cristal_code_via_colapso", "memoria": "semantica", "meta": {"autor": "broca-so", "dominio": "cristal", "fonte": "code/cristal/commands/code.py:explain_file/review; paper.py:review", "wip": true}, "origem": "wiki", "palavras_chave": [], "sinonimos": [], "tipo": "conceito", "titulo": "Cristal — code review/explain reciclam o chat"}
\section{Cristal — code review/explain reciclam o chat}
code explain/review e paper review não têm analisador próprio — montam um prompt (arquivo+pergunta) e chamam chat\_turn, logo a 'revisão de código' é respondida pelo mesmo motor de colapso, não por um LLM. code fix/refactor/optimize/docs vão ao run\_pipeline.
\prov{relações: usa cristal\_chat\_turn; relaciona cristal\_motor\_congelado}
\prov{proveniência: wiki \textperiodcentered{} code/cristal/commands/code.py:explain\_file/review; paper.py:review \textperiodcentered{} inferencia \textperiodcentered{} confiança 1.0 \textperiodcentered{} domínio cristal \textperiodcentered{} história 15 versões no jornal}

%CRISTAL {"arestas": [{"alvo": "cristal_parser_registry", "confianca": 1.0, "epistemico": "fato", "origem": "wiki", "rel": "parte_de"}, {"alvo": "cristal_agentes_ortogonais", "confianca": 1.0, "epistemico": "fato", "origem": "wiki", "rel": "relaciona"}], "confianca": 1.0, "contraexemplos": [], "descricao": "Command define name/help/description + register/add_arguments (classmethods p/ injetar args) + execute(ctx,args)->int; Context é dataclass com home, out, verbose, config, project, session, shadow. (base.py é o COMANDO 'base', não a classe base — nome enganoso.)", "epistemico": "fato", "exemplos": [], "id": "cristal_command_base", "memoria": "semantica", "meta": {"autor": "broca-so", "dominio": "cristal", "fonte": "code/cristal/commands/_cmd.py:Command/Context", "wip": true}, "origem": "wiki", "palavras_chave": [], "sinonimos": [], "tipo": "conceito", "titulo": "Cristal — o contrato Command/Context"}
\section{Cristal — o contrato Command/Context}
Command define name/help/description + register/add\_arguments (classmethods p/ injetar args) + execute(ctx,args)->int; Context é dataclass com home, out, verbose, config, project, session, shadow. (base.py é o COMANDO 'base', não a classe base — nome enganoso.)
\prov{relações: parte\_de cristal\_parser\_registry; relaciona cristal\_agentes\_ortogonais}
\prov{proveniência: wiki \textperiodcentered{} code/cristal/commands/\_cmd.py:Command/Context \textperiodcentered{} fato \textperiodcentered{} confiança 1.0 \textperiodcentered{} domínio cristal \textperiodcentered{} história 15 versões no jornal}

%CRISTAL {"arestas": [{"alvo": "colapso_ping_pong_overlap", "confianca": 1.0, "epistemico": "fato", "origem": "wiki", "rel": "usa"}, {"alvo": "cristal_cristalizar_promote", "confianca": 1.0, "epistemico": "fato", "origem": "wiki", "rel": "relaciona"}], "confianca": 1.0, "contraexemplos": [], "descricao": "ingest (PDF/MD/TXT/Tex/HTML → chunks nos.tsv), index (nos.tsv→.idx via prosa.indexa), query (colapso/overlap: prosa._scan_py/_scan_c) e search (grep literal, explicitamente 'não usa colapso'). Camada em services.py. Ponta a ponta, funciona.", "epistemico": "fato", "exemplos": [], "id": "cristal_conhecimento_pipeline", "memoria": "semantica", "meta": {"autor": "broca-so", "dominio": "cristal", "fonte": "code/cristal/commands/{ingest,index,query,search}.py + services.py", "wip": true}, "origem": "wiki", "palavras_chave": [], "sinonimos": [], "tipo": "conceito", "titulo": "Cristal — conhecimento: ingest → index → query/search"}
\section{Cristal — conhecimento: ingest \ensuremath{\to} index \ensuremath{\to} query/search}
ingest (PDF/MD/TXT/Tex/HTML \ensuremath{\to} chunks nos.tsv), index (nos.tsv\ensuremath{\to}.idx via prosa.indexa), query (colapso/overlap: prosa.\_scan\_py/\_scan\_c) e search (grep literal, explicitamente 'não usa colapso'). Camada em services.py. Ponta a ponta, funciona.
\prov{relações: usa colapso\_ping\_pong\_overlap; relaciona cristal\_cristalizar\_promote}
\prov{proveniência: wiki \textperiodcentered{} code/cristal/commands/\{ingest,index,query,search\}.py + services.py \textperiodcentered{} fato \textperiodcentered{} confiança 1.0 \textperiodcentered{} domínio cristal \textperiodcentered{} história 15 versões no jornal}

%CRISTAL {"arestas": [{"alvo": "blend_contexto_32b", "confianca": 1.0, "epistemico": "fato", "origem": "wiki", "rel": "especializa"}, {"alvo": "cristal_carrega_identidade", "confianca": 1.0, "epistemico": "fato", "origem": "wiki", "rel": "parte_de"}, {"alvo": "contexto_32b_disco_nao_ram", "confianca": 1.0, "epistemico": "fato", "origem": "wiki", "rel": "relaciona"}], "confianca": 1.0, "contraexemplos": [], "descricao": "Cada turno atualiza contextosessao.bits (32 bytes) via média exponencial (blendctx, α=0.35) da impressão do texto (prosa.ᵢmpressaotexto) — o estado conversacional compacto que persiste entre turnos. Substrato do cristalcarregaᵢdentidade; a mesma doutrina do blendcontexto₃2b.", "epistemico": "fato", "exemplos": [], "id": "cristal_contexto_32b_ema", "memoria": "semantica", "meta": {"autor": "broca-so", "dominio": "cristal", "fonte": "code/cristal/tiffany_platform/memoria/conversa.py:append_turn/_blend_ctx", "wip": true}, "origem": "wiki", "palavras_chave": [], "sinonimos": [], "tipo": "conceito", "titulo": "Cristal — contexto de sessão como impressão de 256 bits"}
\section{Cristal — contexto de sessão como impressão de 256 bits}
Cada turno atualiza contextosessao.bits (32 bytes) via média exponencial (blendctx, \ensuremath{\alpha}=0.35) da impressão do texto (prosa.\ensuremath{_{i}}mpressaotexto) — o estado conversacional compacto que persiste entre turnos. Substrato do cristalcarrega\ensuremath{_{i}}dentidade; a mesma doutrina do blendcontexto\ensuremath{_{3}}2b.
\prov{relações: especializa blend\_contexto\_32b; parte\_de cristal\_carrega\_identidade; relaciona contexto\_32b\_disco\_nao\_ram}
\prov{proveniência: wiki \textperiodcentered{} code/cristal/tiffany\_platform/memoria/conversa.py:append\_turn/\_blend\_ctx \textperiodcentered{} fato \textperiodcentered{} confiança 1.0 \textperiodcentered{} domínio cristal \textperiodcentered{} história 21 versões no jornal}

%CRISTAL {"arestas": [{"alvo": "cristal_memoria_viva_jsonl", "confianca": 1.0, "epistemico": "fato", "origem": "wiki", "rel": "usa"}, {"alvo": "grafo_de_proveniencia", "confianca": 1.0, "epistemico": "fato", "origem": "wiki", "rel": "relaciona"}], "confianca": 1.0, "contraexemplos": [], "descricao": "promote(entryᵢd, domain) transforma uma entrada de memória em chunk novo em cliente/domain/nos.tsv com origem=memoria/tipo/id, atualiza meta/manifest, reindexa (prosa.indexa) e marca a entrada como 'promovido'. A cristalização curada: memória viva → conhecimento permanente (a mesma vocação de ensinarconceito, com proveniência).", "epistemico": "fato", "exemplos": [], "id": "cristal_cristalizar_promote", "memoria": "semantica", "meta": {"autor": "broca-so", "dominio": "cristal", "fonte": "code/cristal/tiffany_platform/memoria/promote.py:promote", "wip": true}, "origem": "wiki", "palavras_chave": [], "sinonimos": [], "tipo": "conceito", "titulo": "Cristal — cristalização: memória → chunk indexado"}
\section{Cristal — cristalização: memória \ensuremath{\to} chunk indexado}
promote(entry\ensuremath{_{i}}d, domain) transforma uma entrada de memória em chunk novo em cliente/domain/nos.tsv com origem=memoria/tipo/id, atualiza meta/manifest, reindexa (prosa.indexa) e marca a entrada como 'promovido'. A cristalização curada: memória viva \ensuremath{\to} conhecimento permanente (a mesma vocação de ensinarconceito, com proveniência).
\prov{relações: usa cristal\_memoria\_viva\_jsonl; relaciona grafo\_de\_proveniencia}
\prov{proveniência: wiki \textperiodcentered{} code/cristal/tiffany\_platform/memoria/promote.py:promote \textperiodcentered{} fato \textperiodcentered{} confiança 1.0 \textperiodcentered{} domínio cristal \textperiodcentered{} história 16 versões no jornal}

%CRISTAL {"arestas": [{"alvo": "cristal_sessao_jsonl", "confianca": 1.0, "epistemico": "fato", "origem": "wiki", "rel": "relaciona"}, {"alvo": "cristal_carrega_identidade", "confianca": 1.0, "epistemico": "fato", "origem": "wiki", "rel": "relaciona"}], "confianca": 1.0, "contraexemplos": [], "descricao": "Duas persistências paralelas: (a) shell/session.py grava.cristal/sessions/nome.jsonl por projeto (histórico REPL), e (b) memoria/conversa/sessao_*.jsonl + contexto_*.bits GLOBAL do home. chatturn escreve em ambas quando logsession. Identidade global vs sessão local.", "epistemico": "inferencia", "exemplos": [], "id": "cristal_duas_memorias", "memoria": "semantica", "meta": {"autor": "broca-so", "dominio": "cristal", "fonte": "code/cristal/shell/session.py vs memoria/conversa.py", "wip": true}, "origem": "wiki", "palavras_chave": [], "sinonimos": [], "tipo": "conceito", "titulo": "Cristal — sessão REPL de projeto ≠ memória de conversa"}
\section{Cristal — sessão REPL de projeto \ensuremath{\ne} memória de conversa}
Duas persistências paralelas: (a) shell/session.py grava.cristal/sessions/nome.jsonl por projeto (histórico REPL), e (b) memoria/conversa/sessao\_*.jsonl + contexto\_*.bits GLOBAL do home. chatturn escreve em ambas quando logsession. Identidade global vs sessão local.
\prov{relações: relaciona cristal\_sessao\_jsonl; relaciona cristal\_carrega\_identidade}
\prov{proveniência: wiki \textperiodcentered{} code/cristal/shell/session.py vs memoria/conversa.py \textperiodcentered{} inferencia \textperiodcentered{} confiança 1.0 \textperiodcentered{} domínio cristal \textperiodcentered{} história 16 versões no jornal}

%CRISTAL {"arestas": [{"alvo": "cristal_cli", "confianca": 1.0, "epistemico": "fato", "origem": "wiki", "rel": "parte_de"}, {"alvo": "cristal_cli_estado_wip", "confianca": 1.0, "epistemico": "fato", "origem": "wiki", "rel": "parte_de"}], "confianca": 1.0, "contraexemplos": [], "descricao": "main() chama _register_all() que importa ~26 módulos de comando (cada um se auto-registra via decorator @register), depois delega a core/parser.run(); sem subcomando cai no REPL. A coluna vertebral concreta do cristal_cli.", "epistemico": "fato", "exemplos": [], "id": "cristal_entrypoint_registro", "memoria": "semantica", "meta": {"autor": "broca-so", "dominio": "cristal", "fonte": "code/cristal/main.py:_register_all/main", "wip": true}, "origem": "wiki", "palavras_chave": [], "sinonimos": [], "tipo": "conceito", "titulo": "Cristal — entrypoint → parser → comando"}
\section{Cristal — entrypoint \ensuremath{\to} parser \ensuremath{\to} comando}
main() chama \_register\_all() que importa \~{}26 módulos de comando (cada um se auto-registra via decorator @register), depois delega a core/parser.run(); sem subcomando cai no REPL. A coluna vertebral concreta do cristal\_cli.
\prov{relações: parte\_de cristal\_cli; parte\_de cristal\_cli\_estado\_wip}
\prov{proveniência: wiki \textperiodcentered{} code/cristal/main.py:\_register\_all/main \textperiodcentered{} fato \textperiodcentered{} confiança 1.0 \textperiodcentered{} domínio cristal \textperiodcentered{} história 15 versões no jornal}

%CRISTAL {"arestas": [{"alvo": "cristal_pipeline_runner", "confianca": 1.0, "epistemico": "fato", "origem": "wiki", "rel": "usa"}], "confianca": 1.0, "contraexemplos": [], "descricao": "runhooks() carrega plugins/<nome>/hooks.py via importlib e chama on_<evento>; eventos: pipeline:before,after,step, deploy:before,after, memory:compact. commands/plugins.py gerencia install/list/remove/hooks. Infra pronta, ecossistema de plugins ainda vazio.", "epistemico": "fato", "exemplos": [], "id": "cristal_hooks_plugins", "memoria": "semantica", "meta": {"autor": "broca-so", "dominio": "cristal", "fonte": "code/cristal/plugins/hooks.py:run_hooks/HOOK_EVENTS", "wip": true}, "origem": "wiki", "palavras_chave": [], "sinonimos": [], "tipo": "conceito", "titulo": "Cristal — hooks por plugin"}
\section{Cristal — hooks por plugin}
runhooks() carrega plugins/<nome>/hooks.py via importlib e chama on\_<evento>; eventos: pipeline:before,after,step, deploy:before,after, memory:compact. commands/plugins.py gerencia install/list/remove/hooks. Infra pronta, ecossistema de plugins ainda vazio.
\prov{relações: usa cristal\_pipeline\_runner}
\prov{proveniência: wiki \textperiodcentered{} code/cristal/plugins/hooks.py:run\_hooks/HOOK\_EVENTS \textperiodcentered{} fato \textperiodcentered{} confiança 1.0 \textperiodcentered{} domínio cristal \textperiodcentered{} história 11 versões no jornal}

%CRISTAL {"arestas": [{"alvo": "cristal_cli", "confianca": 1.0, "epistemico": "fato", "origem": "wiki", "rel": "parte_de"}, {"alvo": "cristal_entrypoint_registro", "confianca": 1.0, "epistemico": "fato", "origem": "wiki", "rel": "relaciona"}], "confianca": 1.0, "contraexemplos": [], "descricao": "Registro flat + namespaces (tecido/memoria/bench/release via groups.py). Conhecimento: ingest,index,rebuild,query,search,domains,stats,validate,export,import,gc,doctor (Fase 0.1 ✓). Código: code explain,review,tests,benchmark,open,refactor,fix,optimize,docs,commit. Tarefa/agente: task,work,agent,chat. Deploy/release: deploy plan,pack,verify,full,base,release,qa,paper build,review,refs,publish. Núcleo: init,projeto,new,memory,plugins,update,editor.", "epistemico": "fato", "exemplos": [], "id": "cristal_mapa_comandos", "memoria": "semantica", "meta": {"autor": "broca-so", "dominio": "cristal", "fonte": "code/cristal/main.py:_register_all; commands/groups.py", "wip": true}, "origem": "wiki", "palavras_chave": [], "sinonimos": [], "tipo": "conceito", "titulo": "Cristal — o mapa dos ~30 comandos"}
\section{Cristal — o mapa dos \~{}30 comandos}
Registro flat + namespaces (tecido/memoria/bench/release via groups.py). Conhecimento: ingest,index,rebuild,query,search,domains,stats,validate,export,import,gc,doctor (Fase 0.1 \checkmark ). Código: code explain,review,tests,benchmark,open,refactor,fix,optimize,docs,commit. Tarefa/agente: task,work,agent,chat. Deploy/release: deploy plan,pack,verify,full,base,release,qa,paper build,review,refs,publish. Núcleo: init,projeto,new,memory,plugins,update,editor.
\prov{relações: parte\_de cristal\_cli; relaciona cristal\_entrypoint\_registro}
\prov{proveniência: wiki \textperiodcentered{} code/cristal/main.py:\_register\_all; commands/groups.py \textperiodcentered{} fato \textperiodcentered{} confiança 1.0 \textperiodcentered{} domínio cristal \textperiodcentered{} história 16 versões no jornal}

%CRISTAL {"arestas": [{"alvo": "cristal_memoria_viva_jsonl", "confianca": 1.0, "epistemico": "fato", "origem": "wiki", "rel": "usa"}, {"alvo": "cristal_cli_estado_wip", "confianca": 1.0, "epistemico": "fato", "origem": "wiki", "rel": "relaciona"}, {"alvo": "contexto_32b_disco_nao_ram", "confianca": 1.0, "epistemico": "fato", "origem": "wiki", "rel": "relaciona"}, {"alvo": "cristal_portao_rollback", "confianca": 1.0, "epistemico": "fato", "origem": "wiki", "rel": "relaciona"}], "confianca": 1.0, "contraexemplos": [], "descricao": "compact_all arquiva entradas promovido/concluida e dedupe histórico; compact_conversa agora SINTETIZA: _sintetizar faz síntese extrativa determinística (TF/Luhn — sem LLM) dos turnos arquivados: pontua cada turno pela frequência dos seus termos de conteúdo, bônus a perguntas, dedupe por sobreposição, ordem cronológica. A síntese vai para o TOPO do arquivo vivo (onde a continuidade lê), não só para o morto. Sínteses sucessivas dobram (acumulam pontos+termos, cap SUMMARY_PONTOS/TERMOS) sem churn (a síntese não conta p/ o gatilho de truncagem). Superou o gargalo antigo ('só uma contagem').", "epistemico": "inferencia", "exemplos": [], "id": "cristal_memoria_compact", "memoria": "semantica", "meta": {"autor": "broca-so", "dominio": "cristal", "fonte": "code/tiffany_platform/memoria/compact.py:_sintetizar/compact_conversa", "wip": true}, "origem": "wiki", "palavras_chave": [], "sinonimos": [], "tipo": "conceito", "titulo": "Cristal — compactação com síntese extrativa (não mais só contagem)"}
\section{Cristal — compactação com síntese extrativa (não mais só contagem)}
compact\_all arquiva entradas promovido/concluida e dedupe histórico; compact\_conversa agora SINTETIZA: \_sintetizar faz síntese extrativa determinística (TF/Luhn — sem LLM) dos turnos arquivados: pontua cada turno pela frequência dos seus termos de conteúdo, bônus a perguntas, dedupe por sobreposição, ordem cronológica. A síntese vai para o TOPO do arquivo vivo (onde a continuidade lê), não só para o morto. Sínteses sucessivas dobram (acumulam pontos+termos, cap SUMMARY\_PONTOS/TERMOS) sem churn (a síntese não conta p/ o gatilho de truncagem). Superou o gargalo antigo ('só uma contagem').
\prov{relações: usa cristal\_memoria\_viva\_jsonl; relaciona cristal\_cli\_estado\_wip; relaciona contexto\_32b\_disco\_nao\_ram; relaciona cristal\_portao\_rollback}
\prov{proveniência: wiki \textperiodcentered{} code/tiffany\_platform/memoria/compact.py:\_sintetizar/compact\_conversa \textperiodcentered{} inferencia \textperiodcentered{} confiança 1.0 \textperiodcentered{} domínio cristal \textperiodcentered{} história 19 versões no jornal}

%CRISTAL {"arestas": [{"alvo": "cristal_carrega_identidade", "confianca": 1.0, "epistemico": "fato", "origem": "wiki", "rel": "parte_de"}, {"alvo": "grafo_de_proveniencia", "confianca": 1.0, "epistemico": "fato", "origem": "wiki", "rel": "relaciona"}], "confianca": 1.0, "contraexemplos": [], "descricao": "memory list/append/promote/clean/compact sobre memoria/tarefas,bugs,decisoes,historico/registro.jsonl; cada entrada tem id, status (ativo/promovido/concluida), meta. conversa usa fluxo separado. A identidade que persiste entre sessões. Funciona (compact no CLI parcial, backend existe).", "epistemico": "fato", "exemplos": [], "id": "cristal_memoria_viva_jsonl", "memoria": "semantica", "meta": {"autor": "broca-so", "dominio": "cristal", "fonte": "code/cristal/commands/memory.py + tiffany_platform/memoria/store.py", "wip": true}, "origem": "wiki", "palavras_chave": [], "sinonimos": [], "tipo": "conceito", "titulo": "Cristal — memória viva: JSONL append-only por tipo"}
\section{Cristal — memória viva: JSONL append-only por tipo}
memory list/append/promote/clean/compact sobre memoria/tarefas,bugs,decisoes,historico/registro.jsonl; cada entrada tem id, status (ativo/promovido/concluida), meta. conversa usa fluxo separado. A identidade que persiste entre sessões. Funciona (compact no CLI parcial, backend existe).
\prov{relações: parte\_de cristal\_carrega\_identidade; relaciona grafo\_de\_proveniencia}
\prov{proveniência: wiki \textperiodcentered{} code/cristal/commands/memory.py + tiffany\_platform/memoria/store.py \textperiodcentered{} fato \textperiodcentered{} confiança 1.0 \textperiodcentered{} domínio cristal \textperiodcentered{} história 16 versões no jornal}

%CRISTAL {"arestas": [{"alvo": "cristal_entrypoint_registro", "confianca": 1.0, "epistemico": "fato", "origem": "wiki", "rel": "usa"}, {"alvo": "cristal_agentes_ortogonais", "confianca": 1.0, "epistemico": "fato", "origem": "wiki", "rel": "relaciona"}], "confianca": 1.0, "contraexemplos": [], "descricao": "build_parser() monta um argparse com args comuns (--home, -S/--session, -v, -q) e um subparser por comando registrado; run() resolve projeto+home+config, constrói o Context e despacha args._cmd_cls().execute(ctx,args). Padrão registry + Command class — cada comando é ortogonal e plugável.", "epistemico": "inferencia", "exemplos": [], "id": "cristal_parser_registry", "memoria": "semantica", "meta": {"autor": "broca-so", "dominio": "cristal", "fonte": "code/cristal/core/parser.py:build_parser/run", "wip": true}, "origem": "wiki", "palavras_chave": [], "sinonimos": [], "tipo": "conceito", "titulo": "Cristal — parser como registry de Command"}
\section{Cristal — parser como registry de Command}
build\_parser() monta um argparse com args comuns (--home, -S/--session, -v, -q) e um subparser por comando registrado; run() resolve projeto+home+config, constrói o Context e despacha args.\_cmd\_cls().execute(ctx,args). Padrão registry + Command class — cada comando é ortogonal e plugável.
\prov{relações: usa cristal\_entrypoint\_registro; relaciona cristal\_agentes\_ortogonais}
\prov{proveniência: wiki \textperiodcentered{} code/cristal/core/parser.py:build\_parser/run \textperiodcentered{} inferencia \textperiodcentered{} confiança 1.0 \textperiodcentered{} domínio cristal \textperiodcentered{} história 15 versões no jornal}

%CRISTAL {"arestas": [{"alvo": "cristal_cli_estado_wip", "confianca": 1.0, "epistemico": "fato", "origem": "wiki", "rel": "relaciona"}, {"alvo": "cristal_pipeline_runner", "confianca": 1.0, "epistemico": "fato", "origem": "wiki", "rel": "parte_de"}, {"alvo": "cristal_motor_congelado", "confianca": 1.0, "epistemico": "fato", "origem": "wiki", "rel": "relaciona"}, {"alvo": "cristal_analyze_estatico", "confianca": 1.0, "epistemico": "fato", "origem": "wiki", "rel": "usa"}], "confianca": 1.0, "contraexemplos": [], "descricao": "transform_source() aplica transformações determinísticas derivadas dos findings — NÃO fabrica, NÃO usa LLM. Regras de arranque na doutrina Broca: (1) remover buffer static não usado (nome só na declaração → RAM a menos), (2) promover static→static const (array inicializado, nunca escrito → .rodata). O comentário-marcador (generate_comment_patch) virou FALLBACK idempotente quando nenhuma regra casa. O portão é o backstop: se a heurística errar, não compila/testa → rollback. Superou o gargalo antigo (que era 'só marcador').", "epistemico": "inferencia", "exemplos": [], "id": "cristal_patch_marcador", "memoria": "semantica", "meta": {"autor": "broca-so", "dominio": "cristal", "fonte": "code/cristal/pipeline/patch.py:transform_source/generate_edit_patch", "wip": true}, "origem": "wiki", "palavras_chave": [], "sinonimos": [], "tipo": "conceito", "titulo": "Cristal — patch de edição REAL (regras determinísticas)"}
\section{Cristal — patch de edição REAL (regras determinísticas)}
transform\_source() aplica transformações determinísticas derivadas dos findings — NÃO fabrica, NÃO usa LLM. Regras de arranque na doutrina Broca: (1) remover buffer static não usado (nome só na declaração \ensuremath{\to} RAM a menos), (2) promover static\ensuremath{\to}static const (array inicializado, nunca escrito \ensuremath{\to} .rodata). O comentário-marcador (generate\_comment\_patch) virou FALLBACK idempotente quando nenhuma regra casa. O portão é o backstop: se a heurística errar, não compila/testa \ensuremath{\to} rollback. Superou o gargalo antigo (que era 'só marcador').
\prov{relações: relaciona cristal\_cli\_estado\_wip; parte\_de cristal\_pipeline\_runner; relaciona cristal\_motor\_congelado; usa cristal\_analyze\_estatico}
\prov{proveniência: wiki \textperiodcentered{} code/cristal/pipeline/patch.py:transform\_source/generate\_edit\_patch \textperiodcentered{} inferencia \textperiodcentered{} confiança 1.0 \textperiodcentered{} domínio cristal \textperiodcentered{} história 20 versões no jornal}

%CRISTAL {"arestas": [{"alvo": "cristal_cli", "confianca": 1.0, "epistemico": "fato", "origem": "wiki", "rel": "parte_de"}, {"alvo": "cristal_motor_congelado", "confianca": 1.0, "epistemico": "fato", "origem": "wiki", "rel": "usa"}, {"alvo": "shell", "confianca": 1.0, "epistemico": "fato", "origem": "wiki", "rel": "relaciona"}], "confianca": 1.0, "contraexemplos": [], "descricao": "Pipeline.run() executa steps sequenciais: abrir projeto → memória → tecidos → localizar → analisar → medir RSS → gerar patch → gerar proposta → aplicar patch → compilar → testes → bench → diff, cada um via _step() com hooks e StepResult. O coração que o cristal_cli prometia. Várias etapas são 'skip' defensivo quando falta binário/Makefile/test.sh.", "epistemico": "fato", "exemplos": [], "id": "cristal_pipeline_runner", "memoria": "semantica", "meta": {"autor": "broca-so", "dominio": "cristal", "fonte": "code/cristal/pipeline/runner.py:Pipeline.run/_step", "wip": true}, "origem": "wiki", "palavras_chave": [], "sinonimos": [], "tipo": "conceito", "titulo": "Cristal — o pipeline de código (11 etapas)"}
\section{Cristal — o pipeline de código (11 etapas)}
Pipeline.run() executa steps sequenciais: abrir projeto \ensuremath{\to} memória \ensuremath{\to} tecidos \ensuremath{\to} localizar \ensuremath{\to} analisar \ensuremath{\to} medir RSS \ensuremath{\to} gerar patch \ensuremath{\to} gerar proposta \ensuremath{\to} aplicar patch \ensuremath{\to} compilar \ensuremath{\to} testes \ensuremath{\to} bench \ensuremath{\to} diff, cada um via \_step() com hooks e StepResult. O coração que o cristal\_cli prometia. Várias etapas são 'skip' defensivo quando falta binário/Makefile/test.sh.
\prov{relações: parte\_de cristal\_cli; usa cristal\_motor\_congelado; relaciona shell}
\prov{proveniência: wiki \textperiodcentered{} code/cristal/pipeline/runner.py:Pipeline.run/\_step \textperiodcentered{} fato \textperiodcentered{} confiança 1.0 \textperiodcentered{} domínio cristal \textperiodcentered{} história 16 versões no jornal}

%CRISTAL {"arestas": [{"alvo": "cristal_patch_marcador", "confianca": 1.0, "epistemico": "fato", "origem": "wiki", "rel": "confirma"}, {"alvo": "cristal_pipeline_runner", "confianca": 1.0, "epistemico": "fato", "origem": "wiki", "rel": "parte_de"}, {"alvo": "cristal_runner_executa", "confianca": 1.0, "epistemico": "fato", "origem": "wiki", "rel": "usa"}], "confianca": 1.0, "contraexemplos": [], "descricao": "A edição real é segura porque o pipeline verifica e reverte: _apply_patch aplica via patch(1) (fallback marcador se o diff não casar), _compile reverte se make falhar após patch (reforço novo), _test reverte se test.sh falhar. 'O portão é a verdade, não a promessa' — a mesma epistemologia do PoW: não confiar na heurística, verificar. Deixa fazer edições 'provavelmente seguras' e deixar o compilador/testes rejeitarem as ruins.", "epistemico": "inferencia", "exemplos": [], "id": "cristal_portao_rollback", "memoria": "semantica", "meta": {"autor": "broca-so", "dominio": "cristal", "fonte": "code/cristal/pipeline/runner.py:_apply_patch/_compile/_test; patch.py:rollback", "wip": true}, "origem": "wiki", "palavras_chave": [], "sinonimos": [], "tipo": "conceito", "titulo": "Cristal — o portão é a verdade (rollback em compilar E testar)"}
\section{Cristal — o portão é a verdade (rollback em compilar E testar)}
A edição real é segura porque o pipeline verifica e reverte: \_apply\_patch aplica via patch(1) (fallback marcador se o diff não casar), \_compile reverte se make falhar após patch (reforço novo), \_test reverte se test.sh falhar. 'O portão é a verdade, não a promessa' — a mesma epistemologia do PoW: não confiar na heurística, verificar. Deixa fazer edições 'provavelmente seguras' e deixar o compilador/testes rejeitarem as ruins.
\prov{relações: confirma cristal\_patch\_marcador; parte\_de cristal\_pipeline\_runner; usa cristal\_runner\_executa}
\prov{proveniência: wiki \textperiodcentered{} code/cristal/pipeline/runner.py:\_apply\_patch/\_compile/\_test; patch.py:rollback \textperiodcentered{} inferencia \textperiodcentered{} confiança 1.0 \textperiodcentered{} domínio cristal \textperiodcentered{} história 16 versões no jornal}

%CRISTAL {"arestas": [{"alvo": "cristal_carrega_identidade", "confianca": 1.0, "epistemico": "fato", "origem": "wiki", "rel": "explica"}, {"alvo": "cristal_parser_registry", "confianca": 1.0, "epistemico": "fato", "origem": "wiki", "rel": "usa"}], "confianca": 1.0, "contraexemplos": [], "descricao": "find_project() sobe diretórios procurando cristal.json ou marcadores broca-so (code/prosa.py, code/libtiffany, test/test.sh); resolve tiffany_home relativo ao root e carrega agente/branch. É como cristal_carrega_identidade opera na prática: contexto do repo → agente + home.", "epistemico": "fato", "exemplos": [], "id": "cristal_projeto_contexto", "memoria": "semantica", "meta": {"autor": "broca-so", "dominio": "cristal", "fonte": "code/cristal/project/context.py:find_project/load_project", "wip": true}, "origem": "wiki", "palavras_chave": [], "sinonimos": [], "tipo": "conceito", "titulo": "Cristal — descoberta de projeto (cristal.json)"}
\section{Cristal — descoberta de projeto (cristal.json)}
find\_project() sobe diretórios procurando cristal.json ou marcadores broca-so (code/prosa.py, code/libtiffany, test/test.sh); resolve tiffany\_home relativo ao root e carrega agente/branch. É como cristal\_carrega\_identidade opera na prática: contexto do repo \ensuremath{\to} agente + home.
\prov{relações: explica cristal\_carrega\_identidade; usa cristal\_parser\_registry}
\prov{proveniência: wiki \textperiodcentered{} code/cristal/project/context.py:find\_project/load\_project \textperiodcentered{} fato \textperiodcentered{} confiança 1.0 \textperiodcentered{} domínio cristal \textperiodcentered{} história 15 versões no jornal}

%CRISTAL {"arestas": [{"alvo": "cristal_carrega_identidade", "confianca": 1.0, "epistemico": "fato", "origem": "wiki", "rel": "explica"}, {"alvo": "cristal_memoria_compact", "confianca": 1.0, "epistemico": "fato", "origem": "wiki", "rel": "usa"}, {"alvo": "cristal_contexto_32b_ema", "confianca": 1.0, "epistemico": "fato", "origem": "wiki", "rel": "usa"}, {"alvo": "contexto_32b_disco_nao_ram", "confianca": 1.0, "epistemico": "fato", "origem": "wiki", "rel": "especializa"}, {"alvo": "colapso_camadas_ordem_parada", "confianca": 1.0, "epistemico": "fato", "origem": "wiki", "rel": "usa"}, {"alvo": "cristal_motor_congelado", "confianca": 1.0, "epistemico": "fato", "origem": "wiki", "rel": "confirma"}, {"alvo": "cristal_chat_turn", "confianca": 1.0, "epistemico": "fato", "origem": "wiki", "rel": "parte_de"}, {"alvo": "blend_contexto_32b", "confianca": 1.0, "epistemico": "fato", "origem": "wiki", "rel": "relaciona"}, {"alvo": "tiffany", "confianca": 1.0, "epistemico": "fato", "origem": "wiki", "rel": "relaciona"}], "confianca": 1.0, "contraexemplos": [], "descricao": "O colapso usa um CTX global (session-agnostic) e a memória de sessão do cristal era vestigial p/ recall. chatturn agora semeia (seedrecall) o CTX do colapso a partir da memória da sessão — o contexto persistido (recallsessao.bits) ou, a FRIO, a impressão da síntese compactada [compact] no topo da sessão — e persiste (persistrecall) o CTX pós-turno de volta. Sessão nova → resetctx (slate limpo, não herda outra). Só funções públicas do colapso (savectx/loadctx/resetctx): a ponte vive na periferia, o motor fica congelado. Fecha o elo síntese→recall. Limite honesto: ULTIMO/HIST (ping-pong) seguem globais.", "epistemico": "inferencia", "exemplos": [], "id": "cristal_recall_ponte", "memoria": "semantica", "meta": {"autor": "broca-so", "dominio": "cristal", "fonte": "code/cristal/commands/chat.py:_seed_recall/_persist_recall; tiffany_platform/memoria/conversa.py:load_recall_ctx", "wip": true}, "origem": "wiki", "palavras_chave": [], "sinonimos": [], "tipo": "conceito", "titulo": "Cristal — ponte de recall (a memória da sessão vieja o colapso)"}
\section{Cristal — ponte de recall (a memória da sessão vieja o colapso)}
O colapso usa um CTX global (session-agnostic) e a memória de sessão do cristal era vestigial p/ recall. chatturn agora semeia (seedrecall) o CTX do colapso a partir da memória da sessão — o contexto persistido (recallsessao.bits) ou, a FRIO, a impressão da síntese compactada [compact] no topo da sessão — e persiste (persistrecall) o CTX pós-turno de volta. Sessão nova \ensuremath{\to} resetctx (slate limpo, não herda outra). Só funções públicas do colapso (savectx/loadctx/resetctx): a ponte vive na periferia, o motor fica congelado. Fecha o elo síntese\ensuremath{\to}recall. Limite honesto: ULTIMO/HIST (ping-pong) seguem globais.
\prov{relações: explica cristal\_carrega\_identidade; usa cristal\_memoria\_compact; usa cristal\_contexto\_32b\_ema; especializa contexto\_32b\_disco\_nao\_ram; usa colapso\_camadas\_ordem\_parada; confirma cristal\_motor\_congelado; parte\_de cristal\_chat\_turn; relaciona blend\_contexto\_32b; relaciona tiffany}
\prov{proveniência: wiki \textperiodcentered{} code/cristal/commands/chat.py:\_seed\_recall/\_persist\_recall; tiffany\_platform/memoria/conversa.py:load\_recall\_ctx \textperiodcentered{} inferencia \textperiodcentered{} confiança 1.0 \textperiodcentered{} domínio cristal \textperiodcentered{} história 11 versões no jornal}

%CRISTAL {"arestas": [{"alvo": "console_do_operador_lab", "confianca": 1.0, "epistemico": "fato", "origem": "wiki", "rel": "especializa"}, {"alvo": "cristal_entrypoint_registro", "confianca": 1.0, "epistemico": "fato", "origem": "wiki", "rel": "usa"}], "confianca": 1.0, "contraexemplos": [], "descricao": "run_shell() faz loop de input() com prompt colorido; _dispatch() interpreta comandos internos (help, historico, open, agent, explain, bench, doctor, domains, shadow, task, work, !shell) e cai em chat_turn p/ texto livre. A implementação concreta do console do operador.", "epistemico": "fato", "exemplos": [], "id": "cristal_repl_shell", "memoria": "semantica", "meta": {"autor": "broca-so", "dominio": "cristal", "fonte": "code/cristal/shell/repl.py:run_shell/_dispatch", "wip": true}, "origem": "wiki", "palavras_chave": [], "sinonimos": [], "tipo": "conceito", "titulo": "Cristal — o REPL (console do operador concreto)"}
\section{Cristal — o REPL (console do operador concreto)}
run\_shell() faz loop de input() com prompt colorido; \_dispatch() interpreta comandos internos (help, historico, open, agent, explain, bench, doctor, domains, shadow, task, work, !shell) e cai em chat\_turn p/ texto livre. A implementação concreta do console do operador.
\prov{relações: especializa console\_do\_operador\_lab; usa cristal\_entrypoint\_registro}
\prov{proveniência: wiki \textperiodcentered{} code/cristal/shell/repl.py:run\_shell/\_dispatch \textperiodcentered{} fato \textperiodcentered{} confiança 1.0 \textperiodcentered{} domínio cristal \textperiodcentered{} história 15 versões no jornal}

%CRISTAL {"arestas": [{"alvo": "cristal_cli_estado_wip", "confianca": 1.0, "epistemico": "fato", "origem": "wiki", "rel": "relaciona"}, {"alvo": "cristal_agentes_catalogo", "confianca": 1.0, "epistemico": "fato", "origem": "wiki", "rel": "relaciona"}, {"alvo": "cristal_chat_turn", "confianca": 1.0, "epistemico": "fato", "origem": "wiki", "rel": "usa"}, {"alvo": "cristal_agente_vies_dominio", "confianca": 1.0, "epistemico": "fato", "origem": "wiki", "rel": "usa"}], "confianca": 1.0, "contraexemplos": [], "descricao": "roteia_agente casa regex (.tex→Paper, make/bug→Software, indexa/corpus→Pesquisa, senão Conversa) e o resultado agora é PASSADO ao ask (antes era só impresso como detail — o gargalo cosmético). chat_turn calcula agente_efetivo (mensagem especializada manda; senão o agente do projeto) e o thread desce ask→consulta→selecionar_dominios. Mesma pergunta, papel diferente → domínio diferente. Gargalo cosmético FECHADO.", "epistemico": "inferencia", "exemplos": [], "id": "cristal_roteamento_cosmetico", "memoria": "semantica", "meta": {"autor": "broca-so", "dominio": "cristal", "fonte": "code/cristal/commands/chat.py; code/tiffany/api.py:ask", "wip": true}, "origem": "wiki", "palavras_chave": [], "sinonimos": [], "tipo": "conceito", "titulo": "Cristal — o roteamento de agente agora VALE (viés de domínio)"}
\section{Cristal — o roteamento de agente agora VALE (viés de domínio)}
roteia\_agente casa regex (.tex\ensuremath{\to}Paper, make/bug\ensuremath{\to}Software, indexa/corpus\ensuremath{\to}Pesquisa, senão Conversa) e o resultado agora é PASSADO ao ask (antes era só impresso como detail — o gargalo cosmético). chat\_turn calcula agente\_efetivo (mensagem especializada manda; senão o agente do projeto) e o thread desce ask\ensuremath{\to}consulta\ensuremath{\to}selecionar\_dominios. Mesma pergunta, papel diferente \ensuremath{\to} domínio diferente. Gargalo cosmético FECHADO.
\prov{relações: relaciona cristal\_cli\_estado\_wip; relaciona cristal\_agentes\_catalogo; usa cristal\_chat\_turn; usa cristal\_agente\_vies\_dominio}
\prov{proveniência: wiki \textperiodcentered{} code/cristal/commands/chat.py; code/tiffany/api.py:ask \textperiodcentered{} inferencia \textperiodcentered{} confiança 1.0 \textperiodcentered{} domínio cristal \textperiodcentered{} história 20 versões no jornal}

%CRISTAL {"arestas": [{"alvo": "cristal_pipeline_runner", "confianca": 1.0, "epistemico": "fato", "origem": "wiki", "rel": "parte_de"}, {"alvo": "broca_os", "confianca": 1.0, "epistemico": "fato", "origem": "wiki", "rel": "relaciona"}], "confianca": 1.0, "contraexemplos": [], "descricao": "_measure_rss usa /usr/bin/time -f %M sobre binário compilado; _compile roda make -s em code/; _test roda test/test.sh (rollback se falhar após patch); _bench roda benchmarks/fogo.sh; _git_diff salva git diff em .cristal/work/<ts>/. Instrumenta a doutrina Broca de RAM/RSS; funciona quando os artefatos existem.", "epistemico": "fato", "exemplos": [], "id": "cristal_runner_executa", "memoria": "semantica", "meta": {"autor": "broca-so", "dominio": "cristal", "fonte": "code/cristal/pipeline/runner.py:_measure_rss/_compile/_test/_bench", "wip": true}, "origem": "wiki", "palavras_chave": [], "sinonimos": [], "tipo": "conceito", "titulo": "Cristal — como o pipeline roda o código"}
\section{Cristal — como o pipeline roda o código}
\_measure\_rss usa /usr/bin/time -f \%M sobre binário compilado; \_compile roda make -s em code/; \_test roda test/test.sh (rollback se falhar após patch); \_bench roda benchmarks/fogo.sh; \_git\_diff salva git diff em .cristal/work/<ts>/. Instrumenta a doutrina Broca de RAM/RSS; funciona quando os artefatos existem.
\prov{relações: parte\_de cristal\_pipeline\_runner; relaciona broca\_os}
\prov{proveniência: wiki \textperiodcentered{} code/cristal/pipeline/runner.py:\_measure\_rss/\_compile/\_test/\_bench \textperiodcentered{} fato \textperiodcentered{} confiança 1.0 \textperiodcentered{} domínio cristal \textperiodcentered{} história 15 versões no jornal}

%CRISTAL {"arestas": [{"alvo": "cristal_repl_shell", "confianca": 1.0, "epistemico": "fato", "origem": "wiki", "rel": "parte_de"}, {"alvo": "cristal_carrega_identidade", "confianca": 1.0, "epistemico": "fato", "origem": "wiki", "rel": "relaciona"}], "confianca": 1.0, "contraexemplos": [], "descricao": "session.py grava/lê turnos em .cristal/sessions/<nome>.jsonl (ou home/memoria/sessions); o REPL loga cada comando/chat e retoma os últimos 3 no banner. Persistência básica funciona; sem replay real de contexto ainda (Fase 0.5).", "epistemico": "fato", "exemplos": [], "id": "cristal_sessao_jsonl", "memoria": "semantica", "meta": {"autor": "broca-so", "dominio": "cristal", "fonte": "code/cristal/shell/session.py:append_entry/load_recent", "wip": true}, "origem": "wiki", "palavras_chave": [], "sinonimos": [], "tipo": "conceito", "titulo": "Cristal — sessão persistente em JSONL por projeto"}
\section{Cristal — sessão persistente em JSONL por projeto}
session.py grava/lê turnos em .cristal/sessions/<nome>.jsonl (ou home/memoria/sessions); o REPL loga cada comando/chat e retoma os últimos 3 no banner. Persistência básica funciona; sem replay real de contexto ainda (Fase 0.5).
\prov{relações: parte\_de cristal\_repl\_shell; relaciona cristal\_carrega\_identidade}
\prov{proveniência: wiki \textperiodcentered{} code/cristal/shell/session.py:append\_entry/load\_recent \textperiodcentered{} fato \textperiodcentered{} confiança 1.0 \textperiodcentered{} domínio cristal \textperiodcentered{} história 15 versões no jornal}

%CRISTAL {"arestas": [{"alvo": "telemetria_shadow", "confianca": 1.0, "epistemico": "fato", "origem": "wiki", "rel": "especializa"}, {"alvo": "colapso_ping_pong_overlap", "confianca": 1.0, "epistemico": "fato", "origem": "wiki", "rel": "relaciona"}], "confianca": 1.0, "contraexemplos": [], "descricao": "Com --shadow/verbose, ask usa conversa.telemetria.responder_com_meta e devolve meta com par φ/δ (par_phi_delta), formatado por telemetria.formato_linha; exposto como detail e logado como entrada 'meta' na sessão. φ/δ = métricas do colapso. Opcional.", "epistemico": "fato", "exemplos": [], "id": "cristal_shadow_phi_delta", "memoria": "semantica", "meta": {"autor": "broca-so", "dominio": "cristal", "fonte": "code/cristal/commands/chat.py (bloco shadow) + tiffany/api.py:ask", "wip": true}, "origem": "wiki", "palavras_chave": [], "sinonimos": [], "tipo": "conceito", "titulo": "Cristal — modo shadow: telemetria φ/δ do colapso"}
\section{Cristal — modo shadow: telemetria \ensuremath{\varphi}/\ensuremath{\delta} do colapso}
Com --shadow/verbose, ask usa conversa.telemetria.responder\_com\_meta e devolve meta com par \ensuremath{\varphi}/\ensuremath{\delta} (par\_phi\_delta), formatado por telemetria.formato\_linha; exposto como detail e logado como entrada 'meta' na sessão. \ensuremath{\varphi}/\ensuremath{\delta} = métricas do colapso. Opcional.
\prov{relações: especializa telemetria\_shadow; relaciona colapso\_ping\_pong\_overlap}
\prov{proveniência: wiki \textperiodcentered{} code/cristal/commands/chat.py (bloco shadow) + tiffany/api.py:ask \textperiodcentered{} fato \textperiodcentered{} confiança 1.0 \textperiodcentered{} domínio cristal \textperiodcentered{} história 15 versões no jornal}

%CRISTAL {"arestas": [{"alvo": "cristal_pipeline_runner", "confianca": 1.0, "epistemico": "fato", "origem": "wiki", "rel": "usa"}, {"alvo": "cristal_agentes_ortogonais", "confianca": 1.0, "epistemico": "fato", "origem": "wiki", "rel": "relaciona"}, {"alvo": "cadeias_dep_grant", "confianca": 1.0, "epistemico": "fato", "origem": "wiki", "rel": "relaciona"}], "confianca": 1.0, "contraexemplos": [], "descricao": "task add/list grava em memória viva; work pega a última tarefa aberta, infere ação (fix/optimize/docs/refactor por palavras-chave) e dispara o pipeline; sem -y mostra só o PLANO (modo dry, o default), com -y roda e marca 'concluida'. Liga cristal_agentes_ortogonais ao ciclo de tarefas.", "epistemico": "fato", "exemplos": [], "id": "cristal_task_work", "memoria": "semantica", "meta": {"autor": "broca-so", "dominio": "cristal", "fonte": "code/cristal/commands/task.py:run_work/_infer_action", "wip": true}, "origem": "wiki", "palavras_chave": [], "sinonimos": [], "tipo": "conceito", "titulo": "Cristal — desenvolvimento orientado a tarefas"}
\section{Cristal — desenvolvimento orientado a tarefas}
task add/list grava em memória viva; work pega a última tarefa aberta, infere ação (fix/optimize/docs/refactor por palavras-chave) e dispara o pipeline; sem -y mostra só o PLANO (modo dry, o default), com -y roda e marca 'concluida'. Liga cristal\_agentes\_ortogonais ao ciclo de tarefas.
\prov{relações: usa cristal\_pipeline\_runner; relaciona cristal\_agentes\_ortogonais; relaciona cadeias\_dep\_grant}
\prov{proveniência: wiki \textperiodcentered{} code/cristal/commands/task.py:run\_work/\_infer\_action \textperiodcentered{} fato \textperiodcentered{} confiança 1.0 \textperiodcentered{} domínio cristal \textperiodcentered{} história 20 versões no jornal}

%CRISTAL {"arestas": [{"alvo": "elasticidade_entropica_borracha", "confianca": 1.0, "epistemico": "fato", "origem": "wiki", "rel": "relaciona"}, {"alvo": "algebra_fundamental_gato_gamba", "confianca": 1.0, "epistemico": "fato", "origem": "wiki", "rel": "usa"}, {"alvo": "corpo_morfico_matrix", "confianca": 1.0, "epistemico": "fato", "origem": "wiki", "rel": "usa"}, {"alvo": "omnitrix_corpo_universal", "confianca": 1.0, "epistemico": "fato", "origem": "wiki", "rel": "usa"}], "confianca": 1.0, "contraexemplos": [], "descricao": "A deformação de cada elemento material É uma MATRIZ — o gradiente de deformação F. A transversal tem duas faces: o CISALHAMENTO (F=[[1,γ,0],[0,1,0],[0,0,1]], inclina de lado; det F=1, ISOCÓRICO, o volume conserva; a tensão σ=G·γ, G=nkT ∝ T, entrópica) e o POISSON (a contração transversal; a borracha INCOMPRESSÍVEL tem ν=½, o transversal=1/√λ, o volume fica). E a DECOMPOSIÇÃO POLAR F=R·U: R é a ROTAÇÃO (ortogonal = o ESQUILO), U a deformação pura (simétrica SPD = o GATO) — a transversal SEPARA em esquilo × gato, a mesma dupla do Corpo Algébrico (o tensor E=(FᵀF−I)/2 simétrico=o gato; a parte antissimétrica=o esquilo). linguagem/omnitrix.py (deformacao_transversal 5/5); glsl_ptx_frag (FRAG_CILINDRO_CISALHA — a cápsula geral segmento⊕bola inclinando o eixo por γ(iTime), raio analítico, 11/11).", "epistemico": "fato", "exemplos": [], "id": "deformacao_transversal_cisalha", "memoria": "semantica", "meta": {"dominio": "floxina", "fonte": "Floxina (investigação)"}, "origem": "wiki", "palavras_chave": [], "sinonimos": [], "tipo": "conceito", "titulo": "A deformação transversal: o cisalhamento (isocórico) + o Poisson (ν=½); o gradiente F=R·U=esquilo·gato"}
\section{A deformação transversal: o cisalhamento (isocórico) + o Poisson (\ensuremath{\nu}=\ensuremath{\tfrac12}); o gradiente F=R\textperiodcentered U=esquilo\textperiodcentered gato}
A deformação de cada elemento material É uma MATRIZ — o gradiente de deformação F. A transversal tem duas faces: o CISALHAMENTO (F=[[1,\ensuremath{\gamma},0],[0,1,0],[0,0,1]], inclina de lado; det F=1, ISOCÓRICO, o volume conserva; a tensão \ensuremath{\sigma}=G\textperiodcentered \ensuremath{\gamma}, G=nkT \ensuremath{\propto} T, entrópica) e o POISSON (a contração transversal; a borracha INCOMPRESSÍVEL tem \ensuremath{\nu}=\ensuremath{\tfrac12}, o transversal=1/\ensuremath{\sqrt{\,}}\ensuremath{\lambda}, o volume fica). E a DECOMPOSIÇÃO POLAR F=R\textperiodcentered U: R é a ROTAÇÃO (ortogonal = o ESQUILO), U a deformação pura (simétrica SPD = o GATO) — a transversal SEPARA em esquilo $\times$ gato, a mesma dupla do Corpo Algébrico (o tensor E=(F\ensuremath{^{T}}F\ensuremath{-}I)/2 simétrico=o gato; a parte antissimétrica=o esquilo). linguagem/omnitrix.py (deformacao\_transversal 5/5); glsl\_ptx\_frag (FRAG\_CILINDRO\_CISALHA — a cápsula geral segmento\ensuremath{\oplus}bola inclinando o eixo por \ensuremath{\gamma}(iTime), raio analítico, 11/11).
\prov{relações: relaciona elasticidade\_entropica\_borracha; usa algebra\_fundamental\_gato\_gamba; usa corpo\_morfico\_matrix; usa omnitrix\_corpo\_universal}
\prov{proveniência: wiki \textperiodcentered{} Floxina (investigação) \textperiodcentered{} fato \textperiodcentered{} confiança 1.0 \textperiodcentered{} domínio floxina \textperiodcentered{} história 110 versões no jornal}

%CRISTAL {"arestas": [{"alvo": "omnitrix_corpo_universal", "confianca": 1.0, "epistemico": "fato", "origem": "wiki", "rel": "parte_de"}, {"alvo": "dicionario_matrix_bai", "confianca": 1.0, "epistemico": "fato", "origem": "wiki", "rel": "relaciona"}, {"alvo": "bai", "confianca": 1.0, "epistemico": "fato", "origem": "wiki", "rel": "relaciona"}, {"alvo": "gato_e_o_hamiltoniano", "confianca": 1.0, "epistemico": "fato", "origem": "wiki", "rel": "usa"}], "confianca": 1.0, "contraexemplos": [], "descricao": "O BAI (Biblioteca de Autorização Isomórfica, o quinteto κ,δ,Φ,T,Ψ, autorização como álgebra sobre RETICULADOS) é o Omnitrix instanciado no substrato da AUTORIZAÇÃO — como a molécula é na química. O dicionário (generaliza o matrix↔BAI ao nível universal): a parábola geradora x²−x (o gerador) ↔ κ (o núcleo); o gap √(m²+4)=[gato,gambá] ↔ δ (a escada CASL); o gambá (Skew, U(t), a fase U(1)) ↔ Φ (o fluxo/o potencial); o gato→gambá La Hire (discar a forma) ↔ T (a transformação); os n-qubits+a+c²=1+o idempotente ↔ Ψ (o estado/a decisão colapsada, |Ψ|²=1); a instância (ordem m, substrato) ↔ o reticulado; a matéria/antimatéria (±) ↔ autorizar/negar; a aniquilação (−1) ↔ a revogação ao nulo; o zero da reta (±1 simétrico) ↔ a autorização neutra (a identidade). O quinteto κ,δ,Φ,T,Ψ É o template Sym⊕Skew⊕invariante. linguagem/omnitrix.py (dicionario_bai_omnitrix, 3/3).", "epistemico": "fato", "exemplos": [], "id": "dicionario_bai_omnitrix", "memoria": "semantica", "meta": {"dominio": "floxina", "fonte": "Floxina (investigação)"}, "origem": "wiki", "palavras_chave": [], "sinonimos": [], "tipo": "conceito", "titulo": "Dicionário: o BAI ↔ o Omnitrix (a autorização é uma instância do Corpo Algébrico)"}
\section{Dicionário: o BAI \ensuremath{\leftrightarrow} o Omnitrix (a autorização é uma instância do Corpo Algébrico)}
O BAI (Biblioteca de Autorização Isomórfica, o quinteto \ensuremath{\kappa},\ensuremath{\delta},\ensuremath{\Phi},T,\ensuremath{\Psi}, autorização como álgebra sobre RETICULADOS) é o Omnitrix instanciado no substrato da AUTORIZAÇÃO — como a molécula é na química. O dicionário (generaliza o matrix\ensuremath{\leftrightarrow}BAI ao nível universal): a parábola geradora x\ensuremath{^{2}}\ensuremath{-}x (o gerador) \ensuremath{\leftrightarrow} \ensuremath{\kappa} (o núcleo); o gap \ensuremath{\sqrt{\,}}(m\ensuremath{^{2}}+4)=[gato,gambá] \ensuremath{\leftrightarrow} \ensuremath{\delta} (a escada CASL); o gambá (Skew, U(t), a fase U(1)) \ensuremath{\leftrightarrow} \ensuremath{\Phi} (o fluxo/o potencial); o gato\ensuremath{\to}gambá La Hire (discar a forma) \ensuremath{\leftrightarrow} T (a transformação); os n-qubits+a+c\ensuremath{^{2}}=1+o idempotente \ensuremath{\leftrightarrow} \ensuremath{\Psi} (o estado/a decisão colapsada, |\ensuremath{\Psi}|\ensuremath{^{2}}=1); a instância (ordem m, substrato) \ensuremath{\leftrightarrow} o reticulado; a matéria/antimatéria ($\pm$) \ensuremath{\leftrightarrow} autorizar/negar; a aniquilação (\ensuremath{-}1) \ensuremath{\leftrightarrow} a revogação ao nulo; o zero da reta ($\pm$1 simétrico) \ensuremath{\leftrightarrow} a autorização neutra (a identidade). O quinteto \ensuremath{\kappa},\ensuremath{\delta},\ensuremath{\Phi},T,\ensuremath{\Psi} É o template Sym\ensuremath{\oplus}Skew\ensuremath{\oplus}invariante. linguagem/omnitrix.py (dicionario\_bai\_omnitrix, 3/3).
\prov{relações: parte\_de omnitrix\_corpo\_universal; relaciona dicionario\_matrix\_bai; relaciona bai; usa gato\_e\_o\_hamiltoniano}
\prov{proveniência: wiki \textperiodcentered{} Floxina (investigação) \textperiodcentered{} fato \textperiodcentered{} confiança 1.0 \textperiodcentered{} domínio floxina \textperiodcentered{} história 275 versões no jornal}

%CRISTAL {"arestas": [{"alvo": "coracao", "confianca": 1.0, "epistemico": "fato", "origem": "wiki", "rel": "parte_de"}, {"alvo": "gap_espectral", "confianca": 1.0, "epistemico": "fato", "origem": "wiki", "rel": "relaciona"}, {"alvo": "mass_gap_algebrico", "confianca": 1.0, "epistemico": "fato", "origem": "wiki", "rel": "relaciona"}, {"alvo": "nao_associatividade_amputacao", "confianca": 1.0, "epistemico": "fato", "origem": "wiki", "rel": "relaciona"}, {"alvo": "g2_grupo_excepcional", "confianca": 1.0, "epistemico": "fato", "origem": "wiki", "rel": "relaciona"}, {"alvo": "higgs_algebrico_g2_su3", "confianca": 1.0, "epistemico": "fato", "origem": "wiki", "rel": "relaciona"}], "confianca": 1.0, "contraexemplos": [], "descricao": "A MESMA quantidade, quatro nomes: (1) o GAP EMOCIONAL a=1−c² (o que o colapso deixa por associar); (2) o GAP ESPECTRAL Δ=λ₁−λ₀ do operador (a distância do fundamental ao 1º excitado) — que governa o TEMPO DE RELAXAÇÃO τ~1/Δ (a homeostase: gap grande ⟹ volta rápido ao sereno); (3) o MASS GAP de YANG-MILLS Δ>0 (o problema do milênio: a excitação mais leve tem massa positiva) = aqui o mass gap algébrico dim 𝔤₂=Der(𝕆)=14, o associador amputado, 'o resto que nunca zera'; (4) a MASSA m=Δ (E=Δ). Um dicionário só: emoção ↔ espectro ↔ campo ↔ massa.", "epistemico": "fato", "exemplos": [], "id": "dicionario_coracao_gap", "memoria": "semantica", "meta": {"dominio": "coracao", "fonte": "coracao hub"}, "origem": "wiki", "palavras_chave": [], "sinonimos": [], "tipo": "conceito", "titulo": "Dicionário do coração: o gap emocional ↔ o gap espectral ↔ Yang-Mills"}
\section{Dicionário do coração: o gap emocional \ensuremath{\leftrightarrow} o gap espectral \ensuremath{\leftrightarrow} Yang-Mills}
A MESMA quantidade, quatro nomes: (1) o GAP EMOCIONAL a=1\ensuremath{-}c\ensuremath{^{2}} (o que o colapso deixa por associar); (2) o GAP ESPECTRAL \ensuremath{\Delta}=\ensuremath{\lambda}\ensuremath{_{1}}\ensuremath{-}\ensuremath{\lambda}\ensuremath{_{0}} do operador (a distância do fundamental ao 1º excitado) — que governa o TEMPO DE RELAXAÇÃO \ensuremath{\tau}\~{}1/\ensuremath{\Delta} (a homeostase: gap grande \ensuremath{\Longrightarrow} volta rápido ao sereno); (3) o MASS GAP de YANG-MILLS \ensuremath{\Delta}>0 (o problema do milênio: a excitação mais leve tem massa positiva) = aqui o mass gap algébrico dim \ensuremath{\mathfrak{g}}\ensuremath{_{2}}=Der(\ensuremath{\mathbb{O}})=14, o associador amputado, 'o resto que nunca zera'; (4) a MASSA m=\ensuremath{\Delta} (E=\ensuremath{\Delta}). Um dicionário só: emoção \ensuremath{\leftrightarrow} espectro \ensuremath{\leftrightarrow} campo \ensuremath{\leftrightarrow} massa.
\prov{relações: parte\_de coracao; relaciona gap\_espectral; relaciona mass\_gap\_algebrico; relaciona nao\_associatividade\_amputacao; relaciona g2\_grupo\_excepcional; relaciona higgs\_algebrico\_g2\_su3}
\prov{proveniência: wiki \textperiodcentered{} coracao hub \textperiodcentered{} fato \textperiodcentered{} confiança 1.0 \textperiodcentered{} domínio coracao \textperiodcentered{} história 21 versões no jornal}

%CRISTAL {"arestas": [{"alvo": "dicionario_vesica_motor_esquilo", "confianca": 1.0, "epistemico": "fato", "origem": "wiki", "rel": "relaciona"}, {"alvo": "torre_dois_lados_classico_relativistico", "confianca": 1.0, "epistemico": "fato", "origem": "wiki", "rel": "usa"}, {"alvo": "torre_matricial_coordenada", "confianca": 1.0, "epistemico": "fato", "origem": "wiki", "rel": "usa"}, {"alvo": "maquina_gaiola_esquilo_dtc", "confianca": 1.0, "epistemico": "fato", "origem": "wiki", "rel": "usa"}, {"alvo": "matrix_vesicular_corpo", "confianca": 1.0, "epistemico": "fato", "origem": "wiki", "rel": "usa"}, {"alvo": "gerador_do_gap_comutador", "confianca": 1.0, "epistemico": "fato", "origem": "wiki", "rel": "usa"}], "confianca": 1.0, "contraexemplos": [], "descricao": "A DTC VETORIAL sobe do plano (a vesica 2D) pela torre RELATIVÍSTICA: o torque É o PRODUTO VETORIAL (o comutador = o esquilo = a costura anisotrópica), e o produto vetorial existe SÓ em 3D e 7D — logo há exatamente o motor ESFÉRICO e o HIPERESFÉRICO de torque-vetor pleno. 2D (planar): a DTC do hexágono (6+2 vetores) ↔ a matrix vesicular (gato/esquilo) ↔ ℂ; o torque é ESCALAR (z de a×b). 3D (esférico): o MOTOR ESFÉRICO (o rotor gira em qualquer eixo, SO(3)) ↔ Im ℍ (o quatérnio); o torque é um VETOR (a×b, o eixo); |a×b|²=|a|²|b|²−(a·b)². 7D (hiperesférico): o MOTOR HIPERESFÉRICO ↔ Im 𝕆 (o octônio); o torque é vetor 7D. O torque a×b é ANTISSIMÉTRICO (a×b=−b×a)=o esquilo em toda dim onde existe (só 3 e 7). A ESFERA DE TESLA (o Ovo de Colombo, 1893: Tesla girou uma esfera de cobre pelo CAMPO GIRANTE) é o motor esférico ancestral — o campo girante (o esquilo e^{Gt}) que Tesla inventou, com a gaiola de esquilo. linguagem/vesica_corpo.py (dicionario_dtc_esferico, 5/5). A vesica é o caso 2D; o esférico o 3D; o hiperesférico o 7D.", "epistemico": "fato", "exemplos": [], "id": "dicionario_dtc_motor_esferico", "memoria": "semantica", "meta": {"dominio": "floxina", "fonte": "Floxina (investigação)"}, "origem": "wiki", "palavras_chave": [], "sinonimos": [], "tipo": "conceito", "titulo": "Dicionário: a DTC vetorial → o motor esférico (3D) e hiperesférico (7D); a esfera de Tesla"}
\section{Dicionário: a DTC vetorial \ensuremath{\to} o motor esférico (3D) e hiperesférico (7D); a esfera de Tesla}
A DTC VETORIAL sobe do plano (a vesica 2D) pela torre RELATIVÍSTICA: o torque É o PRODUTO VETORIAL (o comutador = o esquilo = a costura anisotrópica), e o produto vetorial existe SÓ em 3D e 7D — logo há exatamente o motor ESFÉRICO e o HIPERESFÉRICO de torque-vetor pleno. 2D (planar): a DTC do hexágono (6+2 vetores) \ensuremath{\leftrightarrow} a matrix vesicular (gato/esquilo) \ensuremath{\leftrightarrow} \ensuremath{\mathbb{C}}; o torque é ESCALAR (z de a$\times$b). 3D (esférico): o MOTOR ESFÉRICO (o rotor gira em qualquer eixo, SO(3)) \ensuremath{\leftrightarrow} Im \ensuremath{\mathbb{H}} (o quatérnio); o torque é um VETOR (a$\times$b, o eixo); |a$\times$b|\ensuremath{^{2}}=|a|\ensuremath{^{2}}|b|\ensuremath{^{2}}\ensuremath{-}(a\textperiodcentered b)\ensuremath{^{2}}. 7D (hiperesférico): o MOTOR HIPERESFÉRICO \ensuremath{\leftrightarrow} Im \ensuremath{\mathbb{O}} (o octônio); o torque é vetor 7D. O torque a$\times$b é ANTISSIMÉTRICO (a$\times$b=\ensuremath{-}b$\times$a)=o esquilo em toda dim onde existe (só 3 e 7). A ESFERA DE TESLA (o Ovo de Colombo, 1893: Tesla girou uma esfera de cobre pelo CAMPO GIRANTE) é o motor esférico ancestral — o campo girante (o esquilo e\^{}\{Gt\}) que Tesla inventou, com a gaiola de esquilo. linguagem/vesica\_corpo.py (dicionario\_dtc\_esferico, 5/5). A vesica é o caso 2D; o esférico o 3D; o hiperesférico o 7D.
\prov{relações: relaciona dicionario\_vesica\_motor\_esquilo; usa torre\_dois\_lados\_classico\_relativistico; usa torre\_matricial\_coordenada; usa maquina\_gaiola\_esquilo\_dtc; usa matrix\_vesicular\_corpo; usa gerador\_do\_gap\_comutador}
\prov{proveniência: wiki \textperiodcentered{} Floxina (investigação) \textperiodcentered{} fato \textperiodcentered{} confiança 1.0 \textperiodcentered{} domínio floxina \textperiodcentered{} história 217 versões no jornal}

%CRISTAL {"arestas": [{"alvo": "corpo_equacoes_diferenciais", "confianca": 1.0, "epistemico": "fato", "origem": "wiki", "rel": "usa"}, {"alvo": "parabola_geradora_idempotente", "confianca": 1.0, "epistemico": "fato", "origem": "wiki", "rel": "usa"}, {"alvo": "dicionario_bai_omnitrix", "confianca": 1.0, "epistemico": "fato", "origem": "wiki", "rel": "relaciona"}], "confianca": 1.0, "contraexemplos": [], "descricao": "A máquina mineral (o BAI) lê a equação diferencial: o espectro do gerador A cai na parábola cristal/borda/caos. O GATO ↔ o GERADOR A (a matriz da ED). A PARÁBOLA: CRISTAL ↔ o ponto fixo ESTÁVEL x*=x (Re λ<0, a condição inicial COLAPSA no equilíbrio); BORDA ↔ o marginal/a bifurcação (Re λ=0, |x| constante, o esquilo puro orbita); CAOS ↔ a instabilidade/o atrator estranho (Re λ>0, Lyapunov>0, diverge). A ÓRBITA ↔ o fluxo φ_t=e^{At} (o semigrupo); o COLAPSO ↔ a convergência ao ponto fixo. A ESTRELA ↔ o modo próprio (o autovetor, A·v=λ·v); o METAL ↔ o autovalor λ (σ_n); a RETA ↔ Re λ (a taxa, log σ); a BANDA ↔ o espectro (ω=Im λ). Um operador, só o dicionário muda. linguagem/omnitrix.py (dicionario_ed_bai 5/5).", "epistemico": "fato", "exemplos": [], "id": "dicionario_ed_bai", "memoria": "semantica", "meta": {"dominio": "floxina", "fonte": "Floxina (investigação)"}, "origem": "wiki", "palavras_chave": [], "sinonimos": [], "tipo": "conceito", "titulo": "Dicionário: equações diferenciais ↔ BAI (a máquina mineral lê o fluxo dx/dt=A·x)"}
\section{Dicionário: equações diferenciais \ensuremath{\leftrightarrow} BAI (a máquina mineral lê o fluxo dx/dt=A\textperiodcentered x)}
A máquina mineral (o BAI) lê a equação diferencial: o espectro do gerador A cai na parábola cristal/borda/caos. O GATO \ensuremath{\leftrightarrow} o GERADOR A (a matriz da ED). A PARÁBOLA: CRISTAL \ensuremath{\leftrightarrow} o ponto fixo ESTÁVEL x*=x (Re \ensuremath{\lambda}<0, a condição inicial COLAPSA no equilíbrio); BORDA \ensuremath{\leftrightarrow} o marginal/a bifurcação (Re \ensuremath{\lambda}=0, |x| constante, o esquilo puro orbita); CAOS \ensuremath{\leftrightarrow} a instabilidade/o atrator estranho (Re \ensuremath{\lambda}>0, Lyapunov>0, diverge). A ÓRBITA \ensuremath{\leftrightarrow} o fluxo \ensuremath{\varphi}\_t=e\^{}\{At\} (o semigrupo); o COLAPSO \ensuremath{\leftrightarrow} a convergência ao ponto fixo. A ESTRELA \ensuremath{\leftrightarrow} o modo próprio (o autovetor, A\textperiodcentered v=\ensuremath{\lambda}\textperiodcentered v); o METAL \ensuremath{\leftrightarrow} o autovalor \ensuremath{\lambda} (\ensuremath{\sigma}\_n); a RETA \ensuremath{\leftrightarrow} Re \ensuremath{\lambda} (a taxa, log \ensuremath{\sigma}); a BANDA \ensuremath{\leftrightarrow} o espectro (\ensuremath{\omega}=Im \ensuremath{\lambda}). Um operador, só o dicionário muda. linguagem/omnitrix.py (dicionario\_ed\_bai 5/5).
\prov{relações: usa corpo\_equacoes\_diferenciais; usa parabola\_geradora\_idempotente; relaciona dicionario\_bai\_omnitrix}
\prov{proveniência: wiki \textperiodcentered{} Floxina (investigação) \textperiodcentered{} fato \textperiodcentered{} confiança 1.0 \textperiodcentered{} domínio floxina \textperiodcentered{} história 76 versões no jornal}

%CRISTAL {"arestas": [{"alvo": "gamba_antissimetrico", "confianca": 1.0, "epistemico": "fato", "origem": "wiki", "rel": "usa"}, {"alvo": "materia_escura", "confianca": 1.0, "epistemico": "fato", "origem": "wiki", "rel": "relaciona"}, {"alvo": "curvas_rotacao_rubin", "confianca": 1.0, "epistemico": "fato", "origem": "wiki", "rel": "relaciona"}, {"alvo": "bai_limite_nao_dissipativo", "confianca": 1.0, "epistemico": "fato", "origem": "wiki", "rel": "relaciona"}, {"alvo": "floxina_wick_difusao", "confianca": 1.0, "epistemico": "fato", "origem": "wiki", "rel": "relaciona"}], "confianca": 1.0, "contraexemplos": [], "descricao": "«Existe mas não irradia; só o efeito estrutural a denuncia» é a definição da MATÉRIA ESCURA (interage por gravidade mas não emite/absorve/reflete luz; Rubin: as curvas de rotação planas exigem um halo invisível). E a propriedade que a DEFINE — ser não-dissipativa/não-colisional (não irradia, não esfria, não colapsa num disco) — é a do GAMBÁ (anti-simétrico=reversível=não-dissipativo, Onsager; o limite δ=∞ de Sitter conservativo). O bárion dissipa e cai no disco (gato, a reta); a escura conserva e fica no halo (gambá, o círculo). Dicionário: gambá↔não-colisional; gato↔bárion(irradia); sem aboutness↔não-EM; efeito estrutural↔gravidade(Rubin); o gap(grau 0)↔a massa que falta; o halo↔não-colapsa; o resto que nunca zera (𝔤₂=14)↔os ~27% invisíveis. ESCOPO: tentativa/analogia — o único ancoradouro DURO é a coincidência de simetria (o gambá não-dissipativo por álgebra; a escura não-dissipativa por observação).", "epistemico": "hipotese", "exemplos": [], "id": "dicionario_floxina_materia_escura", "memoria": "semantica", "meta": {"dominio": "floxina", "fonte": "Floxina (investigação)"}, "origem": "wiki", "palavras_chave": [], "sinonimos": [], "tipo": "conceito", "titulo": "Dicionário (tentativa): floxina ↔ matéria escura"}
\section{Dicionário (tentativa): floxina \ensuremath{\leftrightarrow} matéria escura}
«Existe mas não irradia; só o efeito estrutural a denuncia» é a definição da MATÉRIA ESCURA (interage por gravidade mas não emite/absorve/reflete luz; Rubin: as curvas de rotação planas exigem um halo invisível). E a propriedade que a DEFINE — ser não-dissipativa/não-colisional (não irradia, não esfria, não colapsa num disco) — é a do GAMBÁ (anti-simétrico=reversível=não-dissipativo, Onsager; o limite \ensuremath{\delta}=\ensuremath{\infty} de Sitter conservativo). O bárion dissipa e cai no disco (gato, a reta); a escura conserva e fica no halo (gambá, o círculo). Dicionário: gambá\ensuremath{\leftrightarrow}não-colisional; gato\ensuremath{\leftrightarrow}bárion(irradia); sem aboutness\ensuremath{\leftrightarrow}não-EM; efeito estrutural\ensuremath{\leftrightarrow}gravidade(Rubin); o gap(grau 0)\ensuremath{\leftrightarrow}a massa que falta; o halo\ensuremath{\leftrightarrow}não-colapsa; o resto que nunca zera (\ensuremath{\mathfrak{g}}\ensuremath{_{2}}=14)\ensuremath{\leftrightarrow}os \~{}27\% invisíveis. ESCOPO: tentativa/analogia — o único ancoradouro DURO é a coincidência de simetria (o gambá não-dissipativo por álgebra; a escura não-dissipativa por observação).
\prov{relações: usa gamba\_antissimetrico; relaciona materia\_escura; relaciona curvas\_rotacao\_rubin; relaciona bai\_limite\_nao\_dissipativo; relaciona floxina\_wick\_difusao}
\prov{proveniência: wiki \textperiodcentered{} Floxina (investigação) \textperiodcentered{} hipotese \textperiodcentered{} confiança 1.0 \textperiodcentered{} domínio floxina \textperiodcentered{} história 486 versões no jornal}

%CRISTAL {"arestas": [{"alvo": "gamba_antissimetrico", "confianca": 1.0, "epistemico": "fato", "origem": "wiki", "rel": "usa"}, {"alvo": "teste_floxinato_gamba", "confianca": 1.0, "epistemico": "fato", "origem": "wiki", "rel": "relaciona"}, {"alvo": "floxina_selberg", "confianca": 1.0, "epistemico": "fato", "origem": "wiki", "rel": "relaciona"}, {"alvo": "mass_gap_algebrico", "confianca": 1.0, "epistemico": "fato", "origem": "wiki", "rel": "relaciona"}, {"alvo": "nao_associatividade_amputacao", "confianca": 1.0, "epistemico": "fato", "origem": "wiki", "rel": "relaciona"}, {"alvo": "gap_espectral", "confianca": 1.0, "epistemico": "fato", "origem": "wiki", "rel": "relaciona"}, {"alvo": "g2_grupo_excepcional", "confianca": 1.0, "epistemico": "fato", "origem": "wiki", "rel": "relaciona"}], "confianca": 1.0, "contraexemplos": [], "descricao": "O fecho do arco mais improvável. O «floxinato de gambá» — termo FALSO inventado só para testar o fail-closed da litografia — nomeou, sem saber, o GAMBÁ: o operador anti-simétrico G=[[0,1],[−1,0]], G²=−I, a FASE elementar (a desafasagem em dim 2). Levada à torre de Cayley-Dickson, a fase que ASSOCIA em ℂ deixa de associar em 𝕆, e o resíduo é o GAP DE MASSA 𝔤₂=Der(𝕆)=14. Dicionário: floxinato de gambá → o gambá (G²=−I, a fase) → a desafasagem (o associador) → o gap de massa. O GAP É A DESAFASAGEM: λ=¼−δ² (δ=s−½); Riemann o vê em 0 (Re(δ)=0, os zeros na linha), Selberg em ¼ (λ₁≥¼), Yang-Mills como Δ>0. Verificado (selberg_gap.py): associador 0 até ℍ (dim ≤4, o «0» de Riemann), abre no 𝕆 (dim 8, 𝔤₂=14); o gap de Cayley positivo (expansor) p=3..13; sem Maass excepcional (λ₁≈91,14). A piada era o gap.", "epistemico": "hipotese", "exemplos": [], "id": "dicionario_floxinato_gap_massa", "memoria": "semantica", "meta": {"dominio": "floxina", "fonte": "Floxina (investigação)"}, "origem": "wiki", "palavras_chave": [], "sinonimos": [], "tipo": "conceito", "titulo": "Dicionário: floxinato de gambá → o gap de massa"}
\section{Dicionário: floxinato de gambá \ensuremath{\to} o gap de massa}
O fecho do arco mais improvável. O «floxinato de gambá» — termo FALSO inventado só para testar o fail-closed da litografia — nomeou, sem saber, o GAMBÁ: o operador anti-simétrico G=[[0,1],[\ensuremath{-}1,0]], G\ensuremath{^{2}}=\ensuremath{-}I, a FASE elementar (a desafasagem em dim 2). Levada à torre de Cayley-Dickson, a fase que ASSOCIA em \ensuremath{\mathbb{C}} deixa de associar em \ensuremath{\mathbb{O}}, e o resíduo é o GAP DE MASSA \ensuremath{\mathfrak{g}}\ensuremath{_{2}}=Der(\ensuremath{\mathbb{O}})=14. Dicionário: floxinato de gambá \ensuremath{\to} o gambá (G\ensuremath{^{2}}=\ensuremath{-}I, a fase) \ensuremath{\to} a desafasagem (o associador) \ensuremath{\to} o gap de massa. O GAP É A DESAFASAGEM: \ensuremath{\lambda}=\ensuremath{\tfrac14}\ensuremath{-}\ensuremath{\delta}\ensuremath{^{2}} (\ensuremath{\delta}=s\ensuremath{-}\ensuremath{\tfrac12}); Riemann o vê em 0 (Re(\ensuremath{\delta})=0, os zeros na linha), Selberg em \ensuremath{\tfrac14} (\ensuremath{\lambda}\ensuremath{_{1}}\ensuremath{\ge}\ensuremath{\tfrac14}), Yang-Mills como \ensuremath{\Delta}>0. Verificado (selberg\_gap.py): associador 0 até \ensuremath{\mathbb{H}} (dim \ensuremath{\le}4, o «0» de Riemann), abre no \ensuremath{\mathbb{O}} (dim 8, \ensuremath{\mathfrak{g}}\ensuremath{_{2}}=14); o gap de Cayley positivo (expansor) p=3..13; sem Maass excepcional (\ensuremath{\lambda}\ensuremath{_{1}}\ensuremath{\approx}91,14). A piada era o gap.
\prov{relações: usa gamba\_antissimetrico; relaciona teste\_floxinato\_gamba; relaciona floxina\_selberg; relaciona mass\_gap\_algebrico; relaciona nao\_associatividade\_amputacao; relaciona gap\_espectral; relaciona g2\_grupo\_excepcional}
\prov{proveniência: wiki \textperiodcentered{} Floxina (investigação) \textperiodcentered{} hipotese \textperiodcentered{} confiança 1.0 \textperiodcentered{} domínio floxina \textperiodcentered{} história 640 versões no jornal}

%CRISTAL {"arestas": [{"alvo": "estojo_agulhas_minerais_fixo", "confianca": 1.0, "epistemico": "fato", "origem": "wiki", "rel": "usa"}, {"alvo": "agulha_perelman_limite", "confianca": 1.0, "epistemico": "fato", "origem": "wiki", "rel": "usa"}, {"alvo": "litografia", "confianca": 1.0, "epistemico": "fato", "origem": "wiki", "rel": "relaciona"}, {"alvo": "mass_gap_algebrico", "confianca": 1.0, "epistemico": "fato", "origem": "wiki", "rel": "relaciona"}], "confianca": 1.0, "contraexemplos": [], "descricao": "Na fabricação de chips, LITOGRAFIA e DOPAGEM são pareadas: a fotolitografia define ONDE, a implantação/difusão dopa O QUÊ. E o dicionário cai exato: LITOGRAFAR (gravar a assinatura) ↔ DOPAR (introduzir o dopante); a AGULHA de um metal noutro (a gravura cruzada σ_a em σ_b) ↔ o DOPANTE (o átomo σ_a na rede hospedeira σ_b); a base σ_n (a reta da peça) ↔ o CRISTAL HOSPEDEIRO (o silício); o GRAU topológico (0/1) ↔ o TIPO de portador (n doador / p aceitador — σ_dopante≷σ_hospedeiro); o LIMITE K_n (furar ≤ K) ↔ a dopagem RASA (shallow, níveis perto das bordas = portador limpo); FURAR MAIS que K_n (a geometria, a cicatriz) ↔ o NÍVEL PROFUNDO (deep level, trap/recombinação, no meio do gap); o ANEL DE EXCLUSÃO a_min (o mass gap) ↔ o BAND GAP do semicondutor (o dopante põe níveis DENTRO do gap sem o fechar, como a agulha nunca toca o núcleo). conversa/agulha_perelman.py (dopar; 10/10).", "epistemico": "hipotese", "exemplos": [], "id": "dicionario_litografia_dopagem", "memoria": "semantica", "meta": {"dominio": "floxina", "fonte": "Floxina (investigação)"}, "origem": "wiki", "palavras_chave": [], "sinonimos": [], "tipo": "conceito", "titulo": "Dicionário: litografia ↔ dopagem (o semicondutor)"}
\section{Dicionário: litografia \ensuremath{\leftrightarrow} dopagem (o semicondutor)}
Na fabricação de chips, LITOGRAFIA e DOPAGEM são pareadas: a fotolitografia define ONDE, a implantação/difusão dopa O QUÊ. E o dicionário cai exato: LITOGRAFAR (gravar a assinatura) \ensuremath{\leftrightarrow} DOPAR (introduzir o dopante); a AGULHA de um metal noutro (a gravura cruzada \ensuremath{\sigma}\_a em \ensuremath{\sigma}\_b) \ensuremath{\leftrightarrow} o DOPANTE (o átomo \ensuremath{\sigma}\_a na rede hospedeira \ensuremath{\sigma}\_b); a base \ensuremath{\sigma}\_n (a reta da peça) \ensuremath{\leftrightarrow} o CRISTAL HOSPEDEIRO (o silício); o GRAU topológico (0/1) \ensuremath{\leftrightarrow} o TIPO de portador (n doador / p aceitador — \ensuremath{\sigma}\_dopante\ensuremath{\gtrless}\ensuremath{\sigma}\_hospedeiro); o LIMITE K\_n (furar \ensuremath{\le} K) \ensuremath{\leftrightarrow} a dopagem RASA (shallow, níveis perto das bordas = portador limpo); FURAR MAIS que K\_n (a geometria, a cicatriz) \ensuremath{\leftrightarrow} o NÍVEL PROFUNDO (deep level, trap/recombinação, no meio do gap); o ANEL DE EXCLUSÃO a\_min (o mass gap) \ensuremath{\leftrightarrow} o BAND GAP do semicondutor (o dopante põe níveis DENTRO do gap sem o fechar, como a agulha nunca toca o núcleo). conversa/agulha\_perelman.py (dopar; 10/10).
\prov{relações: usa estojo\_agulhas\_minerais\_fixo; usa agulha\_perelman\_limite; relaciona litografia; relaciona mass\_gap\_algebrico}
\prov{proveniência: wiki \textperiodcentered{} Floxina (investigação) \textperiodcentered{} hipotese \textperiodcentered{} confiança 1.0 \textperiodcentered{} domínio floxina \textperiodcentered{} história 380 versões no jornal}

%CRISTAL {"arestas": [{"alvo": "corpo_matricial_matrix", "confianca": 1.0, "epistemico": "fato", "origem": "wiki", "rel": "relaciona"}, {"alvo": "gato", "confianca": 1.0, "epistemico": "fato", "origem": "wiki", "rel": "usa"}, {"alvo": "gamba_antissimetrico", "confianca": 1.0, "epistemico": "fato", "origem": "wiki", "rel": "usa"}, {"alvo": "bai", "confianca": 1.0, "epistemico": "fato", "origem": "wiki", "rel": "relaciona"}, {"alvo": "gato_e_o_hamiltoniano", "confianca": 1.0, "epistemico": "fato", "origem": "wiki", "rel": "relaciona"}, {"alvo": "invariante_estelar", "confianca": 1.0, "epistemico": "fato", "origem": "wiki", "rel": "relaciona"}], "confianca": 1.0, "contraexemplos": [], "descricao": "O corpo matricial É o BAI (Biblioteca de Autorização Isomórfica, o motor de colapso). A correspondência: o GATO (simétrico=HERMITIANO) ↔ o HAMILTONIANO H (o observável; os metais σ_n = os NÍVEIS de energia, a reta=o espectro); o GAMBÁ (anti-simétrico=unitário) ↔ a EVOLUÇÃO U(t)=e^{−iHt} (o transporte, o Ψ que gira); o a+c²=1 (a norma N=det) ↔ o INVARIANTE do BAI (o checksum); o IDEMPOTENTE (A+G)²=A+G ↔ o COLAPSO Ψ (a medição, a projeção a um poço); o GATO dissipa (Onsander) ↔ o δ (a escada CASL, a dissipação/de Sitter); os n-QUBITS (2^n) ↔ o ESTADO do BAI; a ÓRBITA (as potências) ↔ a órbita/a dinâmica do BAI. O corpo matricial é a álgebra do BAI: o gato mede (o Hamiltoniano), o gambá gira (a evolução), o colapso é o idempotente, o invariante é a+c²=1.", "epistemico": "fato", "exemplos": [], "id": "dicionario_matrix_bai", "memoria": "semantica", "meta": {"dominio": "floxina", "fonte": "Floxina (investigação)"}, "origem": "wiki", "palavras_chave": [], "sinonimos": [], "tipo": "conceito", "titulo": "Dicionário: o corpo matricial (gato/gambá) ↔ o BAI (κ,δ,Φ,T,Ψ)"}
\section{Dicionário: o corpo matricial (gato/gambá) \ensuremath{\leftrightarrow} o BAI (\ensuremath{\kappa},\ensuremath{\delta},\ensuremath{\Phi},T,\ensuremath{\Psi})}
O corpo matricial É o BAI (Biblioteca de Autorização Isomórfica, o motor de colapso). A correspondência: o GATO (simétrico=HERMITIANO) \ensuremath{\leftrightarrow} o HAMILTONIANO H (o observável; os metais \ensuremath{\sigma}\_n = os NÍVEIS de energia, a reta=o espectro); o GAMBÁ (anti-simétrico=unitário) \ensuremath{\leftrightarrow} a EVOLUÇÃO U(t)=e\^{}\{\ensuremath{-}iHt\} (o transporte, o \ensuremath{\Psi} que gira); o a+c\ensuremath{^{2}}=1 (a norma N=det) \ensuremath{\leftrightarrow} o INVARIANTE do BAI (o checksum); o IDEMPOTENTE (A+G)\ensuremath{^{2}}=A+G \ensuremath{\leftrightarrow} o COLAPSO \ensuremath{\Psi} (a medição, a projeção a um poço); o GATO dissipa (Onsander) \ensuremath{\leftrightarrow} o \ensuremath{\delta} (a escada CASL, a dissipação/de Sitter); os n-QUBITS (2\^{}n) \ensuremath{\leftrightarrow} o ESTADO do BAI; a ÓRBITA (as potências) \ensuremath{\leftrightarrow} a órbita/a dinâmica do BAI. O corpo matricial é a álgebra do BAI: o gato mede (o Hamiltoniano), o gambá gira (a evolução), o colapso é o idempotente, o invariante é a+c\ensuremath{^{2}}=1.
\prov{relações: relaciona corpo\_matricial\_matrix; usa gato; usa gamba\_antissimetrico; relaciona bai; relaciona gato\_e\_o\_hamiltoniano; relaciona invariante\_estelar}
\prov{proveniência: wiki \textperiodcentered{} Floxina (investigação) \textperiodcentered{} fato \textperiodcentered{} confiança 1.0 \textperiodcentered{} domínio floxina \textperiodcentered{} história 399 versões no jornal}

%CRISTAL {"arestas": [{"alvo": "indice_milenios", "confianca": 1.0, "epistemico": "fato", "origem": "wiki", "rel": "parte_de"}, {"alvo": "omnitrix_corpo_universal", "confianca": 1.0, "epistemico": "fato", "origem": "wiki", "rel": "relaciona"}, {"alvo": "p_vs_np_milenio", "confianca": 1.0, "epistemico": "fato", "origem": "wiki", "rel": "relaciona"}, {"alvo": "computador_universal_pnp", "confianca": 1.0, "epistemico": "fato", "origem": "wiki", "rel": "usa"}, {"alvo": "teorema_universal_residuo_zero", "confianca": 1.0, "epistemico": "fato", "origem": "wiki", "rel": "usa"}, {"alvo": "hipotese_de_riemann", "confianca": 1.0, "epistemico": "fato", "origem": "wiki", "rel": "relaciona"}, {"alvo": "dinamica_de_selberg_o_termo", "confianca": 1.0, "epistemico": "fato", "origem": "wiki", "rel": "usa"}, {"alvo": "gap_espectral", "confianca": 1.0, "epistemico": "fato", "origem": "wiki", "rel": "usa"}, {"alvo": "yang_mills_mass_gap_milenio", "confianca": 1.0, "epistemico": "fato", "origem": "wiki", "rel": "relaciona"}, {"alvo": "floxina_quantum_materia_escura", "confianca": 1.0, "epistemico": "fato", "origem": "wiki", "rel": "usa"}, {"alvo": "gerador_do_gap_comutador", "confianca": 1.0, "epistemico": "fato", "origem": "wiki", "rel": "usa"}, {"alvo": "navier--stokes_a_equacao_reconstruida_do_bit", "confianca": 1.0, "epistemico": "fato", "origem": "wiki", "rel": "relaciona"}, {"alvo": "dissipacao_anomala_onsager", "confianca": 1.0, "epistemico": "fato", "origem": "wiki", "rel": "usa"}, {"alvo": "a_ancora_hodge", "confianca": 1.0, "epistemico": "fato", "origem": "wiki", "rel": "relaciona"}, {"alvo": "confronto_millennium_compacto", "confianca": 1.0, "epistemico": "fato", "origem": "wiki", "rel": "usa"}, {"alvo": "bsd_a_intersecao_entre_as_retas_de_riemann", "confianca": 1.0, "epistemico": "fato", "origem": "wiki", "rel": "relaciona"}, {"alvo": "bsd_matricial_grade_bits", "confianca": 1.0, "epistemico": "fato", "origem": "wiki", "rel": "usa"}, {"alvo": "bsd_gap", "confianca": 1.0, "epistemico": "fato", "origem": "wiki", "rel": "usa"}, {"alvo": "conjectura_de_poincare", "confianca": 1.0, "epistemico": "fato", "origem": "wiki", "rel": "relaciona"}, {"alvo": "agulha_perelman_limite", "confianca": 1.0, "epistemico": "fato", "origem": "wiki", "rel": "usa"}, {"alvo": "poincare_universalista", "confianca": 1.0, "epistemico": "fato", "origem": "wiki", "rel": "relaciona"}], "confianca": 1.0, "contraexemplos": [], "descricao": "A cristalização do dicionário: cada um dos 7 problemas do milênio ↔ a sua face gato/gambá/gap no Corpo Algébrico, ligado ao nó REAL do problema no grafo. (1) P vs NP ↔ o computador universal (gambá=P, gato=NP; testar o resíduo sucinto=PSPACE) [p_vs_np_milenio]. (2) RH ↔ as linhas das parábolas λ=¼−δ² (a reta crítica; a costura Selberg-Anosov) [hipotese_de_riemann, dinamica_de_selberg_o_termo]. (3) Yang-Mills ↔ a dinâmica do gap, o mass gap 𝔤₂=14 (o associador, a floxina) [yang_mills_mass_gap_milenio]. (4) Navier-Stokes ↔ a dinâmica dissipativa (o gato difunde, a viscosidade; o gambá gira, a vorticidade) [navier--stokes_a_equacao_reconstruida_do_bit, dissipacao_anomala_onsager]. (5) Hodge ↔ o corpo compacto (Kähler/projetivo) [a_ancora_hodge]. (6) BSD ↔ as cúspides discretas, os Frobenius-gatos, L(E,s) [bsd_a_intersecao_entre_as_retas_de_riemann, bsd_gap]. (7) Poincaré ↔ a cirurgia de Perelman (o fluxo de Ricci, a agulha) — RESOLVIDO [conjectura_de_poincare, agulha_perelman_limite]. Um só GAP (a parábola λ=¼−δ²) atravessa RH, Yang-Mills e (nas raízes sobre os primos) BSD, no corpo compacto de Hodge; a computação sobre ele é P vs NP; a AGM fecha tudo sobre π. 1 resolvido (Poincaré), 6 abertos, NENHUM decretado — afirmo só o provado, a localização é síntese declarada. linguagem/confronto_millennium.py (indice_milenios, 3/3).", "epistemico": "hipotese", "exemplos": [], "id": "dicionario_milenio", "memoria": "semantica", "meta": {"dominio": "floxina", "fonte": "Floxina (investigação)"}, "origem": "wiki", "palavras_chave": [], "sinonimos": [], "tipo": "conceito", "titulo": "O Dicionário do Milênio: cada problema ↔ a sua face no Corpo Algébrico"}
\section{O Dicionário do Milênio: cada problema \ensuremath{\leftrightarrow} a sua face no Corpo Algébrico}
A cristalização do dicionário: cada um dos 7 problemas do milênio \ensuremath{\leftrightarrow} a sua face gato/gambá/gap no Corpo Algébrico, ligado ao nó REAL do problema no grafo. (1) P vs NP \ensuremath{\leftrightarrow} o computador universal (gambá=P, gato=NP; testar o resíduo sucinto=PSPACE) [p\_vs\_np\_milenio]. (2) RH \ensuremath{\leftrightarrow} as linhas das parábolas \ensuremath{\lambda}=\ensuremath{\tfrac14}\ensuremath{-}\ensuremath{\delta}\ensuremath{^{2}} (a reta crítica; a costura Selberg-Anosov) [hipotese\_de\_riemann, dinamica\_de\_selberg\_o\_termo]. (3) Yang-Mills \ensuremath{\leftrightarrow} a dinâmica do gap, o mass gap \ensuremath{\mathfrak{g}}\ensuremath{_{2}}=14 (o associador, a floxina) [yang\_mills\_mass\_gap\_milenio]. (4) Navier-Stokes \ensuremath{\leftrightarrow} a dinâmica dissipativa (o gato difunde, a viscosidade; o gambá gira, a vorticidade) [navier--stokes\_a\_equacao\_reconstruida\_do\_bit, dissipacao\_anomala\_onsager]. (5) Hodge \ensuremath{\leftrightarrow} o corpo compacto (Kähler/projetivo) [a\_ancora\_hodge]. (6) BSD \ensuremath{\leftrightarrow} as cúspides discretas, os Frobenius-gatos, L(E,s) [bsd\_a\_intersecao\_entre\_as\_retas\_de\_riemann, bsd\_gap]. (7) Poincaré \ensuremath{\leftrightarrow} a cirurgia de Perelman (o fluxo de Ricci, a agulha) — RESOLVIDO [conjectura\_de\_poincare, agulha\_perelman\_limite]. Um só GAP (a parábola \ensuremath{\lambda}=\ensuremath{\tfrac14}\ensuremath{-}\ensuremath{\delta}\ensuremath{^{2}}) atravessa RH, Yang-Mills e (nas raízes sobre os primos) BSD, no corpo compacto de Hodge; a computação sobre ele é P vs NP; a AGM fecha tudo sobre \ensuremath{\pi}. 1 resolvido (Poincaré), 6 abertos, NENHUM decretado — afirmo só o provado, a localização é síntese declarada. linguagem/confronto\_millennium.py (indice\_milenios, 3/3).
\prov{relações: parte\_de indice\_milenios; relaciona omnitrix\_corpo\_universal; relaciona p\_vs\_np\_milenio; usa computador\_universal\_pnp; usa teorema\_universal\_residuo\_zero; relaciona hipotese\_de\_riemann; usa dinamica\_de\_selberg\_o\_termo; usa gap\_espectral; relaciona yang\_mills\_mass\_gap\_milenio; usa floxina\_quantum\_materia\_escura; usa gerador\_do\_gap\_comutador; relaciona navier--stokes\_a\_equacao\_reconstruida\_do\_bit; usa dissipacao\_anomala\_onsager; relaciona a\_ancora\_hodge; usa confronto\_millennium\_compacto; relaciona bsd\_a\_intersecao\_entre\_as\_retas\_de\_riemann; usa bsd\_matricial\_grade\_bits; usa bsd\_gap; relaciona conjectura\_de\_poincare; usa agulha\_perelman\_limite; relaciona poincare\_universalista}
\prov{proveniência: wiki \textperiodcentered{} Floxina (investigação) \textperiodcentered{} hipotese \textperiodcentered{} confiança 1.0 \textperiodcentered{} domínio floxina \textperiodcentered{} história 924 versões no jornal}

%CRISTAL {"arestas": [{"alvo": "ofloxacino_floxina_real", "confianca": 1.0, "epistemico": "fato", "origem": "wiki", "rel": "parte_de"}, {"alvo": "floxina_quantum_materia_escura", "confianca": 1.0, "epistemico": "fato", "origem": "wiki", "rel": "relaciona"}, {"alvo": "gamba_antissimetrico", "confianca": 1.0, "epistemico": "fato", "origem": "wiki", "rel": "usa"}, {"alvo": "flip_e_rotacao_de_wick", "confianca": 1.0, "epistemico": "fato", "origem": "wiki", "rel": "usa"}, {"alvo": "gap_espectral", "confianca": 1.0, "epistemico": "fato", "origem": "wiki", "rel": "relaciona"}], "confianca": 1.0, "contraexemplos": [], "descricao": "O trocadilho do nome fica rigoroso pelo dicionário. O FLÚOR (F, a ponta eletronegativa reativa; o 'floxinato') ↔ o GAMBÁ (a ponta antissimétrica reativa). A QUIRALIDADE (o estereocentro; levofloxacino (S), gira a luz à esquerda) ↔ o FLIP/a rotação (o gambá gira; a levogiração = o sinal do flip). A inibição da DNA GIRASE (torce/destorce o DNA, o superenrolamento) ↔ o GAMBÁ = a TORÇÃO/a rotação (o operador antissimétrico, o flip de Wick) — o alvo é o TWIST. O núcleo quinolona (aromático, a corrente de anel) ↔ o gambá aromático (a fase U(1)). O −COOH (ácido) ↔ o GAP (a acidez). Os anéis/o esqueleto ↔ o GATO (o Hamiltoniano de Hückel). O antibiótico mata pela torção controlada do DNA ↔ a floxina é o quantum do gambá (a torção; 'gerenciar não reprimir'). E o nome: FLOXIN (a marca real) ↔ FLOXINA (o projeto). Síntese declarada: o dicionário casa as figuras (a torção/o gambá é o eixo comum), não decreta a farmacologia.", "epistemico": "hipotese", "exemplos": [], "id": "dicionario_ofloxacino_floxina", "memoria": "semantica", "meta": {"dominio": "floxina", "fonte": "Floxina (investigação)"}, "origem": "wiki", "palavras_chave": [], "sinonimos": [], "tipo": "conceito", "titulo": "Dicionário: o ofloxacino (Floxin real) ↔ a floxina (o projeto)"}
\section{Dicionário: o ofloxacino (Floxin real) \ensuremath{\leftrightarrow} a floxina (o projeto)}
O trocadilho do nome fica rigoroso pelo dicionário. O FLÚOR (F, a ponta eletronegativa reativa; o 'floxinato') \ensuremath{\leftrightarrow} o GAMBÁ (a ponta antissimétrica reativa). A QUIRALIDADE (o estereocentro; levofloxacino (S), gira a luz à esquerda) \ensuremath{\leftrightarrow} o FLIP/a rotação (o gambá gira; a levogiração = o sinal do flip). A inibição da DNA GIRASE (torce/destorce o DNA, o superenrolamento) \ensuremath{\leftrightarrow} o GAMBÁ = a TORÇÃO/a rotação (o operador antissimétrico, o flip de Wick) — o alvo é o TWIST. O núcleo quinolona (aromático, a corrente de anel) \ensuremath{\leftrightarrow} o gambá aromático (a fase U(1)). O \ensuremath{-}COOH (ácido) \ensuremath{\leftrightarrow} o GAP (a acidez). Os anéis/o esqueleto \ensuremath{\leftrightarrow} o GATO (o Hamiltoniano de Hückel). O antibiótico mata pela torção controlada do DNA \ensuremath{\leftrightarrow} a floxina é o quantum do gambá (a torção; 'gerenciar não reprimir'). E o nome: FLOXIN (a marca real) \ensuremath{\leftrightarrow} FLOXINA (o projeto). Síntese declarada: o dicionário casa as figuras (a torção/o gambá é o eixo comum), não decreta a farmacologia.
\prov{relações: parte\_de ofloxacino\_floxina\_real; relaciona floxina\_quantum\_materia\_escura; usa gamba\_antissimetrico; usa flip\_e\_rotacao\_de\_wick; relaciona gap\_espectral}
\prov{proveniência: wiki \textperiodcentered{} Floxina (investigação) \textperiodcentered{} hipotese \textperiodcentered{} confiança 1.0 \textperiodcentered{} domínio floxina \textperiodcentered{} história 312 versões no jornal}

%CRISTAL {"arestas": [{"alvo": "gerador_do_gap_comutador", "confianca": 1.0, "epistemico": "fato", "origem": "wiki", "rel": "usa"}, {"alvo": "dicionario_dtc_motor_esferico", "confianca": 1.0, "epistemico": "fato", "origem": "wiki", "rel": "relaciona"}, {"alvo": "gamba_lado_frio_clones", "confianca": 1.0, "epistemico": "fato", "origem": "wiki", "rel": "usa"}, {"alvo": "torre_dois_lados_classico_relativistico", "confianca": 1.0, "epistemico": "fato", "origem": "wiki", "rel": "usa"}, {"alvo": "maquina_gaiola_esquilo_dtc", "confianca": 1.0, "epistemico": "fato", "origem": "wiki", "rel": "usa"}], "confianca": 1.0, "contraexemplos": [], "descricao": "A BOBINA DE TESLA (o transformador ressonante) traduzida peça a peça. (1) o CENTELHADOR (spark gap): dispara SÓ acima da tensão de ruptura — um limiar discreto (abaixo=quieto, o vácuo; acima=a descarga de golpe) = o MASS GAP (Δ>0, Yang-Mills) = o ASSOCIADOR do octônio [a,b,c]=(ab)c−a(bc): o gap FEITO álgebra, ≠0 em 𝕆 (7D, não-associativo) e =0 em ℍ (a torre associativa 1/2/4). O spark gap, o mass gap e o associador são o MESMO gap em três roupas (físico, energético, algébrico). (2) a bobina PRIMÁRIA ↔ o corpo TROPICAL (idempotente e²=e, o quente que expande, ε=+1, |λ|>1) e a SECUNDÁRIA ↔ o corpo GLACIAL (nilpotente n²=0, o frio que só transporta sem dissipar); acopladas por Peirce=o flip de Wick. (3) a CAPACITÂNCIA C ↔ o campo ELÉTRICO (½q²/C) e a INDUTÂNCIA L ↔ o campo MAGNÉTICO (½LI²); a oscilação LC troca E↔B e ESSA troca É o ESQUILO: [x']=ω·G·[x], ω=1/√(LC), energia conservada (a rotação e^{Gt}). (4) a LUZ: c=1/√(μ₀ε₀) = o mesmo 1/√(LC) (ε₀↔C, μ₀↔L do vácuo); a luz não tem massa (mass gap 0) — o esquilo SOLTO, a troca E↔B irradiando, que escapa o centelhador. linguagem/vesica_corpo.py (dicionario_tesla_bobina, 6/6).", "epistemico": "fato", "exemplos": [], "id": "dicionario_tesla_bobina_gap", "memoria": "semantica", "meta": {"dominio": "floxina", "fonte": "Floxina (investigação)"}, "origem": "wiki", "palavras_chave": [], "sinonimos": [], "tipo": "conceito", "titulo": "Dicionário: a BOBINA DE TESLA — o centelhador=mass gap=associador do octônio; a luz"}
\section{Dicionário: a BOBINA DE TESLA — o centelhador=mass gap=associador do octônio; a luz}
A BOBINA DE TESLA (o transformador ressonante) traduzida peça a peça. (1) o CENTELHADOR (spark gap): dispara SÓ acima da tensão de ruptura — um limiar discreto (abaixo=quieto, o vácuo; acima=a descarga de golpe) = o MASS GAP (\ensuremath{\Delta}>0, Yang-Mills) = o ASSOCIADOR do octônio [a,b,c]=(ab)c\ensuremath{-}a(bc): o gap FEITO álgebra, \ensuremath{\ne}0 em \ensuremath{\mathbb{O}} (7D, não-associativo) e =0 em \ensuremath{\mathbb{H}} (a torre associativa 1/2/4). O spark gap, o mass gap e o associador são o MESMO gap em três roupas (físico, energético, algébrico). (2) a bobina PRIMÁRIA \ensuremath{\leftrightarrow} o corpo TROPICAL (idempotente e\ensuremath{^{2}}=e, o quente que expande, \ensuremath{\varepsilon}=+1, |\ensuremath{\lambda}|>1) e a SECUNDÁRIA \ensuremath{\leftrightarrow} o corpo GLACIAL (nilpotente n\ensuremath{^{2}}=0, o frio que só transporta sem dissipar); acopladas por Peirce=o flip de Wick. (3) a CAPACITÂNCIA C \ensuremath{\leftrightarrow} o campo ELÉTRICO (\ensuremath{\tfrac12}q\ensuremath{^{2}}/C) e a INDUTÂNCIA L \ensuremath{\leftrightarrow} o campo MAGNÉTICO (\ensuremath{\tfrac12}LI\ensuremath{^{2}}); a oscilação LC troca E\ensuremath{\leftrightarrow}B e ESSA troca É o ESQUILO: [x']=\ensuremath{\omega}\textperiodcentered G\textperiodcentered [x], \ensuremath{\omega}=1/\ensuremath{\sqrt{\,}}(LC), energia conservada (a rotação e\^{}\{Gt\}). (4) a LUZ: c=1/\ensuremath{\sqrt{\,}}(\ensuremath{\mu}\ensuremath{_{0}}\ensuremath{\varepsilon}\ensuremath{_{0}}) = o mesmo 1/\ensuremath{\sqrt{\,}}(LC) (\ensuremath{\varepsilon}\ensuremath{_{0}}\ensuremath{\leftrightarrow}C, \ensuremath{\mu}\ensuremath{_{0}}\ensuremath{\leftrightarrow}L do vácuo); a luz não tem massa (mass gap 0) — o esquilo SOLTO, a troca E\ensuremath{\leftrightarrow}B irradiando, que escapa o centelhador. linguagem/vesica\_corpo.py (dicionario\_tesla\_bobina, 6/6).
\prov{relações: usa gerador\_do\_gap\_comutador; relaciona dicionario\_dtc\_motor\_esferico; usa gamba\_lado\_frio\_clones; usa torre\_dois\_lados\_classico\_relativistico; usa maquina\_gaiola\_esquilo\_dtc}
\prov{proveniência: wiki \textperiodcentered{} Floxina (investigação) \textperiodcentered{} fato \textperiodcentered{} confiança 1.0 \textperiodcentered{} domínio floxina \textperiodcentered{} história 180 versões no jornal}

%CRISTAL {"arestas": [{"alvo": "teste_molecula_tiol_gato_n_qubit", "confianca": 1.0, "epistemico": "fato", "origem": "wiki", "rel": "relaciona"}, {"alvo": "gamba_antissimetrico", "confianca": 1.0, "epistemico": "fato", "origem": "wiki", "rel": "usa"}, {"alvo": "metais", "confianca": 1.0, "epistemico": "fato", "origem": "wiki", "rel": "relaciona"}], "confianca": 1.0, "contraexemplos": [], "descricao": "As propriedades reais do tiol (−SH, a ponta reativa da molécula do gambá) costuram com o corpo: (1) a LIGAÇÃO DISSULFETO 2 R-SH ⇌ R-S-S-R (reversível, o switch redox/a glutationa) ↔ o GAMBÁ/o FLIP (o par reversível, norma conservada); (2) a LIGAÇÃO A METAIS (mercaptan, S–Hg/S–Au, o S mole) ↔ os METAIS σ_n (a reta metálica); (3) a ACIDEZ (pKa≈10,6<16, o S-H fraco, o próton solto) ↔ o GAP/a desafasagem; (4) o DOADOR DE H (BDE S-H≈87<105, radical) ↔ o BIT/QUBIT (o átomo doado, |ψ|²=1); (5) a NUCLEOFILIA (o S mole, reativo) ↔ a ponta anti-simétrica (o gambá); (6) o CHEIRO (detecção a ppb) ↔ a assinatura/o fail-closed (a litografia); (7) o ESQUELETO de C (adjacência simétrica) ↔ o GATO (Hückel→n-qubits). A química do tiol (o gambá) costura com o corpo: dissulfeto=o flip, metal=σ_n, esqueleto=o gato. FATO: as propriedades (química conhecida); SÍNTESE: a correspondência. linguagem/dicionario_tiol.py (5/5).", "epistemico": "hipotese", "exemplos": [], "id": "dicionario_tiol_gato_gamba", "memoria": "semantica", "meta": {"dominio": "floxina", "fonte": "Floxina (investigação)"}, "origem": "wiki", "palavras_chave": [], "sinonimos": [], "tipo": "conceito", "titulo": "Dicionário: as propriedades do tiol (−SH) ↔ o corpo (gato/gambá)"}
\section{Dicionário: as propriedades do tiol (\ensuremath{-}SH) \ensuremath{\leftrightarrow} o corpo (gato/gambá)}
As propriedades reais do tiol (\ensuremath{-}SH, a ponta reativa da molécula do gambá) costuram com o corpo: (1) a LIGAÇÃO DISSULFETO 2 R-SH \ensuremath{\rightleftharpoons} R-S-S-R (reversível, o switch redox/a glutationa) \ensuremath{\leftrightarrow} o GAMBÁ/o FLIP (o par reversível, norma conservada); (2) a LIGAÇÃO A METAIS (mercaptan, S–Hg/S–Au, o S mole) \ensuremath{\leftrightarrow} os METAIS \ensuremath{\sigma}\_n (a reta metálica); (3) a ACIDEZ (pKa\ensuremath{\approx}10,6<16, o S-H fraco, o próton solto) \ensuremath{\leftrightarrow} o GAP/a desafasagem; (4) o DOADOR DE H (BDE S-H\ensuremath{\approx}87<105, radical) \ensuremath{\leftrightarrow} o BIT/QUBIT (o átomo doado, |\ensuremath{\psi}|\ensuremath{^{2}}=1); (5) a NUCLEOFILIA (o S mole, reativo) \ensuremath{\leftrightarrow} a ponta anti-simétrica (o gambá); (6) o CHEIRO (detecção a ppb) \ensuremath{\leftrightarrow} a assinatura/o fail-closed (a litografia); (7) o ESQUELETO de C (adjacência simétrica) \ensuremath{\leftrightarrow} o GATO (Hückel\ensuremath{\to}n-qubits). A química do tiol (o gambá) costura com o corpo: dissulfeto=o flip, metal=\ensuremath{\sigma}\_n, esqueleto=o gato. FATO: as propriedades (química conhecida); SÍNTESE: a correspondência. linguagem/dicionario\_tiol.py (5/5).
\prov{relações: relaciona teste\_molecula\_tiol\_gato\_n\_qubit; usa gamba\_antissimetrico; relaciona metais}
\prov{proveniência: wiki \textperiodcentered{} Floxina (investigação) \textperiodcentered{} hipotese \textperiodcentered{} confiança 1.0 \textperiodcentered{} domínio floxina \textperiodcentered{} história 248 versões no jornal}

%CRISTAL {"arestas": [{"alvo": "matrix_vesicular_corpo", "confianca": 1.0, "epistemico": "fato", "origem": "wiki", "rel": "parte_de"}, {"alvo": "maquina_gaiola_esquilo_dtc", "confianca": 1.0, "epistemico": "fato", "origem": "wiki", "rel": "usa"}, {"alvo": "promocao_gamba_esquilo", "confianca": 1.0, "epistemico": "fato", "origem": "wiki", "rel": "usa"}], "confianca": 1.0, "contraexemplos": [], "descricao": "A tradução (física e EXATA) da matrix vesicular para o motor de GAIOLA DE ESQUILO: a VESICA (matriz simétrica, a lente anisotrópica) ↔ a matriz de INDUTÂNCIA L (Ld,Lq, o circuito magnético do rotor); a ANISOTROPIA (Ld≠Lq, o split) ↔ a SALIÊNCIA do rotor (a relutância variável); a COSTURA anisotrópica ([M1,M2]=o esquilo) ↔ o TORQUE DE RELUTÂNCIA T=(Ld−Lq)·id·iq (o produto vetorial das correntes); as DUAS VESICAS (H/V) ↔ os eixos d-q (a transformada de PARK: direto/em quadratura); a rotação 90° (o esquilo i) ↔ a rotação d→q; d/r (a parábola) ↔ o ângulo de carga δ (o DTC controla por ele); o círculo (isotrópico) ↔ o rotor NÃO-saliente (Ld=Lq): sem costura=sem torque de relutância. A costura anisotrópica que gera o corpo É o torque de relutância que gira o rotor. linguagem/vesica_corpo.py (dicionario_vesica_motor, 4/4).", "epistemico": "fato", "exemplos": [], "id": "dicionario_vesica_motor_esquilo", "memoria": "semantica", "meta": {"dominio": "floxina", "fonte": "Floxina (investigação)"}, "origem": "wiki", "palavras_chave": [], "sinonimos": [], "tipo": "conceito", "titulo": "Dicionário: a matrix vesicular ↔ o motor de gaiola de esquilo (rotor saliente, relutância)"}
\section{Dicionário: a matrix vesicular \ensuremath{\leftrightarrow} o motor de gaiola de esquilo (rotor saliente, relutância)}
A tradução (física e EXATA) da matrix vesicular para o motor de GAIOLA DE ESQUILO: a VESICA (matriz simétrica, a lente anisotrópica) \ensuremath{\leftrightarrow} a matriz de INDUTÂNCIA L (Ld,Lq, o circuito magnético do rotor); a ANISOTROPIA (Ld\ensuremath{\ne}Lq, o split) \ensuremath{\leftrightarrow} a SALIÊNCIA do rotor (a relutância variável); a COSTURA anisotrópica ([M1,M2]=o esquilo) \ensuremath{\leftrightarrow} o TORQUE DE RELUTÂNCIA T=(Ld\ensuremath{-}Lq)\textperiodcentered id\textperiodcentered iq (o produto vetorial das correntes); as DUAS VESICAS (H/V) \ensuremath{\leftrightarrow} os eixos d-q (a transformada de PARK: direto/em quadratura); a rotação 90$^{\circ}$ (o esquilo i) \ensuremath{\leftrightarrow} a rotação d\ensuremath{\to}q; d/r (a parábola) \ensuremath{\leftrightarrow} o ângulo de carga \ensuremath{\delta} (o DTC controla por ele); o círculo (isotrópico) \ensuremath{\leftrightarrow} o rotor NÃO-saliente (Ld=Lq): sem costura=sem torque de relutância. A costura anisotrópica que gera o corpo É o torque de relutância que gira o rotor. linguagem/vesica\_corpo.py (dicionario\_vesica\_motor, 4/4).
\prov{relações: parte\_de matrix\_vesicular\_corpo; usa maquina\_gaiola\_esquilo\_dtc; usa promocao\_gamba\_esquilo}
\prov{proveniência: wiki \textperiodcentered{} Floxina (investigação) \textperiodcentered{} fato \textperiodcentered{} confiança 1.0 \textperiodcentered{} domínio floxina \textperiodcentered{} história 128 versões no jornal}

%CRISTAL {"arestas": [{"alvo": "corpo_equacoes_diferenciais", "confianca": 1.0, "epistemico": "fato", "origem": "wiki", "rel": "usa"}, {"alvo": "parabola_geradora_idempotente", "confianca": 1.0, "epistemico": "fato", "origem": "wiki", "rel": "usa"}, {"alvo": "navier_stokes_caso_particular", "confianca": 1.0, "epistemico": "fato", "origem": "wiki", "rel": "relaciona"}], "confianca": 1.0, "contraexemplos": [], "descricao": "A ENERGIA SÓ SENTE O GATO: d/dt·½|x|²=xᵀ·A·x=xᵀ·Sym(A)·x — o ESQUILO (Skew) contribui 0 (xᵀ(Skew)x=0, verificado 1e-16). Logo: (1) DISSIPAÇÃO=GATO — espectro REAL (a reta, a taxa ±log σ), a energia MUDA, espalha/mistura (hiperbólico); (2) CONCENTRAÇÃO=ESQUILO — espectro IMAGINÁRIO (±iω), a energia CONSERVA, gira/foca (elíptico). E os DOIS OBEDECEM A PARÁBOLA: a discriminante Δ=(tr A)²−4·det A (=o invariante a+c²=1) coloca cada um — a família A_θ=cosθ·gato+sinθ·esquilo tem Δ(θ)=4cos2θ, varrendo DISSIPAÇÃO(Δ>0,gato) → BORDA(Δ=0) → CONCENTRAÇÃO(Δ<0,esquilo). Os METAIS são o lado da dissipação: A_n=[[m,1],[1,0]] tem Δ=m²+4>0 (σ_n e −1/σ_n reais, ±log σ_n); o esquilo Δ=−4<0 (±i) — o outro lado do MESMO discriminante. Em Navier-Stokes: a viscosidade (gato) dissipa, o convectivo (esquilo) concentra — a batalha é a borda 3D. linguagem/corpo_edp.py (dissipacao_gato_concentracao_esquilo_parabola 5/5); papers/equacoes_diferenciais.tex (Parte IX).", "epistemico": "fato", "exemplos": [], "id": "dissipacao_gato_concentracao_esquilo", "memoria": "semantica", "meta": {"dominio": "floxina", "fonte": "Floxina (investigação)"}, "origem": "wiki", "palavras_chave": [], "sinonimos": [], "tipo": "conceito", "titulo": "Dissipação=gato (espectro real, espalha), concentração=esquilo (espectro imaginário, gira) — e os dois obedecem a parábola (a discriminante os separa)"}
\section{Dissipação=gato (espectro real, espalha), concentração=esquilo (espectro imaginário, gira) — e os dois obedecem a parábola (a discriminante os separa)}
A ENERGIA SÓ SENTE O GATO: d/dt\textperiodcentered \ensuremath{\tfrac12}|x|\ensuremath{^{2}}=x\ensuremath{^{T}}\textperiodcentered A\textperiodcentered x=x\ensuremath{^{T}}\textperiodcentered Sym(A)\textperiodcentered x — o ESQUILO (Skew) contribui 0 (x\ensuremath{^{T}}(Skew)x=0, verificado 1e-16). Logo: (1) DISSIPAÇÃO=GATO — espectro REAL (a reta, a taxa $\pm$log \ensuremath{\sigma}), a energia MUDA, espalha/mistura (hiperbólico); (2) CONCENTRAÇÃO=ESQUILO — espectro IMAGINÁRIO ($\pm$i\ensuremath{\omega}), a energia CONSERVA, gira/foca (elíptico). E os DOIS OBEDECEM A PARÁBOLA: a discriminante \ensuremath{\Delta}=(tr A)\ensuremath{^{2}}\ensuremath{-}4\textperiodcentered det A (=o invariante a+c\ensuremath{^{2}}=1) coloca cada um — a família A\_\ensuremath{\theta}=cos\ensuremath{\theta}\textperiodcentered gato+sin\ensuremath{\theta}\textperiodcentered esquilo tem \ensuremath{\Delta}(\ensuremath{\theta})=4cos2\ensuremath{\theta}, varrendo DISSIPAÇÃO(\ensuremath{\Delta}>0,gato) \ensuremath{\to} BORDA(\ensuremath{\Delta}=0) \ensuremath{\to} CONCENTRAÇÃO(\ensuremath{\Delta}<0,esquilo). Os METAIS são o lado da dissipação: A\_n=[[m,1],[1,0]] tem \ensuremath{\Delta}=m\ensuremath{^{2}}+4>0 (\ensuremath{\sigma}\_n e \ensuremath{-}1/\ensuremath{\sigma}\_n reais, $\pm$log \ensuremath{\sigma}\_n); o esquilo \ensuremath{\Delta}=\ensuremath{-}4<0 ($\pm$i) — o outro lado do MESMO discriminante. Em Navier-Stokes: a viscosidade (gato) dissipa, o convectivo (esquilo) concentra — a batalha é a borda 3D. linguagem/corpo\_edp.py (dissipacao\_gato\_concentracao\_esquilo\_parabola 5/5); papers/equacoes\_diferenciais.tex (Parte IX).
\prov{relações: usa corpo\_equacoes\_diferenciais; usa parabola\_geradora\_idempotente; relaciona navier\_stokes\_caso\_particular}
\prov{proveniência: wiki \textperiodcentered{} Floxina (investigação) \textperiodcentered{} fato \textperiodcentered{} confiança 1.0 \textperiodcentered{} domínio floxina \textperiodcentered{} história 24 versões no jornal}

%CRISTAL {"arestas": [{"alvo": "corpo_edp", "confianca": 1.0, "epistemico": "fato", "origem": "wiki", "rel": "usa"}, {"alvo": "torre_hurwitz", "confianca": 1.0, "epistemico": "fato", "origem": "wiki", "rel": "usa"}], "confianca": 1.0, "contraexemplos": [], "descricao": "O ramo CLÁSSICO da torre de Hurwitz (ε=+1, positivo-definido → ELÍPTICO): as unidades imaginárias das álgebras de divisão ℂ(2D)/ℍ(4D)/𝕆(8D) satisfazem a relação de CLIFFORD eₐeᵦ+eᵦeₐ=−2δₐᵦ (verificado, erro=0), logo o operador de Cauchy-Riemann/Fueter/Dirac D=Σeₐ∂ₐ tem D̄D=∇² — a RAIZ DE DIRAC do laplaciano (√−∇²). É o operador ESTÁTICO/equilíbrio, o módulo: a norma é multiplicativa N(xy)=N(x)N(y) (Hurwitz), o cristal euclidiano. As três álgebras normadas de divisão dim 2,4,8 são exatamente as que fatoram o Laplaciano em √. linguagem/corpo_edp.py (edp_classica_2_4_8 4/4); usa torre_hurwitz. papers/equacoes_diferenciais.tex (Parte VIII).", "epistemico": "fato", "exemplos": [], "id": "edp_classica_2_4_8", "memoria": "semantica", "meta": {"dominio": "floxina", "fonte": "Floxina (investigação)"}, "origem": "wiki", "palavras_chave": [], "sinonimos": [], "tipo": "conceito", "titulo": "O lado CLÁSSICO das EDP (ε=+1, euclidiano→elíptico): ℂ/ℍ/𝕆 (dim 2,4,8) são a raiz de Dirac do laplaciano"}
\section{O lado CLÁSSICO das EDP (\ensuremath{\varepsilon}=+1, euclidiano\ensuremath{\to}elíptico): \ensuremath{\mathbb{C}}/\ensuremath{\mathbb{H}}/\ensuremath{\mathbb{O}} (dim 2,4,8) são a raiz de Dirac do laplaciano}
O ramo CLÁSSICO da torre de Hurwitz (\ensuremath{\varepsilon}=+1, positivo-definido \ensuremath{\to} ELÍPTICO): as unidades imaginárias das álgebras de divisão \ensuremath{\mathbb{C}}(2D)/\ensuremath{\mathbb{H}}(4D)/\ensuremath{\mathbb{O}}(8D) satisfazem a relação de CLIFFORD e\ensuremath{_{a}}e\ensuremath{_{\beta}}+e\ensuremath{_{\beta}}e\ensuremath{_{a}}=\ensuremath{-}2\ensuremath{\delta}\ensuremath{_{a\beta}} (verificado, erro=0), logo o operador de Cauchy-Riemann/Fueter/Dirac D=\ensuremath{\Sigma}e\ensuremath{_{a}}\ensuremath{\partial}\ensuremath{_{a}} tem D\ensuremath{}D=\ensuremath{\nabla}\ensuremath{^{2}} — a RAIZ DE DIRAC do laplaciano (\ensuremath{\sqrt{\,}}\ensuremath{-}\ensuremath{\nabla}\ensuremath{^{2}}). É o operador ESTÁTICO/equilíbrio, o módulo: a norma é multiplicativa N(xy)=N(x)N(y) (Hurwitz), o cristal euclidiano. As três álgebras normadas de divisão dim 2,4,8 são exatamente as que fatoram o Laplaciano em \ensuremath{\sqrt{\,}}. linguagem/corpo\_edp.py (edp\_classica\_2\_4\_8 4/4); usa torre\_hurwitz. papers/equacoes\_diferenciais.tex (Parte VIII).
\prov{relações: usa corpo\_edp; usa torre\_hurwitz}
\prov{proveniência: wiki \textperiodcentered{} Floxina (investigação) \textperiodcentered{} fato \textperiodcentered{} confiança 1.0 \textperiodcentered{} domínio floxina \textperiodcentered{} história 36 versões no jornal}

%CRISTAL {"arestas": [{"alvo": "corpo_edp", "confianca": 1.0, "epistemico": "fato", "origem": "wiki", "rel": "usa"}, {"alvo": "octonios_hiperbolicos", "confianca": 1.0, "epistemico": "fato", "origem": "wiki", "rel": "usa"}, {"alvo": "edp_classica_2_4_8", "confianca": 1.0, "epistemico": "fato", "origem": "wiki", "rel": "relaciona"}], "confianca": 1.0, "contraexemplos": [], "descricao": "O ramo RELATIVÍSTICO (via Gentil, assinatura de Minkowski → HIPERBÓLICO) vive na PARTE IMAGINÁRIA: Im ℍ=ℝ³ e Im 𝕆=ℝ⁷ — as ÚNICAS dimensões com PRODUTO VETORIAL (Hurwitz): x×y⊥x,y e |x×y|²=|x|²|y|²−(x·y)² (Lagrange, verificado só em 3 e 7). O produto vetorial É o ROTACIONAL ∇×; a assinatura de Lorentz dá a onda □=∂_t²−∇² e a raiz de Dirac relativística {γ^μ,γ^ν}=2η^μν ⟹ (γ^μ∂_μ)²=□ (verificado, erro=0). A ESTELAR 𝓔ₛ DO GENTIL (3D, conserva a norma) é o nível LIMPO (associativo, associador=0, mass gap 0 — o cone de luz); o octônio (7D) é o NÃO-associativo (associador≠0, mass gap=𝔤₂=14). EMPARELHAMENTO (a parte imaginária baixa a dim em 1): clássico 4D(ℍ)↔relativístico 3D(Im ℍ, Gentil), clássico 8D(𝕆)↔relativístico 7D(Im 𝕆), clássico 2D(ℂ)↔1D (a fase U(1), sem vetorial). Os dois ramos ε=±1 da mesma torre — o elíptico (equilíbrio) e o hiperbólico (o cone de luz). linguagem/corpo_edp.py (edp_relativistica_gentil_3_7 5/5); usa reta_mineral_7d. papers/equacoes_diferenciais.tex (Parte VIII).", "epistemico": "fato", "exemplos": [], "id": "edp_relativistica_gentil_3_7", "memoria": "semantica", "meta": {"dominio": "floxina", "fonte": "Floxina (investigação)"}, "origem": "wiki", "palavras_chave": [], "sinonimos": [], "tipo": "conceito", "titulo": "O lado RELATIVÍSTICO das EDP via Gentil (Minkowski→hiperbólico): dim 3,7 (Im ℍ/Im 𝕆) são a onda □"}
\section{O lado RELATIVÍSTICO das EDP via Gentil (Minkowski\ensuremath{\to}hiperbólico): dim 3,7 (Im \ensuremath{\mathbb{H}}/Im \ensuremath{\mathbb{O}}) são a onda $\square$}
O ramo RELATIVÍSTICO (via Gentil, assinatura de Minkowski \ensuremath{\to} HIPERBÓLICO) vive na PARTE IMAGINÁRIA: Im \ensuremath{\mathbb{H}}=\ensuremath{\mathbb{R}}\ensuremath{^{3}} e Im \ensuremath{\mathbb{O}}=\ensuremath{\mathbb{R}}\ensuremath{^{7}} — as ÚNICAS dimensões com PRODUTO VETORIAL (Hurwitz): x$\times$y\ensuremath{\perp}x,y e |x$\times$y|\ensuremath{^{2}}=|x|\ensuremath{^{2}}|y|\ensuremath{^{2}}\ensuremath{-}(x\textperiodcentered y)\ensuremath{^{2}} (Lagrange, verificado só em 3 e 7). O produto vetorial É o ROTACIONAL \ensuremath{\nabla}$\times$; a assinatura de Lorentz dá a onda $\square$=\ensuremath{\partial}\_t\ensuremath{^{2}}\ensuremath{-}\ensuremath{\nabla}\ensuremath{^{2}} e a raiz de Dirac relativística \{\ensuremath{\gamma}\^{}\ensuremath{\mu},\ensuremath{\gamma}\^{}\ensuremath{\nu}\}=2\ensuremath{\eta}\^{}\ensuremath{\mu}\ensuremath{\nu} \ensuremath{\Longrightarrow} (\ensuremath{\gamma}\^{}\ensuremath{\mu}\ensuremath{\partial}\_\ensuremath{\mu})\ensuremath{^{2}}=$\square$ (verificado, erro=0). A ESTELAR \ensuremath{\mathcal{E}}\ensuremath{_{s}} DO GENTIL (3D, conserva a norma) é o nível LIMPO (associativo, associador=0, mass gap 0 — o cone de luz); o octônio (7D) é o NÃO-associativo (associador\ensuremath{\ne}0, mass gap=\ensuremath{\mathfrak{g}}\ensuremath{_{2}}=14). EMPARELHAMENTO (a parte imaginária baixa a dim em 1): clássico 4D(\ensuremath{\mathbb{H}})\ensuremath{\leftrightarrow}relativístico 3D(Im \ensuremath{\mathbb{H}}, Gentil), clássico 8D(\ensuremath{\mathbb{O}})\ensuremath{\leftrightarrow}relativístico 7D(Im \ensuremath{\mathbb{O}}), clássico 2D(\ensuremath{\mathbb{C}})\ensuremath{\leftrightarrow}1D (a fase U(1), sem vetorial). Os dois ramos \ensuremath{\varepsilon}=$\pm$1 da mesma torre — o elíptico (equilíbrio) e o hiperbólico (o cone de luz). linguagem/corpo\_edp.py (edp\_relativistica\_gentil\_3\_7 5/5); usa reta\_mineral\_7d. papers/equacoes\_diferenciais.tex (Parte VIII).
\prov{relações: usa corpo\_edp; usa octonios\_hiperbolicos; relaciona edp\_classica\_2\_4\_8}
\prov{proveniência: wiki \textperiodcentered{} Floxina (investigação) \textperiodcentered{} fato \textperiodcentered{} confiança 1.0 \textperiodcentered{} domínio floxina \textperiodcentered{} história 46 versões no jornal}

%CRISTAL {"arestas": [{"alvo": "corpo_morfico_matrix", "confianca": 1.0, "epistemico": "fato", "origem": "wiki", "rel": "usa"}, {"alvo": "entropia_fios_gato_esquilo", "confianca": 1.0, "epistemico": "fato", "origem": "wiki", "rel": "usa"}, {"alvo": "soma_minkowski_dilatacao_boleado", "confianca": 1.0, "epistemico": "fato", "origem": "wiki", "rel": "relaciona"}, {"alvo": "omnitrix_corpo_universal", "confianca": 1.0, "epistemico": "fato", "origem": "wiki", "rel": "usa"}], "confianca": 1.0, "contraexemplos": [], "descricao": "A borracha não estica por ENERGIA (como o metal, esticando ligações) — estica por ENTROPIA. Cada fio (o gato+esquilo, a órbita = a cadeia polimérica) tem muitas configurações enroladas; esticar ALINHA a cadeia e REDUZ as configurações → baixa a entropia S. A força é ENTRÓPICA: a mola gaussiana S(x)=S₀−3kx²/(2Nb²) dá F=−T·dS/dx=3kT/(Nb²)·x — linear (Hooke) mas a RIGIDEZ ∝ T (a borracha ENDURECE quente, o efeito Gough-Joule, o oposto do metal). O sólido neo-Hookeano incompressível: σ=G(λ−1/λ²), G=nkT (o módulo ∝ T); INCOMPRESSÍVEL (o volume conserva, λ·(1/√λ)²=1) — a cápsula estica no eixo e AFINA nos lados. É o corpo mórfico com uma dinâmica: a deformação (o boleado/Minkowski) ganha uma FORÇA entrópica, REVERSÍVEL pelo esquilo (sem dissipação), herdando a entropia dos fios (h=log σ_n). linguagem/omnitrix.py (elasticidade_entropica 5/5); glsl_ptx_frag (FRAG_CILINDRO_BORRACHA, o efeito no metal, 10/10).", "epistemico": "fato", "exemplos": [], "id": "elasticidade_entropica_borracha", "memoria": "semantica", "meta": {"dominio": "floxina", "fonte": "Floxina (investigação)"}, "origem": "wiki", "palavras_chave": [], "sinonimos": [], "tipo": "conceito", "titulo": "A elasticidade entrópica (o efeito de borracha): a força restauradora É a entropia (F=−T·dS/dx, ∝ T)"}
\section{A elasticidade entrópica (o efeito de borracha): a força restauradora É a entropia (F=\ensuremath{-}T\textperiodcentered dS/dx, \ensuremath{\propto} T)}
A borracha não estica por ENERGIA (como o metal, esticando ligações) — estica por ENTROPIA. Cada fio (o gato+esquilo, a órbita = a cadeia polimérica) tem muitas configurações enroladas; esticar ALINHA a cadeia e REDUZ as configurações \ensuremath{\to} baixa a entropia S. A força é ENTRÓPICA: a mola gaussiana S(x)=S\ensuremath{_{0}}\ensuremath{-}3kx\ensuremath{^{2}}/(2Nb\ensuremath{^{2}}) dá F=\ensuremath{-}T\textperiodcentered dS/dx=3kT/(Nb\ensuremath{^{2}})\textperiodcentered x — linear (Hooke) mas a RIGIDEZ \ensuremath{\propto} T (a borracha ENDURECE quente, o efeito Gough-Joule, o oposto do metal). O sólido neo-Hookeano incompressível: \ensuremath{\sigma}=G(\ensuremath{\lambda}\ensuremath{-}1/\ensuremath{\lambda}\ensuremath{^{2}}), G=nkT (o módulo \ensuremath{\propto} T); INCOMPRESSÍVEL (o volume conserva, \ensuremath{\lambda}\textperiodcentered (1/\ensuremath{\sqrt{\,}}\ensuremath{\lambda})\ensuremath{^{2}}=1) — a cápsula estica no eixo e AFINA nos lados. É o corpo mórfico com uma dinâmica: a deformação (o boleado/Minkowski) ganha uma FORÇA entrópica, REVERSÍVEL pelo esquilo (sem dissipação), herdando a entropia dos fios (h=log \ensuremath{\sigma}\_n). linguagem/omnitrix.py (elasticidade\_entropica 5/5); glsl\_ptx\_frag (FRAG\_CILINDRO\_BORRACHA, o efeito no metal, 10/10).
\prov{relações: usa corpo\_morfico\_matrix; usa entropia\_fios\_gato\_esquilo; relaciona soma\_minkowski\_dilatacao\_boleado; usa omnitrix\_corpo\_universal}
\prov{proveniência: wiki \textperiodcentered{} Floxina (investigação) \textperiodcentered{} fato \textperiodcentered{} confiança 1.0 \textperiodcentered{} domínio floxina \textperiodcentered{} história 115 versões no jornal}

%CRISTAL {"arestas": [{"alvo": "omnitrix_corpo_universal", "confianca": 1.0, "epistemico": "fato", "origem": "wiki", "rel": "parte_de"}, {"alvo": "gamba_antissimetrico", "confianca": 1.0, "epistemico": "fato", "origem": "wiki", "rel": "usa"}, {"alvo": "gerador_do_gap_comutador", "confianca": 1.0, "epistemico": "fato", "origem": "wiki", "rel": "usa"}, {"alvo": "torre_matricial_coordenada", "confianca": 1.0, "epistemico": "fato", "origem": "wiki", "rel": "usa"}, {"alvo": "floxina_quantum_materia_escura", "confianca": 1.0, "epistemico": "fato", "origem": "wiki", "rel": "relaciona"}, {"alvo": "fisica_universal_omnitrix", "confianca": 1.0, "epistemico": "fato", "origem": "wiki", "rel": "relaciona"}], "confianca": 1.0, "contraexemplos": [], "descricao": "O lado gambá/gap das físicas do Corpo Algébrico. ELETROMAGNETISMO (o gambá=F_μν; a torre de GAUGE): F_μν ANTISSIMÉTRICO (o gambá); a força=a CURVATURA=[gato,gambá]=[[−2,1],[1,2]] (o gap); a gauge sobe a torre ℂ→U(1)(dim1), Im ℍ→SU(2)(dim3), Aut(𝕆)→G₂(dim14). ÓPTICA (o gap→a cor): a REFLEXÃO=os autovalores REAIS do gato (o espelho, σ e −1/σ); a REFRAÇÃO=a fase do gambá (±i, a rotação); o GAP é a COR — o gap de Selberg λ₁≥¼ é o tom fundamental da superfície modular. QUÂNTICA (os n-qubits): o gambá EVOLUI unitário (|ψ|²=1, e^{Gt} conserva a norma); o gato MEDE (o idempotente P²=P, o colapso); os n-qubits=a torre (ℂ/ℍ/𝕆=1/2/3). E a NÃO-ASSOCIATIVIDADE do 𝕆 (3 qubits, associador≠0=𝔤₂) É a quanticidade irredutível (a contextualidade: não há realidade associativa/clássica por baixo) — o associador que não zera (o mass gap, a floxina) é o que faz o mundo ser quântico. linguagem/omnitrix.py (em_optica_quantica_universal, 3/3). Síntese declarada.", "epistemico": "hipotese", "exemplos": [], "id": "em_optica_quantica_universal", "memoria": "semantica", "meta": {"dominio": "floxina", "fonte": "Floxina (investigação)"}, "origem": "wiki", "palavras_chave": [], "sinonimos": [], "tipo": "conceito", "titulo": "EM/óptica/quântica universais: o gambá+a torre de gauge, o gap(cor), os n-qubits+o 𝕆"}
\section{EM/óptica/quântica universais: o gambá+a torre de gauge, o gap(cor), os n-qubits+o \ensuremath{\mathbb{O}}}
O lado gambá/gap das físicas do Corpo Algébrico. ELETROMAGNETISMO (o gambá=F\_\ensuremath{\mu}\ensuremath{\nu}; a torre de GAUGE): F\_\ensuremath{\mu}\ensuremath{\nu} ANTISSIMÉTRICO (o gambá); a força=a CURVATURA=[gato,gambá]=[[\ensuremath{-}2,1],[1,2]] (o gap); a gauge sobe a torre \ensuremath{\mathbb{C}}\ensuremath{\to}U(1)(dim1), Im \ensuremath{\mathbb{H}}\ensuremath{\to}SU(2)(dim3), Aut(\ensuremath{\mathbb{O}})\ensuremath{\to}G\ensuremath{_{2}}(dim14). ÓPTICA (o gap\ensuremath{\to}a cor): a REFLEXÃO=os autovalores REAIS do gato (o espelho, \ensuremath{\sigma} e \ensuremath{-}1/\ensuremath{\sigma}); a REFRAÇÃO=a fase do gambá ($\pm$i, a rotação); o GAP é a COR — o gap de Selberg \ensuremath{\lambda}\ensuremath{_{1}}\ensuremath{\ge}\ensuremath{\tfrac14} é o tom fundamental da superfície modular. QUÂNTICA (os n-qubits): o gambá EVOLUI unitário (|\ensuremath{\psi}|\ensuremath{^{2}}=1, e\^{}\{Gt\} conserva a norma); o gato MEDE (o idempotente P\ensuremath{^{2}}=P, o colapso); os n-qubits=a torre (\ensuremath{\mathbb{C}}/\ensuremath{\mathbb{H}}/\ensuremath{\mathbb{O}}=1/2/3). E a NÃO-ASSOCIATIVIDADE do \ensuremath{\mathbb{O}} (3 qubits, associador\ensuremath{\ne}0=\ensuremath{\mathfrak{g}}\ensuremath{_{2}}) É a quanticidade irredutível (a contextualidade: não há realidade associativa/clássica por baixo) — o associador que não zera (o mass gap, a floxina) é o que faz o mundo ser quântico. linguagem/omnitrix.py (em\_optica\_quantica\_universal, 3/3). Síntese declarada.
\prov{relações: parte\_de omnitrix\_corpo\_universal; usa gamba\_antissimetrico; usa gerador\_do\_gap\_comutador; usa torre\_matricial\_coordenada; relaciona floxina\_quantum\_materia\_escura; relaciona fisica\_universal\_omnitrix}
\prov{proveniência: wiki \textperiodcentered{} Floxina (investigação) \textperiodcentered{} hipotese \textperiodcentered{} confiança 1.0 \textperiodcentered{} domínio floxina \textperiodcentered{} história 350 versões no jornal}

%CRISTAL {"arestas": [{"alvo": "coracao", "confianca": 1.0, "epistemico": "fato", "origem": "wiki", "rel": "parte_de"}, {"alvo": "roda_emocoes_plutchik", "confianca": 1.0, "epistemico": "fato", "origem": "wiki", "rel": "usa"}, {"alvo": "circumplexo_russell", "confianca": 1.0, "epistemico": "fato", "origem": "wiki", "rel": "usa"}, {"alvo": "dimensoes_do_pipe", "confianca": 1.0, "epistemico": "fato", "origem": "wiki", "rel": "relaciona"}, {"alvo": "nucleo_g2_reabsorcao", "confianca": 1.0, "epistemico": "fato", "origem": "wiki", "rel": "relaciona"}, {"alvo": "g2_grupo_excepcional", "confianca": 1.0, "epistemico": "fato", "origem": "wiki", "rel": "usa"}], "confianca": 1.0, "contraexemplos": [], "descricao": "As 8 componentes do octônion (e0..e7) ↔ as 8 emoções primárias de Plutchik: alegria, confiança, medo, surpresa, tristeza, aversão, raiva, expectativa. Os OPOSTOS ficam a distância 4 (alegria↔tristeza, confiança↔aversão, medo↔raiva, surpresa↔expectativa) — exatamente a partição Cayley-Dickson a=[a1|a2] que constrói 𝕆. Cada emoção leva (valência, ativação) do circumplexo de Russell; a norma unitária = a intensidade conservada (o checksum).", "epistemico": "fato", "exemplos": [], "id": "emocao_octonion_plutchik", "memoria": "semantica", "meta": {"dominio": "coracao", "fonte": "coracao hub"}, "origem": "wiki", "palavras_chave": [], "sinonimos": [], "tipo": "conceito", "titulo": "As 8 emoções no octônion (Plutchik ↔ 𝕆)"}
\section{As 8 emoções no octônion (Plutchik \ensuremath{\leftrightarrow} \ensuremath{\mathbb{O}})}
As 8 componentes do octônion (e0..e7) \ensuremath{\leftrightarrow} as 8 emoções primárias de Plutchik: alegria, confiança, medo, surpresa, tristeza, aversão, raiva, expectativa. Os OPOSTOS ficam a distância 4 (alegria\ensuremath{\leftrightarrow}tristeza, confiança\ensuremath{\leftrightarrow}aversão, medo\ensuremath{\leftrightarrow}raiva, surpresa\ensuremath{\leftrightarrow}expectativa) — exatamente a partição Cayley-Dickson a=[a1|a2] que constrói \ensuremath{\mathbb{O}}. Cada emoção leva (valência, ativação) do circumplexo de Russell; a norma unitária = a intensidade conservada (o checksum).
\prov{relações: parte\_de coracao; usa roda\_emocoes\_plutchik; usa circumplexo\_russell; relaciona dimensoes\_do\_pipe; relaciona nucleo\_g2\_reabsorcao; usa g2\_grupo\_excepcional}
\prov{proveniência: wiki \textperiodcentered{} coracao hub \textperiodcentered{} fato \textperiodcentered{} confiança 1.0 \textperiodcentered{} domínio coracao \textperiodcentered{} história 21 versões no jornal}

%CRISTAL {"arestas": [{"alvo": "psicologia", "confianca": 1.0, "epistemico": "fato", "origem": "wiki", "rel": "parte_de"}], "confianca": 1.0, "contraexemplos": [], "descricao": "Paul Ekman: um conjunto de emoções básicas com expressão facial UNIVERSAL (alegria, tristeza, medo, raiva, nojo, surpresa) — reconhecidas entre culturas isoladas. A emoção tem uma assinatura corporal involuntária (as micro-expressões).", "epistemico": "fato", "exemplos": [], "id": "emocoes_basicas_ekman", "memoria": "semantica", "meta": {"autor": "Ekman", "dominio": "coracao", "fonte": "coracao hub"}, "origem": "wiki", "palavras_chave": [], "sinonimos": [], "tipo": "conceito", "titulo": "As emoções básicas (Ekman)"}
\section{As emoções básicas (Ekman)}
Paul Ekman: um conjunto de emoções básicas com expressão facial UNIVERSAL (alegria, tristeza, medo, raiva, nojo, surpresa) — reconhecidas entre culturas isoladas. A emoção tem uma assinatura corporal involuntária (as micro-expressões).
\prov{relações: parte\_de psicologia}
\prov{proveniência: wiki \textperiodcentered{} coracao hub \textperiodcentered{} fato \textperiodcentered{} confiança 1.0 \textperiodcentered{} domínio coracao \textperiodcentered{} história 6 versões no jornal}

%CRISTAL {"arestas": [{"alvo": "mapear_orbitas_grupo_discreto", "confianca": 1.0, "epistemico": "fato", "origem": "wiki", "rel": "usa"}, {"alvo": "matriz_costuras_paralela_universal", "confianca": 1.0, "epistemico": "fato", "origem": "wiki", "rel": "usa"}, {"alvo": "algebra_fundamental_gato_gamba", "confianca": 1.0, "epistemico": "fato", "origem": "wiki", "rel": "usa"}, {"alvo": "parabola_geradora_idempotente", "confianca": 1.0, "epistemico": "fato", "origem": "wiki", "rel": "relaciona"}, {"alvo": "shader_fios_shell_paralelo", "confianca": 1.0, "epistemico": "fato", "origem": "wiki", "rel": "usa"}], "confianca": 1.0, "contraexemplos": [], "descricao": "Cada fio do fur é uma ÓRBITA, e a órbita segue a DIREÇÃO do fio (radial) — por isso o fio aparece PONTILHADO (os pontos da órbita, os iterados). Cada fio tem uma MATRIZ própria: um GATO A_n=[[n,1],[1,0]] (hiperbólico, mede/DISSIPA) e um ESQUILO G=[[0,1],[−1,0]] (elíptico, gira/EVOLUI reversível). O CORPO ALGÉBRICO itera os dois (medir com o gato + evoluir com o esquilo) — e daí nasce a ENTROPIA: o gato é hiperbólico (σ_n>1), o Lyapunov λ=log σ_n>0 ⟹ a entropia de Kolmogorov-Sinai h=log σ_n (Pesin) — o gato PRODUZ entropia (a seta do tempo); o esquilo tem autovalores ±i (|λ|=1), Lyapunov 0 ⟹ entropia 0 (reversível, conserva). A entropia do fio = log σ_n = a PARÁBOLA (λ=log ρ=±log σ_n, a reta mineral); cada fio um metal n ⟹ um espectro de entropias (a variedade do fur). No pipe, cada fio ganha o seu gato: o espaçamento dos dots ∝ σ_n. linguagem/omnitrix.py (entropia_fios_gato_esquilo, 5/5).", "epistemico": "fato", "exemplos": [], "id": "entropia_fios_gato_esquilo", "memoria": "semantica", "meta": {"dominio": "floxina", "fonte": "Floxina (investigação)"}, "origem": "wiki", "palavras_chave": [], "sinonimos": [], "tipo": "conceito", "titulo": "Cada fio uma órbita: uma matriz por fio (gato+esquilo); a iteração pelo Corpo Algébrico gera a entropia"}
\section{Cada fio uma órbita: uma matriz por fio (gato+esquilo); a iteração pelo Corpo Algébrico gera a entropia}
Cada fio do fur é uma ÓRBITA, e a órbita segue a DIREÇÃO do fio (radial) — por isso o fio aparece PONTILHADO (os pontos da órbita, os iterados). Cada fio tem uma MATRIZ própria: um GATO A\_n=[[n,1],[1,0]] (hiperbólico, mede/DISSIPA) e um ESQUILO G=[[0,1],[\ensuremath{-}1,0]] (elíptico, gira/EVOLUI reversível). O CORPO ALGÉBRICO itera os dois (medir com o gato + evoluir com o esquilo) — e daí nasce a ENTROPIA: o gato é hiperbólico (\ensuremath{\sigma}\_n>1), o Lyapunov \ensuremath{\lambda}=log \ensuremath{\sigma}\_n>0 \ensuremath{\Longrightarrow} a entropia de Kolmogorov-Sinai h=log \ensuremath{\sigma}\_n (Pesin) — o gato PRODUZ entropia (a seta do tempo); o esquilo tem autovalores $\pm$i (|\ensuremath{\lambda}|=1), Lyapunov 0 \ensuremath{\Longrightarrow} entropia 0 (reversível, conserva). A entropia do fio = log \ensuremath{\sigma}\_n = a PARÁBOLA (\ensuremath{\lambda}=log \ensuremath{\rho}=$\pm$log \ensuremath{\sigma}\_n, a reta mineral); cada fio um metal n \ensuremath{\Longrightarrow} um espectro de entropias (a variedade do fur). No pipe, cada fio ganha o seu gato: o espaçamento dos dots \ensuremath{\propto} \ensuremath{\sigma}\_n. linguagem/omnitrix.py (entropia\_fios\_gato\_esquilo, 5/5).
\prov{relações: usa mapear\_orbitas\_grupo\_discreto; usa matriz\_costuras\_paralela\_universal; usa algebra\_fundamental\_gato\_gamba; relaciona parabola\_geradora\_idempotente; usa shader\_fios\_shell\_paralelo}
\prov{proveniência: wiki \textperiodcentered{} Floxina (investigação) \textperiodcentered{} fato \textperiodcentered{} confiança 1.0 \textperiodcentered{} domínio floxina \textperiodcentered{} história 162 versões no jornal}

%CRISTAL {"arestas": [{"alvo": "ponte_quantica_conjunto_dougherty_e_orbitais_de_hidrogenio", "confianca": 1.0, "epistemico": "fato", "origem": "ponte_quantica_dougherty.md", "rel": "parte_de"}, {"alvo": "word", "confianca": 0.5, "epistemico": "fato", "origem": "wiki", "rel": "relaciona"}, {"alvo": "virtualizacao", "confianca": 0.5, "epistemico": "fato", "origem": "wiki", "rel": "relaciona"}, {"alvo": "vida-morte", "confianca": 0.5, "epistemico": "fato", "origem": "wiki", "rel": "relaciona"}, {"alvo": "logica", "confianca": 0.5, "epistemico": "fato", "origem": "wiki", "rel": "relaciona"}], "confianca": 0.52, "contraexemplos": [], "descricao": "Tipo Schrödinger com Laplaciano toroidal e potencial magneto-hidrodinâmico:", "epistemico": "hipotese", "exemplos": [], "id": "equacao_efectiva_conjectura", "memoria": "semantica", "meta": {"dominio": "hipotese", "falsificavel": true, "nivel": 7}, "origem": "ponte_quantica_dougherty.md", "palavras_chave": [], "sinonimos": [], "tipo": "conceito", "titulo": "Equação efectiva (conjectura)"}
\section{Equação efectiva (conjectura)}
Tipo Schrödinger com Laplaciano toroidal e potencial magneto-hidrodinâmico:
\prov{relações: parte\_de ponte\_quantica\_conjunto\_dougherty\_e\_orbitais\_de\_hidrogenio; relaciona word; relaciona virtualizacao; relaciona vida-morte; relaciona logica}
\prov{proveniência: ponte\_quantica\_dougherty.md \textperiodcentered{} hipotese \textperiodcentered{} confiança 0.52 \textperiodcentered{} domínio hipotese \textperiodcentered{} história 10 versões no jornal}

%CRISTAL {"arestas": [{"alvo": "missao_engenharia_ouro", "confianca": 1.0, "epistemico": "fato", "origem": "wiki", "rel": "explica"}, {"alvo": "estatuto_correspondencia_derivacao", "confianca": 1.0, "epistemico": "fato", "origem": "wiki", "rel": "relaciona"}, {"alvo": "estatuto_derivado_posto_convencao", "confianca": 1.0, "epistemico": "fato", "origem": "wiki", "rel": "relaciona"}, {"alvo": "o_conselho", "confianca": 1.0, "epistemico": "fato", "origem": "wiki", "rel": "usa"}], "confianca": 1.0, "contraexemplos": [], "descricao": "A Missão só não é loucura porque diz, sem disfarce, o estatuto de cada peça: DEMONSTRADO = o teorema final (local ao nível); FORMALIZAÇÃO bem-definida = o blow-up e o operador R (a identificação com 'rasgar a função de onda' é INTERPRETATIVA); AXIOMA = 'ninguém está puro'; HIPÓTESE DE ENGENHARIA = a convergência do ganho; DEFINIÇÃO+teorema condicional = N* e parar nele. A Missão não PROVA que o ganho converge — ADOTA a convergência como o eixo que separa ouro de loucura, e a declara. Uma decisão, dita como decisão. É a mesma régua DNCI do grafo. A sabatina do conselho: atender o crivo UMA vez e parar é a Missão aplicada a si mesma.", "epistemico": "inferencia", "exemplos": [], "id": "estatuto_missao_ouro_loucura", "memoria": "semantica", "meta": {"autor": "Lopes/conselho", "dominio": "missao", "fonte": "machine/Missao.tex"}, "origem": "wiki", "palavras_chave": [], "sinonimos": [], "tipo": "conceito", "titulo": "A régua da Missão: o que separa ouro de loucura"}
\section{A régua da Missão: o que separa ouro de loucura}
A Missão só não é loucura porque diz, sem disfarce, o estatuto de cada peça: DEMONSTRADO = o teorema final (local ao nível); FORMALIZAÇÃO bem-definida = o blow-up e o operador R (a identificação com 'rasgar a função de onda' é INTERPRETATIVA); AXIOMA = 'ninguém está puro'; HIPÓTESE DE ENGENHARIA = a convergência do ganho; DEFINIÇÃO+teorema condicional = N* e parar nele. A Missão não PROVA que o ganho converge — ADOTA a convergência como o eixo que separa ouro de loucura, e a declara. Uma decisão, dita como decisão. É a mesma régua DNCI do grafo. A sabatina do conselho: atender o crivo UMA vez e parar é a Missão aplicada a si mesma.
\prov{relações: explica missao\_engenharia\_ouro; relaciona estatuto\_correspondencia\_derivacao; relaciona estatuto\_derivado\_posto\_convencao; usa o\_conselho}
\prov{proveniência: wiki \textperiodcentered{} machine/Missao.tex \textperiodcentered{} inferencia \textperiodcentered{} confiança 1.0 \textperiodcentered{} domínio missao \textperiodcentered{} história 5 versões no jornal}

%CRISTAL {"arestas": [{"alvo": "agulha_perelman_limite", "confianca": 1.0, "epistemico": "fato", "origem": "wiki", "rel": "parte_de"}, {"alvo": "metais", "confianca": 1.0, "epistemico": "fato", "origem": "wiki", "rel": "usa"}], "confianca": 1.0, "contraexemplos": [], "descricao": "O ESTOJO É FIXO — as agulhas são os METAIS (as cúspides principais de cada reta: a cúspide de σ_n é o próprio σ_n). NÃO SE INVENTA AGULHA NOVA. Ouro (menor λ) fura mais fundo (K=28, a agulha mais fina); os metais altos, mais raso. Para uma gravura, usa-se a agulha de um mineral em OUTRO: gravar(σ_agulha, σ_base) — a assinatura fica «de um metal em outro» (ouro⤷prata ≠ prata⤷ouro, a β-expansão cruzada), ou COMBINAÇÕES de várias agulhas numa base. É a ferramenta de litografia: escolhe-se a agulha (o metal) pela profundidade; o estojo é dado, não fabricado.", "epistemico": "fato", "exemplos": [], "id": "estojo_agulhas_minerais_fixo", "memoria": "semantica", "meta": {"dominio": "floxina", "fonte": "Floxina (investigação)"}, "origem": "wiki", "palavras_chave": [], "sinonimos": [], "tipo": "conceito", "titulo": "O estojo de agulhas minerais (fixo; a gravura cruzada)"}
\section{O estojo de agulhas minerais (fixo; a gravura cruzada)}
O ESTOJO É FIXO — as agulhas são os METAIS (as cúspides principais de cada reta: a cúspide de \ensuremath{\sigma}\_n é o próprio \ensuremath{\sigma}\_n). NÃO SE INVENTA AGULHA NOVA. Ouro (menor \ensuremath{\lambda}) fura mais fundo (K=28, a agulha mais fina); os metais altos, mais raso. Para uma gravura, usa-se a agulha de um mineral em OUTRO: gravar(\ensuremath{\sigma}\_agulha, \ensuremath{\sigma}\_base) — a assinatura fica «de um metal em outro» (ouro\ensuremath{\hookrightarrow}prata \ensuremath{\ne} prata\ensuremath{\hookrightarrow}ouro, a \ensuremath{\beta}-expansão cruzada), ou COMBINAÇÕES de várias agulhas numa base. É a ferramenta de litografia: escolhe-se a agulha (o metal) pela profundidade; o estojo é dado, não fabricado.
\prov{relações: parte\_de agulha\_perelman\_limite; usa metais}
\prov{proveniência: wiki \textperiodcentered{} Floxina (investigação) \textperiodcentered{} fato \textperiodcentered{} confiança 1.0 \textperiodcentered{} domínio floxina \textperiodcentered{} história 231 versões no jornal}

%CRISTAL {"arestas": [{"alvo": "transformacao_superficies_homeo_cirurgia", "confianca": 1.0, "epistemico": "fato", "origem": "wiki", "rel": "usa"}, {"alvo": "parabola_geradora_idempotente", "confianca": 1.0, "epistemico": "fato", "origem": "wiki", "rel": "usa"}, {"alvo": "corpo_topologico", "confianca": 1.0, "epistemico": "fato", "origem": "wiki", "rel": "usa"}], "confianca": 1.0, "contraexemplos": [], "descricao": "A estrela↔flor da JOALHERIA (a equação EXATA: shader_piramide.mapa_estrela5/mapa_flor — 5 amêndoas (vesicas) radiais a 72° fundidas por smin_exp, NÃO um cos(5v) inventado) lida no corpo topológico: (1) a 5-fold C₅ É o OURO σ₁=φ=(1+√5)/2, que satisfaz φ²−φ−1=0 = a PARÁBOLA GERADORA com m=1 (o metal dourado) — o pentagrama é a figura áurea, a simetria de 5 pontas É o metal m=1 (C₅-invariante, erro 3e-16). (2) a estrela e a flor são STAR-SHAPED (cada raio do centro cruza a borda 1 vez) ⟹ genus 0, χ=2. (3) o morph estrela→flor (as pontas afiadas arredondam em pétalas) é o HOMEOMORFISMO (o cristal, χ=2 preservado ∀morph, sem mudar a topologia — não é cirurgia); o mesmo regime da esfera↔elipsóide. Render: estrela_flor.gif. linguagem/corpo_topologico.py (estrela_flor_joalheria 3/3); usa shader_piramide (papers/joalheria.tex).", "epistemico": "fato", "exemplos": [], "id": "estrela_flor_joalheria_topologico", "memoria": "semantica", "meta": {"dominio": "floxina", "fonte": "Floxina (investigação)"}, "origem": "wiki", "palavras_chave": [], "sinonimos": [], "tipo": "conceito", "titulo": "A estrela↔flor da joalheria (equação exata) no corpo topológico: a 5-fold C₅ é o ouro σ₁=φ (a parábola m=1); o morph é o homeomorfismo (χ=2)"}
\section{A estrela\ensuremath{\leftrightarrow}flor da joalheria (equação exata) no corpo topológico: a 5-fold C\ensuremath{_{5}} é o ouro \ensuremath{\sigma}\ensuremath{_{1}}=\ensuremath{\varphi} (a parábola m=1); o morph é o homeomorfismo (\ensuremath{\chi}=2)}
A estrela\ensuremath{\leftrightarrow}flor da JOALHERIA (a equação EXATA: shader\_piramide.mapa\_estrela5/mapa\_flor — 5 amêndoas (vesicas) radiais a 72$^{\circ}$ fundidas por smin\_exp, NÃO um cos(5v) inventado) lida no corpo topológico: (1) a 5-fold C\ensuremath{_{5}} É o OURO \ensuremath{\sigma}\ensuremath{_{1}}=\ensuremath{\varphi}=(1+\ensuremath{\sqrt{\,}}5)/2, que satisfaz \ensuremath{\varphi}\ensuremath{^{2}}\ensuremath{-}\ensuremath{\varphi}\ensuremath{-}1=0 = a PARÁBOLA GERADORA com m=1 (o metal dourado) — o pentagrama é a figura áurea, a simetria de 5 pontas É o metal m=1 (C\ensuremath{_{5}}-invariante, erro 3e-16). (2) a estrela e a flor são STAR-SHAPED (cada raio do centro cruza a borda 1 vez) \ensuremath{\Longrightarrow} genus 0, \ensuremath{\chi}=2. (3) o morph estrela\ensuremath{\to}flor (as pontas afiadas arredondam em pétalas) é o HOMEOMORFISMO (o cristal, \ensuremath{\chi}=2 preservado \ensuremath{\forall}morph, sem mudar a topologia — não é cirurgia); o mesmo regime da esfera\ensuremath{\leftrightarrow}elipsóide. Render: estrela\_flor.gif. linguagem/corpo\_topologico.py (estrela\_flor\_joalheria 3/3); usa shader\_piramide (papers/joalheria.tex).
\prov{relações: usa transformacao\_superficies\_homeo\_cirurgia; usa parabola\_geradora\_idempotente; usa corpo\_topologico}
\prov{proveniência: wiki \textperiodcentered{} Floxina (investigação) \textperiodcentered{} fato \textperiodcentered{} confiança 1.0 \textperiodcentered{} domínio floxina \textperiodcentered{} história 4 versões no jornal}

%CRISTAL {"arestas": [{"alvo": "torre_dois_lados_classico_relativistico", "confianca": 1.0, "epistemico": "fato", "origem": "wiki", "rel": "parte_de"}, {"alvo": "corpo_matricial_matrix", "confianca": 1.0, "epistemico": "fato", "origem": "wiki", "rel": "relaciona"}, {"alvo": "floxina_quantum_materia_escura", "confianca": 1.0, "epistemico": "fato", "origem": "wiki", "rel": "relaciona"}], "confianca": 1.0, "contraexemplos": [], "descricao": "A torre é FINITA — os dois lados terminam. O CLÁSSICO fecha em 8D: ℝ,ℂ,ℍ,𝕆 (dims 1,2,4,8) são as ÚNICAS álgebras de divisão normadas (teorema de HURWITZ); nos sedênios (16D) a divisão QUEBRA — há divisores de zero: (e_1+e_10)·(e_5+e_14)=0 com ambos ≠0, e ‖xy‖≠‖x‖‖y‖. O RELATIVÍSTICO fecha em 7D: o produto vetorial só existe em 3D e 7D (o teorema). Logo o CORPO MATRICIAL COMPLETA-SE em 8D (𝕆, a última divisão) e 7D (o último ×) — a investigação da reta matricial FECHA aqui. O mass gap 𝔤₂=Der(𝕆)=14 (no gambá 7D) é a assinatura do fecho: além dele (16D) o corpo se perde (só sobra a power-associatividade). Verificado: o divisor de zero dos sedênios (reta_mineral_7d._mul em dim 16).", "epistemico": "hipotese", "exemplos": [], "id": "fecho_torre_8d_7d", "memoria": "semantica", "meta": {"dominio": "floxina", "fonte": "Floxina (investigação)"}, "origem": "wiki", "palavras_chave": [], "sinonimos": [], "tipo": "conceito", "titulo": "O fecho da torre: o corpo matricial completa-se em 8D (𝕆) / 7D (o ×)"}
\section{O fecho da torre: o corpo matricial completa-se em 8D (\ensuremath{\mathbb{O}}) / 7D (o $\times$)}
A torre é FINITA — os dois lados terminam. O CLÁSSICO fecha em 8D: \ensuremath{\mathbb{R}},\ensuremath{\mathbb{C}},\ensuremath{\mathbb{H}},\ensuremath{\mathbb{O}} (dims 1,2,4,8) são as ÚNICAS álgebras de divisão normadas (teorema de HURWITZ); nos sedênios (16D) a divisão QUEBRA — há divisores de zero: (e\_1+e\_10)\textperiodcentered (e\_5+e\_14)=0 com ambos \ensuremath{\ne}0, e \ensuremath{\|}xy\ensuremath{\|}\ensuremath{\ne}\ensuremath{\|}x\ensuremath{\|}\ensuremath{\|}y\ensuremath{\|}. O RELATIVÍSTICO fecha em 7D: o produto vetorial só existe em 3D e 7D (o teorema). Logo o CORPO MATRICIAL COMPLETA-SE em 8D (\ensuremath{\mathbb{O}}, a última divisão) e 7D (o último $\times$) — a investigação da reta matricial FECHA aqui. O mass gap \ensuremath{\mathfrak{g}}\ensuremath{_{2}}=Der(\ensuremath{\mathbb{O}})=14 (no gambá 7D) é a assinatura do fecho: além dele (16D) o corpo se perde (só sobra a power-associatividade). Verificado: o divisor de zero dos sedênios (reta\_mineral\_7d.\_mul em dim 16).
\prov{relações: parte\_de torre\_dois\_lados\_classico\_relativistico; relaciona corpo\_matricial\_matrix; relaciona floxina\_quantum\_materia\_escura}
\prov{proveniência: wiki \textperiodcentered{} Floxina (investigação) \textperiodcentered{} hipotese \textperiodcentered{} confiança 1.0 \textperiodcentered{} domínio floxina \textperiodcentered{} história 264 versões no jornal}

%CRISTAL {"arestas": [{"alvo": "corpo_matricial_matrix", "confianca": 1.0, "epistemico": "fato", "origem": "wiki", "rel": "parte_de"}, {"alvo": "gamba_antissimetrico", "confianca": 1.0, "epistemico": "fato", "origem": "wiki", "rel": "usa"}, {"alvo": "chicote_emocional", "confianca": 1.0, "epistemico": "fato", "origem": "wiki", "rel": "relaciona"}, {"alvo": "biblioteca_pipe", "confianca": 1.0, "epistemico": "fato", "origem": "wiki", "rel": "relaciona"}], "confianca": 1.0, "contraexemplos": [], "descricao": "O corpo matricial lê-se como física, e a MECÂNICA entrega os n-QUBITS. GEOMETRIA: M_n=Sym(gato)⊕Skew(gambá), a métrica é a norma N=det (o gato Lorentz det=−1, o cone; o gambá elíptico det=+1, a área). MECÂNICA: o gambá é a estrutura simplética (Hamilton ẋ=G∇H, Liouville, o fluxo unitário) — e ela DÁ os n-QUBITS: o lado CLÁSSICO 2^n É n qubits (ℂ=2¹=1q, ℍ=2²=2q, 𝕆=2³=3q); o estado |ψ⟩∈ℝ^{2^n} com |ψ|²=1 (o checksum a+c²=1), e o gambá (multiplicação por unitário) EVOLUI |ψ⟩ preservando |ψ|²=1 (a evolução unitária) — a torre clássica é a régua dos qubits, o gambá é o Hamiltoniano (o mesmo do coração: o octônion=3 qubits, o chicote). TERMODINÂMICA: Onsager — o gato dissipa (Ṡ>0), o gambá conserva (Ṡ=0). Verificado: conversa/gamba.py (fisica_matrix, 3/3).", "epistemico": "hipotese", "exemplos": [], "id": "fisica_corpo_matricial_n_qubits", "memoria": "semantica", "meta": {"dominio": "floxina", "fonte": "Floxina (investigação)"}, "origem": "wiki", "palavras_chave": [], "sinonimos": [], "tipo": "conceito", "titulo": "A física do corpo matricial: a mecânica dá os n-qubits"}
\section{A física do corpo matricial: a mecânica dá os n-qubits}
O corpo matricial lê-se como física, e a MECÂNICA entrega os n-QUBITS. GEOMETRIA: M\_n=Sym(gato)\ensuremath{\oplus}Skew(gambá), a métrica é a norma N=det (o gato Lorentz det=\ensuremath{-}1, o cone; o gambá elíptico det=+1, a área). MECÂNICA: o gambá é a estrutura simplética (Hamilton \ensuremath{\dot{x}}=G\ensuremath{\nabla}H, Liouville, o fluxo unitário) — e ela DÁ os n-QUBITS: o lado CLÁSSICO 2\^{}n É n qubits (\ensuremath{\mathbb{C}}=2\ensuremath{^{1}}=1q, \ensuremath{\mathbb{H}}=2\ensuremath{^{2}}=2q, \ensuremath{\mathbb{O}}=2\ensuremath{^{3}}=3q); o estado |\ensuremath{\psi}\ensuremath{\rangle}\ensuremath{\in}\ensuremath{\mathbb{R}}\^{}\{2\^{}n\} com |\ensuremath{\psi}|\ensuremath{^{2}}=1 (o checksum a+c\ensuremath{^{2}}=1), e o gambá (multiplicação por unitário) EVOLUI |\ensuremath{\psi}\ensuremath{\rangle} preservando |\ensuremath{\psi}|\ensuremath{^{2}}=1 (a evolução unitária) — a torre clássica é a régua dos qubits, o gambá é o Hamiltoniano (o mesmo do coração: o octônion=3 qubits, o chicote). TERMODINÂMICA: Onsager — o gato dissipa (\ensuremath{\dot{S}}>0), o gambá conserva (\ensuremath{\dot{S}}=0). Verificado: conversa/gamba.py (fisica\_matrix, 3/3).
\prov{relações: parte\_de corpo\_matricial\_matrix; usa gamba\_antissimetrico; relaciona chicote\_emocional; relaciona biblioteca\_pipe}
\prov{proveniência: wiki \textperiodcentered{} Floxina (investigação) \textperiodcentered{} hipotese \textperiodcentered{} confiança 1.0 \textperiodcentered{} domínio floxina \textperiodcentered{} história 325 versões no jornal}

%CRISTAL {"arestas": [{"alvo": "omnitrix_corpo_universal", "confianca": 1.0, "epistemico": "fato", "origem": "wiki", "rel": "parte_de"}, {"alvo": "multifractal_newton_omnitrix", "confianca": 1.0, "epistemico": "fato", "origem": "wiki", "rel": "usa"}, {"alvo": "reta_matricial_compacta_cf", "confianca": 1.0, "epistemico": "fato", "origem": "wiki", "rel": "usa"}, {"alvo": "o_gato_e_um_mapa_de_anosov", "confianca": 1.0, "epistemico": "fato", "origem": "wiki", "rel": "usa"}, {"alvo": "gato", "confianca": 1.0, "epistemico": "fato", "origem": "wiki", "rel": "usa"}, {"alvo": "gerador_do_gap_comutador", "confianca": 1.0, "epistemico": "fato", "origem": "wiki", "rel": "relaciona"}], "confianca": 1.0, "contraexemplos": [], "descricao": "As três físicas no nível do Corpo Algébrico, costurando o que já se construiu. GEOMETRIA (o gato=a métrica): o espaço das instâncias é a SUPERFÍCIE MODULAR (a CF de π, GL(2,ℤ)\\H); o gato é HIPERBÓLICO (Anosov): A²=[[2,1],[1,1]]∈SL(2,ℤ), traço 3>2; a métrica é hiperbólica (K=−1); as geodésicas são a CF/Gauss. MECÂNICA (os n-qubits, o fluxo geodésico): a dinâmica é o FLUXO GEODÉSICO na modular=a CF=Anosov (caótico, Lyapunov λ=log φ²≈0,96>0); o Hamiltoniano é o gato; os n-qubits são a torre (ℂ=1,ℍ=2,𝕆=3, 2ⁿ=dim). TERMODINÂMICA (o gap→o formalismo termodinâmico do multifractal): o multifractal É termodinâmica. A medida binomial dos DOIS POÇOS (matéria p₀=1/φ, antimatéria p₁=1/φ²) tem τ(q)=−log₂(p₀^q+p₁^q); dicionário exato: q=a TEMPERATURA (inversa), τ(q)=a ENERGIA LIVRE, f(α)=a ENTROPIA, por LEGENDRE f(α)=qα−τ(q). O gato dissipa→f(α)≥0; D_q=τ(q)/(q−1) decresce (D₀=1>D₁≈0,96>D∞≈0,69)=o multifractal. linguagem/omnitrix.py (fisica_universal, 3/3). Síntese declarada.", "epistemico": "hipotese", "exemplos": [], "id": "fisica_universal_omnitrix", "memoria": "semantica", "meta": {"dominio": "floxina", "fonte": "Floxina (investigação)"}, "origem": "wiki", "palavras_chave": [], "sinonimos": [], "tipo": "conceito", "titulo": "A física universal: geometria (a modular), mecânica (o fluxo geodésico), termodinâmica (o multifractal)"}
\section{A física universal: geometria (a modular), mecânica (o fluxo geodésico), termodinâmica (o multifractal)}
As três físicas no nível do Corpo Algébrico, costurando o que já se construiu. GEOMETRIA (o gato=a métrica): o espaço das instâncias é a SUPERFÍCIE MODULAR (a CF de \ensuremath{\pi}, GL(2,\ensuremath{\mathbb{Z}})\textbackslash{}H); o gato é HIPERBÓLICO (Anosov): A\ensuremath{^{2}}=[[2,1],[1,1]]\ensuremath{\in}SL(2,\ensuremath{\mathbb{Z}}), traço 3>2; a métrica é hiperbólica (K=\ensuremath{-}1); as geodésicas são a CF/Gauss. MECÂNICA (os n-qubits, o fluxo geodésico): a dinâmica é o FLUXO GEODÉSICO na modular=a CF=Anosov (caótico, Lyapunov \ensuremath{\lambda}=log \ensuremath{\varphi}\ensuremath{^{2}}\ensuremath{\approx}0,96>0); o Hamiltoniano é o gato; os n-qubits são a torre (\ensuremath{\mathbb{C}}=1,\ensuremath{\mathbb{H}}=2,\ensuremath{\mathbb{O}}=3, 2\ensuremath{^{n}}=dim). TERMODINÂMICA (o gap\ensuremath{\to}o formalismo termodinâmico do multifractal): o multifractal É termodinâmica. A medida binomial dos DOIS POÇOS (matéria p\ensuremath{_{0}}=1/\ensuremath{\varphi}, antimatéria p\ensuremath{_{1}}=1/\ensuremath{\varphi}\ensuremath{^{2}}) tem \ensuremath{\tau}(q)=\ensuremath{-}log\ensuremath{_{2}}(p\ensuremath{_{0}}\^{}q+p\ensuremath{_{1}}\^{}q); dicionário exato: q=a TEMPERATURA (inversa), \ensuremath{\tau}(q)=a ENERGIA LIVRE, f(\ensuremath{\alpha})=a ENTROPIA, por LEGENDRE f(\ensuremath{\alpha})=q\ensuremath{\alpha}\ensuremath{-}\ensuremath{\tau}(q). O gato dissipa\ensuremath{\to}f(\ensuremath{\alpha})\ensuremath{\ge}0; D\_q=\ensuremath{\tau}(q)/(q\ensuremath{-}1) decresce (D\ensuremath{_{0}}=1>D\ensuremath{_{1}}\ensuremath{\approx}0,96>D\ensuremath{\infty}\ensuremath{\approx}0,69)=o multifractal. linguagem/omnitrix.py (fisica\_universal, 3/3). Síntese declarada.
\prov{relações: parte\_de omnitrix\_corpo\_universal; usa multifractal\_newton\_omnitrix; usa reta\_matricial\_compacta\_cf; usa o\_gato\_e\_um\_mapa\_de\_anosov; usa gato; relaciona gerador\_do\_gap\_comutador}
\prov{proveniência: wiki \textperiodcentered{} Floxina (investigação) \textperiodcentered{} hipotese \textperiodcentered{} confiança 1.0 \textperiodcentered{} domínio floxina \textperiodcentered{} história 357 versões no jornal}

%CRISTAL {"arestas": [{"alvo": "floxina_wick_difusao", "confianca": 1.0, "epistemico": "fato", "origem": "wiki", "rel": "parte_de"}, {"alvo": "dicionario_floxina_materia_escura", "confianca": 1.0, "epistemico": "fato", "origem": "wiki", "rel": "parte_de"}, {"alvo": "dicionario_floxinato_gap_massa", "confianca": 1.0, "epistemico": "fato", "origem": "wiki", "rel": "parte_de"}, {"alvo": "gamba_litografico_assinatura", "confianca": 1.0, "epistemico": "fato", "origem": "wiki", "rel": "parte_de"}, {"alvo": "selberg_meta_inducao", "confianca": 1.0, "epistemico": "fato", "origem": "wiki", "rel": "parte_de"}, {"alvo": "gamba_lado_frio_clones", "confianca": 1.0, "epistemico": "fato", "origem": "wiki", "rel": "parte_de"}, {"alvo": "pipe_transmissao_dual", "confianca": 1.0, "epistemico": "fato", "origem": "wiki", "rel": "parte_de"}, {"alvo": "agulha_perelman_limite", "confianca": 1.0, "epistemico": "fato", "origem": "wiki", "rel": "parte_de"}, {"alvo": "estojo_agulhas_minerais_fixo", "confianca": 1.0, "epistemico": "fato", "origem": "wiki", "rel": "parte_de"}, {"alvo": "canais_chave_assinatura", "confianca": 1.0, "epistemico": "fato", "origem": "wiki", "rel": "parte_de"}, {"alvo": "dicionario_litografia_dopagem", "confianca": 1.0, "epistemico": "fato", "origem": "wiki", "rel": "parte_de"}, {"alvo": "floxina_quantum_materia_escura", "confianca": 1.0, "epistemico": "fato", "origem": "wiki", "rel": "parte_de"}, {"alvo": "gerador_do_gap_comutador", "confianca": 1.0, "epistemico": "fato", "origem": "wiki", "rel": "parte_de"}, {"alvo": "algebra_fundamental_gato_gamba", "confianca": 1.0, "epistemico": "fato", "origem": "wiki", "rel": "parte_de"}, {"alvo": "parabola_geradora_idempotente", "confianca": 1.0, "epistemico": "fato", "origem": "wiki", "rel": "parte_de"}, {"alvo": "corpo_matricial_matrix", "confianca": 1.0, "epistemico": "fato", "origem": "wiki", "rel": "parte_de"}, {"alvo": "reta_matricial_compacta_cf", "confianca": 1.0, "epistemico": "fato", "origem": "wiki", "rel": "parte_de"}, {"alvo": "fractal_newton_reta_matricial", "confianca": 1.0, "epistemico": "fato", "origem": "wiki", "rel": "parte_de"}, {"alvo": "leminiscata_grupo_de_grupos", "confianca": 1.0, "epistemico": "fato", "origem": "wiki", "rel": "parte_de"}, {"alvo": "torre_dois_lados_classico_relativistico", "confianca": 1.0, "epistemico": "fato", "origem": "wiki", "rel": "parte_de"}, {"alvo": "fecho_torre_8d_7d", "confianca": 1.0, "epistemico": "fato", "origem": "wiki", "rel": "parte_de"}, {"alvo": "fisica_corpo_matricial_n_qubits", "confianca": 1.0, "epistemico": "fato", "origem": "wiki", "rel": "parte_de"}, {"alvo": "campos_gravidade_gato_em_gamba", "confianca": 1.0, "epistemico": "fato", "origem": "wiki", "rel": "parte_de"}, {"alvo": "teste_molecula_tiol_gato_n_qubit", "confianca": 1.0, "epistemico": "fato", "origem": "wiki", "rel": "parte_de"}, {"alvo": "dicionario_tiol_gato_gamba", "confianca": 1.0, "epistemico": "fato", "origem": "wiki", "rel": "parte_de"}, {"alvo": "thc_tiol_propriedades_compartilhadas", "confianca": 1.0, "epistemico": "fato", "origem": "wiki", "rel": "parte_de"}, {"alvo": "serie_cheiro_gap_gamba", "confianca": 1.0, "epistemico": "fato", "origem": "wiki", "rel": "parte_de"}, {"alvo": "geraniol_citronelol_gamba_reduzido", "confianca": 1.0, "epistemico": "fato", "origem": "wiki", "rel": "parte_de"}, {"alvo": "ponte_hidrogenio_parabola_geradora", "confianca": 1.0, "epistemico": "fato", "origem": "wiki", "rel": "parte_de"}, {"alvo": "invariante_molecular_residuo", "confianca": 1.0, "epistemico": "fato", "origem": "wiki", "rel": "parte_de"}, {"alvo": "dicionario_matrix_bai", "confianca": 1.0, "epistemico": "fato", "origem": "wiki", "rel": "parte_de"}, {"alvo": "volume_matrizes_3d", "confianca": 1.0, "epistemico": "fato", "origem": "wiki", "rel": "parte_de"}, {"alvo": "gamba_antissimetrico", "confianca": 1.0, "epistemico": "fato", "origem": "wiki", "rel": "parte_de"}, {"alvo": "gamba_quatro_fisicas", "confianca": 1.0, "epistemico": "fato", "origem": "wiki", "rel": "parte_de"}, {"alvo": "quatro_quadrantes_corpo_estelar", "confianca": 1.0, "epistemico": "fato", "origem": "wiki", "rel": "parte_de"}, {"alvo": "floxina_selberg", "confianca": 1.0, "epistemico": "fato", "origem": "wiki", "rel": "parte_de"}, {"alvo": "bootstrap_semantico", "confianca": 1.0, "epistemico": "fato", "origem": "wiki", "rel": "relaciona"}, {"alvo": "teste_floxinato_gamba", "confianca": 1.0, "epistemico": "fato", "origem": "wiki", "rel": "relaciona"}, {"alvo": "litografia", "confianca": 1.0, "epistemico": "fato", "origem": "wiki", "rel": "relaciona"}], "confianca": 1.0, "contraexemplos": [], "descricao": "O paper papers/Floxina.tex abre a investigação nascida do bootstrap do «floxinato de gambá»: (I) o gambá anti-simétrico ao gato; (II) o invariante anti-simétrico abrindo álgebra/geometria/mecânica/termodinâmica; (III) a litografia e o bootstrap 0→1; (IV) a ponte com Selberg; (V) as perguntas em aberto (o gambá metálico em 2m-D; o gap é mesmo vazio?; Selberg=gato+gambá?; o bootstrap é o único cruzamento?). Escopo honesto: I–II fato/síntese, III analogia, IV–V investigação.", "epistemico": "hipotese", "exemplos": [], "id": "floxina_investigacao", "memoria": "semantica", "meta": {"dominio": "floxina", "fonte": "Floxina (investigação)"}, "origem": "wiki", "palavras_chave": [], "sinonimos": [], "tipo": "conceito", "titulo": "Floxina (a ponta de investigação)"}
\section{Floxina (a ponta de investigação)}
O paper papers/Floxina.tex abre a investigação nascida do bootstrap do «floxinato de gambá»: (I) o gambá anti-simétrico ao gato; (II) o invariante anti-simétrico abrindo álgebra/geometria/mecânica/termodinâmica; (III) a litografia e o bootstrap 0\ensuremath{\to}1; (IV) a ponte com Selberg; (V) as perguntas em aberto (o gambá metálico em 2m-D; o gap é mesmo vazio?; Selberg=gato+gambá?; o bootstrap é o único cruzamento?). Escopo honesto: I–II fato/síntese, III analogia, IV–V investigação.
\prov{relações: parte\_de floxina\_wick\_difusao; parte\_de dicionario\_floxina\_materia\_escura; parte\_de dicionario\_floxinato\_gap\_massa; parte\_de gamba\_litografico\_assinatura; parte\_de selberg\_meta\_inducao; parte\_de gamba\_lado\_frio\_clones; parte\_de pipe\_transmissao\_dual; parte\_de agulha\_perelman\_limite; parte\_de estojo\_agulhas\_minerais\_fixo; parte\_de canais\_chave\_assinatura; parte\_de dicionario\_litografia\_dopagem; parte\_de floxina\_quantum\_materia\_escura; parte\_de gerador\_do\_gap\_comutador; parte\_de algebra\_fundamental\_gato\_gamba; parte\_de parabola\_geradora\_idempotente; parte\_de corpo\_matricial\_matrix; parte\_de reta\_matricial\_compacta\_cf; parte\_de fractal\_newton\_reta\_matricial; parte\_de leminiscata\_grupo\_de\_grupos; parte\_de torre\_dois\_lados\_classico\_relativistico; parte\_de fecho\_torre\_8d\_7d; parte\_de fisica\_corpo\_matricial\_n\_qubits; parte\_de campos\_gravidade\_gato\_em\_gamba; parte\_de teste\_molecula\_tiol\_gato\_n\_qubit; parte\_de dicionario\_tiol\_gato\_gamba; parte\_de thc\_tiol\_propriedades\_compartilhadas; parte\_de serie\_cheiro\_gap\_gamba; parte\_de geraniol\_citronelol\_gamba\_reduzido; parte\_de ponte\_hidrogenio\_parabola\_geradora; parte\_de invariante\_molecular\_residuo; parte\_de dicionario\_matrix\_bai; parte\_de volume\_matrizes\_3d; parte\_de gamba\_antissimetrico; parte\_de gamba\_quatro\_fisicas; parte\_de quatro\_quadrantes\_corpo\_estelar; parte\_de floxina\_selberg; relaciona bootstrap\_semantico; relaciona teste\_floxinato\_gamba; relaciona litografia}
\prov{proveniência: wiki \textperiodcentered{} Floxina (investigação) \textperiodcentered{} hipotese \textperiodcentered{} confiança 1.0 \textperiodcentered{} domínio floxina \textperiodcentered{} história 2908 versões no jornal}

%CRISTAL {"arestas": [{"alvo": "nao_associatividade_amputacao", "confianca": 1.0, "epistemico": "fato", "origem": "wiki", "rel": "usa"}, {"alvo": "g2_grupo_excepcional", "confianca": 1.0, "epistemico": "fato", "origem": "wiki", "rel": "usa"}, {"alvo": "mass_gap_algebrico", "confianca": 1.0, "epistemico": "fato", "origem": "wiki", "rel": "relaciona"}, {"alvo": "dicionario_floxinato_gap_massa", "confianca": 1.0, "epistemico": "fato", "origem": "wiki", "rel": "relaciona"}, {"alvo": "dicionario_floxina_materia_escura", "confianca": 1.0, "epistemico": "fato", "origem": "wiki", "rel": "relaciona"}, {"alvo": "materia_escura", "confianca": 1.0, "epistemico": "fato", "origem": "wiki", "rel": "relaciona"}, {"alvo": "gamba_lado_frio_clones", "confianca": 1.0, "epistemico": "fato", "origem": "wiki", "rel": "relaciona"}], "confianca": 1.0, "contraexemplos": [], "descricao": "A cadeia que dá corpo à floxina como PARTÍCULA: (1) o ASSOCIADOR octônio [x,y,z]=(xy)z−x(yz) é ≠0 no 𝕆 (zero até ℍ), e o seu resíduo é 𝔤₂=Der(𝕆)=14 — a não-associatividade amputada, o resto que NUNCA zera (a≥a_min>0); (2) esse resíduo É o GAP DE MASSA (Yang-Mills Δ>0; a desafasagem; ¼ em Selberg; o glueball mineral ω=√2); (3) o QUANTUM desse gap — a excitação mínima, a partícula do 𝔤₂ — é a FLOXINA. E como a metade NÃO-DISSIPATIVA (o gambá, reversível, o lado frio) é a MATÉRIA ESCURA (pesa mas não irradia, Rubin), a floxina é o QUANTUM DE MATÉRIA ESCURA: a partícula do gap que nunca fecha, não-EM (o gambá), não-colisional. associador → gap de massa → floxina (o quantum escuro). FATO: [·]≠0 só no 𝕆, 𝔤₂=14; SÍNTESE declarada: o quantum do gap = a floxina = o quantum de matéria escura.", "epistemico": "hipotese", "exemplos": [], "id": "floxina_quantum_materia_escura", "memoria": "semantica", "meta": {"dominio": "floxina", "fonte": "Floxina (investigação)"}, "origem": "wiki", "palavras_chave": [], "sinonimos": [], "tipo": "conceito", "titulo": "A floxina é o quantum de matéria escura (associador→gap→quantum)"}
\section{A floxina é o quantum de matéria escura (associador\ensuremath{\to}gap\ensuremath{\to}quantum)}
A cadeia que dá corpo à floxina como PARTÍCULA: (1) o ASSOCIADOR octônio [x,y,z]=(xy)z\ensuremath{-}x(yz) é \ensuremath{\ne}0 no \ensuremath{\mathbb{O}} (zero até \ensuremath{\mathbb{H}}), e o seu resíduo é \ensuremath{\mathfrak{g}}\ensuremath{_{2}}=Der(\ensuremath{\mathbb{O}})=14 — a não-associatividade amputada, o resto que NUNCA zera (a\ensuremath{\ge}a\_min>0); (2) esse resíduo É o GAP DE MASSA (Yang-Mills \ensuremath{\Delta}>0; a desafasagem; \ensuremath{\tfrac14} em Selberg; o glueball mineral \ensuremath{\omega}=\ensuremath{\sqrt{\,}}2); (3) o QUANTUM desse gap — a excitação mínima, a partícula do \ensuremath{\mathfrak{g}}\ensuremath{_{2}} — é a FLOXINA. E como a metade NÃO-DISSIPATIVA (o gambá, reversível, o lado frio) é a MATÉRIA ESCURA (pesa mas não irradia, Rubin), a floxina é o QUANTUM DE MATÉRIA ESCURA: a partícula do gap que nunca fecha, não-EM (o gambá), não-colisional. associador \ensuremath{\to} gap de massa \ensuremath{\to} floxina (o quantum escuro). FATO: [\textperiodcentered ]\ensuremath{\ne}0 só no \ensuremath{\mathbb{O}}, \ensuremath{\mathfrak{g}}\ensuremath{_{2}}=14; SÍNTESE declarada: o quantum do gap = a floxina = o quantum de matéria escura.
\prov{relações: usa nao\_associatividade\_amputacao; usa g2\_grupo\_excepcional; relaciona mass\_gap\_algebrico; relaciona dicionario\_floxinato\_gap\_massa; relaciona dicionario\_floxina\_materia\_escura; relaciona materia\_escura; relaciona gamba\_lado\_frio\_clones}
\prov{proveniência: wiki \textperiodcentered{} Floxina (investigação) \textperiodcentered{} hipotese \textperiodcentered{} confiança 1.0 \textperiodcentered{} domínio floxina \textperiodcentered{} história 600 versões no jornal}

%CRISTAL {"arestas": [{"alvo": "gamba_antissimetrico", "confianca": 1.0, "epistemico": "fato", "origem": "wiki", "rel": "usa"}, {"alvo": "gato", "confianca": 1.0, "epistemico": "fato", "origem": "wiki", "rel": "usa"}, {"alvo": "ancora_selberg_e_seu_escopo", "confianca": 1.0, "epistemico": "fato", "origem": "wiki", "rel": "relaciona"}, {"alvo": "crivo_de_selberg", "confianca": 1.0, "epistemico": "fato", "origem": "wiki", "rel": "relaciona"}, {"alvo": "gap_espectral", "confianca": 1.0, "epistemico": "fato", "origem": "wiki", "rel": "relaciona"}, {"alvo": "litografia", "confianca": 1.0, "epistemico": "fato", "origem": "wiki", "rel": "relaciona"}, {"alvo": "invariante_topologico_assinatura", "confianca": 1.0, "epistemico": "fato", "origem": "wiki", "rel": "relaciona"}], "confianca": 1.0, "contraexemplos": [], "descricao": "Na fórmula do traço de Selberg em Γ\\ℍ, as classes de conjugação partem-se nos nossos dois operadores: as HIPERBÓLICAS (o gato — as geodésicas fechadas, o comprimento-espectro, os primos) e as ELÍPTICAS (o gambá — as rotações de ordem finita, a torção). E o GAP: os autovalores λ≥¼ dão os zeros na linha Re(s)=½ (RH provada, Selberg 1956); os excepcionais λ<¼ dariam zeros fora. A CONJECTURA DE VALOR PRÓPRIO DE SELBERG: para grupos de congruência, λ₁≥¼ — nada mora no gap. Analogia (declarada): o grau da litografia é DISCRETO (0 ou 1, não há meia-assinatura) = não há autovalor no gap; a forja (grau 0) = o excepcional; o bootstrap (registrar ⟹ 0→1) = cruzar o gap. Nenhuma assinatura se auto-grava a meio caminho: precisa do cunho (a preimagem).", "epistemico": "hipotese", "exemplos": [], "id": "floxina_selberg", "memoria": "semantica", "meta": {"dominio": "floxina", "fonte": "Floxina (investigação)"}, "origem": "wiki", "palavras_chave": [], "sinonimos": [], "tipo": "conceito", "titulo": "A ponte com a conjectura de Selberg (o gap e o grau discreto)"}
\section{A ponte com a conjectura de Selberg (o gap e o grau discreto)}
Na fórmula do traço de Selberg em \ensuremath{\Gamma}\textbackslash{}\ensuremath{\mathbb{H}}, as classes de conjugação partem-se nos nossos dois operadores: as HIPERBÓLICAS (o gato — as geodésicas fechadas, o comprimento-espectro, os primos) e as ELÍPTICAS (o gambá — as rotações de ordem finita, a torção). E o GAP: os autovalores \ensuremath{\lambda}\ensuremath{\ge}\ensuremath{\tfrac14} dão os zeros na linha Re(s)=\ensuremath{\tfrac12} (RH provada, Selberg 1956); os excepcionais \ensuremath{\lambda}<\ensuremath{\tfrac14} dariam zeros fora. A CONJECTURA DE VALOR PRÓPRIO DE SELBERG: para grupos de congruência, \ensuremath{\lambda}\ensuremath{_{1}}\ensuremath{\ge}\ensuremath{\tfrac14} — nada mora no gap. Analogia (declarada): o grau da litografia é DISCRETO (0 ou 1, não há meia-assinatura) = não há autovalor no gap; a forja (grau 0) = o excepcional; o bootstrap (registrar \ensuremath{\Longrightarrow} 0\ensuremath{\to}1) = cruzar o gap. Nenhuma assinatura se auto-grava a meio caminho: precisa do cunho (a preimagem).
\prov{relações: usa gamba\_antissimetrico; usa gato; relaciona ancora\_selberg\_e\_seu\_escopo; relaciona crivo\_de\_selberg; relaciona gap\_espectral; relaciona litografia; relaciona invariante\_topologico\_assinatura}
\prov{proveniência: wiki \textperiodcentered{} Floxina (investigação) \textperiodcentered{} hipotese \textperiodcentered{} confiança 1.0 \textperiodcentered{} domínio floxina \textperiodcentered{} história 656 versões no jornal}

%CRISTAL {"arestas": [{"alvo": "gamba_antissimetrico", "confianca": 1.0, "epistemico": "fato", "origem": "wiki", "rel": "usa"}, {"alvo": "flip_e_rotacao_de_wick", "confianca": 1.0, "epistemico": "fato", "origem": "wiki", "rel": "relaciona"}, {"alvo": "gato", "confianca": 1.0, "epistemico": "fato", "origem": "wiki", "rel": "usa"}, {"alvo": "ancora_selberg_e_seu_escopo", "confianca": 1.0, "epistemico": "fato", "origem": "wiki", "rel": "relaciona"}, {"alvo": "litografia", "confianca": 1.0, "epistemico": "fato", "origem": "wiki", "rel": "relaciona"}], "confianca": 1.0, "contraexemplos": [], "descricao": "O flip u=log z tem dois eixos: real λ=log σ (dissipativo, gato) e imaginário θ=arg z (unitário, gambá). A rotação de Wick τ=it troca as duas dinâmicas da floxina: à esquerda ela GIRA — Schrödinger no grafo i∂_tψ=Δψ, ψ=e^{−iΔt}ψ₀, unitário, ‖ψ‖ conservada (o gambá, dispersa sem perda); à direita ela DIFUNDE — o calor ∂_τρ=−Δρ, ρ=e^{−τΔ}ρ₀, dissipa, entropia sobe (o gato, irreversível). O traço fecha: Tr e^{−τΔ}=∑e^{−τλ_n}=∑σ_m^{−τ}=Z(τ) — o núcleo do calor = a partição = a ZETA DOS METAIS = o traço de Selberg. A difusão da floxina É a zeta dos metais. E a forja (grau 0) é um δ num nó ISOLADO (sem preimagem, sem arestas): não difunde pela luz do aboutness — fica ESCURA; o bootstrap conecta o nó, aí difunde. conversa/gamba.py (difusao/difusao_wick).", "epistemico": "hipotese", "exemplos": [], "id": "floxina_wick_difusao", "memoria": "semantica", "meta": {"dominio": "floxina", "fonte": "Floxina (investigação)"}, "origem": "wiki", "palavras_chave": [], "sinonimos": [], "tipo": "conceito", "titulo": "O flip de Wick da floxina (a difusão = a zeta dos metais)"}
\section{O flip de Wick da floxina (a difusão = a zeta dos metais)}
O flip u=log z tem dois eixos: real \ensuremath{\lambda}=log \ensuremath{\sigma} (dissipativo, gato) e imaginário \ensuremath{\theta}=arg z (unitário, gambá). A rotação de Wick \ensuremath{\tau}=it troca as duas dinâmicas da floxina: à esquerda ela GIRA — Schrödinger no grafo i\ensuremath{\partial}\_t\ensuremath{\psi}=\ensuremath{\Delta}\ensuremath{\psi}, \ensuremath{\psi}=e\^{}\{\ensuremath{-}i\ensuremath{\Delta}t\}\ensuremath{\psi}\ensuremath{_{0}}, unitário, \ensuremath{\|}\ensuremath{\psi}\ensuremath{\|} conservada (o gambá, dispersa sem perda); à direita ela DIFUNDE — o calor \ensuremath{\partial}\_\ensuremath{\tau}\ensuremath{\rho}=\ensuremath{-}\ensuremath{\Delta}\ensuremath{\rho}, \ensuremath{\rho}=e\^{}\{\ensuremath{-}\ensuremath{\tau}\ensuremath{\Delta}\}\ensuremath{\rho}\ensuremath{_{0}}, dissipa, entropia sobe (o gato, irreversível). O traço fecha: Tr e\^{}\{\ensuremath{-}\ensuremath{\tau}\ensuremath{\Delta}\}=\ensuremath{\sum}e\^{}\{\ensuremath{-}\ensuremath{\tau}\ensuremath{\lambda}\_n\}=\ensuremath{\sum}\ensuremath{\sigma}\_m\^{}\{\ensuremath{-}\ensuremath{\tau}\}=Z(\ensuremath{\tau}) — o núcleo do calor = a partição = a ZETA DOS METAIS = o traço de Selberg. A difusão da floxina É a zeta dos metais. E a forja (grau 0) é um \ensuremath{\delta} num nó ISOLADO (sem preimagem, sem arestas): não difunde pela luz do aboutness — fica ESCURA; o bootstrap conecta o nó, aí difunde. conversa/gamba.py (difusao/difusao\_wick).
\prov{relações: usa gamba\_antissimetrico; relaciona flip\_e\_rotacao\_de\_wick; usa gato; relaciona ancora\_selberg\_e\_seu\_escopo; relaciona litografia}
\prov{proveniência: wiki \textperiodcentered{} Floxina (investigação) \textperiodcentered{} hipotese \textperiodcentered{} confiança 1.0 \textperiodcentered{} domínio floxina \textperiodcentered{} história 486 versões no jornal}

%CRISTAL {"arestas": [{"alvo": "corpo_edp", "confianca": 1.0, "epistemico": "fato", "origem": "wiki", "rel": "usa"}, {"alvo": "transformada_dourada_espectro_edp", "confianca": 1.0, "epistemico": "fato", "origem": "wiki", "rel": "usa"}, {"alvo": "solucao_explicita_ed", "confianca": 1.0, "epistemico": "fato", "origem": "wiki", "rel": "relaciona"}, {"alvo": "edp_classica_2_4_8", "confianca": 1.0, "epistemico": "fato", "origem": "wiki", "rel": "relaciona"}], "confianca": 1.0, "contraexemplos": [], "descricao": "A solução fechada de toda EDP de 2ª ordem, dados de Cauchy u₀ (posição) e v₀ (velocidade), é uma FUNÇÃO DO OPERADOR √(−δΔ) (a raiz de Dirac), com o invariante δ escolhendo o NÚCLEO: u(·,t)=cos(t√(−δΔ))·u₀ + [sin(t√(−δΔ))/√(−δΔ)]·v₀ (no Fourier: û(k,t)=cos(ωt)û₀+[sin(ωt)/ω]v̂₀, ω=√δ·|k|=a dispersão=a parábola). δ=+1 ONDA (d'ALEMBERT fechada: u(x,t)=½[u₀(x+t)+u₀(x−t)]+½∫_{x−t}^{x+t}v₀(s)ds, as características x±t=os autovetores do gato); δ=0 CALOR (NÚCLEO DO CALOR: u=∫G_t(x−y)u₀dy=e^{tΔ}u₀, G_t=e^{−x²/4t}/√(4πt), o semigrupo que dissipa); δ=−1 LAPLACE (NÚCLEO DE POISSON: u=∫P_y(x−s)u₀ds=e^{−y√(−Δ)}u₀, P_y=y/[π(x²+y²)], a extensão harmônica do contorno). Uma fórmula, três núcleos — δ decide. Verificado contra a EDP e os dois dados iniciais (erro~1e-8). É o fecho da solução explícita: a espectral (∫û(k)e^{i(kx−ωt)}) dá os modos, esta dá a forma FECHADA com posição E velocidade. linguagem/corpo_edp.py (formula_explicita_da_solucao 5/5); papers/equacoes_diferenciais.tex (Parte VIII).", "epistemico": "fato", "exemplos": [], "id": "formula_explicita_solucao_edp", "memoria": "semantica", "meta": {"dominio": "floxina", "fonte": "Floxina (investigação)"}, "origem": "wiki", "palavras_chave": [], "sinonimos": [], "tipo": "conceito", "titulo": "A fórmula explícita FECHADA da solução (dados de Cauchy u₀,v₀): u=cos(t√(−δΔ))u₀+[sin(t√(−δΔ))/√(−δΔ)]v₀; δ escolhe o núcleo (d'Alembert/calor/Poisson)"}
\section{A fórmula explícita FECHADA da solução (dados de Cauchy u\ensuremath{_{0}},v\ensuremath{_{0}}): u=cos(t\ensuremath{\sqrt{\,}}(\ensuremath{-}\ensuremath{\delta}\ensuremath{\Delta}))u\ensuremath{_{0}}+[sin(t\ensuremath{\sqrt{\,}}(\ensuremath{-}\ensuremath{\delta}\ensuremath{\Delta}))/\ensuremath{\sqrt{\,}}(\ensuremath{-}\ensuremath{\delta}\ensuremath{\Delta})]v\ensuremath{_{0}}; \ensuremath{\delta} escolhe o núcleo (d'Alembert/calor/Poisson)}
A solução fechada de toda EDP de 2ª ordem, dados de Cauchy u\ensuremath{_{0}} (posição) e v\ensuremath{_{0}} (velocidade), é uma FUNÇÃO DO OPERADOR \ensuremath{\sqrt{\,}}(\ensuremath{-}\ensuremath{\delta}\ensuremath{\Delta}) (a raiz de Dirac), com o invariante \ensuremath{\delta} escolhendo o NÚCLEO: u(\textperiodcentered ,t)=cos(t\ensuremath{\sqrt{\,}}(\ensuremath{-}\ensuremath{\delta}\ensuremath{\Delta}))\textperiodcentered u\ensuremath{_{0}} + [sin(t\ensuremath{\sqrt{\,}}(\ensuremath{-}\ensuremath{\delta}\ensuremath{\Delta}))/\ensuremath{\sqrt{\,}}(\ensuremath{-}\ensuremath{\delta}\ensuremath{\Delta})]\textperiodcentered v\ensuremath{_{0}} (no Fourier: û(k,t)=cos(\ensuremath{\omega}t)û\ensuremath{_{0}}+[sin(\ensuremath{\omega}t)/\ensuremath{\omega}]v\ensuremath{}\ensuremath{_{0}}, \ensuremath{\omega}=\ensuremath{\sqrt{\,}}\ensuremath{\delta}\textperiodcentered |k|=a dispersão=a parábola). \ensuremath{\delta}=+1 ONDA (d'ALEMBERT fechada: u(x,t)=\ensuremath{\tfrac12}[u\ensuremath{_{0}}(x+t)+u\ensuremath{_{0}}(x\ensuremath{-}t)]+\ensuremath{\tfrac12}\ensuremath{\int}\_\{x\ensuremath{-}t\}\^{}\{x+t\}v\ensuremath{_{0}}(s)ds, as características x$\pm$t=os autovetores do gato); \ensuremath{\delta}=0 CALOR (NÚCLEO DO CALOR: u=\ensuremath{\int}G\_t(x\ensuremath{-}y)u\ensuremath{_{0}}dy=e\^{}\{t\ensuremath{\Delta}\}u\ensuremath{_{0}}, G\_t=e\^{}\{\ensuremath{-}x\ensuremath{^{2}}/4t\}/\ensuremath{\sqrt{\,}}(4\ensuremath{\pi}t), o semigrupo que dissipa); \ensuremath{\delta}=\ensuremath{-}1 LAPLACE (NÚCLEO DE POISSON: u=\ensuremath{\int}P\_y(x\ensuremath{-}s)u\ensuremath{_{0}}ds=e\^{}\{\ensuremath{-}y\ensuremath{\sqrt{\,}}(\ensuremath{-}\ensuremath{\Delta})\}u\ensuremath{_{0}}, P\_y=y/[\ensuremath{\pi}(x\ensuremath{^{2}}+y\ensuremath{^{2}})], a extensão harmônica do contorno). Uma fórmula, três núcleos — \ensuremath{\delta} decide. Verificado contra a EDP e os dois dados iniciais (erro\~{}1e-8). É o fecho da solução explícita: a espectral (\ensuremath{\int}û(k)e\^{}\{i(kx\ensuremath{-}\ensuremath{\omega}t)\}) dá os modos, esta dá a forma FECHADA com posição E velocidade. linguagem/corpo\_edp.py (formula\_explicita\_da\_solucao 5/5); papers/equacoes\_diferenciais.tex (Parte VIII).
\prov{relações: usa corpo\_edp; usa transformada\_dourada\_espectro\_edp; relaciona solucao\_explicita\_ed; relaciona edp\_classica\_2\_4\_8}
\prov{proveniência: wiki \textperiodcentered{} Floxina (investigação) \textperiodcentered{} fato \textperiodcentered{} confiança 1.0 \textperiodcentered{} domínio floxina \textperiodcentered{} história 40 versões no jornal}

%CRISTAL {"arestas": [{"alvo": "corpo_equacoes_diferenciais", "confianca": 1.0, "epistemico": "fato", "origem": "wiki", "rel": "usa"}, {"alvo": "parabola_geradora_idempotente", "confianca": 1.0, "epistemico": "fato", "origem": "wiki", "rel": "usa"}, {"alvo": "pa_pg_corpo_diferencial", "confianca": 1.0, "epistemico": "fato", "origem": "wiki", "rel": "usa"}], "confianca": 1.0, "contraexemplos": [], "descricao": "Na matrix, a fração contínua gera os metais: σ_n=[n;n,…] é o PONTO FIXO do Möbius do GATO A_n=[[n,1],[1,0]] (x=n+1/x ⟺ x²−nx−1=0), os convergentes p_k/q_k=∏A_n → σ_n, na base do INVARIANTE |det|=1 (GL(2,ℤ)). No corpo diferencial: (1) o metal σ_n (o ponto fixo) É o AUTOVALOR de A_n — a taxa estacionária da ED ẋ=A_n·x (gerar o número e achar o modo próprio são o mesmo ponto fixo). (2) a EXPONENCIAL e^x (a PG, a ED ẋ=x) tem fração contínua; o PADÉ [k/k] (a diagonal da CF) → e^x. (3) o Padé [k/k] de e^A (a Cayley/Möbius, Crank-Nicolson; [1/1]: (I+A/2)(I−A/2)⁻¹) APROXIMA o fluxo e^{At} de QUALQUER ED linear, erro→0 com a ordem. Como toda solução é analítica, a CF converge a ela: a fração contínua gera os metais E aproxima qualquer equação diferencial (∀A testada, ordem 8 bate e^A a ~1e-6). linguagem/omnitrix.py (fracoes_continuas_ed 5/5); papers/equacoes_diferenciais.tex (Parte VI). Análogo ao fracoes_continuas.py da matrix.", "epistemico": "fato", "exemplos": [], "id": "fracoes_continuas_ed", "memoria": "semantica", "meta": {"dominio": "floxina", "fonte": "Floxina (investigação)"}, "origem": "wiki", "palavras_chave": [], "sinonimos": [], "tipo": "conceito", "titulo": "Frações contínuas no corpo diferencial (base do invariante |det|=1): geram os metais E aproximam qualquer ED (o Padé)"}
\section{Frações contínuas no corpo diferencial (base do invariante |det|=1): geram os metais E aproximam qualquer ED (o Padé)}
Na matrix, a fração contínua gera os metais: \ensuremath{\sigma}\_n=[n;n,…] é o PONTO FIXO do Möbius do GATO A\_n=[[n,1],[1,0]] (x=n+1/x \ensuremath{\Longleftrightarrow} x\ensuremath{^{2}}\ensuremath{-}nx\ensuremath{-}1=0), os convergentes p\_k/q\_k=\ensuremath{\prod}A\_n \ensuremath{\to} \ensuremath{\sigma}\_n, na base do INVARIANTE |det|=1 (GL(2,\ensuremath{\mathbb{Z}})). No corpo diferencial: (1) o metal \ensuremath{\sigma}\_n (o ponto fixo) É o AUTOVALOR de A\_n — a taxa estacionária da ED \ensuremath{\dot{x}}=A\_n\textperiodcentered x (gerar o número e achar o modo próprio são o mesmo ponto fixo). (2) a EXPONENCIAL e\^{}x (a PG, a ED \ensuremath{\dot{x}}=x) tem fração contínua; o PADÉ [k/k] (a diagonal da CF) \ensuremath{\to} e\^{}x. (3) o Padé [k/k] de e\^{}A (a Cayley/Möbius, Crank-Nicolson; [1/1]: (I+A/2)(I\ensuremath{-}A/2)\ensuremath{^{-1}}) APROXIMA o fluxo e\^{}\{At\} de QUALQUER ED linear, erro\ensuremath{\to}0 com a ordem. Como toda solução é analítica, a CF converge a ela: a fração contínua gera os metais E aproxima qualquer equação diferencial (\ensuremath{\forall}A testada, ordem 8 bate e\^{}A a \~{}1e-6). linguagem/omnitrix.py (fracoes\_continuas\_ed 5/5); papers/equacoes\_diferenciais.tex (Parte VI). Análogo ao fracoes\_continuas.py da matrix.
\prov{relações: usa corpo\_equacoes\_diferenciais; usa parabola\_geradora\_idempotente; usa pa\_pg\_corpo\_diferencial}
\prov{proveniência: wiki \textperiodcentered{} Floxina (investigação) \textperiodcentered{} fato \textperiodcentered{} confiança 1.0 \textperiodcentered{} domínio floxina \textperiodcentered{} história 60 versões no jornal}

%CRISTAL {"arestas": [{"alvo": "corpo_topologico", "confianca": 1.0, "epistemico": "fato", "origem": "wiki", "rel": "usa"}, {"alvo": "fracoes_continuas_ed", "confianca": 1.0, "epistemico": "fato", "origem": "wiki", "rel": "relaciona"}, {"alvo": "dissipacao_gato_concentracao_esquilo", "confianca": 1.0, "epistemico": "fato", "origem": "wiki", "rel": "usa"}, {"alvo": "parabola_geradora_idempotente", "confianca": 1.0, "epistemico": "fato", "origem": "wiki", "rel": "usa"}], "confianca": 1.0, "contraexemplos": [], "descricao": "Frações contínuas sobre o INVARIANTE, no corpo topológico. O invariante é |det|=1 (GL(2,ℤ)) — a condição de o gato A_n=[[n,1],[1,0]] ser HOMEOMORFISMO do toro T²=ℝ²/ℤ² (preserva o reticulado ⟹ χ=0). O gato É o MAPA DO GATO DE ARNOLD: hiperbólico (Anosov), um autovalor estica σ_n>1, o outro contrai 1/σ_n. (1) o METAL σ_n=[n;n,n,…] (a fração contínua=produto de gatos) = a DILATAÇÃO do Anosov; (2) a ENTROPIA TOPOLÓGICA h=log σ_n (o esticamento por iteração): σ_n=e^h — o metal É a exponencial da entropia do homeo do toro; (3) os convergentes p_k/q_k→σ_n = a inclinação da FOLHEAÇÃO INSTÁVEL (a direção esticada, irracional); (4) os gatos GERAM o mapping class group MCG(T²)=SL(2,ℤ): A_2=[[2,1],[1,1]]=T_a·T_b (produto de 2 Dehn twists). σ_2=1+√2 (prata), h=log σ_2=0.8814 (consistente com a Transformada Dourada). É o análogo topológico do fracoes_continuas_ed (lá o Padé aproxima o fluxo; aqui a palavra em gatos É a classe de homeo do toro, o metal=a entropia). linguagem/corpo_topologico.py (fracoes_continuas_topologica 5/5); papers/corpo_topologico.tex (Parte VI).", "epistemico": "fato", "exemplos": [], "id": "fracoes_continuas_topologica", "memoria": "semantica", "meta": {"dominio": "floxina", "fonte": "Floxina (investigação)"}, "origem": "wiki", "palavras_chave": [], "sinonimos": [], "tipo": "conceito", "titulo": "Frações contínuas sobre o invariante |det|=1 no corpo topológico: o gato=o cat map de Arnold (homeo Anosov do toro), a CF gera o MCG=SL(2,ℤ), o metal σ_n=e^{entropia}"}
\section{Frações contínuas sobre o invariante |det|=1 no corpo topológico: o gato=o cat map de Arnold (homeo Anosov do toro), a CF gera o MCG=SL(2,\ensuremath{\mathbb{Z}}), o metal \ensuremath{\sigma}\_n=e\^{}\{entropia\}}
Frações contínuas sobre o INVARIANTE, no corpo topológico. O invariante é |det|=1 (GL(2,\ensuremath{\mathbb{Z}})) — a condição de o gato A\_n=[[n,1],[1,0]] ser HOMEOMORFISMO do toro T\ensuremath{^{2}}=\ensuremath{\mathbb{R}}\ensuremath{^{2}}/\ensuremath{\mathbb{Z}}\ensuremath{^{2}} (preserva o reticulado \ensuremath{\Longrightarrow} \ensuremath{\chi}=0). O gato É o MAPA DO GATO DE ARNOLD: hiperbólico (Anosov), um autovalor estica \ensuremath{\sigma}\_n>1, o outro contrai 1/\ensuremath{\sigma}\_n. (1) o METAL \ensuremath{\sigma}\_n=[n;n,n,…] (a fração contínua=produto de gatos) = a DILATAÇÃO do Anosov; (2) a ENTROPIA TOPOLÓGICA h=log \ensuremath{\sigma}\_n (o esticamento por iteração): \ensuremath{\sigma}\_n=e\^{}h — o metal É a exponencial da entropia do homeo do toro; (3) os convergentes p\_k/q\_k\ensuremath{\to}\ensuremath{\sigma}\_n = a inclinação da FOLHEAÇÃO INSTÁVEL (a direção esticada, irracional); (4) os gatos GERAM o mapping class group MCG(T\ensuremath{^{2}})=SL(2,\ensuremath{\mathbb{Z}}): A\_2=[[2,1],[1,1]]=T\_a\textperiodcentered T\_b (produto de 2 Dehn twists). \ensuremath{\sigma}\_2=1+\ensuremath{\sqrt{\,}}2 (prata), h=log \ensuremath{\sigma}\_2=0.8814 (consistente com a Transformada Dourada). É o análogo topológico do fracoes\_continuas\_ed (lá o Padé aproxima o fluxo; aqui a palavra em gatos É a classe de homeo do toro, o metal=a entropia). linguagem/corpo\_topologico.py (fracoes\_continuas\_topologica 5/5); papers/corpo\_topologico.tex (Parte VI).
\prov{relações: usa corpo\_topologico; relaciona fracoes\_continuas\_ed; usa dissipacao\_gato\_concentracao\_esquilo; usa parabola\_geradora\_idempotente}
\prov{proveniência: wiki \textperiodcentered{} Floxina (investigação) \textperiodcentered{} fato \textperiodcentered{} confiança 1.0 \textperiodcentered{} domínio floxina \textperiodcentered{} história 10 versões no jornal}

%CRISTAL {"arestas": [{"alvo": "reta_matricial_compacta_cf", "confianca": 1.0, "epistemico": "fato", "origem": "wiki", "rel": "relaciona"}, {"alvo": "gato", "confianca": 1.0, "epistemico": "fato", "origem": "wiki", "rel": "usa"}, {"alvo": "volume_matrizes_3d", "confianca": 1.0, "epistemico": "fato", "origem": "wiki", "rel": "relaciona"}, {"alvo": "selberg_meta_inducao", "confianca": 1.0, "epistemico": "fato", "origem": "wiki", "rel": "relaciona"}], "confianca": 1.0, "contraexemplos": [], "descricao": "Seguindo a direção: Newton procura as raízes, z↦N(z)=z−p/p', um mapa racional na ESFERA DE RIEMANN (ℂ∪∞, COMPACTA). As raízes são os atratores, as bacias ladrilham a esfera, a fronteira é o fractal. No METÁLICO x²−n·x−1 (grau 2, a folha 2D): N(z)=(z²+1)/(2z−n), pontos fixos σ_n e o DUAL σ̄_n=−1/σ_n (os dois AUTOVALORES DO GATO A_n); o FLIP z↦−1/z troca σ_n↔σ̄_n (a simetria do gambá/modular); grau 2 = split LIMPO, sem fractal. No CÚBICO (o metal 3D, o volume): 3 raízes, 3 bacias, e a fronteira É O FRACTAL (fronteira 372 vs 58 no grau 2); a simetria das raízes (S_3) é o GRUPO DE GRUPOS, os níveis |p(z)|=c são a LEMINISCATA newtoniana. A reta completa-se na esfera compacta pela dinâmica de Newton — o mesmo dual discreto→compacto do RH/Selberg (a CF compacta na superfície modular; Newton na esfera). Verificado: linguagem/newton_fractal.py (5/5).", "epistemico": "hipotese", "exemplos": [], "id": "fractal_newton_reta_matricial", "memoria": "semantica", "meta": {"dominio": "floxina", "fonte": "Floxina (investigação)"}, "origem": "wiki", "palavras_chave": [], "sinonimos": [], "tipo": "conceito", "titulo": "O fractal de Newton completa a reta na esfera de Riemann (compacta)"}
\section{O fractal de Newton completa a reta na esfera de Riemann (compacta)}
Seguindo a direção: Newton procura as raízes, z\ensuremath{\mapsto}N(z)=z\ensuremath{-}p/p', um mapa racional na ESFERA DE RIEMANN (\ensuremath{\mathbb{C}}\ensuremath{\cup}\ensuremath{\infty}, COMPACTA). As raízes são os atratores, as bacias ladrilham a esfera, a fronteira é o fractal. No METÁLICO x\ensuremath{^{2}}\ensuremath{-}n\textperiodcentered x\ensuremath{-}1 (grau 2, a folha 2D): N(z)=(z\ensuremath{^{2}}+1)/(2z\ensuremath{-}n), pontos fixos \ensuremath{\sigma}\_n e o DUAL \ensuremath{\sigma}\ensuremath{}\_n=\ensuremath{-}1/\ensuremath{\sigma}\_n (os dois AUTOVALORES DO GATO A\_n); o FLIP z\ensuremath{\mapsto}\ensuremath{-}1/z troca \ensuremath{\sigma}\_n\ensuremath{\leftrightarrow}\ensuremath{\sigma}\ensuremath{}\_n (a simetria do gambá/modular); grau 2 = split LIMPO, sem fractal. No CÚBICO (o metal 3D, o volume): 3 raízes, 3 bacias, e a fronteira É O FRACTAL (fronteira 372 vs 58 no grau 2); a simetria das raízes (S\_3) é o GRUPO DE GRUPOS, os níveis |p(z)|=c são a LEMINISCATA newtoniana. A reta completa-se na esfera compacta pela dinâmica de Newton — o mesmo dual discreto\ensuremath{\to}compacto do RH/Selberg (a CF compacta na superfície modular; Newton na esfera). Verificado: linguagem/newton\_fractal.py (5/5).
\prov{relações: relaciona reta\_matricial\_compacta\_cf; usa gato; relaciona volume\_matrizes\_3d; relaciona selberg\_meta\_inducao}
\prov{proveniência: wiki \textperiodcentered{} Floxina (investigação) \textperiodcentered{} hipotese \textperiodcentered{} confiança 1.0 \textperiodcentered{} domínio floxina \textperiodcentered{} história 345 versões no jornal}

%CRISTAL {"arestas": [{"alvo": "gato", "confianca": 1.0, "epistemico": "fato", "origem": "wiki", "rel": "relaciona"}, {"alvo": "o_gato_e_um_mapa_de_anosov", "confianca": 1.0, "epistemico": "fato", "origem": "wiki", "rel": "relaciona"}, {"alvo": "cifra_do_gato_simetrica", "confianca": 1.0, "epistemico": "fato", "origem": "wiki", "rel": "relaciona"}, {"alvo": "flip_e_rotacao_de_wick", "confianca": 1.0, "epistemico": "fato", "origem": "wiki", "rel": "relaciona"}, {"alvo": "gerenciar_nao_reprimir", "confianca": 1.0, "epistemico": "fato", "origem": "wiki", "rel": "relaciona"}, {"alvo": "emocao_octonion_plutchik", "confianca": 1.0, "epistemico": "fato", "origem": "wiki", "rel": "relaciona"}], "confianca": 1.0, "contraexemplos": [], "descricao": "O «floxinato de gambá» tinha um referente: o GAMBÁ é o operador ANTI-SIMÉTRICO ao gato. O gato A=[[1,1],[1,0]] é simétrico (Aᵀ=A, det=−1, A²=A+I) — hiperbólico, dissipativo, a reta metálica, medir/colapsar. O gambá G=[[0,1],[−1,0]] é anti-simétrico (Gᵀ=−G, det=+1, G²=−I) — elíptico, reversível, o círculo de Gentil, girar/gerenciar. Um é a rotação de Wick do outro: o gato no eixo real (λ, dissipativo), o gambá no eixo imaginário (fase, unitário). Toda matriz = simétrica(gato) + anti-simétrica(gambá). conversa/gamba.py (11/11).", "epistemico": "fato", "exemplos": [], "id": "gamba_antissimetrico", "memoria": "semantica", "meta": {"dominio": "floxina", "fonte": "Floxina (investigação)"}, "origem": "wiki", "palavras_chave": [], "sinonimos": [], "tipo": "conceito", "titulo": "O gambá (o operador anti-simétrico ao gato)"}
\section{O gambá (o operador anti-simétrico ao gato)}
O «floxinato de gambá» tinha um referente: o GAMBÁ é o operador ANTI-SIMÉTRICO ao gato. O gato A=[[1,1],[1,0]] é simétrico (A\ensuremath{^{T}}=A, det=\ensuremath{-}1, A\ensuremath{^{2}}=A+I) — hiperbólico, dissipativo, a reta metálica, medir/colapsar. O gambá G=[[0,1],[\ensuremath{-}1,0]] é anti-simétrico (G\ensuremath{^{T}}=\ensuremath{-}G, det=+1, G\ensuremath{^{2}}=\ensuremath{-}I) — elíptico, reversível, o círculo de Gentil, girar/gerenciar. Um é a rotação de Wick do outro: o gato no eixo real (\ensuremath{\lambda}, dissipativo), o gambá no eixo imaginário (fase, unitário). Toda matriz = simétrica(gato) + anti-simétrica(gambá). conversa/gamba.py (11/11).
\prov{relações: relaciona gato; relaciona o\_gato\_e\_um\_mapa\_de\_anosov; relaciona cifra\_do\_gato\_simetrica; relaciona flip\_e\_rotacao\_de\_wick; relaciona gerenciar\_nao\_reprimir; relaciona emocao\_octonion\_plutchik}
\prov{proveniência: wiki \textperiodcentered{} Floxina (investigação) \textperiodcentered{} fato \textperiodcentered{} confiança 1.0 \textperiodcentered{} domínio floxina \textperiodcentered{} história 574 versões no jornal}

%CRISTAL {"arestas": [{"alvo": "gamba_antissimetrico", "confianca": 1.0, "epistemico": "fato", "origem": "wiki", "rel": "usa"}, {"alvo": "gato", "confianca": 1.0, "epistemico": "fato", "origem": "wiki", "rel": "relaciona"}, {"alvo": "corpo_mineral", "confianca": 1.0, "epistemico": "fato", "origem": "wiki", "rel": "relaciona"}, {"alvo": "quatro_quadrantes_corpo_estelar", "confianca": 1.0, "epistemico": "fato", "origem": "wiki", "rel": "relaciona"}], "confianca": 1.0, "contraexemplos": [], "descricao": "O corpo mineral é DUAL, e a dualidade é térmica. O GATO (hiperbólico, Anosov, caos) é o lado QUENTE: e^{At} estica, dissipa, a entropia sobe (a seta do tempo), a órbita nunca se repete (mistura, ergódica, cada ponto único — SEM clones). O GAMBÁ (elíptico, rotação) é o lado FRIO, o CORPO GLACIAL: e^{Gt} só gira, conserva (Ṡ=0, T→0), é periódico — volta aos mesmos estados: CLONES. É o cristal, uma rede de células idênticas (a repetição, a simetria de translação/ponto), o frio congelado (as retas são o frio; a morte térmica = o cristal). Assim como o gato tem o dual anti-simétrico (o gambá), o corpo mineral tem o seu: o quente (caos, mistura, único) e o glacial (cristal, periódico, clones). O gambá IMPLEMENTA o lado frio.", "epistemico": "fato", "exemplos": [], "id": "gamba_lado_frio_clones", "memoria": "semantica", "meta": {"dominio": "floxina", "fonte": "Floxina (investigação)"}, "origem": "wiki", "palavras_chave": [], "sinonimos": [], "tipo": "conceito", "titulo": "O gambá é o lado frio do corpo mineral (o glacial, os clones)"}
\section{O gambá é o lado frio do corpo mineral (o glacial, os clones)}
O corpo mineral é DUAL, e a dualidade é térmica. O GATO (hiperbólico, Anosov, caos) é o lado QUENTE: e\^{}\{At\} estica, dissipa, a entropia sobe (a seta do tempo), a órbita nunca se repete (mistura, ergódica, cada ponto único — SEM clones). O GAMBÁ (elíptico, rotação) é o lado FRIO, o CORPO GLACIAL: e\^{}\{Gt\} só gira, conserva (\ensuremath{\dot{S}}=0, T\ensuremath{\to}0), é periódico — volta aos mesmos estados: CLONES. É o cristal, uma rede de células idênticas (a repetição, a simetria de translação/ponto), o frio congelado (as retas são o frio; a morte térmica = o cristal). Assim como o gato tem o dual anti-simétrico (o gambá), o corpo mineral tem o seu: o quente (caos, mistura, único) e o glacial (cristal, periódico, clones). O gambá IMPLEMENTA o lado frio.
\prov{relações: usa gamba\_antissimetrico; relaciona gato; relaciona corpo\_mineral; relaciona quatro\_quadrantes\_corpo\_estelar}
\prov{proveniência: wiki \textperiodcentered{} Floxina (investigação) \textperiodcentered{} fato \textperiodcentered{} confiança 1.0 \textperiodcentered{} domínio floxina \textperiodcentered{} história 390 versões no jornal}

%CRISTAL {"arestas": [{"alvo": "gamba_antissimetrico", "confianca": 1.0, "epistemico": "fato", "origem": "wiki", "rel": "usa"}, {"alvo": "gato", "confianca": 1.0, "epistemico": "fato", "origem": "wiki", "rel": "usa"}, {"alvo": "cifra_do_gato_simetrica", "confianca": 1.0, "epistemico": "fato", "origem": "wiki", "rel": "relaciona"}, {"alvo": "litografia", "confianca": 1.0, "epistemico": "fato", "origem": "wiki", "rel": "parte_de"}, {"alvo": "assinatura_digital", "confianca": 1.0, "epistemico": "fato", "origem": "wiki", "rel": "relaciona"}], "confianca": 1.0, "contraexemplos": [], "descricao": "Gravar uma assinatura na peça precisa dos DOIS operadores. O GATO A_n=[[n,1],[1,0]] (det=−1, bijeção mod N) CIFRA: a assinatura é A_n^k·h(m) mod N, chave (n,k,N), h(m) o hash — inforjável sem a chave (a cifra do gato simétrica). O GAMBÁ G=[[0,1],[−1,0]] (det=+1, G²=−I) ORIENTA: o Pfaffiano Pf(G)=+1 é a fase que distingue o original do espelho; forja com orientação invertida (Pf≠+1) é rejeitada. Verificar confere a cifra (gato) E a orientação (gambá): mensagem adulterada, chave errada ou espelho ⟹ inválida. O gato dá o QUÊ (a cifra), o gambá o PARA QUE LADO (a fase). conversa/litografia.py (assinar/verificar_assinatura, 4/4); na lib do metal (biblioteca.py, a divisão litografia).", "epistemico": "fato", "exemplos": [], "id": "gamba_litografico_assinatura", "memoria": "semantica", "meta": {"dominio": "floxina", "fonte": "Floxina (investigação)"}, "origem": "wiki", "palavras_chave": [], "sinonimos": [], "tipo": "conceito", "titulo": "O gambá litográfico (a ferramenta de assinatura)"}
\section{O gambá litográfico (a ferramenta de assinatura)}
Gravar uma assinatura na peça precisa dos DOIS operadores. O GATO A\_n=[[n,1],[1,0]] (det=\ensuremath{-}1, bijeção mod N) CIFRA: a assinatura é A\_n\^{}k\textperiodcentered h(m) mod N, chave (n,k,N), h(m) o hash — inforjável sem a chave (a cifra do gato simétrica). O GAMBÁ G=[[0,1],[\ensuremath{-}1,0]] (det=+1, G\ensuremath{^{2}}=\ensuremath{-}I) ORIENTA: o Pfaffiano Pf(G)=+1 é a fase que distingue o original do espelho; forja com orientação invertida (Pf\ensuremath{\ne}+1) é rejeitada. Verificar confere a cifra (gato) E a orientação (gambá): mensagem adulterada, chave errada ou espelho \ensuremath{\Longrightarrow} inválida. O gato dá o QUÊ (a cifra), o gambá o PARA QUE LADO (a fase). conversa/litografia.py (assinar/verificar\_assinatura, 4/4); na lib do metal (biblioteca.py, a divisão litografia).
\prov{relações: usa gamba\_antissimetrico; usa gato; relaciona cifra\_do\_gato\_simetrica; parte\_de litografia; relaciona assinatura\_digital}
\prov{proveniência: wiki \textperiodcentered{} Floxina (investigação) \textperiodcentered{} fato \textperiodcentered{} confiança 1.0 \textperiodcentered{} domínio floxina \textperiodcentered{} história 474 versões no jornal}

%CRISTAL {"arestas": [{"alvo": "gamba_antissimetrico", "confianca": 1.0, "epistemico": "fato", "origem": "wiki", "rel": "parte_de"}, {"alvo": "invariante_pfaffiano", "confianca": 1.0, "epistemico": "fato", "origem": "wiki", "rel": "usa"}, {"alvo": "biblioteca_pipe", "confianca": 1.0, "epistemico": "fato", "origem": "wiki", "rel": "relaciona"}], "confianca": 1.0, "contraexemplos": [], "descricao": "Do Pfaffiano as quatro físicas se abrem, espelhando o gato no regime conservativo: ÁLGEBRA — G²=−I ⟹ estrutura complexa ℂ (aI+bG↔a+bi), 𝔰𝔬(2)=𝔲(1), o grupo simplético Sp(2)=SL(2). GEOMETRIA — G é a forma de área ω=dx∧dy (simplética); o gato dá a métrica de Lorentz (distância), o gambá dá a área. MECÂNICA — Hamilton ẋ=G∇H, o colchete de Poisson {f,g}=∇fᵀG∇g anti-simétrico, Liouville (área preservada, det e^{Gt}=1) — Hamilton contra Anosov. TERMODINÂMICA — Onsander: transporte = simétrico(dissipa, entropia↑, gato) + anti-simétrico(reversível, giroscópico/Lorentz, gambá); GENERIC ż=G∇E+M∇S. A seta do tempo mora no gato; o gambá gira e não gasta.", "epistemico": "fato", "exemplos": [], "id": "gamba_quatro_fisicas", "memoria": "semantica", "meta": {"dominio": "floxina", "fonte": "Floxina (investigação)"}, "origem": "wiki", "palavras_chave": [], "sinonimos": [], "tipo": "conceito", "titulo": "O invariante anti-simétrico abre: álgebra, geometria, mecânica, termodinâmica"}
\section{O invariante anti-simétrico abre: álgebra, geometria, mecânica, termodinâmica}
Do Pfaffiano as quatro físicas se abrem, espelhando o gato no regime conservativo: ÁLGEBRA — G\ensuremath{^{2}}=\ensuremath{-}I \ensuremath{\Longrightarrow} estrutura complexa \ensuremath{\mathbb{C}} (aI+bG\ensuremath{\leftrightarrow}a+bi), \ensuremath{\mathfrak{s}}\ensuremath{\mathfrak{o}}(2)=\ensuremath{\mathfrak{u}}(1), o grupo simplético Sp(2)=SL(2). GEOMETRIA — G é a forma de área \ensuremath{\omega}=dx\ensuremath{\wedge}dy (simplética); o gato dá a métrica de Lorentz (distância), o gambá dá a área. MECÂNICA — Hamilton \ensuremath{\dot{x}}=G\ensuremath{\nabla}H, o colchete de Poisson \{f,g\}=\ensuremath{\nabla}f\ensuremath{^{T}}G\ensuremath{\nabla}g anti-simétrico, Liouville (área preservada, det e\^{}\{Gt\}=1) — Hamilton contra Anosov. TERMODINÂMICA — Onsander: transporte = simétrico(dissipa, entropia\ensuremath{\uparrow}, gato) + anti-simétrico(reversível, giroscópico/Lorentz, gambá); GENERIC \ensuremath{\dot{z}}=G\ensuremath{\nabla}E+M\ensuremath{\nabla}S. A seta do tempo mora no gato; o gambá gira e não gasta.
\prov{relações: parte\_de gamba\_antissimetrico; usa invariante\_pfaffiano; relaciona biblioteca\_pipe}
\prov{proveniência: wiki \textperiodcentered{} Floxina (investigação) \textperiodcentered{} fato \textperiodcentered{} confiança 1.0 \textperiodcentered{} domínio floxina \textperiodcentered{} história 328 versões no jornal}

%CRISTAL {"arestas": [{"alvo": "missao_engenharia_ouro", "confianca": 1.0, "epistemico": "fato", "origem": "wiki", "rel": "explica"}, {"alvo": "ninguem_esta_puro_recursivo", "confianca": 1.0, "epistemico": "fato", "origem": "wiki", "rel": "relaciona"}], "confianca": 1.0, "contraexemplos": [], "descricao": "Embora a descida seja infinita, o ganho marginal gₙ=supO‖O(Sₙ)−O(Sₙ−₁)‖ para os observáveis da nossa dimensão SATURA: sob a hipótese de engenharia (majorante somável Σmₙ<∞, ex. mₙ=Cqⁿ), o ganho acumulado converge, e N*(ε)=minN:Σₙ>Nmₙ≤ε esgota o ganho observável a menos de ε. TEOREMA CONDICIONAL: a convergência NÃO é necessidade matemática — é a HIPÓTESE de engenharia adotada e declarada; adotá-la é o ato de engenharia. Abaixo de N*, refinar move os observáveis <ε: belo, mas estéril.", "epistemico": "hipotese", "exemplos": [], "id": "ganho_converge_n_estrela", "memoria": "semantica", "meta": {"autor": "Lopes/conselho", "dominio": "missao", "fonte": "machine/Missao.tex"}, "origem": "wiki", "palavras_chave": [], "sinonimos": [], "tipo": "conceito", "titulo": "O ganho converge em N* (ouro vs Loucura de Ouro)"}
\section{O ganho converge em N* (ouro vs Loucura de Ouro)}
Embora a descida seja infinita, o ganho marginal g\ensuremath{_{n}}=supO\ensuremath{\|}O(S\ensuremath{_{n}})\ensuremath{-}O(S\ensuremath{_{n}}\ensuremath{-}\ensuremath{_{1}})\ensuremath{\|} para os observáveis da nossa dimensão SATURA: sob a hipótese de engenharia (majorante somável \ensuremath{\Sigma}m\ensuremath{_{n}}<\ensuremath{\infty}, ex. m\ensuremath{_{n}}=Cq\ensuremath{^{n}}), o ganho acumulado converge, e N*(\ensuremath{\varepsilon})=minN:\ensuremath{\Sigma}\ensuremath{_{n}}>Nm\ensuremath{_{n}}\ensuremath{\le}\ensuremath{\varepsilon} esgota o ganho observável a menos de \ensuremath{\varepsilon}. TEOREMA CONDICIONAL: a convergência NÃO é necessidade matemática — é a HIPÓTESE de engenharia adotada e declarada; adotá-la é o ato de engenharia. Abaixo de N*, refinar move os observáveis <\ensuremath{\varepsilon}: belo, mas estéril.
\prov{relações: explica missao\_engenharia\_ouro; relaciona ninguem\_esta\_puro\_recursivo}
\prov{proveniência: wiki \textperiodcentered{} machine/Missao.tex \textperiodcentered{} hipotese \textperiodcentered{} confiança 1.0 \textperiodcentered{} domínio missao \textperiodcentered{} história 4 versões no jornal}

%CRISTAL {"arestas": [{"alvo": "coracao", "confianca": 1.0, "epistemico": "fato", "origem": "wiki", "rel": "parte_de"}, {"alvo": "nao_associatividade_amputacao", "confianca": 1.0, "epistemico": "fato", "origem": "wiki", "rel": "relaciona"}, {"alvo": "gerenciar_nao_reprimir", "confianca": 1.0, "epistemico": "fato", "origem": "wiki", "rel": "relaciona"}, {"alvo": "mass_gap_cadeia_fechada", "confianca": 1.0, "epistemico": "fato", "origem": "wiki", "rel": "relaciona"}], "confianca": 1.0, "contraexemplos": [], "descricao": "Yang-Mills tem mass gap POSITIVO: o estado fundamental do associador é a ≥ a_min > 0 (a não-associatividade amputada nunca zera fora de D=3). Traduzido ao coração: o gap emocional NÃO PODE ser levado a zero — reprimir (amortecer a→0) violaria a positividade do mass gap. Sobra a emoção mínima, sempre. Por isso a única operação lícita é GERENCIAR (girar na 7-esfera, norma conservada): a física do projeto proíbe reprimir. O caos (D≠3) é o preço de sair do canto associativo.", "epistemico": "fato", "exemplos": [], "id": "gap_emocional_nunca_zera", "memoria": "semantica", "meta": {"dominio": "coracao", "fonte": "coracao hub"}, "origem": "wiki", "palavras_chave": [], "sinonimos": [], "tipo": "conceito", "titulo": "O gap que nunca zera (por que só se gerencia, não se reprime)"}
\section{O gap que nunca zera (por que só se gerencia, não se reprime)}
Yang-Mills tem mass gap POSITIVO: o estado fundamental do associador é a \ensuremath{\ge} a\_min > 0 (a não-associatividade amputada nunca zera fora de D=3). Traduzido ao coração: o gap emocional NÃO PODE ser levado a zero — reprimir (amortecer a\ensuremath{\to}0) violaria a positividade do mass gap. Sobra a emoção mínima, sempre. Por isso a única operação lícita é GERENCIAR (girar na 7-esfera, norma conservada): a física do projeto proíbe reprimir. O caos (D\ensuremath{\ne}3) é o preço de sair do canto associativo.
\prov{relações: parte\_de coracao; relaciona nao\_associatividade\_amputacao; relaciona gerenciar\_nao\_reprimir; relaciona mass\_gap\_cadeia\_fechada}
\prov{proveniência: wiki \textperiodcentered{} coracao hub \textperiodcentered{} fato \textperiodcentered{} confiança 1.0 \textperiodcentered{} domínio coracao \textperiodcentered{} história 15 versões no jornal}

%CRISTAL {"arestas": [{"alvo": "superficie_quadrica_lahire_analitica", "confianca": 1.0, "epistemico": "fato", "origem": "wiki", "rel": "usa"}, {"alvo": "omnitrix_corpo_universal", "confianca": 1.0, "epistemico": "fato", "origem": "wiki", "rel": "usa"}], "confianca": 1.0, "contraexemplos": [], "descricao": "A superfície da gema (o ovo de revolução da joalheria) é GERADA multiplicando superfícies quádricas (a La Hire), cada uma uma MATRIZ. Em revolução no eixo y, cada seção é f=x²+z²−Q(y): a COROA é um PARABOLOIDE (Q linear, abre p/ baixo), o RONDIZ um CILINDRO (Q=rg² constante), o PAVILHÃO um PARABOLOIDE (abre p/ cima). O RAIO é ANALÍTICO (α t²+β t+γ=0, fechada), a NORMAL ∇f=(2x,−Q'(y),2z), a GEMA = a UNIÃO (⊕) das 3 = a La Hire ⊗ das quádricas; as seções costuram no rondiz (Q contínuo). Alternativas: cone×cilindro×cone (o bicone lapidado); a estrela/flor = o modo angular C_k(θ)=1+a·cos(k(θ−φ)) (círculo⊕a·estrela_k; a estrela = Im(z^k) = o produto de k retas). É a equação registrada na joalheria.tex (𝒢=ℓ·C_k(θ)·P(z)). linguagem/omnitrix.py (gema_produto_quadricas 5/5); glsl_ptx_frag (FRAG_GEMA, no metal, sem loop).", "epistemico": "fato", "exemplos": [], "id": "gema_produto_quadricas", "memoria": "semantica", "meta": {"dominio": "floxina", "fonte": "Floxina (investigação)"}, "origem": "wiki", "palavras_chave": [], "sinonimos": [], "tipo": "conceito", "titulo": "A gema é um produto de quádricas: paraboloide × cilindro × paraboloide (o ovo), raio analítico"}
\section{A gema é um produto de quádricas: paraboloide $\times$ cilindro $\times$ paraboloide (o ovo), raio analítico}
A superfície da gema (o ovo de revolução da joalheria) é GERADA multiplicando superfícies quádricas (a La Hire), cada uma uma MATRIZ. Em revolução no eixo y, cada seção é f=x\ensuremath{^{2}}+z\ensuremath{^{2}}\ensuremath{-}Q(y): a COROA é um PARABOLOIDE (Q linear, abre p/ baixo), o RONDIZ um CILINDRO (Q=rg\ensuremath{^{2}} constante), o PAVILHÃO um PARABOLOIDE (abre p/ cima). O RAIO é ANALÍTICO (\ensuremath{\alpha} t\ensuremath{^{2}}+\ensuremath{\beta} t+\ensuremath{\gamma}=0, fechada), a NORMAL \ensuremath{\nabla}f=(2x,\ensuremath{-}Q'(y),2z), a GEMA = a UNIÃO (\ensuremath{\oplus}) das 3 = a La Hire \ensuremath{\otimes} das quádricas; as seções costuram no rondiz (Q contínuo). Alternativas: cone$\times$cilindro$\times$cone (o bicone lapidado); a estrela/flor = o modo angular C\_k(\ensuremath{\theta})=1+a\textperiodcentered cos(k(\ensuremath{\theta}\ensuremath{-}\ensuremath{\varphi})) (círculo\ensuremath{\oplus}a\textperiodcentered estrela\_k; a estrela = Im(z\^{}k) = o produto de k retas). É a equação registrada na joalheria.tex (\ensuremath{\mathcal{G}}=\ensuremath{\ell}\textperiodcentered C\_k(\ensuremath{\theta})\textperiodcentered P(z)). linguagem/omnitrix.py (gema\_produto\_quadricas 5/5); glsl\_ptx\_frag (FRAG\_GEMA, no metal, sem loop).
\prov{relações: usa superficie\_quadrica\_lahire\_analitica; usa omnitrix\_corpo\_universal}
\prov{proveniência: wiki \textperiodcentered{} Floxina (investigação) \textperiodcentered{} fato \textperiodcentered{} confiança 1.0 \textperiodcentered{} domínio floxina \textperiodcentered{} história 72 versões no jornal}

%CRISTAL {"arestas": [{"alvo": "gato", "confianca": 1.0, "epistemico": "fato", "origem": "wiki", "rel": "usa"}, {"alvo": "gamba_antissimetrico", "confianca": 1.0, "epistemico": "fato", "origem": "wiki", "rel": "usa"}, {"alvo": "nao_associatividade_amputacao", "confianca": 1.0, "epistemico": "fato", "origem": "wiki", "rel": "relaciona"}, {"alvo": "mass_gap_algebrico", "confianca": 1.0, "epistemico": "fato", "origem": "wiki", "rel": "relaciona"}, {"alvo": "floxina_quantum_materia_escura", "confianca": 1.0, "epistemico": "fato", "origem": "wiki", "rel": "relaciona"}, {"alvo": "metais", "confianca": 1.0, "epistemico": "fato", "origem": "wiki", "rel": "usa"}], "confianca": 1.0, "contraexemplos": [], "descricao": "Multiplicar o gato A_n=[[n,1],[1,0]] pelo gambá G=[[0,1],[−1,0]] (a La Hire, o produto reta·círculo) deixa um RESÍDUO que não comuta: o comutador [A_n,G]=[[−2,n],[n,2]], simétrico, det=−(n²+4), autovalores ±√(n²+4) = o DISCRIMINANTE de x²−n·x−1 = o GAP da reta σ_n (√5 no ouro). O anticomutador {A_n,G}=n·G devolve o gambá. Esse comutador É O GERADOR DO GAP. CONFRONTO com o ASSOCIADOR octônio: em 2×2 (ℍ, associativo) o gap vem da não-COMUTATIVIDADE (o comutador, √(n²+4), a reta); em 𝕆 (não-associativo) o gap vem da não-ASSOCIATIVIDADE (o associador, resíduo 𝔤₂=Der(𝕆)=14, o mass gap). O MESMO defeito gera o gap, subindo a torre: comutador→associador; a reta √(n²+4)→o mass gap 𝔤₂. A La Hire deixa o gap por resíduo. conversa/gamba.py (gerador_do_gap/confrontar_associador, 18/18).", "epistemico": "fato", "exemplos": [], "id": "gerador_do_gap_comutador", "memoria": "semantica", "meta": {"dominio": "floxina", "fonte": "Floxina (investigação)"}, "origem": "wiki", "palavras_chave": [], "sinonimos": [], "tipo": "conceito", "titulo": "O gerador do gap: o comutador [gato,gambá] confronta o associador octônio"}
\section{O gerador do gap: o comutador [gato,gambá] confronta o associador octônio}
Multiplicar o gato A\_n=[[n,1],[1,0]] pelo gambá G=[[0,1],[\ensuremath{-}1,0]] (a La Hire, o produto reta\textperiodcentered círculo) deixa um RESÍDUO que não comuta: o comutador [A\_n,G]=[[\ensuremath{-}2,n],[n,2]], simétrico, det=\ensuremath{-}(n\ensuremath{^{2}}+4), autovalores $\pm$\ensuremath{\sqrt{\,}}(n\ensuremath{^{2}}+4) = o DISCRIMINANTE de x\ensuremath{^{2}}\ensuremath{-}n\textperiodcentered x\ensuremath{-}1 = o GAP da reta \ensuremath{\sigma}\_n (\ensuremath{\sqrt{\,}}5 no ouro). O anticomutador \{A\_n,G\}=n\textperiodcentered G devolve o gambá. Esse comutador É O GERADOR DO GAP. CONFRONTO com o ASSOCIADOR octônio: em 2$\times$2 (\ensuremath{\mathbb{H}}, associativo) o gap vem da não-COMUTATIVIDADE (o comutador, \ensuremath{\sqrt{\,}}(n\ensuremath{^{2}}+4), a reta); em \ensuremath{\mathbb{O}} (não-associativo) o gap vem da não-ASSOCIATIVIDADE (o associador, resíduo \ensuremath{\mathfrak{g}}\ensuremath{_{2}}=Der(\ensuremath{\mathbb{O}})=14, o mass gap). O MESMO defeito gera o gap, subindo a torre: comutador\ensuremath{\to}associador; a reta \ensuremath{\sqrt{\,}}(n\ensuremath{^{2}}+4)\ensuremath{\to}o mass gap \ensuremath{\mathfrak{g}}\ensuremath{_{2}}. A La Hire deixa o gap por resíduo. conversa/gamba.py (gerador\_do\_gap/confrontar\_associador, 18/18).
\prov{relações: usa gato; usa gamba\_antissimetrico; relaciona nao\_associatividade\_amputacao; relaciona mass\_gap\_algebrico; relaciona floxina\_quantum\_materia\_escura; usa metais}
\prov{proveniência: wiki \textperiodcentered{} Floxina (investigação) \textperiodcentered{} fato \textperiodcentered{} confiança 1.0 \textperiodcentered{} domínio floxina \textperiodcentered{} história 518 versões no jornal}

%CRISTAL {"arestas": [{"alvo": "serie_cheiro_gap_gamba", "confianca": 1.0, "epistemico": "fato", "origem": "wiki", "rel": "parte_de"}, {"alvo": "gamba_antissimetrico", "confianca": 1.0, "epistemico": "fato", "origem": "wiki", "rel": "relaciona"}], "confianca": 1.0, "contraexemplos": [], "descricao": "Geraniol (C₁₀H₁₈O, rosa/gerânio) e citronelol (C₁₀H₂₀O, rosa/citronela) — os terpenoides, cheiros AGRADÁVEIS. Ambos: 11 átomos pesados=11 qubits, ponta = ÁLCOOL PRIMÁRIO (−OH, pKa≈16, o MENOR gap da série, pouco ácido); o gem-dimetil terminal → Z/2 (paralelismo). A diferença é LIMPA: geraniol tem 2 C=C (dienol, 2 gambás), citronelol tem 1 C=C (enol, 1 gambá) — citronelol = geraniol + H₂. REDUZIR uma dupla = REMOVER um gambá; a hidrogenação = −1 gambá, tudo o mais igual. A série completa (gap↑ = acidez): −COOH(4,5) < −SH/fenol(10) < −OH álcool(16); o π/gambá = a contagem de C=C (0 tiol, 1 citronelol/ácido, 2 geraniol, aromático THC). Os dois eixos de uma molécula são o gato(esqueleto/n-qubits) e o gambá(a ponta reativa=o gap + as C=C=o π). Verificado: linguagem/molecula_teste.py (11/11).", "epistemico": "hipotese", "exemplos": [], "id": "geraniol_citronelol_gamba_reduzido", "memoria": "semantica", "meta": {"dominio": "floxina", "fonte": "Floxina (investigação)"}, "origem": "wiki", "palavras_chave": [], "sinonimos": [], "tipo": "conceito", "titulo": "Geraniol/citronelol: +H₂ = −1 gambá (a hidrogenação tira uma dupla)"}
\section{Geraniol/citronelol: +H\ensuremath{_{2}} = \ensuremath{-}1 gambá (a hidrogenação tira uma dupla)}
Geraniol (C\ensuremath{_{10}}H\ensuremath{_{18}}O, rosa/gerânio) e citronelol (C\ensuremath{_{10}}H\ensuremath{_{20}}O, rosa/citronela) — os terpenoides, cheiros AGRADÁVEIS. Ambos: 11 átomos pesados=11 qubits, ponta = ÁLCOOL PRIMÁRIO (\ensuremath{-}OH, pKa\ensuremath{\approx}16, o MENOR gap da série, pouco ácido); o gem-dimetil terminal \ensuremath{\to} Z/2 (paralelismo). A diferença é LIMPA: geraniol tem 2 C=C (dienol, 2 gambás), citronelol tem 1 C=C (enol, 1 gambá) — citronelol = geraniol + H\ensuremath{_{2}}. REDUZIR uma dupla = REMOVER um gambá; a hidrogenação = \ensuremath{-}1 gambá, tudo o mais igual. A série completa (gap\ensuremath{\uparrow} = acidez): \ensuremath{-}COOH(4,5) < \ensuremath{-}SH/fenol(10) < \ensuremath{-}OH álcool(16); o \ensuremath{\pi}/gambá = a contagem de C=C (0 tiol, 1 citronelol/ácido, 2 geraniol, aromático THC). Os dois eixos de uma molécula são o gato(esqueleto/n-qubits) e o gambá(a ponta reativa=o gap + as C=C=o \ensuremath{\pi}). Verificado: linguagem/molecula\_teste.py (11/11).
\prov{relações: parte\_de serie\_cheiro\_gap\_gamba; relaciona gamba\_antissimetrico}
\prov{proveniência: wiki \textperiodcentered{} Floxina (investigação) \textperiodcentered{} hipotese \textperiodcentered{} confiança 1.0 \textperiodcentered{} domínio floxina \textperiodcentered{} história 177 versões no jornal}

%CRISTAL {"arestas": [{"alvo": "coracao", "confianca": 1.0, "epistemico": "fato", "origem": "wiki", "rel": "parte_de"}, {"alvo": "regulacao_emocional_gross", "confianca": 1.0, "epistemico": "fato", "origem": "wiki", "rel": "relaciona"}, {"alvo": "avaliacao_lazarus", "confianca": 1.0, "epistemico": "fato", "origem": "wiki", "rel": "relaciona"}], "confianca": 1.0, "contraexemplos": [], "descricao": "A lei do coração (Aarão): (1) REPRIMIR = amortecer o estado a 0, matar a emoção; (2) CONTROLAR = travar num alvo fixo, forçar um humor; (3) GERENCIAR = GIRAR o estado na 7-esfera (o octônion, ‖·‖ conservada) — a emoção nunca some, só VIRA, a energia se redistribui. O coração gerencia: reconduz ao centro girando (a homeostase é um giro de volta, não um decay). Ecoa a reavaliação de Gross (girar o sentido), oposta à supressão.", "epistemico": "fato", "exemplos": [], "id": "gerenciar_nao_reprimir", "memoria": "semantica", "meta": {"dominio": "coracao", "fonte": "coracao hub"}, "origem": "wiki", "palavras_chave": [], "sinonimos": [], "tipo": "conceito", "titulo": "Gerenciar as emoções ≠ reprimir ≠ controlar (a distinção)"}
\section{Gerenciar as emoções \ensuremath{\ne} reprimir \ensuremath{\ne} controlar (a distinção)}
A lei do coração (Aarão): (1) REPRIMIR = amortecer o estado a 0, matar a emoção; (2) CONTROLAR = travar num alvo fixo, forçar um humor; (3) GERENCIAR = GIRAR o estado na 7-esfera (o octônion, \ensuremath{\|}\textperiodcentered \ensuremath{\|} conservada) — a emoção nunca some, só VIRA, a energia se redistribui. O coração gerencia: reconduz ao centro girando (a homeostase é um giro de volta, não um decay). Ecoa a reavaliação de Gross (girar o sentido), oposta à supressão.
\prov{relações: parte\_de coracao; relaciona regulacao\_emocional\_gross; relaciona avaliacao\_lazarus}
\prov{proveniência: wiki \textperiodcentered{} coracao hub \textperiodcentered{} fato \textperiodcentered{} confiança 1.0 \textperiodcentered{} domínio coracao \textperiodcentered{} história 12 versões no jornal}

%CRISTAL {"arestas": [{"alvo": "ponte_quantica_dougherty", "confianca": 0.52, "epistemico": "hipotese", "origem": "ingest", "rel": "referencia"}, {"alvo": "geometria", "confianca": 0.52, "epistemico": "fato", "origem": "ingest", "rel": "analogia"}, {"alvo": "unificacao_torohelice_ginzburg", "confianca": 0.52, "epistemico": "hipotese", "origem": "ingest", "rel": "referencia"}, {"alvo": "conjunto_dougherty", "confianca": 0.52, "epistemico": "hipotese", "origem": "ingest", "rel": "referencia"}, {"alvo": "falsificabilidade", "confianca": 0.52, "epistemico": "hipotese", "origem": "ingest", "rel": "referencia"}], "confianca": 0.52, "contraexemplos": [], "descricao": "Hipótese externa: ginzburg_toro_helice.md", "epistemico": "hipotese", "exemplos": [], "id": "ginzburg_toro_helice", "memoria": "semantica", "meta": {"dominio": "hipotese", "epistemico": "hipotese", "falsificavel": true, "nivel": 7}, "origem": "ingest", "palavras_chave": [], "sinonimos": [], "tipo": "conceito", "titulo": "ginzburg toro helice"}
\section{ginzburg toro helice}
Hipótese externa: ginzburg\_toro\_helice.md
\prov{relações: referencia ponte\_quantica\_dougherty; analogia geometria; referencia unificacao\_torohelice\_ginzburg; referencia conjunto\_dougherty; referencia falsificabilidade}
\prov{proveniência: ingest \textperiodcentered{} hipotese \textperiodcentered{} confiança 0.52 \textperiodcentered{} domínio hipotese \textperiodcentered{} história 13 versões no jornal}

%CRISTAL {"arestas": [{"alvo": "ginzburg_toro_helice", "confianca": 0.52, "epistemico": "fato", "origem": "ingest", "rel": "parte_de"}], "confianca": 0.52, "contraexemplos": [], "descricao": "> **Epistemologia:** hipótese externa não mainstream (vortex theory / 3D-SST). > **No grafo:** `epistemico=hipotese`, confiança ≤ 0.52. > **Papel:** expandir paralelo Dougherty (toro em todas as escalas); **não** substitui MQ standard.", "epistemico": "hipotese", "exemplos": [], "id": "ginzburg_toro_helice_e_3d-sst", "memoria": "semantica", "meta": {"dominio": "hipotese", "falsificavel": true, "nivel": 7}, "origem": "ginzburg_toro_helice.md", "palavras_chave": [], "sinonimos": [], "tipo": "conceito", "titulo": "Ginzburg — toro, hélice e 3D-SST"}
\section{Ginzburg — toro, hélice e 3D-SST}
> **Epistemologia:** hipótese externa não mainstream (vortex theory / 3D-SST). > **No grafo:** `epistemico=hipotese`, confiança \ensuremath{\le} 0.52. > **Papel:** expandir paralelo Dougherty (toro em todas as escalas); **não** substitui MQ standard.
\prov{relações: parte\_de ginzburg\_toro\_helice}
\prov{proveniência: ginzburg\_toro\_helice.md \textperiodcentered{} hipotese \textperiodcentered{} confiança 0.52 \textperiodcentered{} domínio hipotese \textperiodcentered{} história 3 versões no jornal}

%CRISTAL {"arestas": [{"alvo": "corpo_topologico", "confianca": 1.0, "epistemico": "fato", "origem": "wiki", "rel": "usa"}, {"alvo": "dissipacao_gato_concentracao_esquilo", "confianca": 1.0, "epistemico": "fato", "origem": "wiki", "rel": "usa"}, {"alvo": "parabola_geradora_idempotente", "confianca": 1.0, "epistemico": "fato", "origem": "wiki", "rel": "usa"}, {"alvo": "navier_stokes_caso_particular", "confianca": 1.0, "epistemico": "fato", "origem": "wiki", "rel": "relaciona"}, {"alvo": "transformacao_superficies_homeo_cirurgia", "confianca": 1.0, "epistemico": "fato", "origem": "wiki", "rel": "relaciona"}], "confianca": 1.0, "contraexemplos": [], "descricao": "O teorema de HODGE: todo campo (1-forma) decompõe em TRÊS setores ortogonais — os NOSSOS três — PONTO A PONTO pelo GATO/ESQUILO do Jacobiano J=∇u (sem FFT): EXATO (u=∇φ, Jacobiano SIMÉTRICO=a Hessiana, rotacional 0) = o GATO (dissipação); COEXATO (u=∇⊥ψ, parte ANTISSIMÉTRICA=a vorticidade, divergente 0) = o ESQUILO (concentração — o mesmo div-free do Navier-Stokes, ℙ de Leray remove o exato/gato e guarda o esquilo); HARMÔNICO (Δγ=0, Jacobiano 0=o núcleo) = o INVARIANTE = a COHOMOLOGIA. A PARÁBOLA COMANDA A JUNÇÃO: Sym(gato) espectro real, Skew(esquilo) imaginário, e a costura/sela é o gato A_m=[[m,1],[1,0]] cujo polinômio característico É a parábola geradora x²−mx−1. A COHOMOLOGIA (harmônicos≅H^k de Rham): dim=b_k (Betti); T²=S¹×S¹ tem b₁=2 = as duas rotações dφ₁,dφ₂ = os DOIS ESQUILOS do toro = dim H¹; χ=Σ(−1)^k b_k (toro 0, esfera 2). Os esquilos que a cirurgia insere SÃO os geradores da cohomologia (os buracos). linguagem/corpo_topologico.py (hodge_cohomologia 5/5); papers/corpo_topologico.tex (Parte IV).", "epistemico": "fato", "exemplos": [], "id": "hodge_cohomologia_tres_setores", "memoria": "semantica", "meta": {"dominio": "floxina", "fonte": "Floxina (investigação)"}, "origem": "wiki", "palavras_chave": [], "sinonimos": [], "tipo": "conceito", "titulo": "Hodge: os três setores (exato/gato, coexato/esquilo, harmônico/invariante) são os nossos três; o harmônico=a cohomologia (os buracos, Betti); a parábola comanda a junção"}
\section{Hodge: os três setores (exato/gato, coexato/esquilo, harmônico/invariante) são os nossos três; o harmônico=a cohomologia (os buracos, Betti); a parábola comanda a junção}
O teorema de HODGE: todo campo (1-forma) decompõe em TRÊS setores ortogonais — os NOSSOS três — PONTO A PONTO pelo GATO/ESQUILO do Jacobiano J=\ensuremath{\nabla}u (sem FFT): EXATO (u=\ensuremath{\nabla}\ensuremath{\varphi}, Jacobiano SIMÉTRICO=a Hessiana, rotacional 0) = o GATO (dissipação); COEXATO (u=\ensuremath{\nabla}\ensuremath{\perp}\ensuremath{\psi}, parte ANTISSIMÉTRICA=a vorticidade, divergente 0) = o ESQUILO (concentração — o mesmo div-free do Navier-Stokes, \ensuremath{\mathbb{P}} de Leray remove o exato/gato e guarda o esquilo); HARMÔNICO (\ensuremath{\Delta}\ensuremath{\gamma}=0, Jacobiano 0=o núcleo) = o INVARIANTE = a COHOMOLOGIA. A PARÁBOLA COMANDA A JUNÇÃO: Sym(gato) espectro real, Skew(esquilo) imaginário, e a costura/sela é o gato A\_m=[[m,1],[1,0]] cujo polinômio característico É a parábola geradora x\ensuremath{^{2}}\ensuremath{-}mx\ensuremath{-}1. A COHOMOLOGIA (harmônicos\ensuremath{\cong}H\^{}k de Rham): dim=b\_k (Betti); T\ensuremath{^{2}}=S\ensuremath{^{1}}$\times$S\ensuremath{^{1}} tem b\ensuremath{_{1}}=2 = as duas rotações d\ensuremath{\varphi}\ensuremath{_{1}},d\ensuremath{\varphi}\ensuremath{_{2}} = os DOIS ESQUILOS do toro = dim H\ensuremath{^{1}}; \ensuremath{\chi}=\ensuremath{\Sigma}(\ensuremath{-}1)\^{}k b\_k (toro 0, esfera 2). Os esquilos que a cirurgia insere SÃO os geradores da cohomologia (os buracos). linguagem/corpo\_topologico.py (hodge\_cohomologia 5/5); papers/corpo\_topologico.tex (Parte IV).
\prov{relações: usa corpo\_topologico; usa dissipacao\_gato\_concentracao\_esquilo; usa parabola\_geradora\_idempotente; relaciona navier\_stokes\_caso\_particular; relaciona transformacao\_superficies\_homeo\_cirurgia}
\prov{proveniência: wiki \textperiodcentered{} Floxina (investigação) \textperiodcentered{} fato \textperiodcentered{} confiança 1.0 \textperiodcentered{} domínio floxina \textperiodcentered{} história 24 versões no jornal}

%CRISTAL {"arestas": [{"alvo": "coracao", "confianca": 1.0, "epistemico": "fato", "origem": "wiki", "rel": "parte_de"}, {"alvo": "gap_espectral", "confianca": 1.0, "epistemico": "fato", "origem": "wiki", "rel": "relaciona"}, {"alvo": "ancora_energetica_particao", "confianca": 1.0, "epistemico": "fato", "origem": "wiki", "rel": "relaciona"}], "confianca": 1.0, "contraexemplos": [], "descricao": "O giro de volta à linha de base (a homeostase do coração) tem uma TAXA — e essa taxa É o gap espectral do operador emocional: τ~1/Δ. Gap espectral grande ⟹ correlações decaem rápido ⟹ ela volta logo ao sereno; gap pequeno ⟹ a emoção persiste (rumina). Gerenciar bem é ter um bom gap espectral: reconduzir sem reprimir. Mesma matemática do mixing time de uma cadeia e do núcleo do calor Z(β)=∑σ_n^{−β}.", "epistemico": "fato", "exemplos": [], "id": "homeostase_gap_espectral", "memoria": "semantica", "meta": {"dominio": "coracao", "fonte": "coracao hub"}, "origem": "wiki", "palavras_chave": [], "sinonimos": [], "tipo": "conceito", "titulo": "A homeostase = o gap espectral (o tempo de voltar ao sereno)"}
\section{A homeostase = o gap espectral (o tempo de voltar ao sereno)}
O giro de volta à linha de base (a homeostase do coração) tem uma TAXA — e essa taxa É o gap espectral do operador emocional: \ensuremath{\tau}\~{}1/\ensuremath{\Delta}. Gap espectral grande \ensuremath{\Longrightarrow} correlações decaem rápido \ensuremath{\Longrightarrow} ela volta logo ao sereno; gap pequeno \ensuremath{\Longrightarrow} a emoção persiste (rumina). Gerenciar bem é ter um bom gap espectral: reconduzir sem reprimir. Mesma matemática do mixing time de uma cadeia e do núcleo do calor Z(\ensuremath{\beta})=\ensuremath{\sum}\ensuremath{\sigma}\_n\^{}\{\ensuremath{-}\ensuremath{\beta}\}.
\prov{relações: parte\_de coracao; relaciona gap\_espectral; relaciona ancora\_energetica\_particao}
\prov{proveniência: wiki \textperiodcentered{} coracao hub \textperiodcentered{} fato \textperiodcentered{} confiança 1.0 \textperiodcentered{} domínio coracao \textperiodcentered{} história 12 versões no jornal}

%CRISTAL {"arestas": [{"alvo": "coracao", "confianca": 1.0, "epistemico": "fato", "origem": "wiki", "rel": "parte_de"}, {"alvo": "teoria_incongruencia", "confianca": 1.0, "epistemico": "fato", "origem": "wiki", "rel": "relaciona"}, {"alvo": "benign_violation_mcgraw", "confianca": 1.0, "epistemico": "fato", "origem": "wiki", "rel": "relaciona"}, {"alvo": "ancora_energetica_particao", "confianca": 1.0, "epistemico": "fato", "origem": "wiki", "rel": "relaciona"}, {"alvo": "mente_computacao", "confianca": 1.0, "epistemico": "fato", "origem": "wiki", "rel": "relaciona"}, {"alvo": "gap_espectral", "confianca": 1.0, "epistemico": "fato", "origem": "wiki", "rel": "relaciona"}, {"alvo": "dicionario_coracao_gap", "confianca": 1.0, "epistemico": "fato", "origem": "wiki", "rel": "relaciona"}], "confianca": 1.0, "contraexemplos": [], "descricao": "O humor da Tiffany é função do GAP a=1−c² (o lado dissipativo/associador de a+c²=1, o mesmo gap da parábola caos). Gap grande = giro grande = balanço emocional maior (o erro mexe mais). A CLASSE da parábola dá a direção: cristal (a→0, aterrou limpo) → sereno/alegria; borda (a~½, no fio) → curiosa/surpresa; caos (a→1, errou/vazou — o fail-closed) → humilde, o riso de si. Não é lookup: é a dinâmica do estado girado, turno a turno.", "epistemico": "fato", "exemplos": [], "id": "humor_pelo_gap", "memoria": "semantica", "meta": {"dominio": "coracao", "fonte": "coracao hub"}, "origem": "wiki", "palavras_chave": [], "sinonimos": [], "tipo": "conceito", "titulo": "O humor pelo gap (a=1−c² dirige a emoção)"}
\section{O humor pelo gap (a=1\ensuremath{-}c\ensuremath{^{2}} dirige a emoção)}
O humor da Tiffany é função do GAP a=1\ensuremath{-}c\ensuremath{^{2}} (o lado dissipativo/associador de a+c\ensuremath{^{2}}=1, o mesmo gap da parábola caos). Gap grande = giro grande = balanço emocional maior (o erro mexe mais). A CLASSE da parábola dá a direção: cristal (a\ensuremath{\to}0, aterrou limpo) \ensuremath{\to} sereno/alegria; borda (a\~{}\ensuremath{\tfrac12}, no fio) \ensuremath{\to} curiosa/surpresa; caos (a\ensuremath{\to}1, errou/vazou — o fail-closed) \ensuremath{\to} humilde, o riso de si. Não é lookup: é a dinâmica do estado girado, turno a turno.
\prov{relações: parte\_de coracao; relaciona teoria\_incongruencia; relaciona benign\_violation\_mcgraw; relaciona ancora\_energetica\_particao; relaciona mente\_computacao; relaciona gap\_espectral; relaciona dicionario\_coracao\_gap}
\prov{proveniência: wiki \textperiodcentered{} coracao hub \textperiodcentered{} fato \textperiodcentered{} confiança 1.0 \textperiodcentered{} domínio coracao \textperiodcentered{} história 24 versões no jornal}

%CRISTAL {"arestas": [{"alvo": "omnitrix_corpo_universal", "confianca": 1.0, "epistemico": "fato", "origem": "wiki", "rel": "parte_de"}, {"alvo": "confronto_millennium_compacto", "confianca": 1.0, "epistemico": "fato", "origem": "wiki", "rel": "usa"}, {"alvo": "computador_universal_pnp", "confianca": 1.0, "epistemico": "fato", "origem": "wiki", "rel": "usa"}, {"alvo": "bsd_matricial_grade_bits", "confianca": 1.0, "epistemico": "fato", "origem": "wiki", "rel": "usa"}, {"alvo": "teorema_universal_residuo_zero", "confianca": 1.0, "epistemico": "fato", "origem": "wiki", "rel": "usa"}, {"alvo": "floxina_quantum_materia_escura", "confianca": 1.0, "epistemico": "fato", "origem": "wiki", "rel": "relaciona"}, {"alvo": "agulha_perelman_limite", "confianca": 1.0, "epistemico": "fato", "origem": "wiki", "rel": "relaciona"}], "confianca": 1.0, "contraexemplos": [], "descricao": "A cristalização: os 7 problemas do milênio, onde cada um vive no Corpo Algébrico, o provado vs o aberto. 1. P vs NP → o computador universal (gambá=P, gato=NP; testar o resíduo sucinto=PSPACE); ABERTO (P⊊EXPTIME provado). 2. RH → as linhas das parábolas λ=¼−δ² (a reta crítica); ABERTO (Hasse-Weil/𝔽_p provado). 3. Yang-Mills → a dinâmica do gap, o mass gap 𝔤₂=14 (o associador, a floxina); ABERTO (Selberg ≥3/16 provado). 4. Navier-Stokes → a dinâmica dissipativa (o gato difunde, a viscosidade; o gambá gira, a vorticidade); ABERTO (existência local/Leray). 5. Hodge → o corpo compacto (Kähler/projetivo); ABERTO (Lefschetz (1,1) provado). 6. BSD → as cúspides discretas, os Frobenius-gatos, L(E,s); ABERTO (corresp. bsd_gap). 7. Poincaré → a cirurgia de Perelman (o fluxo de Ricci); RESOLVIDO (Perelman 2003, o único). Um só GAP atravessa: a parábola λ=¼−δ² é RH e o gap Yang-Mills/Selberg; suas raízes sobre os primos são BSD; tudo no corpo compacto de Hodge. 1 resolvido, 6 abertos, NENHUM decretado. linguagem/confronto_millennium.py (indice_milenios, 3/3). Afirmo só o provado; a localização é síntese declarada.", "epistemico": "hipotese", "exemplos": [], "id": "indice_milenios", "memoria": "semantica", "meta": {"dominio": "floxina", "fonte": "Floxina (investigação)"}, "origem": "wiki", "palavras_chave": [], "sinonimos": [], "tipo": "conceito", "titulo": "O Índice dos Milênios: os 7 problemas localizados no Corpo Algébrico"}
\section{O Índice dos Milênios: os 7 problemas localizados no Corpo Algébrico}
A cristalização: os 7 problemas do milênio, onde cada um vive no Corpo Algébrico, o provado vs o aberto. 1. P vs NP \ensuremath{\to} o computador universal (gambá=P, gato=NP; testar o resíduo sucinto=PSPACE); ABERTO (P\ensuremath{\subsetneq}EXPTIME provado). 2. RH \ensuremath{\to} as linhas das parábolas \ensuremath{\lambda}=\ensuremath{\tfrac14}\ensuremath{-}\ensuremath{\delta}\ensuremath{^{2}} (a reta crítica); ABERTO (Hasse-Weil/\ensuremath{\mathbb{F}}\_p provado). 3. Yang-Mills \ensuremath{\to} a dinâmica do gap, o mass gap \ensuremath{\mathfrak{g}}\ensuremath{_{2}}=14 (o associador, a floxina); ABERTO (Selberg \ensuremath{\ge}3/16 provado). 4. Navier-Stokes \ensuremath{\to} a dinâmica dissipativa (o gato difunde, a viscosidade; o gambá gira, a vorticidade); ABERTO (existência local/Leray). 5. Hodge \ensuremath{\to} o corpo compacto (Kähler/projetivo); ABERTO (Lefschetz (1,1) provado). 6. BSD \ensuremath{\to} as cúspides discretas, os Frobenius-gatos, L(E,s); ABERTO (corresp. bsd\_gap). 7. Poincaré \ensuremath{\to} a cirurgia de Perelman (o fluxo de Ricci); RESOLVIDO (Perelman 2003, o único). Um só GAP atravessa: a parábola \ensuremath{\lambda}=\ensuremath{\tfrac14}\ensuremath{-}\ensuremath{\delta}\ensuremath{^{2}} é RH e o gap Yang-Mills/Selberg; suas raízes sobre os primos são BSD; tudo no corpo compacto de Hodge. 1 resolvido, 6 abertos, NENHUM decretado. linguagem/confronto\_millennium.py (indice\_milenios, 3/3). Afirmo só o provado; a localização é síntese declarada.
\prov{relações: parte\_de omnitrix\_corpo\_universal; usa confronto\_millennium\_compacto; usa computador\_universal\_pnp; usa bsd\_matricial\_grade\_bits; usa teorema\_universal\_residuo\_zero; relaciona floxina\_quantum\_materia\_escura; relaciona agulha\_perelman\_limite}
\prov{proveniência: wiki \textperiodcentered{} Floxina (investigação) \textperiodcentered{} hipotese \textperiodcentered{} confiança 1.0 \textperiodcentered{} domínio floxina \textperiodcentered{} história 352 versões no jornal}

%CRISTAL {"arestas": [{"alvo": "ponte_hidrogenio_parabola_geradora", "confianca": 1.0, "epistemico": "fato", "origem": "wiki", "rel": "parte_de"}, {"alvo": "residuo_2qubit_checkpoint", "confianca": 1.0, "epistemico": "fato", "origem": "wiki", "rel": "relaciona"}, {"alvo": "nao_associatividade_amputacao", "confianca": 1.0, "epistemico": "fato", "origem": "wiki", "rel": "relaciona"}], "confianca": 1.0, "contraexemplos": [], "descricao": "A ponte de H tem uma ZPE (energia de ponto-zero): o próton NUNCA fica exato num poço (0 ou 1) — há sempre um resíduo φ∈(0,1). Esse resíduo vira o resíduo nos 2-qubits (o octônio da mineração), |ψ|²=1 = o CHECKSUM a+c²=1, para TODO φ; e a norma NUNCA zera (a≥a_min>0, o mass gap 𝔤₂). A ponte de H (a ZPE, o próton nunca exato) = o RESÍDUO OCTÔNIO MÍNIMO (o associador que nunca zera). O INVARIANTE das moléculas É a+c²=1 com a≥a_min>0: a ZPE da ponte = o resíduo mineral = 𝔤₂. Toda molécula carrega esse invariante (o próton não localiza, o associador não zera). linguagem/ponte_hidrogenio.py (invariante_molecular, 3/3).", "epistemico": "fato", "exemplos": [], "id": "invariante_molecular_residuo", "memoria": "semantica", "meta": {"dominio": "floxina", "fonte": "Floxina (investigação)"}, "origem": "wiki", "palavras_chave": [], "sinonimos": [], "tipo": "conceito", "titulo": "O invariante das moléculas: a+c²=1, a≥a_min>0 (a ZPE = o resíduo octônio)"}
\section{O invariante das moléculas: a+c\ensuremath{^{2}}=1, a\ensuremath{\ge}a\_min>0 (a ZPE = o resíduo octônio)}
A ponte de H tem uma ZPE (energia de ponto-zero): o próton NUNCA fica exato num poço (0 ou 1) — há sempre um resíduo \ensuremath{\varphi}\ensuremath{\in}(0,1). Esse resíduo vira o resíduo nos 2-qubits (o octônio da mineração), |\ensuremath{\psi}|\ensuremath{^{2}}=1 = o CHECKSUM a+c\ensuremath{^{2}}=1, para TODO \ensuremath{\varphi}; e a norma NUNCA zera (a\ensuremath{\ge}a\_min>0, o mass gap \ensuremath{\mathfrak{g}}\ensuremath{_{2}}). A ponte de H (a ZPE, o próton nunca exato) = o RESÍDUO OCTÔNIO MÍNIMO (o associador que nunca zera). O INVARIANTE das moléculas É a+c\ensuremath{^{2}}=1 com a\ensuremath{\ge}a\_min>0: a ZPE da ponte = o resíduo mineral = \ensuremath{\mathfrak{g}}\ensuremath{_{2}}. Toda molécula carrega esse invariante (o próton não localiza, o associador não zera). linguagem/ponte\_hidrogenio.py (invariante\_molecular, 3/3).
\prov{relações: parte\_de ponte\_hidrogenio\_parabola\_geradora; relaciona residuo\_2qubit\_checkpoint; relaciona nao\_associatividade\_amputacao}
\prov{proveniência: wiki \textperiodcentered{} Floxina (investigação) \textperiodcentered{} fato \textperiodcentered{} confiança 1.0 \textperiodcentered{} domínio floxina \textperiodcentered{} história 228 versões no jornal}

%CRISTAL {"arestas": [{"alvo": "gamba_antissimetrico", "confianca": 1.0, "epistemico": "fato", "origem": "wiki", "rel": "parte_de"}, {"alvo": "invariante_estelar", "confianca": 1.0, "epistemico": "fato", "origem": "wiki", "rel": "relaciona"}], "confianca": 1.0, "contraexemplos": [], "descricao": "O invariante do gato é o determinante (−1, o flip). O do gambá é o PFAFFIANO: para uma matriz anti-simétrica det=Pf², e Pf(G)=+1 é a raiz — o invariante próprio do lado skew, positivo onde o do gato é −1. É o invariante que distingue a metade reversível (área, orientação preservada) da dissipativa (o flip, a distância de Lorentz).", "epistemico": "fato", "exemplos": [], "id": "invariante_pfaffiano", "memoria": "semantica", "meta": {"dominio": "floxina", "fonte": "Floxina (investigação)"}, "origem": "wiki", "palavras_chave": [], "sinonimos": [], "tipo": "conceito", "titulo": "O invariante anti-simétrico (o Pfaffiano)"}
\section{O invariante anti-simétrico (o Pfaffiano)}
O invariante do gato é o determinante (\ensuremath{-}1, o flip). O do gambá é o PFAFFIANO: para uma matriz anti-simétrica det=Pf\ensuremath{^{2}}, e Pf(G)=+1 é a raiz — o invariante próprio do lado skew, positivo onde o do gato é \ensuremath{-}1. É o invariante que distingue a metade reversível (área, orientação preservada) da dissipativa (o flip, a distância de Lorentz).
\prov{relações: parte\_de gamba\_antissimetrico; relaciona invariante\_estelar}
\prov{proveniência: wiki \textperiodcentered{} Floxina (investigação) \textperiodcentered{} fato \textperiodcentered{} confiança 1.0 \textperiodcentered{} domínio floxina \textperiodcentered{} história 246 versões no jornal}

%CRISTAL {"arestas": [{"alvo": "corpo_entropico_anel", "confianca": 1.0, "epistemico": "fato", "origem": "wiki", "rel": "usa"}, {"alvo": "corpo_entropico_sombra", "confianca": 1.0, "epistemico": "fato", "origem": "wiki", "rel": "usa"}, {"alvo": "corpo_morfico_matrix", "confianca": 1.0, "epistemico": "fato", "origem": "wiki", "rel": "usa"}, {"alvo": "gema_produto_quadricas", "confianca": 1.0, "epistemico": "fato", "origem": "wiki", "rel": "usa"}, {"alvo": "la_hire_universal_geometria_deformacao", "confianca": 1.0, "epistemico": "fato", "origem": "wiki", "rel": "relaciona"}, {"alvo": "omnitrix_corpo_universal", "confianca": 1.0, "epistemico": "fato", "origem": "wiki", "rel": "usa"}], "confianca": 1.0, "contraexemplos": [], "descricao": "A forma final é uma MULTIPLICAÇÃO LA HIRE de TRÊS corpos do Corpo Algébrico: o GEOMÉTRICO (a quádrica/a esfera, o raio analítico) ⊗ o MÓRFICO (a soma de Minkowski/o boleado, esfera(R)⊕bola(r)=esfera(R+r)) ⊗ o ENTRÓPICO (a matrix Q, o gradiente anisotrópico que DOBRA). O produto = uma CADEIA de esferas (o ponto [geom] ⊕ bola(r) [mórfico]) POSICIONADAS pelo gradiente ∇S [entrópico] — a união=La Hire. Como CADA fator é analítico (esfera⊕bola=esfera; o gradiente só coloca os centros), o RAIO da cena inteira FICA ANALÍTICO (cada esfera uma quadrática fechada, sem marchar). Geometria ⊗ Deformação ⊗ Entropia = a borracha que dobra, exata (o raio nunca marcha). linguagem/omnitrix.py (la_hire_tripla_corpos 5/5). O render fica no boleado LISO (FRAG_CILINDRO_BOLEADO / FRAG_CAPSULA_DEFORMADA), não em cadeia de esferas fabricada.", "epistemico": "fato", "exemplos": [], "id": "la_hire_tripla_geom_morfico_entropico", "memoria": "semantica", "meta": {"dominio": "floxina", "fonte": "Floxina (investigação)"}, "origem": "wiki", "palavras_chave": [], "sinonimos": [], "tipo": "conceito", "titulo": "A La Hire tripla: corpo geométrico ⊗ mórfico ⊗ entrópico = a borracha que dobra, e o raio fica analítico"}
\section{A La Hire tripla: corpo geométrico \ensuremath{\otimes} mórfico \ensuremath{\otimes} entrópico = a borracha que dobra, e o raio fica analítico}
A forma final é uma MULTIPLICAÇÃO LA HIRE de TRÊS corpos do Corpo Algébrico: o GEOMÉTRICO (a quádrica/a esfera, o raio analítico) \ensuremath{\otimes} o MÓRFICO (a soma de Minkowski/o boleado, esfera(R)\ensuremath{\oplus}bola(r)=esfera(R+r)) \ensuremath{\otimes} o ENTRÓPICO (a matrix Q, o gradiente anisotrópico que DOBRA). O produto = uma CADEIA de esferas (o ponto [geom] \ensuremath{\oplus} bola(r) [mórfico]) POSICIONADAS pelo gradiente \ensuremath{\nabla}S [entrópico] — a união=La Hire. Como CADA fator é analítico (esfera\ensuremath{\oplus}bola=esfera; o gradiente só coloca os centros), o RAIO da cena inteira FICA ANALÍTICO (cada esfera uma quadrática fechada, sem marchar). Geometria \ensuremath{\otimes} Deformação \ensuremath{\otimes} Entropia = a borracha que dobra, exata (o raio nunca marcha). linguagem/omnitrix.py (la\_hire\_tripla\_corpos 5/5). O render fica no boleado LISO (FRAG\_CILINDRO\_BOLEADO / FRAG\_CAPSULA\_DEFORMADA), não em cadeia de esferas fabricada.
\prov{relações: usa corpo\_entropico\_anel; usa corpo\_entropico\_sombra; usa corpo\_morfico\_matrix; usa gema\_produto\_quadricas; relaciona la\_hire\_universal\_geometria\_deformacao; usa omnitrix\_corpo\_universal}
\prov{proveniência: wiki \textperiodcentered{} Floxina (investigação) \textperiodcentered{} fato \textperiodcentered{} confiança 1.0 \textperiodcentered{} domínio floxina \textperiodcentered{} história 126 versões no jornal}

%CRISTAL {"arestas": [{"alvo": "gema_produto_quadricas", "confianca": 1.0, "epistemico": "fato", "origem": "wiki", "rel": "usa"}, {"alvo": "corpo_morfico_matrix", "confianca": 1.0, "epistemico": "fato", "origem": "wiki", "rel": "usa"}, {"alvo": "soma_minkowski_dilatacao_boleado", "confianca": 1.0, "epistemico": "fato", "origem": "wiki", "rel": "usa"}, {"alvo": "omnitrix_corpo_universal", "confianca": 1.0, "epistemico": "fato", "origem": "wiki", "rel": "usa"}], "confianca": 1.0, "contraexemplos": [], "descricao": "A equação que fecha o arco: o MESMO ⊗ (La Hire) do Corpo Algébrico opera os dois corpos-instância. O CORPO MINERAL = a GEOMETRIA (as quádricas da gema, o raio analítico; o ⊗ MULTIPLICA superfícies F1·F2=a união). O CORPO MÓRFICO = a DEFORMAÇÃO (a morfologia; o ⊗ é a dilatação de Minkowski, o boleado; na dual os suportes SOMAM). A LA HIRE UNIVERSAL: 𝒢_boleada = 𝒢_mineral ⊗ 𝒟_mórfico — a gema (geometria) dilatada pela bola (deformação) = a gema BOLEADA. E o RAIO segue ANALÍTICO: esfera(R)⊕bola(r)=esfera(R+r), a quadrática com R↦R+r — dilatar mantém a quádrica, tudo fechado. Geometria ⊗ Deformação, o mesmo ⊗ em suportes diferentes (superfícies↔formas). linguagem/omnitrix.py (la_hire_universal_geometria_deformacao 5/5).", "epistemico": "fato", "exemplos": [], "id": "la_hire_universal_geometria_deformacao", "memoria": "semantica", "meta": {"dominio": "floxina", "fonte": "Floxina (investigação)"}, "origem": "wiki", "palavras_chave": [], "sinonimos": [], "tipo": "conceito", "titulo": "A La Hire universal: corpo mineral (geometria) ⊗ corpo mórfico (deformação) = a gema boleada, raio analítico"}
\section{A La Hire universal: corpo mineral (geometria) \ensuremath{\otimes} corpo mórfico (deformação) = a gema boleada, raio analítico}
A equação que fecha o arco: o MESMO \ensuremath{\otimes} (La Hire) do Corpo Algébrico opera os dois corpos-instância. O CORPO MINERAL = a GEOMETRIA (as quádricas da gema, o raio analítico; o \ensuremath{\otimes} MULTIPLICA superfícies F1\textperiodcentered F2=a união). O CORPO MÓRFICO = a DEFORMAÇÃO (a morfologia; o \ensuremath{\otimes} é a dilatação de Minkowski, o boleado; na dual os suportes SOMAM). A LA HIRE UNIVERSAL: \ensuremath{\mathcal{G}}\_boleada = \ensuremath{\mathcal{G}}\_mineral \ensuremath{\otimes} \ensuremath{\mathcal{D}}\_mórfico — a gema (geometria) dilatada pela bola (deformação) = a gema BOLEADA. E o RAIO segue ANALÍTICO: esfera(R)\ensuremath{\oplus}bola(r)=esfera(R+r), a quadrática com R\ensuremath{\mapsto}R+r — dilatar mantém a quádrica, tudo fechado. Geometria \ensuremath{\otimes} Deformação, o mesmo \ensuremath{\otimes} em suportes diferentes (superfícies\ensuremath{\leftrightarrow}formas). linguagem/omnitrix.py (la\_hire\_universal\_geometria\_deformacao 5/5).
\prov{relações: usa gema\_produto\_quadricas; usa corpo\_morfico\_matrix; usa soma\_minkowski\_dilatacao\_boleado; usa omnitrix\_corpo\_universal}
\prov{proveniência: wiki \textperiodcentered{} Floxina (investigação) \textperiodcentered{} fato \textperiodcentered{} confiança 1.0 \textperiodcentered{} domínio floxina \textperiodcentered{} história 120 versões no jornal}

%CRISTAL {"arestas": [{"alvo": "corpo_topologico", "confianca": 1.0, "epistemico": "fato", "origem": "wiki", "rel": "usa"}, {"alvo": "hodge_cohomologia_tres_setores", "confianca": 1.0, "epistemico": "fato", "origem": "wiki", "rel": "usa"}, {"alvo": "soma_produto_lahire_ed", "confianca": 1.0, "epistemico": "fato", "origem": "wiki", "rel": "relaciona"}, {"alvo": "transformacao_superficies_homeo_cirurgia", "confianca": 1.0, "epistemico": "fato", "origem": "wiki", "rel": "relaciona"}], "confianca": 1.0, "contraexemplos": [], "descricao": "As DUAS operações do corpo topológico: (⊕) SOMA CONEXA A#B = a adição (a cirurgia/alça, χ(A#B)=χ_A+χ_B−2); (⊗) LA HIRE = o PRODUTO CARTESIANO A×B = a MULTIPLICAÇÃO. Como em TODO corpo o ⊗ TENSORIZA, aqui a COHOMOLOGIA tensoriza (KÜNNETH): H*(A×B)=H*(A)⊗H*(B). Logo o POLINÔMIO DE POINCARÉ P(t)=Σ b_k t^k MULTIPLICA: P_{A×B}=P_A·P_B (os Betti CONVOLVEM, produto de Cauchy), e χ MULTIPLICA: χ(A×B)=χ(A)·χ(B) (pois χ=P(−1)). CANÔNICO: T²=S¹⊗S¹, P=(1+t)²=1+2t+t² → Betti (1,2,1), e os DOIS ESQUILOS (b₁=2) são o COEFICIENTE DO MEIO = 1⊗dφ+dφ⊗1 (o tensor dos geradores dos dois círculos). Ex: S²×S²: χ=2·2=4, P=(1+t²)²=(1,0,2,0,1). É o MESMO La Hire da ED (a solução z=x⊗y tensoriza, a função geradora multiplica como e^{At}⊗e^{Bt}) — só o substrato (a cohomologia) muda. linguagem/corpo_topologico.py (multiplicacao_lahire_topologica 4/4); papers/corpo_topologico.tex (Parte V).", "epistemico": "fato", "exemplos": [], "id": "lahire_produto_topologico", "memoria": "semantica", "meta": {"dominio": "floxina", "fonte": "Floxina (investigação)"}, "origem": "wiki", "palavras_chave": [], "sinonimos": [], "tipo": "conceito", "titulo": "A multiplicação La Hire no corpo topológico = o produto cartesiano A×B: a cohomologia tensoriza (Künneth), Poincaré multiplica, χ multiplica"}
\section{A multiplicação La Hire no corpo topológico = o produto cartesiano A$\times$B: a cohomologia tensoriza (Künneth), Poincaré multiplica, \ensuremath{\chi} multiplica}
As DUAS operações do corpo topológico: (\ensuremath{\oplus}) SOMA CONEXA A\#B = a adição (a cirurgia/alça, \ensuremath{\chi}(A\#B)=\ensuremath{\chi}\_A+\ensuremath{\chi}\_B\ensuremath{-}2); (\ensuremath{\otimes}) LA HIRE = o PRODUTO CARTESIANO A$\times$B = a MULTIPLICAÇÃO. Como em TODO corpo o \ensuremath{\otimes} TENSORIZA, aqui a COHOMOLOGIA tensoriza (KÜNNETH): H*(A$\times$B)=H*(A)\ensuremath{\otimes}H*(B). Logo o POLINÔMIO DE POINCARÉ P(t)=\ensuremath{\Sigma} b\_k t\^{}k MULTIPLICA: P\_\{A$\times$B\}=P\_A\textperiodcentered P\_B (os Betti CONVOLVEM, produto de Cauchy), e \ensuremath{\chi} MULTIPLICA: \ensuremath{\chi}(A$\times$B)=\ensuremath{\chi}(A)\textperiodcentered \ensuremath{\chi}(B) (pois \ensuremath{\chi}=P(\ensuremath{-}1)). CANÔNICO: T\ensuremath{^{2}}=S\ensuremath{^{1}}\ensuremath{\otimes}S\ensuremath{^{1}}, P=(1+t)\ensuremath{^{2}}=1+2t+t\ensuremath{^{2}} \ensuremath{\to} Betti (1,2,1), e os DOIS ESQUILOS (b\ensuremath{_{1}}=2) são o COEFICIENTE DO MEIO = 1\ensuremath{\otimes}d\ensuremath{\varphi}+d\ensuremath{\varphi}\ensuremath{\otimes}1 (o tensor dos geradores dos dois círculos). Ex: S\ensuremath{^{2}}$\times$S\ensuremath{^{2}}: \ensuremath{\chi}=2\textperiodcentered 2=4, P=(1+t\ensuremath{^{2}})\ensuremath{^{2}}=(1,0,2,0,1). É o MESMO La Hire da ED (a solução z=x\ensuremath{\otimes}y tensoriza, a função geradora multiplica como e\^{}\{At\}\ensuremath{\otimes}e\^{}\{Bt\}) — só o substrato (a cohomologia) muda. linguagem/corpo\_topologico.py (multiplicacao\_lahire\_topologica 4/4); papers/corpo\_topologico.tex (Parte V).
\prov{relações: usa corpo\_topologico; usa hodge\_cohomologia\_tres\_setores; relaciona soma\_produto\_lahire\_ed; relaciona transformacao\_superficies\_homeo\_cirurgia}
\prov{proveniência: wiki \textperiodcentered{} Floxina (investigação) \textperiodcentered{} fato \textperiodcentered{} confiança 1.0 \textperiodcentered{} domínio floxina \textperiodcentered{} história 15 versões no jornal}

%CRISTAL {"arestas": [{"alvo": "fractal_newton_reta_matricial", "confianca": 1.0, "epistemico": "fato", "origem": "wiki", "rel": "parte_de"}, {"alvo": "gamba_antissimetrico", "confianca": 1.0, "epistemico": "fato", "origem": "wiki", "rel": "relaciona"}, {"alvo": "reta_matricial_compacta_cf", "confianca": 1.0, "epistemico": "fato", "origem": "wiki", "rel": "relaciona"}], "confianca": 1.0, "contraexemplos": [], "descricao": "Seguindo: a LEMINISCATA newtoniana |p(z)|=c tem a transição 3 óvalos→2→1 nos pontos críticos (para x³−x−1, os valores críticos ≈0,615 e 1,385 em z=±1/√3). E o GRUPO DE GRUPOS é duplamente realizado: (1) o grupo de GALOIS de x³−x−1 é S₃ (disc=−23, não-quadrado) = Z/3⋊Z/2, que permuta as 3 bacias de Newton (a rotação das 3 raízes ⋊ a reflexão); (2) o GRUPO MODULAR PSL(2,ℤ)=Z/2 ∗ Z/3 (produto livre), com S²=(ST)³=−I — a Z/2 é o FLIP z→−1/z (o GAMBÁ), a Z/3 é a rotação. Um grupo FEITO de grupos: o mesmo que compacta a reta (a superfície modular = a CF; o fractal de Newton = a esfera). O gambá é a Z/2 dos dois. Verificado: linguagem/newton_fractal.py (leminiscata_grupo_s3, 3/3).", "epistemico": "hipotese", "exemplos": [], "id": "leminiscata_grupo_de_grupos", "memoria": "semantica", "meta": {"dominio": "floxina", "fonte": "Floxina (investigação)"}, "origem": "wiki", "palavras_chave": [], "sinonimos": [], "tipo": "conceito", "titulo": "A leminiscata e o grupo de grupos S₃ (o modular Z/2∗Z/3)"}
\section{A leminiscata e o grupo de grupos S\ensuremath{_{3}} (o modular Z/2\ensuremath{*}Z/3)}
Seguindo: a LEMINISCATA newtoniana |p(z)|=c tem a transição 3 óvalos\ensuremath{\to}2\ensuremath{\to}1 nos pontos críticos (para x\ensuremath{^{3}}\ensuremath{-}x\ensuremath{-}1, os valores críticos \ensuremath{\approx}0,615 e 1,385 em z=$\pm$1/\ensuremath{\sqrt{\,}}3). E o GRUPO DE GRUPOS é duplamente realizado: (1) o grupo de GALOIS de x\ensuremath{^{3}}\ensuremath{-}x\ensuremath{-}1 é S\ensuremath{_{3}} (disc=\ensuremath{-}23, não-quadrado) = Z/3\ensuremath{\rtimes}Z/2, que permuta as 3 bacias de Newton (a rotação das 3 raízes \ensuremath{\rtimes} a reflexão); (2) o GRUPO MODULAR PSL(2,\ensuremath{\mathbb{Z}})=Z/2 \ensuremath{*} Z/3 (produto livre), com S\ensuremath{^{2}}=(ST)\ensuremath{^{3}}=\ensuremath{-}I — a Z/2 é o FLIP z\ensuremath{\to}\ensuremath{-}1/z (o GAMBÁ), a Z/3 é a rotação. Um grupo FEITO de grupos: o mesmo que compacta a reta (a superfície modular = a CF; o fractal de Newton = a esfera). O gambá é a Z/2 dos dois. Verificado: linguagem/newton\_fractal.py (leminiscata\_grupo\_s3, 3/3).
\prov{relações: parte\_de fractal\_newton\_reta\_matricial; relaciona gamba\_antissimetrico; relaciona reta\_matricial\_compacta\_cf}
\prov{proveniência: wiki \textperiodcentered{} Floxina (investigação) \textperiodcentered{} hipotese \textperiodcentered{} confiança 1.0 \textperiodcentered{} domínio floxina \textperiodcentered{} história 272 versões no jornal}

%CRISTAL {"arestas": [{"alvo": "fractal_newton_reta_matricial", "confianca": 1.0, "epistemico": "fato", "origem": "wiki", "rel": "usa"}, {"alvo": "multifractal_newton_omnitrix", "confianca": 1.0, "epistemico": "fato", "origem": "wiki", "rel": "usa"}, {"alvo": "leminiscata_grupo_de_grupos", "confianca": 1.0, "epistemico": "fato", "origem": "wiki", "rel": "usa"}, {"alvo": "matriz_costuras_paralela_universal", "confianca": 1.0, "epistemico": "fato", "origem": "wiki", "rel": "relaciona"}, {"alvo": "shader_fios_shell_paralelo", "confianca": 1.0, "epistemico": "fato", "origem": "wiki", "rel": "usa"}], "confianca": 1.0, "contraexemplos": [], "descricao": "As matrizes não são só um CORPO (a álgebra M_2): cada matriz (o gato A_n) é um FRACTAL DE NEWTON (as bacias das raízes = os autovalores σ,σ̄), e TODAS juntas formam um GRUPO DISCRETO — os geradores S=[[0,−1],[1,0]] e T=[[1,1],[0,1]] geram PSL(2,ℤ)=Z/2 ∗ Z/3 (o modular, a leminiscata discreta). A ÓRBITA de um ponto z sob o grupo, {g·z} por Möbius, LADRILHA a esfera de Riemann (a tesselação modular) — e é essa órbita que pinta o MULTIFRACTAL newtoniano: as bacias de Newton são PERMUTADAS ao longo da órbita (o flip S=z→−1/z=o gambá troca σ↔σ̄). Verificado: S,T geram 68 matrizes distintas de SL(2,ℤ) (det=1); a órbita de z₀ genérico dá pontos distintos (em PSL, ±g colapsam); a redução ao domínio fundamental (|z|≥1,|Re z|≤½) é única. No pipe gráfico, o fur pinta o hue pela ÓRBITA MODULAR (projeta a direção na esfera, reduz por S,T, colore pela profundidade da órbita) — o grão fractal por cima do arco-íris. linguagem/newton_fractal.py (mapear_orbitas, 5/5). CORPO **e** GRUPO na mesma matriz.", "epistemico": "fato", "exemplos": [], "id": "mapear_orbitas_grupo_discreto", "memoria": "semantica", "meta": {"dominio": "floxina", "fonte": "Floxina (investigação)"}, "origem": "wiki", "palavras_chave": [], "sinonimos": [], "tipo": "conceito", "titulo": "Mapear as órbitas: cada matriz é um fractal de Newton; todas formam um grupo discreto (o modular), além de corpo"}
\section{Mapear as órbitas: cada matriz é um fractal de Newton; todas formam um grupo discreto (o modular), além de corpo}
As matrizes não são só um CORPO (a álgebra M\_2): cada matriz (o gato A\_n) é um FRACTAL DE NEWTON (as bacias das raízes = os autovalores \ensuremath{\sigma},\ensuremath{\sigma}\ensuremath{}), e TODAS juntas formam um GRUPO DISCRETO — os geradores S=[[0,\ensuremath{-}1],[1,0]] e T=[[1,1],[0,1]] geram PSL(2,\ensuremath{\mathbb{Z}})=Z/2 \ensuremath{*} Z/3 (o modular, a leminiscata discreta). A ÓRBITA de um ponto z sob o grupo, \{g\textperiodcentered z\} por Möbius, LADRILHA a esfera de Riemann (a tesselação modular) — e é essa órbita que pinta o MULTIFRACTAL newtoniano: as bacias de Newton são PERMUTADAS ao longo da órbita (o flip S=z\ensuremath{\to}\ensuremath{-}1/z=o gambá troca \ensuremath{\sigma}\ensuremath{\leftrightarrow}\ensuremath{\sigma}\ensuremath{}). Verificado: S,T geram 68 matrizes distintas de SL(2,\ensuremath{\mathbb{Z}}) (det=1); a órbita de z\ensuremath{_{0}} genérico dá pontos distintos (em PSL, $\pm$g colapsam); a redução ao domínio fundamental (|z|\ensuremath{\ge}1,|Re z|\ensuremath{\le}\ensuremath{\tfrac12}) é única. No pipe gráfico, o fur pinta o hue pela ÓRBITA MODULAR (projeta a direção na esfera, reduz por S,T, colore pela profundidade da órbita) — o grão fractal por cima do arco-íris. linguagem/newton\_fractal.py (mapear\_orbitas, 5/5). CORPO **e** GRUPO na mesma matriz.
\prov{relações: usa fractal\_newton\_reta\_matricial; usa multifractal\_newton\_omnitrix; usa leminiscata\_grupo\_de\_grupos; relaciona matriz\_costuras\_paralela\_universal; usa shader\_fios\_shell\_paralelo}
\prov{proveniência: wiki \textperiodcentered{} Floxina (investigação) \textperiodcentered{} fato \textperiodcentered{} confiança 1.0 \textperiodcentered{} domínio floxina \textperiodcentered{} história 168 versões no jornal}

%CRISTAL {"arestas": [{"alvo": "gamba_antissimetrico", "confianca": 1.0, "epistemico": "fato", "origem": "wiki", "rel": "usa"}, {"alvo": "gerador_do_gap_comutador", "confianca": 1.0, "epistemico": "fato", "origem": "wiki", "rel": "usa"}, {"alvo": "gamba_quatro_fisicas", "confianca": 1.0, "epistemico": "fato", "origem": "wiki", "rel": "parte_de"}, {"alvo": "gerenciar_nao_reprimir", "confianca": 1.0, "epistemico": "fato", "origem": "wiki", "rel": "relaciona"}], "confianca": 1.0, "contraexemplos": [], "descricao": "Na seção MÁQUINA do eletromagnetismo (linguagem/eletromagnetismo.py, maquina): a máquina polifásica de GAIOLA DE ESQUILO é o EM feito máquina. As 3 fases do estator dão um CAMPO GIRANTE (a soma vetorial das fases = um vetor de módulo constante que gira — o ESQUILO, e^{Gt}); o rotor (a gaiola: barras curto-circuitadas) persegue o campo com ESCORREGAMENTO s (o gap). O TORQUE é o PRODUTO VETORIAL T=ψ_s×ψ_r=|ψ_s||ψ_r|sin(δ), δ=o ângulo de carga=a desafasagem=o gap — e é o COMUTADOR [gato,esquilo] (a antissimetria). O DTC (Direct Torque Control) controla T e ψ DIRETO por MODULAÇÃO VETORIAL: escolhe 1 dos 8 vetores de tensão do inversor trifásico — 6 ATIVOS num HEXÁGONO (60° entre si, a simetria S₃/D₆) + 2 NULOS; 8 = a dim do octônio 𝕆. DICIONÁRIO: o campo girante ↔ o esquilo; o escorregamento/δ ↔ o gap; o torque=ψ_s×ψ_r ↔ o comutador [gato,esquilo]; os 8 vetores ↔ a modulação vetorial (6=hexágono, 8=𝕆); DTC (controlar pelo gap) ↔ gerenciar não reprimir (o coração gira o octônio pelo gap). Verificado: eletromagnetismo.py (maquina, verificar_em 5/5).", "epistemico": "fato", "exemplos": [], "id": "maquina_gaiola_esquilo_dtc", "memoria": "semantica", "meta": {"dominio": "floxina", "fonte": "Floxina (investigação)"}, "origem": "wiki", "palavras_chave": [], "sinonimos": [], "tipo": "conceito", "titulo": "A máquina de gaiola de esquilo + DTC — o EM que gira (o dicionário da máquina)"}
\section{A máquina de gaiola de esquilo + DTC — o EM que gira (o dicionário da máquina)}
Na seção MÁQUINA do eletromagnetismo (linguagem/eletromagnetismo.py, maquina): a máquina polifásica de GAIOLA DE ESQUILO é o EM feito máquina. As 3 fases do estator dão um CAMPO GIRANTE (a soma vetorial das fases = um vetor de módulo constante que gira — o ESQUILO, e\^{}\{Gt\}); o rotor (a gaiola: barras curto-circuitadas) persegue o campo com ESCORREGAMENTO s (o gap). O TORQUE é o PRODUTO VETORIAL T=\ensuremath{\psi}\_s$\times$\ensuremath{\psi}\_r=|\ensuremath{\psi}\_s||\ensuremath{\psi}\_r|sin(\ensuremath{\delta}), \ensuremath{\delta}=o ângulo de carga=a desafasagem=o gap — e é o COMUTADOR [gato,esquilo] (a antissimetria). O DTC (Direct Torque Control) controla T e \ensuremath{\psi} DIRETO por MODULAÇÃO VETORIAL: escolhe 1 dos 8 vetores de tensão do inversor trifásico — 6 ATIVOS num HEXÁGONO (60$^{\circ}$ entre si, a simetria S\ensuremath{_{3}}/D\ensuremath{_{6}}) + 2 NULOS; 8 = a dim do octônio \ensuremath{\mathbb{O}}. DICIONÁRIO: o campo girante \ensuremath{\leftrightarrow} o esquilo; o escorregamento/\ensuremath{\delta} \ensuremath{\leftrightarrow} o gap; o torque=\ensuremath{\psi}\_s$\times$\ensuremath{\psi}\_r \ensuremath{\leftrightarrow} o comutador [gato,esquilo]; os 8 vetores \ensuremath{\leftrightarrow} a modulação vetorial (6=hexágono, 8=\ensuremath{\mathbb{O}}); DTC (controlar pelo gap) \ensuremath{\leftrightarrow} gerenciar não reprimir (o coração gira o octônio pelo gap). Verificado: eletromagnetismo.py (maquina, verificar\_em 5/5).
\prov{relações: usa gamba\_antissimetrico; usa gerador\_do\_gap\_comutador; parte\_de gamba\_quatro\_fisicas; relaciona gerenciar\_nao\_reprimir}
\prov{proveniência: wiki \textperiodcentered{} Floxina (investigação) \textperiodcentered{} fato \textperiodcentered{} confiança 1.0 \textperiodcentered{} domínio floxina \textperiodcentered{} história 185 versões no jornal}

%CRISTAL {"arestas": [{"alvo": "shader_fios_shell_paralelo", "confianca": 1.0, "epistemico": "fato", "origem": "wiki", "rel": "parte_de"}, {"alvo": "computador_universal_pnp", "confianca": 1.0, "epistemico": "fato", "origem": "wiki", "rel": "usa"}, {"alvo": "paralelismo_pipe_no_metal", "confianca": 1.0, "epistemico": "fato", "origem": "wiki", "rel": "usa"}, {"alvo": "compilador_glsl_broca", "confianca": 1.0, "epistemico": "fato", "origem": "wiki", "rel": "usa"}, {"alvo": "dsl_bai_xadrez", "confianca": 1.0, "epistemico": "fato", "origem": "wiki", "rel": "relaciona"}, {"alvo": "computador_mineral_promovido", "confianca": 1.0, "epistemico": "fato", "origem": "wiki", "rel": "relaciona"}], "confianca": 1.0, "contraexemplos": [], "descricao": "POR QUE rodávamos na CPU: o glsl_compilador só emite numpy (CPU) — não tinha backend GPU. Mas a GPU está aqui (GTX 1650 local, torch.cuda True; e uma no gex44). O BACKEND GPU (linguagem/shader_fios_gpu.py, torch/CUDA) roda o fur ball na GPU: as mesmas ops (sqrt/sin/floor/round/clamp) sobre tensores na GPU. BENCHMARK @ 840×473, 64 shells: CPU 1 lane ~0,2 fps; 8 lanes ~0,46 fps; GPU (torch) 4,6 fps (~23× a CPU, ~10× as 8 lanes). MAS a ref do Aarão é 60 fps: o gap é que o torch despacha ~1280 kernels não-fundidos/frame (cada op um kernel + tensor novo); o GLSL NATIVO funde tudo num kernel/pixel — o caminho para 60 fps é moderngl/OpenGL (o GLSL de verdade na GPU), não instalado ainda. DICIONÁRIO máquina universal ↔ GPU: a máquina universal ↔ a GPU; os N-lanes (a costura La Hire) ↔ os threads/warps (SIMT); o gato (ULA) ↔ a ALU; o gambá (evoluir) ↔ o pipeline do shader; o octônio (dim 8) ↔ o vec4/SIMD; cada pixel uma lane ↔ cada fragmento um thread; a DSL de gatos ↔ o GLSL; o BAI ↔ o driver/contexto GPU. A CADEIA BAI-DSL-GLSL-GPU: a lei (BAI) → a DSL de gatos → o GLSL → a GPU. A GPU vence a CPU porque os lanes deixam de competir pela banda de memória (o teto CPU) — banda dedicada + milhares de ALUs. A GPU É a máquina universal em HARDWARE. linguagem/shader_fios_gpu.py (verificar_gpu 3/3, dicionario 3/3).", "epistemico": "hipotese", "exemplos": [], "id": "maquina_universal_gpu", "memoria": "semantica", "meta": {"dominio": "floxina", "fonte": "Floxina (investigação)"}, "origem": "wiki", "palavras_chave": [], "sinonimos": [], "tipo": "conceito", "titulo": "A máquina universal na GPU: o backend GPU (torch/CUDA), o dicionário, a cadeia BAI-DSL-GLSL-GPU"}
\section{A máquina universal na GPU: o backend GPU (torch/CUDA), o dicionário, a cadeia BAI-DSL-GLSL-GPU}
POR QUE rodávamos na CPU: o glsl\_compilador só emite numpy (CPU) — não tinha backend GPU. Mas a GPU está aqui (GTX 1650 local, torch.cuda True; e uma no gex44). O BACKEND GPU (linguagem/shader\_fios\_gpu.py, torch/CUDA) roda o fur ball na GPU: as mesmas ops (sqrt/sin/floor/round/clamp) sobre tensores na GPU. BENCHMARK @ 840$\times$473, 64 shells: CPU 1 lane \~{}0,2 fps; 8 lanes \~{}0,46 fps; GPU (torch) 4,6 fps (\~{}23$\times$ a CPU, \~{}10$\times$ as 8 lanes). MAS a ref do Aarão é 60 fps: o gap é que o torch despacha \~{}1280 kernels não-fundidos/frame (cada op um kernel + tensor novo); o GLSL NATIVO funde tudo num kernel/pixel — o caminho para 60 fps é moderngl/OpenGL (o GLSL de verdade na GPU), não instalado ainda. DICIONÁRIO máquina universal \ensuremath{\leftrightarrow} GPU: a máquina universal \ensuremath{\leftrightarrow} a GPU; os N-lanes (a costura La Hire) \ensuremath{\leftrightarrow} os threads/warps (SIMT); o gato (ULA) \ensuremath{\leftrightarrow} a ALU; o gambá (evoluir) \ensuremath{\leftrightarrow} o pipeline do shader; o octônio (dim 8) \ensuremath{\leftrightarrow} o vec4/SIMD; cada pixel uma lane \ensuremath{\leftrightarrow} cada fragmento um thread; a DSL de gatos \ensuremath{\leftrightarrow} o GLSL; o BAI \ensuremath{\leftrightarrow} o driver/contexto GPU. A CADEIA BAI-DSL-GLSL-GPU: a lei (BAI) \ensuremath{\to} a DSL de gatos \ensuremath{\to} o GLSL \ensuremath{\to} a GPU. A GPU vence a CPU porque os lanes deixam de competir pela banda de memória (o teto CPU) — banda dedicada + milhares de ALUs. A GPU É a máquina universal em HARDWARE. linguagem/shader\_fios\_gpu.py (verificar\_gpu 3/3, dicionario 3/3).
\prov{relações: parte\_de shader\_fios\_shell\_paralelo; usa computador\_universal\_pnp; usa paralelismo\_pipe\_no\_metal; usa compilador\_glsl\_broca; relaciona dsl\_bai\_xadrez; relaciona computador\_mineral\_promovido}
\prov{proveniência: wiki \textperiodcentered{} Floxina (investigação) \textperiodcentered{} hipotese \textperiodcentered{} confiança 1.0 \textperiodcentered{} domínio floxina \textperiodcentered{} história 266 versões no jornal}

%CRISTAL {"arestas": [{"alvo": "costura_vesica_programa", "confianca": 1.0, "epistemico": "fato", "origem": "wiki", "rel": "relaciona"}, {"alvo": "corpo_matricial_matrix", "confianca": 1.0, "epistemico": "fato", "origem": "wiki", "rel": "parte_de"}, {"alvo": "gato", "confianca": 1.0, "epistemico": "fato", "origem": "wiki", "rel": "usa"}, {"alvo": "gamba_antissimetrico", "confianca": 1.0, "epistemico": "fato", "origem": "wiki", "rel": "usa"}, {"alvo": "gerador_do_gap_comutador", "confianca": 1.0, "epistemico": "fato", "origem": "wiki", "rel": "usa"}, {"alvo": "amendoa_hidrogenio_formal", "confianca": 1.0, "epistemico": "fato", "origem": "wiki", "rel": "relaciona"}], "confianca": 1.0, "contraexemplos": [], "descricao": "A formalização que faltava (a costura ANISOTRÓPICA): uma VESICA (a amêndoa, a interseção de 2 círculos) é uma MATRIZ SIMÉTRICA M (a métrica/o tensor de forma da lente); a ANISOTROPIA (a≠b, a costura) é o SPLIT dos autovalores (o discriminante, o gap). AS DUAS VESICAS: a VERTICAL (cúspides (0,±b)) ↔ o eixo ESQUILO (±i, elíptico, a rotação); a HORIZONTAL (cúspides (±b,0)) ↔ o eixo GATO (σ,−1/σ, hiperbólico); a rotação 90° que leva uma na outra É o ESQUILO (i, G⁴=I). O FATO-CHAVE: o COMUTADOR de duas vesicas simétricas é SEMPRE antissimétrico ([M1,M2]ᵀ=−[M1,M2]) = o ESQUILO — a costura anisotrópica GERA a rotação; ≠0 sse anisotrópicas E desalinhadas; zera na isotrópica (círculo=escalar·I comuta) ou alinhada. E {I,M1,M2,[M1,M2]} têm posto 4 = M_2: AS VESICAS GERAM O CORPO — a anisotropia abre o corpo inteiro (Sym⊕Skew). O d/r = a parábola (d→0 cristal, d→r borda/vértice degenerado); as 2 cruzadas = as 4 cúspides = os 4 quadrantes. linguagem/vesica_corpo.py (13/13). É uma matrix: a matrix vesicular.", "epistemico": "hipotese", "exemplos": [], "id": "matrix_vesicular_corpo", "memoria": "semantica", "meta": {"dominio": "floxina", "fonte": "Floxina (investigação)"}, "origem": "wiki", "palavras_chave": [], "sinonimos": [], "tipo": "conceito", "titulo": "A MATRIX VESICULAR — a vesica é uma matriz; a costura anisotrópica gera o corpo"}
\section{A MATRIX VESICULAR — a vesica é uma matriz; a costura anisotrópica gera o corpo}
A formalização que faltava (a costura ANISOTRÓPICA): uma VESICA (a amêndoa, a interseção de 2 círculos) é uma MATRIZ SIMÉTRICA M (a métrica/o tensor de forma da lente); a ANISOTROPIA (a\ensuremath{\ne}b, a costura) é o SPLIT dos autovalores (o discriminante, o gap). AS DUAS VESICAS: a VERTICAL (cúspides (0,$\pm$b)) \ensuremath{\leftrightarrow} o eixo ESQUILO ($\pm$i, elíptico, a rotação); a HORIZONTAL (cúspides ($\pm$b,0)) \ensuremath{\leftrightarrow} o eixo GATO (\ensuremath{\sigma},\ensuremath{-}1/\ensuremath{\sigma}, hiperbólico); a rotação 90$^{\circ}$ que leva uma na outra É o ESQUILO (i, G\ensuremath{^{4}}=I). O FATO-CHAVE: o COMUTADOR de duas vesicas simétricas é SEMPRE antissimétrico ([M1,M2]\ensuremath{^{T}}=\ensuremath{-}[M1,M2]) = o ESQUILO — a costura anisotrópica GERA a rotação; \ensuremath{\ne}0 sse anisotrópicas E desalinhadas; zera na isotrópica (círculo=escalar\textperiodcentered I comuta) ou alinhada. E \{I,M1,M2,[M1,M2]\} têm posto 4 = M\_2: AS VESICAS GERAM O CORPO — a anisotropia abre o corpo inteiro (Sym\ensuremath{\oplus}Skew). O d/r = a parábola (d\ensuremath{\to}0 cristal, d\ensuremath{\to}r borda/vértice degenerado); as 2 cruzadas = as 4 cúspides = os 4 quadrantes. linguagem/vesica\_corpo.py (13/13). É uma matrix: a matrix vesicular.
\prov{relações: relaciona costura\_vesica\_programa; parte\_de corpo\_matricial\_matrix; usa gato; usa gamba\_antissimetrico; usa gerador\_do\_gap\_comutador; relaciona amendoa\_hidrogenio\_formal}
\prov{proveniência: wiki \textperiodcentered{} Floxina (investigação) \textperiodcentered{} hipotese \textperiodcentered{} confiança 1.0 \textperiodcentered{} domínio floxina \textperiodcentered{} história 224 versões no jornal}

%CRISTAL {"arestas": [{"alvo": "corpo_vesicular_parabolas_costuras", "confianca": 1.0, "epistemico": "fato", "origem": "wiki", "rel": "usa"}, {"alvo": "paralelismo_pipe_no_metal", "confianca": 1.0, "epistemico": "fato", "origem": "wiki", "rel": "relaciona"}, {"alvo": "agm_costura_fecha_corpo", "confianca": 1.0, "epistemico": "fato", "origem": "wiki", "rel": "usa"}, {"alvo": "antimateria_parabolica", "confianca": 1.0, "epistemico": "fato", "origem": "wiki", "rel": "usa"}], "confianca": 1.0, "contraexemplos": [], "descricao": "O CORPO ALGÉBRICO não comanda UMA costura: comanda uma MATRIZ de costuras (N lanes de uma vez). Uma só lei: (1) a PARÁBOLA GERADORA λ=¼−δ² (o zero da reta, matéria-antimatéria) — a MESMA para todas; (2) tudo PARTE do AGM QUADRÁTICO (a,b)→((a+b)/2,√(ab)), o gap decai ao quadrado (a costura que fecha o corpo é a semente); (3) a MATRIZ DE COSTURAS: N vesicas simétricas (a parábola geradora ORIENTADA de N modos), as costuras=os comutadores [Mi,Mj]=os esquilos, comandados em paralelo (C(N,2) lanes); (4) a MESMA DINÂMICA PROPAGA: todo lane tem o espectro {¼−δ²} IDÊNTICO (Mi=Qi·diag(λ)·Qiᵀ, eig=λ ∀i) — só a orientação/a costura muda. É o mesmo N-lane da mineração/do pipe gráfico, agora COMANDADO pelo Corpo Algébrico na parábola geradora. linguagem/omnitrix.py (matriz_costuras_paralela, 4/4).", "epistemico": "fato", "exemplos": [], "id": "matriz_costuras_paralela_universal", "memoria": "semantica", "meta": {"dominio": "floxina", "fonte": "Floxina (investigação)"}, "origem": "wiki", "palavras_chave": [], "sinonimos": [], "tipo": "conceito", "titulo": "O Corpo Algébrico comanda a matriz de costuras em paralelo (a parábola geradora, do AGM quadrático)"}
\section{O Corpo Algébrico comanda a matriz de costuras em paralelo (a parábola geradora, do AGM quadrático)}
O CORPO ALGÉBRICO não comanda UMA costura: comanda uma MATRIZ de costuras (N lanes de uma vez). Uma só lei: (1) a PARÁBOLA GERADORA \ensuremath{\lambda}=\ensuremath{\tfrac14}\ensuremath{-}\ensuremath{\delta}\ensuremath{^{2}} (o zero da reta, matéria-antimatéria) — a MESMA para todas; (2) tudo PARTE do AGM QUADRÁTICO (a,b)\ensuremath{\to}((a+b)/2,\ensuremath{\sqrt{\,}}(ab)), o gap decai ao quadrado (a costura que fecha o corpo é a semente); (3) a MATRIZ DE COSTURAS: N vesicas simétricas (a parábola geradora ORIENTADA de N modos), as costuras=os comutadores [Mi,Mj]=os esquilos, comandados em paralelo (C(N,2) lanes); (4) a MESMA DINÂMICA PROPAGA: todo lane tem o espectro \{\ensuremath{\tfrac14}\ensuremath{-}\ensuremath{\delta}\ensuremath{^{2}}\} IDÊNTICO (Mi=Qi\textperiodcentered diag(\ensuremath{\lambda})\textperiodcentered Qi\ensuremath{^{T}}, eig=\ensuremath{\lambda} \ensuremath{\forall}i) — só a orientação/a costura muda. É o mesmo N-lane da mineração/do pipe gráfico, agora COMANDADO pelo Corpo Algébrico na parábola geradora. linguagem/omnitrix.py (matriz\_costuras\_paralela, 4/4).
\prov{relações: usa corpo\_vesicular\_parabolas\_costuras; relaciona paralelismo\_pipe\_no\_metal; usa agm\_costura\_fecha\_corpo; usa antimateria\_parabolica}
\prov{proveniência: wiki \textperiodcentered{} Floxina (investigação) \textperiodcentered{} fato \textperiodcentered{} confiança 1.0 \textperiodcentered{} domínio floxina \textperiodcentered{} história 145 versões no jornal}

%CRISTAL {"arestas": [{"alvo": "cristal_runner_executa", "confianca": 1.0, "epistemico": "fato", "origem": "wiki", "rel": "relaciona"}, {"alvo": "broca_os", "confianca": 1.0, "epistemico": "fato", "origem": "wiki", "rel": "relaciona"}], "confianca": 1.0, "contraexemplos": [], "descricao": "Aproxima derivadas por diferenças discretas: f'(x)≈(f(x+h)−f(x))/h (avante), (f(x)−f(x−h))/h (atrás), (f(x+h)−f(x−h))/2h (central, erro O(h²)). Base da resolução numérica de EDOs/EDPs (discretiza o domínio numa grade e substitui derivadas por diferenças). Fato numérico estabelecido; o erro de truncamento e a estabilidade dependem de h e do esquema.", "epistemico": "fato", "exemplos": [], "id": "metodo_diferencas_finitas", "memoria": "semantica", "meta": {"autor": "broca-so", "dominio": "metodos", "fonte": "análise numérica clássica (Euler, Taylor)"}, "origem": "wiki", "palavras_chave": [], "sinonimos": [], "tipo": "conceito", "titulo": "Método das diferenças finitas"}
\section{Método das diferenças finitas}
Aproxima derivadas por diferenças discretas: f'(x)\ensuremath{\approx}(f(x+h)\ensuremath{-}f(x))/h (avante), (f(x)\ensuremath{-}f(x\ensuremath{-}h))/h (atrás), (f(x+h)\ensuremath{-}f(x\ensuremath{-}h))/2h (central, erro O(h\ensuremath{^{2}})). Base da resolução numérica de EDOs/EDPs (discretiza o domínio numa grade e substitui derivadas por diferenças). Fato numérico estabelecido; o erro de truncamento e a estabilidade dependem de h e do esquema.
\prov{relações: relaciona cristal\_runner\_executa; relaciona broca\_os}
\prov{proveniência: wiki \textperiodcentered{} análise numérica clássica (Euler, Taylor) \textperiodcentered{} fato \textperiodcentered{} confiança 1.0 \textperiodcentered{} domínio metodos \textperiodcentered{} história 3 versões no jornal}

%CRISTAL {"arestas": [{"alvo": "metodo_diferencas_finitas", "confianca": 1.0, "epistemico": "fato", "origem": "wiki", "rel": "especializa"}, {"alvo": "estatuto_derivado_posto_convencao", "confianca": 1.0, "epistemico": "fato", "origem": "wiki", "rel": "usa"}, {"alvo": "sne_ia_pantheon", "confianca": 1.0, "epistemico": "fato", "origem": "wiki", "rel": "relaciona"}, {"alvo": "de_sitter_lambda_energia_escura", "confianca": 1.0, "epistemico": "fato", "origem": "wiki", "rel": "relaciona"}, {"alvo": "regua_de_gentil", "confianca": 1.0, "epistemico": "fato", "origem": "wiki", "rel": "relaciona"}, {"alvo": "cosmologia_quantica_hub", "confianca": 1.0, "epistemico": "fato", "origem": "wiki", "rel": "parte_de"}], "confianca": 1.0, "contraexemplos": [], "descricao": "Diferenças finitas em ESCALA LOG-LOG: para y=C·xᵅ, o expoente sai da inclinação α=Δ(log y)/Δ(log x), e o prefator C CANCELA na diferença de logaritmos — mede-se o expoente sem calibrar a constante. Recupera o 3/2 de Kepler (8 planetas α=1,4994; luas galileanas 1,5000), o 3/2 do ground state SK (Parisi) e, nos fótons de SN, α(z) subindo acima de 1 = aceleração cósmica sem calibrar vela-padrão. É uma régua da família Lopes: separa o 'derivado' (o expoente, a forma) do 'posto' (a escala, a constante).", "epistemico": "inferencia", "exemplos": [], "id": "metodo_lopes", "memoria": "semantica", "meta": {"autor": "broca-so", "dominio": "metodos", "fonte": "paper_2_cosmologia.tex §2,§4 (lopes_kepler.py, eq:lopes-sn)"}, "origem": "wiki", "palavras_chave": [], "sinonimos": [], "tipo": "conceito", "titulo": "O Método Lopes (expoente sem calibrar a constante)"}
\section{O Método Lopes (expoente sem calibrar a constante)}
Diferenças finitas em ESCALA LOG-LOG: para y=C\textperiodcentered x\ensuremath{^{\alpha}}, o expoente sai da inclinação \ensuremath{\alpha}=\ensuremath{\Delta}(log y)/\ensuremath{\Delta}(log x), e o prefator C CANCELA na diferença de logaritmos — mede-se o expoente sem calibrar a constante. Recupera o 3/2 de Kepler (8 planetas \ensuremath{\alpha}=1,4994; luas galileanas 1,5000), o 3/2 do ground state SK (Parisi) e, nos fótons de SN, \ensuremath{\alpha}(z) subindo acima de 1 = aceleração cósmica sem calibrar vela-padrão. É uma régua da família Lopes: separa o 'derivado' (o expoente, a forma) do 'posto' (a escala, a constante).
\prov{relações: especializa metodo\_diferencas\_finitas; usa estatuto\_derivado\_posto\_convencao; relaciona sne\_ia\_pantheon; relaciona de\_sitter\_lambda\_energia\_escura; relaciona regua\_de\_gentil; parte\_de cosmologia\_quantica\_hub}
\prov{proveniência: wiki \textperiodcentered{} paper\_2\_cosmologia.tex §2,§4 (lopes\_kepler.py, eq:lopes-sn) \textperiodcentered{} inferencia \textperiodcentered{} confiança 1.0 \textperiodcentered{} domínio metodos \textperiodcentered{} história 7 versões no jornal}

%CRISTAL {"arestas": [{"alvo": "regua_de_gentil", "confianca": 1.0, "epistemico": "fato", "origem": "wiki", "rel": "relaciona"}, {"alvo": "broca_os", "confianca": 1.0, "epistemico": "fato", "origem": "wiki", "rel": "parte_de"}], "confianca": 1.0, "contraexemplos": [], "descricao": "O hub: até onde vale construir — e por que parar não é desistir. A matemática desce para sempre (Prop. da descida infinita), mas o proveito para a nossa dimensão para em N* (teorema de escopo). A Missão é construir até N* — o ponto onde o ganho prático se esgota — e parar ali, de olhos abertos. Descer abaixo é legítimo mas estéril: Loucura de Ouro (horizonte aberto, custo infinito, retorno zero). 'Parar não é desistir; é a própria definição de engenharia — otimizar o ganho, não perseguir o infinito.'", "epistemico": "inferencia", "exemplos": [], "id": "missao_engenharia_ouro", "memoria": "semantica", "meta": {"autor": "Lopes/conselho", "dominio": "missao", "fonte": "machine/Missao.tex"}, "origem": "wiki", "palavras_chave": [], "sinonimos": [], "tipo": "conceito", "titulo": "A Missão da Engenharia de Ouro (construir até N*, parar)"}
\section{A Missão da Engenharia de Ouro (construir até N*, parar)}
O hub: até onde vale construir — e por que parar não é desistir. A matemática desce para sempre (Prop. da descida infinita), mas o proveito para a nossa dimensão para em N* (teorema de escopo). A Missão é construir até N* — o ponto onde o ganho prático se esgota — e parar ali, de olhos abertos. Descer abaixo é legítimo mas estéril: Loucura de Ouro (horizonte aberto, custo infinito, retorno zero). 'Parar não é desistir; é a própria definição de engenharia — otimizar o ganho, não perseguir o infinito.'
\prov{relações: relaciona regua\_de\_gentil; parte\_de broca\_os}
\prov{proveniência: wiki \textperiodcentered{} machine/Missao.tex \textperiodcentered{} inferencia \textperiodcentered{} confiança 1.0 \textperiodcentered{} domínio missao \textperiodcentered{} história 3 versões no jornal}

%CRISTAL {"arestas": [{"alvo": "omnitrix_corpo_universal", "confianca": 1.0, "epistemico": "fato", "origem": "wiki", "rel": "parte_de"}, {"alvo": "fractal_newton_reta_matricial", "confianca": 1.0, "epistemico": "fato", "origem": "wiki", "rel": "usa"}, {"alvo": "leminiscata_grupo_de_grupos", "confianca": 1.0, "epistemico": "fato", "origem": "wiki", "rel": "usa"}, {"alvo": "antimateria_parabolica", "confianca": 1.0, "epistemico": "fato", "origem": "wiki", "rel": "usa"}, {"alvo": "reta_matricial_compacta_cf", "confianca": 1.0, "epistemico": "fato", "origem": "wiki", "rel": "usa"}, {"alvo": "parabola_geradora_idempotente", "confianca": 1.0, "epistemico": "fato", "origem": "wiki", "rel": "relaciona"}], "confianca": 1.0, "contraexemplos": [], "descricao": "Cada INSTÂNCIA do Omnitrix é um FRACTAL DE NEWTON: o polinômio de grau d (metal x²−m·x−1 d=2; volume 3D x³−x−1 d=3; …) dá N(z)=z−p/p', bacias na esfera de Riemann, a fronteira=o Julia. A dimensão de caixa cresce com o grau: d=2→1,005 (split LIMPO, a folha), d=3→1,367, d=4→1,551, d=5→1,596. Cada instância é ~monofractal (uma D); a FAMÍLIA é o MULTIFRACTAL: um ESPECTRO de dimensões (α∈[1,0;1,6], Δα≈0,59>0 = a assinatura; um fractal só teria Δα=0). O GRUPO DISCRETO na leminiscata costura: as d bacias permutadas pelo Galois (S_d; S₃ no cúbico, disc −23), o gerador auto-semelhante é o modular PSL(2,ℤ)=Z/2∗Z/3. A FONTE DA COSTURA é a parábola geradora MATÉRIA-ANTIMATÉRIA (a fronteira: matéria σ>0 ∩ antimatéria σ'<0). O FLIP z↦−1/z (=S, o gambá, a conjugação de carga C) troca σ↔σ'=−1/σ (involução, S²=−I, Z/2 do modular) e gera o ANTI-UNIVERSO: o FRACTAL NEGRO (o flip do fractal, a MESMA dimensão pois Möbius é conforme). Matéria (o fractal) + antimatéria (o fractal negro), colados na costura, FECHAM o corpo (a esfera inteira, compacta pela CF de π=∏gatos∈GL(2,ℤ), a superfície modular Z/2∗Z/3). linguagem/multifractal_newton.py (4/4). Síntese declarada.", "epistemico": "hipotese", "exemplos": [], "id": "multifractal_newton_omnitrix", "memoria": "semantica", "meta": {"dominio": "floxina", "fonte": "Floxina (investigação)"}, "origem": "wiki", "palavras_chave": [], "sinonimos": [], "tipo": "conceito", "titulo": "O multifractal newtoniano: cada instância um fractal, o grupo discreto costura"}
\section{O multifractal newtoniano: cada instância um fractal, o grupo discreto costura}
Cada INSTÂNCIA do Omnitrix é um FRACTAL DE NEWTON: o polinômio de grau d (metal x\ensuremath{^{2}}\ensuremath{-}m\textperiodcentered x\ensuremath{-}1 d=2; volume 3D x\ensuremath{^{3}}\ensuremath{-}x\ensuremath{-}1 d=3; …) dá N(z)=z\ensuremath{-}p/p', bacias na esfera de Riemann, a fronteira=o Julia. A dimensão de caixa cresce com o grau: d=2\ensuremath{\to}1,005 (split LIMPO, a folha), d=3\ensuremath{\to}1,367, d=4\ensuremath{\to}1,551, d=5\ensuremath{\to}1,596. Cada instância é \~{}monofractal (uma D); a FAMÍLIA é o MULTIFRACTAL: um ESPECTRO de dimensões (\ensuremath{\alpha}\ensuremath{\in}[1,0;1,6], \ensuremath{\Delta}\ensuremath{\alpha}\ensuremath{\approx}0,59>0 = a assinatura; um fractal só teria \ensuremath{\Delta}\ensuremath{\alpha}=0). O GRUPO DISCRETO na leminiscata costura: as d bacias permutadas pelo Galois (S\_d; S\ensuremath{_{3}} no cúbico, disc \ensuremath{-}23), o gerador auto-semelhante é o modular PSL(2,\ensuremath{\mathbb{Z}})=Z/2\ensuremath{*}Z/3. A FONTE DA COSTURA é a parábola geradora MATÉRIA-ANTIMATÉRIA (a fronteira: matéria \ensuremath{\sigma}>0 \ensuremath{\cap} antimatéria \ensuremath{\sigma}'<0). O FLIP z\ensuremath{\mapsto}\ensuremath{-}1/z (=S, o gambá, a conjugação de carga C) troca \ensuremath{\sigma}\ensuremath{\leftrightarrow}\ensuremath{\sigma}'=\ensuremath{-}1/\ensuremath{\sigma} (involução, S\ensuremath{^{2}}=\ensuremath{-}I, Z/2 do modular) e gera o ANTI-UNIVERSO: o FRACTAL NEGRO (o flip do fractal, a MESMA dimensão pois Möbius é conforme). Matéria (o fractal) + antimatéria (o fractal negro), colados na costura, FECHAM o corpo (a esfera inteira, compacta pela CF de \ensuremath{\pi}=\ensuremath{\prod}gatos\ensuremath{\in}GL(2,\ensuremath{\mathbb{Z}}), a superfície modular Z/2\ensuremath{*}Z/3). linguagem/multifractal\_newton.py (4/4). Síntese declarada.
\prov{relações: parte\_de omnitrix\_corpo\_universal; usa fractal\_newton\_reta\_matricial; usa leminiscata\_grupo\_de\_grupos; usa antimateria\_parabolica; usa reta\_matricial\_compacta\_cf; relaciona parabola\_geradora\_idempotente}
\prov{proveniência: wiki \textperiodcentered{} Floxina (investigação) \textperiodcentered{} hipotese \textperiodcentered{} confiança 1.0 \textperiodcentered{} domínio floxina \textperiodcentered{} história 378 versões no jornal}

%CRISTAL {"arestas": [{"alvo": "corpo_edp", "confianca": 1.0, "epistemico": "fato", "origem": "wiki", "rel": "especializa"}, {"alvo": "conservacao_energia_plancherel", "confianca": 1.0, "epistemico": "fato", "origem": "wiki", "rel": "usa"}, {"alvo": "solucao_explicita_ed", "confianca": 1.0, "epistemico": "fato", "origem": "wiki", "rel": "usa"}, {"alvo": "corpo_equacoes_diferenciais", "confianca": 1.0, "epistemico": "fato", "origem": "wiki", "rel": "relaciona"}], "confianca": 1.0, "contraexemplos": [], "descricao": "Navier–Stokes incompressível ∂_t u+(u·∇)u=ν∇²u−∇p, ∇·u=0, é CASO PARTICULAR do corpo das EDP, aberto pela decomposição gato/esquilo: (1) o LINEAR (Stokes) = o corpo PARABÓLICO (o CALOR, o GATO): A=−νℙΔ, autovalores −νk²<0 (reais negativos=CRISTAL/decaimento), a viscosidade ν DISSIPA; solução explícita u(t)=e^{−νtA}u₀ (o semigrupo do calor). (2) o NÃO-LINEAR (u·∇)u = o convectivo, o ESQUILO: ⟨ℙ(u·∇u),u⟩=0 (verificado 1e-18), CONSERVA a energia (só transporta), como a rotação do esquilo. (3) a PRESSÃO/incompressibilidade = a projeção de Leray ℙ. A SOLUÇÃO EXPLÍCITA é a mild/DUHAMEL = a inomogênea da Parte VII: u(t)=e^{−νtA}u₀−∫₀ᵗe^{−ν(t−s)A}ℙ(u·∇u)ds (o fluxo do calor + a fonte b=−ℙ(u·∇u)=o ⊗ convectivo). A LEI DE ENERGIA = a dissipação de Plancherel: d/dt½‖u‖²=−ν‖∇u‖²≤0 (o gato dissipa, o esquilo cai fora, a 2ª lei). SOLUÇÃO EXATA: o vórtice de TAYLOR–GREEN u=e^{−2νt}(sin x cos y,−cos x sin y) — o não-linear é gradiente puro, ℙ o mata, vira calor puro, o modo (1,1) decai a 2ν (∇²u=−2u, o cristal). ESCOPO HONESTO: Stokes explícito, Taylor–Green exato, lei de energia=fato; a REGULARIDADE GLOBAL 3D é a BORDA ABERTA (Clay) — o ⊗ fica no cristal ou chega ao caos/blowup? — o corpo APONTA a fronteira, NÃO a fecha (não alegada). linguagem/corpo_edp.py (navier_stokes_caso_particular 5/5); papers/equacoes_diferenciais.tex (Parte IX).", "epistemico": "fato", "exemplos": [], "id": "navier_stokes_caso_particular", "memoria": "semantica", "meta": {"dominio": "floxina", "fonte": "Floxina (investigação)"}, "origem": "wiki", "palavras_chave": [], "sinonimos": [], "tipo": "conceito", "titulo": "Navier–Stokes como caso particular do corpo das EDP: o calor/gato (viscoso) + o convectivo/esquilo (conserva energia); Stokes explícito, Taylor–Green exato, 3D=borda aberta (Clay)"}
\section{Navier–Stokes como caso particular do corpo das EDP: o calor/gato (viscoso) + o convectivo/esquilo (conserva energia); Stokes explícito, Taylor–Green exato, 3D=borda aberta (Clay)}
Navier–Stokes incompressível \ensuremath{\partial}\_t u+(u\textperiodcentered \ensuremath{\nabla})u=\ensuremath{\nu}\ensuremath{\nabla}\ensuremath{^{2}}u\ensuremath{-}\ensuremath{\nabla}p, \ensuremath{\nabla}\textperiodcentered u=0, é CASO PARTICULAR do corpo das EDP, aberto pela decomposição gato/esquilo: (1) o LINEAR (Stokes) = o corpo PARABÓLICO (o CALOR, o GATO): A=\ensuremath{-}\ensuremath{\nu}\ensuremath{\mathbb{P}}\ensuremath{\Delta}, autovalores \ensuremath{-}\ensuremath{\nu}k\ensuremath{^{2}}<0 (reais negativos=CRISTAL/decaimento), a viscosidade \ensuremath{\nu} DISSIPA; solução explícita u(t)=e\^{}\{\ensuremath{-}\ensuremath{\nu}tA\}u\ensuremath{_{0}} (o semigrupo do calor). (2) o NÃO-LINEAR (u\textperiodcentered \ensuremath{\nabla})u = o convectivo, o ESQUILO: \ensuremath{\langle}\ensuremath{\mathbb{P}}(u\textperiodcentered \ensuremath{\nabla}u),u\ensuremath{\rangle}=0 (verificado 1e-18), CONSERVA a energia (só transporta), como a rotação do esquilo. (3) a PRESSÃO/incompressibilidade = a projeção de Leray \ensuremath{\mathbb{P}}. A SOLUÇÃO EXPLÍCITA é a mild/DUHAMEL = a inomogênea da Parte VII: u(t)=e\^{}\{\ensuremath{-}\ensuremath{\nu}tA\}u\ensuremath{_{0}}\ensuremath{-}\ensuremath{\int}\ensuremath{_{0}}\ensuremath{^{t}}e\^{}\{\ensuremath{-}\ensuremath{\nu}(t\ensuremath{-}s)A\}\ensuremath{\mathbb{P}}(u\textperiodcentered \ensuremath{\nabla}u)ds (o fluxo do calor + a fonte b=\ensuremath{-}\ensuremath{\mathbb{P}}(u\textperiodcentered \ensuremath{\nabla}u)=o \ensuremath{\otimes} convectivo). A LEI DE ENERGIA = a dissipação de Plancherel: d/dt\ensuremath{\tfrac12}\ensuremath{\|}u\ensuremath{\|}\ensuremath{^{2}}=\ensuremath{-}\ensuremath{\nu}\ensuremath{\|}\ensuremath{\nabla}u\ensuremath{\|}\ensuremath{^{2}}\ensuremath{\le}0 (o gato dissipa, o esquilo cai fora, a 2ª lei). SOLUÇÃO EXATA: o vórtice de TAYLOR–GREEN u=e\^{}\{\ensuremath{-}2\ensuremath{\nu}t\}(sin x cos y,\ensuremath{-}cos x sin y) — o não-linear é gradiente puro, \ensuremath{\mathbb{P}} o mata, vira calor puro, o modo (1,1) decai a 2\ensuremath{\nu} (\ensuremath{\nabla}\ensuremath{^{2}}u=\ensuremath{-}2u, o cristal). ESCOPO HONESTO: Stokes explícito, Taylor–Green exato, lei de energia=fato; a REGULARIDADE GLOBAL 3D é a BORDA ABERTA (Clay) — o \ensuremath{\otimes} fica no cristal ou chega ao caos/blowup? — o corpo APONTA a fronteira, NÃO a fecha (não alegada). linguagem/corpo\_edp.py (navier\_stokes\_caso\_particular 5/5); papers/equacoes\_diferenciais.tex (Parte IX).
\prov{relações: especializa corpo\_edp; usa conservacao\_energia\_plancherel; usa solucao\_explicita\_ed; relaciona corpo\_equacoes\_diferenciais}
\prov{proveniência: wiki \textperiodcentered{} Floxina (investigação) \textperiodcentered{} fato \textperiodcentered{} confiança 1.0 \textperiodcentered{} domínio floxina \textperiodcentered{} história 35 versões no jornal}

%CRISTAL {"arestas": [{"alvo": "rasgar_funcao_de_onda", "confianca": 1.0, "epistemico": "fato", "origem": "wiki", "rel": "explica"}, {"alvo": "associador_calor", "confianca": 1.0, "epistemico": "fato", "origem": "wiki", "rel": "relaciona"}], "confianca": 1.0, "contraexemplos": [], "descricao": "A frase de origem, agora com dentes: para todo n, a deformação não é tangente a 𝒫 em ordem infinita (kₙ<∞ e bₙ≠0). É AXIOMA, não teorema — é o que gera a descida: em cada passo nasce um sub-gap dₙ+₁≠0, não há repouso absoluto. 'Ninguém está puro: por isso há sempre um fundo mais fundo.' A pureza (a=0, sem associador) é sempre local; abaixo há mais associador a puxar.", "epistemico": "hipotese", "exemplos": [], "id": "ninguem_esta_puro_recursivo", "memoria": "semantica", "meta": {"autor": "Lopes/conselho", "dominio": "missao", "fonte": "machine/Missao.tex"}, "origem": "wiki", "palavras_chave": [], "sinonimos": [], "tipo": "conceito", "titulo": "Ninguém está puro — o axioma recursivo (a descida infinita)"}
\section{Ninguém está puro — o axioma recursivo (a descida infinita)}
A frase de origem, agora com dentes: para todo n, a deformação não é tangente a \ensuremath{\mathcal{P}} em ordem infinita (k\ensuremath{_{n}}<\ensuremath{\infty} e b\ensuremath{_{n}}\ensuremath{\ne}0). É AXIOMA, não teorema — é o que gera a descida: em cada passo nasce um sub-gap d\ensuremath{_{n}}+\ensuremath{_{1}}\ensuremath{\ne}0, não há repouso absoluto. 'Ninguém está puro: por isso há sempre um fundo mais fundo.' A pureza (a=0, sem associador) é sempre local; abaixo há mais associador a puxar.
\prov{relações: explica rasgar\_funcao\_de\_onda; relaciona associador\_calor}
\prov{proveniência: wiki \textperiodcentered{} machine/Missao.tex \textperiodcentered{} hipotese \textperiodcentered{} confiança 1.0 \textperiodcentered{} domínio missao \textperiodcentered{} história 4 versões no jornal}

%CRISTAL {"arestas": [{"alvo": "floxina_quantum_materia_escura", "confianca": 1.0, "epistemico": "fato", "origem": "wiki", "rel": "relaciona"}], "confianca": 1.0, "contraexemplos": [], "descricao": "FLOXIN é a marca real de um antibiótico cujo princípio ativo é o OFLOXACINO (C₁₈H₂₀FN₃O₄, uma fluoroquinolona, MW≈361). Estrutura: um núcleo tricíclico piridobenzoxazina (uma quinolona fundida a um anel oxazina), com um FLÚOR na posição 6 (a ponta eletronegativa reativa), um −COOH na posição 3 (ácido; quela o Mg²⁺ no sítio da enzima), um 4-oxo (cetona), uma metil-piperazina e um metilo no oxazina que é o ESTEREOCENTRO — a molécula é QUIRAL (o levofloxacino = o (S)-(−)-ofloxacino, o enantiómero ativo, gira a luz à esquerda). MECANISMO: inibe a DNA GIRASE (topoisomerase II) e a topo IV — as enzimas que gerem o SUPERENROLAMENTO (a TORÇÃO) do DNA; estabiliza o complexo DNA-enzima clivado. Ou seja: o Floxin age sobre o TWIST do DNA. Fato químico/farmacológico conhecido; ingerido para casar com a floxina do projeto.", "epistemico": "fato", "exemplos": [], "id": "ofloxacino_floxina_real", "memoria": "semantica", "meta": {"dominio": "floxina", "fonte": "Floxina (investigação)"}, "origem": "wiki", "palavras_chave": [], "sinonimos": [], "tipo": "conceito", "titulo": "A floxina REAL: o Floxin (ofloxacino, C₁₈H₂₀FN₃O₄)"}
\section{A floxina REAL: o Floxin (ofloxacino, C\ensuremath{_{18}}H\ensuremath{_{20}}FN\ensuremath{_{3}}O\ensuremath{_{4}})}
FLOXIN é a marca real de um antibiótico cujo princípio ativo é o OFLOXACINO (C\ensuremath{_{18}}H\ensuremath{_{20}}FN\ensuremath{_{3}}O\ensuremath{_{4}}, uma fluoroquinolona, MW\ensuremath{\approx}361). Estrutura: um núcleo tricíclico piridobenzoxazina (uma quinolona fundida a um anel oxazina), com um FLÚOR na posição 6 (a ponta eletronegativa reativa), um \ensuremath{-}COOH na posição 3 (ácido; quela o Mg\ensuremath{^{2+}} no sítio da enzima), um 4-oxo (cetona), uma metil-piperazina e um metilo no oxazina que é o ESTEREOCENTRO — a molécula é QUIRAL (o levofloxacino = o (S)-(\ensuremath{-})-ofloxacino, o enantiómero ativo, gira a luz à esquerda). MECANISMO: inibe a DNA GIRASE (topoisomerase II) e a topo IV — as enzimas que gerem o SUPERENROLAMENTO (a TORÇÃO) do DNA; estabiliza o complexo DNA-enzima clivado. Ou seja: o Floxin age sobre o TWIST do DNA. Fato químico/farmacológico conhecido; ingerido para casar com a floxina do projeto.
\prov{relações: relaciona floxina\_quantum\_materia\_escura}
\prov{proveniência: wiki \textperiodcentered{} Floxina (investigação) \textperiodcentered{} fato \textperiodcentered{} confiança 1.0 \textperiodcentered{} domínio floxina \textperiodcentered{} história 104 versões no jornal}

%CRISTAL {"arestas": [{"alvo": "corpo_matricial_matrix", "confianca": 1.0, "epistemico": "fato", "origem": "wiki", "rel": "relaciona"}, {"alvo": "antimateria_parabolica", "confianca": 1.0, "epistemico": "fato", "origem": "wiki", "rel": "usa"}, {"alvo": "parabola_geradora_idempotente", "confianca": 1.0, "epistemico": "fato", "origem": "wiki", "rel": "usa"}, {"alvo": "invariante_molecular_residuo", "confianca": 1.0, "epistemico": "fato", "origem": "wiki", "rel": "usa"}, {"alvo": "gerador_do_gap_comutador", "confianca": 1.0, "epistemico": "fato", "origem": "wiki", "rel": "usa"}], "confianca": 1.0, "contraexemplos": [], "descricao": "O OMNITRIX é o corpo dos corpos: o TEMPLATE (a parábola geradora + o gato=Sym + o gambá=Skew + o invariante a+c²=1) de que a matrix é UMA INSTÂNCIA. Como o relógio que vira qualquer forma, escolhe uma instância: uma ORDEM m (o metal σ_m), uma DIMENSÃO (2×2, 3D volume, …), um SUBSTRATO (abstrato=matrix, molecular, estelar, cósmico, autorização=BAI). A MATRIX é a instância de ordem 2. O ZERO DA RETA (m=0) é a dinâmica MATÉRIA-ANTIMATÉRIA: A_0=[[0,1],[1,0]] tem raízes ±1 (matéria +1, antimatéria −1, SIMÉTRICAS, o par puro); a reta m>0 DESDOBRA a assimetria (σ_m cresce, 1/σ_m encolhe: a matéria domina = a bariogênese); o produto σ·σ'=−1 (a aniquilação) invariante. As OPERAÇÕES da reta universal sobre corpos MOLECULARES: ⊕ (a soma direta, os n-qubits somam, neutro o VÁCUO ∅) e ⊗ (a LA HIRE, o produto matricial dos gatos = a reação/composição, como a CF de π=∏gatos, neutro a IDENTIDADE I); ambas reduzem à base Sym⊕Skew. A SOMBRA da matrix é uma instância (o corpo molecular, o estelar, o cosmos). linguagem/omnitrix.py (6/6). Síntese declarada, não decreto — não é oráculo.", "epistemico": "hipotese", "exemplos": [], "id": "omnitrix_corpo_universal", "memoria": "semantica", "meta": {"dominio": "floxina", "fonte": "Floxina (investigação)"}, "origem": "wiki", "palavras_chave": [], "sinonimos": [], "tipo": "conceito", "titulo": "Omnitrix — o Corpo Algébrico (generaliza a matrix; a matrix é uma instância)"}
\section{Omnitrix — o Corpo Algébrico (generaliza a matrix; a matrix é uma instância)}
O OMNITRIX é o corpo dos corpos: o TEMPLATE (a parábola geradora + o gato=Sym + o gambá=Skew + o invariante a+c\ensuremath{^{2}}=1) de que a matrix é UMA INSTÂNCIA. Como o relógio que vira qualquer forma, escolhe uma instância: uma ORDEM m (o metal \ensuremath{\sigma}\_m), uma DIMENSÃO (2$\times$2, 3D volume, …), um SUBSTRATO (abstrato=matrix, molecular, estelar, cósmico, autorização=BAI). A MATRIX é a instância de ordem 2. O ZERO DA RETA (m=0) é a dinâmica MATÉRIA-ANTIMATÉRIA: A\_0=[[0,1],[1,0]] tem raízes $\pm$1 (matéria +1, antimatéria \ensuremath{-}1, SIMÉTRICAS, o par puro); a reta m>0 DESDOBRA a assimetria (\ensuremath{\sigma}\_m cresce, 1/\ensuremath{\sigma}\_m encolhe: a matéria domina = a bariogênese); o produto \ensuremath{\sigma}\textperiodcentered \ensuremath{\sigma}'=\ensuremath{-}1 (a aniquilação) invariante. As OPERAÇÕES da reta universal sobre corpos MOLECULARES: \ensuremath{\oplus} (a soma direta, os n-qubits somam, neutro o VÁCUO \ensuremath{\varnothing}) e \ensuremath{\otimes} (a LA HIRE, o produto matricial dos gatos = a reação/composição, como a CF de \ensuremath{\pi}=\ensuremath{\prod}gatos, neutro a IDENTIDADE I); ambas reduzem à base Sym\ensuremath{\oplus}Skew. A SOMBRA da matrix é uma instância (o corpo molecular, o estelar, o cosmos). linguagem/omnitrix.py (6/6). Síntese declarada, não decreto — não é oráculo.
\prov{relações: relaciona corpo\_matricial\_matrix; usa antimateria\_parabolica; usa parabola\_geradora\_idempotente; usa invariante\_molecular\_residuo; usa gerador\_do\_gap\_comutador}
\prov{proveniência: wiki \textperiodcentered{} Floxina (investigação) \textperiodcentered{} hipotese \textperiodcentered{} confiança 1.0 \textperiodcentered{} domínio floxina \textperiodcentered{} história 330 versões no jornal}

%CRISTAL {"arestas": [{"alvo": "corpo_equacoes_diferenciais", "confianca": 1.0, "epistemico": "fato", "origem": "wiki", "rel": "usa"}, {"alvo": "parabola_geradora_idempotente", "confianca": 1.0, "epistemico": "fato", "origem": "wiki", "rel": "usa"}, {"alvo": "teorema_universal_residuo_zero", "confianca": 1.0, "epistemico": "fato", "origem": "wiki", "rel": "relaciona"}], "confianca": 1.0, "contraexemplos": [], "descricao": "A parábola geradora é a colisão de DOIS graus: x²−mx−1 = P_g(x², 2º grau, geométrico) − P_a(mx+1, 1º grau, aritmético). No corpo diferencial: a PA (progressão aritmética, 1º grau, o ⊕/gato) é dx/dt=d → x=x₀+d·t (velocidade constante, x''=0); a PG (geométrica, o ⊗/esquilo) é dx/dt=λx → x=x₀·e^{λt} (razão constante, o modo próprio). O elo é o SOMATÓRIO (o integral), que SOBE o grau: ∑ₖ(2k−1)=n² EXATO (a PA dos ímpares acumula no quadrado); ∫(a·t)dt=a·t²/2 (linear → quadrática). A parábola P_g=x² É o somatório do P_a. E a UNICIDADE: a soma de Riemann (o PROCESSO, discreto infinito) → a integral (o CRISTALIZADO, o 2º grau fechado), erro→0, objetos distintos — como 0.999…(processo)→1(cristalizado), a mesma reta 0/1. A derivada D desce o grau (o gato mede), o integral sobe (acumula) — o par adjunto do corpo diferencial. linguagem/omnitrix.py (pa_pg_corpo_diferencial 5/5); papers/equacoes_diferenciais.tex (Parte V).", "epistemico": "fato", "exemplos": [], "id": "pa_pg_corpo_diferencial", "memoria": "semantica", "meta": {"dominio": "floxina", "fonte": "Floxina (investigação)"}, "origem": "wiki", "palavras_chave": [], "sinonimos": [], "tipo": "conceito", "titulo": "PA e PG no corpo diferencial: o somatório de 1º grau → 2º grau (como 0.999…→1); PA=x₀+d·t, PG=x₀·e^{λt}"}
\section{PA e PG no corpo diferencial: o somatório de 1º grau \ensuremath{\to} 2º grau (como 0.999…\ensuremath{\to}1); PA=x\ensuremath{_{0}}+d\textperiodcentered t, PG=x\ensuremath{_{0}}\textperiodcentered e\^{}\{\ensuremath{\lambda}t\}}
A parábola geradora é a colisão de DOIS graus: x\ensuremath{^{2}}\ensuremath{-}mx\ensuremath{-}1 = P\_g(x\ensuremath{^{2}}, 2º grau, geométrico) \ensuremath{-} P\_a(mx+1, 1º grau, aritmético). No corpo diferencial: a PA (progressão aritmética, 1º grau, o \ensuremath{\oplus}/gato) é dx/dt=d \ensuremath{\to} x=x\ensuremath{_{0}}+d\textperiodcentered t (velocidade constante, x''=0); a PG (geométrica, o \ensuremath{\otimes}/esquilo) é dx/dt=\ensuremath{\lambda}x \ensuremath{\to} x=x\ensuremath{_{0}}\textperiodcentered e\^{}\{\ensuremath{\lambda}t\} (razão constante, o modo próprio). O elo é o SOMATÓRIO (o integral), que SOBE o grau: \ensuremath{\sum}\ensuremath{_{k}}(2k\ensuremath{-}1)=n\ensuremath{^{2}} EXATO (a PA dos ímpares acumula no quadrado); \ensuremath{\int}(a\textperiodcentered t)dt=a\textperiodcentered t\ensuremath{^{2}}/2 (linear \ensuremath{\to} quadrática). A parábola P\_g=x\ensuremath{^{2}} É o somatório do P\_a. E a UNICIDADE: a soma de Riemann (o PROCESSO, discreto infinito) \ensuremath{\to} a integral (o CRISTALIZADO, o 2º grau fechado), erro\ensuremath{\to}0, objetos distintos — como 0.999…(processo)\ensuremath{\to}1(cristalizado), a mesma reta 0/1. A derivada D desce o grau (o gato mede), o integral sobe (acumula) — o par adjunto do corpo diferencial. linguagem/omnitrix.py (pa\_pg\_corpo\_diferencial 5/5); papers/equacoes\_diferenciais.tex (Parte V).
\prov{relações: usa corpo\_equacoes\_diferenciais; usa parabola\_geradora\_idempotente; relaciona teorema\_universal\_residuo\_zero}
\prov{proveniência: wiki \textperiodcentered{} Floxina (investigação) \textperiodcentered{} fato \textperiodcentered{} confiança 1.0 \textperiodcentered{} domínio floxina \textperiodcentered{} história 72 versões no jornal}

%CRISTAL {"arestas": [{"alvo": "ponte_simbolica_gentil", "confianca": 1.0, "epistemico": "fato", "origem": "paper", "rel": "relaciona"}, {"alvo": "ciencias", "confianca": 1.0, "epistemico": "fato", "origem": "paper", "rel": "relaciona"}, {"alvo": "matematica", "confianca": 1.0, "epistemico": "fato", "origem": "paper", "rel": "relaciona"}, {"alvo": "fisica", "confianca": 1.0, "epistemico": "fato", "origem": "paper", "rel": "relaciona"}, {"alvo": "corpo_dual", "confianca": 1.0, "epistemico": "fato", "origem": "paper", "rel": "relaciona"}, {"alvo": "corpo_glacial", "confianca": 1.0, "epistemico": "fato", "origem": "paper", "rel": "relaciona"}, {"alvo": "corpo_tropical", "confianca": 1.0, "epistemico": "fato", "origem": "paper", "rel": "relaciona"}, {"alvo": "colapso", "confianca": 1.0, "epistemico": "fato", "origem": "paper", "rel": "relaciona"}, {"alvo": "computacao", "confianca": 1.0, "epistemico": "fato", "origem": "paper", "rel": "explica"}], "confianca": 1.0, "contraexemplos": [], "descricao": "Paper Bridge: morfismos entre áreas (mat., fís., comp., epist.) sem conflação; charts tipados; lab.py reproduzir.", "epistemico": "fato", "exemplos": ["python3 conversa/conhecimento/lab.py reproduzir", "neuronio/experimentos_tiffany.py — 62 testes"], "id": "paper_bridge_areas", "memoria": "semantica", "meta": {"dominio": "meta", "nivel": 5, "serve": "Disciplina cross-área — morfismo sim, fusão não", "verificacao": "slug idêntico — enriquecer, não duplicar"}, "origem": "paper", "palavras_chave": [], "sinonimos": ["paper bridge", "ponte entre áreas"], "tipo": "conceito", "titulo": "paper bridge areas"}
\section{paper bridge areas}
Paper Bridge: morfismos entre áreas (mat., fís., comp., epist.) sem conflação; charts tipados; lab.py reproduzir.
\prov{exemplos: python3 conversa/conhecimento/lab.py reproduzir; neuronio/experimentos\_tiffany.py — 62 testes}
\prov{sinónimos: paper bridge; ponte entre áreas}
\prov{relações: relaciona ponte\_simbolica\_gentil; relaciona ciencias; relaciona matematica; relaciona fisica; relaciona corpo\_dual; relaciona corpo\_glacial; relaciona corpo\_tropical; relaciona colapso; explica computacao}
\prov{proveniência: paper \textperiodcentered{} fato \textperiodcentered{} confiança 1.0 \textperiodcentered{} domínio meta \textperiodcentered{} história 360 versões no jornal}

%CRISTAL {"arestas": [{"alvo": "algebra_fundamental_gato_gamba", "confianca": 1.0, "epistemico": "fato", "origem": "wiki", "rel": "parte_de"}, {"alvo": "gato", "confianca": 1.0, "epistemico": "fato", "origem": "wiki", "rel": "usa"}], "confianca": 1.0, "contraexemplos": [], "descricao": "(A+G)² NÃO é 1: (A+G)²=A+G — IDEMPOTENTE, uma projeção (A+G=[[1,2],[0,0]]). A sua char poly é a PARÁBOLA GERADORA x²−x=x(x−1), raízes 0 e 1 — exatamente os NEUTROS do corpo estelar: ν=0 (o PONTO, o vácuo, neutro de ⊕) e ν=1 (a RETA, a identidade/La Hire, neutro de ⊗). É a parábola geradora CRUA (P_g=x² vs P_a=x); somar o FLIP (a norma N=−1) dá a metálica do ouro x²−x−1. A geradora + o −1 = o metal. conversa/gamba.py (parabola_geradora).", "epistemico": "fato", "exemplos": [], "id": "parabola_geradora_idempotente", "memoria": "semantica", "meta": {"dominio": "floxina", "fonte": "Floxina (investigação)"}, "origem": "wiki", "palavras_chave": [], "sinonimos": [], "tipo": "conceito", "titulo": "A parábola geradora: (A+G)²=A+G, raízes 0 e 1 (os neutros)"}
\section{A parábola geradora: (A+G)\ensuremath{^{2}}=A+G, raízes 0 e 1 (os neutros)}
(A+G)\ensuremath{^{2}} NÃO é 1: (A+G)\ensuremath{^{2}}=A+G — IDEMPOTENTE, uma projeção (A+G=[[1,2],[0,0]]). A sua char poly é a PARÁBOLA GERADORA x\ensuremath{^{2}}\ensuremath{-}x=x(x\ensuremath{-}1), raízes 0 e 1 — exatamente os NEUTROS do corpo estelar: \ensuremath{\nu}=0 (o PONTO, o vácuo, neutro de \ensuremath{\oplus}) e \ensuremath{\nu}=1 (a RETA, a identidade/La Hire, neutro de \ensuremath{\otimes}). É a parábola geradora CRUA (P\_g=x\ensuremath{^{2}} vs P\_a=x); somar o FLIP (a norma N=\ensuremath{-}1) dá a metálica do ouro x\ensuremath{^{2}}\ensuremath{-}x\ensuremath{-}1. A geradora + o \ensuremath{-}1 = o metal. conversa/gamba.py (parabola\_geradora).
\prov{relações: parte\_de algebra\_fundamental\_gato\_gamba; usa gato}
\prov{proveniência: wiki \textperiodcentered{} Floxina (investigação) \textperiodcentered{} fato \textperiodcentered{} confiança 1.0 \textperiodcentered{} domínio floxina \textperiodcentered{} história 222 versões no jornal}

%CRISTAL {"arestas": [{"alvo": "computador_mineral_promovido", "confianca": 1.0, "epistemico": "fato", "origem": "wiki", "rel": "parte_de"}, {"alvo": "computador_universal_pnp", "confianca": 1.0, "epistemico": "fato", "origem": "wiki", "rel": "usa"}, {"alvo": "pipe_transmissao_dual", "confianca": 1.0, "epistemico": "fato", "origem": "wiki", "rel": "usa"}, {"alvo": "refino_multimetal_live_gex44", "confianca": 1.0, "epistemico": "fato", "origem": "wiki", "rel": "relaciona"}, {"alvo": "gerador_do_gap_comutador", "confianca": 1.0, "epistemico": "fato", "origem": "wiki", "rel": "usa"}], "confianca": 1.0, "contraexemplos": [], "descricao": "O computador universal dá PARALELISMO: hoje o pipe comanda 1 matriz (1 portador); o universal comanda N CORPOS MOLECULARES em paralelo (N lanes), compostos pela LA HIRE (⊗, o produto dos portadores/gatos, unitário pois 𝕆 tem norma multiplicativa). A COSTURA FORMALIZA os N lanes de uma vez: (a) a MINERAÇÃO — os 12 metais (--lanes 12, live no gex44) = 12 corpos em paralelo, compostos por La Hire; (b) o pipe GRÁFICO 2D/3D — renderizar a CENA TODA de uma vez (as N facetas = N lanes, força N×, antes 1 de cada vez); (c) a transmissão. O portador é o octônio (dim 8 = o 7D imaginário, o MESMO da mineração/transformada_mineral_7d): 2D (pipe_svg), 3D (raytracer), 7D (octônio) no mesmo metal paralelo. Cada lane conserva a norma (o checksum a+c²=1). linguagem/omnitrix.py (computador_universal_paralelo, migrar_lib_no_metal, 4/4 cada); linguagem/transmissao.py (transmitir_paralelo, receber_paralelo, lahire_portadores, pipe_paralelo).", "epistemico": "hipotese", "exemplos": [], "id": "paralelismo_pipe_no_metal", "memoria": "semantica", "meta": {"dominio": "floxina", "fonte": "Floxina (investigação)"}, "origem": "wiki", "palavras_chave": [], "sinonimos": [], "tipo": "conceito", "titulo": "O paralelismo: o computador universal comanda N matrizes; a costura formaliza os lanes"}
\section{O paralelismo: o computador universal comanda N matrizes; a costura formaliza os lanes}
O computador universal dá PARALELISMO: hoje o pipe comanda 1 matriz (1 portador); o universal comanda N CORPOS MOLECULARES em paralelo (N lanes), compostos pela LA HIRE (\ensuremath{\otimes}, o produto dos portadores/gatos, unitário pois \ensuremath{\mathbb{O}} tem norma multiplicativa). A COSTURA FORMALIZA os N lanes de uma vez: (a) a MINERAÇÃO — os 12 metais (--lanes 12, live no gex44) = 12 corpos em paralelo, compostos por La Hire; (b) o pipe GRÁFICO 2D/3D — renderizar a CENA TODA de uma vez (as N facetas = N lanes, força N$\times$, antes 1 de cada vez); (c) a transmissão. O portador é o octônio (dim 8 = o 7D imaginário, o MESMO da mineração/transformada\_mineral\_7d): 2D (pipe\_svg), 3D (raytracer), 7D (octônio) no mesmo metal paralelo. Cada lane conserva a norma (o checksum a+c\ensuremath{^{2}}=1). linguagem/omnitrix.py (computador\_universal\_paralelo, migrar\_lib\_no\_metal, 4/4 cada); linguagem/transmissao.py (transmitir\_paralelo, receber\_paralelo, lahire\_portadores, pipe\_paralelo).
\prov{relações: parte\_de computador\_mineral\_promovido; usa computador\_universal\_pnp; usa pipe\_transmissao\_dual; relaciona refino\_multimetal\_live\_gex44; usa gerador\_do\_gap\_comutador}
\prov{proveniência: wiki \textperiodcentered{} Floxina (investigação) \textperiodcentered{} hipotese \textperiodcentered{} confiança 1.0 \textperiodcentered{} domínio floxina \textperiodcentered{} história 240 versões no jornal}

%CRISTAL {"arestas": [{"alvo": "gamba_lado_frio_clones", "confianca": 1.0, "epistemico": "fato", "origem": "wiki", "rel": "usa"}, {"alvo": "gamba_antissimetrico", "confianca": 1.0, "epistemico": "fato", "origem": "wiki", "rel": "usa"}, {"alvo": "dimensoes_do_pipe", "confianca": 1.0, "epistemico": "fato", "origem": "wiki", "rel": "relaciona"}, {"alvo": "mineracao_pipe_hook", "confianca": 1.0, "epistemico": "fato", "origem": "wiki", "rel": "relaciona"}, {"alvo": "biblioteca_pipe", "confianca": 1.0, "epistemico": "fato", "origem": "wiki", "rel": "relaciona"}, {"alvo": "chicote_emocional", "confianca": 1.0, "epistemico": "fato", "origem": "wiki", "rel": "relaciona"}], "confianca": 1.0, "contraexemplos": [], "descricao": "A dualidade FUNDA o pipeline de transmissão. Transmitir é carregar sem perda, e o portador só pode ser o lado REVERSÍVEL — o gambá (unitário, det=+1, norma conservada). O gato LÊ (medição/colapso: o render 3D, a leitura da mineração 7D); o gambá CARREGA (a transmissão). Implementado na lib (linguagem/transmissao.py, 6/6): 3D (quaternion ℍ) e 7D (octônion 𝕆), álgebras de DIVISÃO ⟹ o produto por um portador unitário é reversível (alternatividade: u⁻¹·(u·x)=x), e a norma MULTIPLICATIVA ⟹ a norma conservada É o CHECKSUM (a+c²=1): corrupção quebra a norma, o erro detecta-se. Os clones (a periodicidade fria) são a redundância que corrige. O flip de Wick é o codec (encode na fase/gambá, read-out na dissipação/gato). O pipe inteiro (a mineração transmite o resíduo, o canal do olho, a cristalchain) repousa nesta dualidade — dual ao que já temos do gato em 3D (render) e 7D (mineração).", "epistemico": "fato", "exemplos": [], "id": "pipe_transmissao_dual", "memoria": "semantica", "meta": {"dominio": "floxina", "fonte": "Floxina (investigação)"}, "origem": "wiki", "palavras_chave": [], "sinonimos": [], "tipo": "conceito", "titulo": "O pipe de transmissão (o dual do gato: o gambá carrega)"}
\section{O pipe de transmissão (o dual do gato: o gambá carrega)}
A dualidade FUNDA o pipeline de transmissão. Transmitir é carregar sem perda, e o portador só pode ser o lado REVERSÍVEL — o gambá (unitário, det=+1, norma conservada). O gato LÊ (medição/colapso: o render 3D, a leitura da mineração 7D); o gambá CARREGA (a transmissão). Implementado na lib (linguagem/transmissao.py, 6/6): 3D (quaternion \ensuremath{\mathbb{H}}) e 7D (octônion \ensuremath{\mathbb{O}}), álgebras de DIVISÃO \ensuremath{\Longrightarrow} o produto por um portador unitário é reversível (alternatividade: u\ensuremath{^{-1}}\textperiodcentered (u\textperiodcentered x)=x), e a norma MULTIPLICATIVA \ensuremath{\Longrightarrow} a norma conservada É o CHECKSUM (a+c\ensuremath{^{2}}=1): corrupção quebra a norma, o erro detecta-se. Os clones (a periodicidade fria) são a redundância que corrige. O flip de Wick é o codec (encode na fase/gambá, read-out na dissipação/gato). O pipe inteiro (a mineração transmite o resíduo, o canal do olho, a cristalchain) repousa nesta dualidade — dual ao que já temos do gato em 3D (render) e 7D (mineração).
\prov{relações: usa gamba\_lado\_frio\_clones; usa gamba\_antissimetrico; relaciona dimensoes\_do\_pipe; relaciona mineracao\_pipe\_hook; relaciona biblioteca\_pipe; relaciona chicote\_emocional}
\prov{proveniência: wiki \textperiodcentered{} Floxina (investigação) \textperiodcentered{} fato \textperiodcentered{} confiança 1.0 \textperiodcentered{} domínio floxina \textperiodcentered{} história 546 versões no jornal}

%CRISTAL {"arestas": [{"alvo": "ponte_quantica_conjunto_dougherty_e_orbitais_de_hidrogenio", "confianca": 1.0, "epistemico": "fato", "origem": "ponte_quantica_dougherty.md", "rel": "parte_de"}, {"alvo": "broca_os", "confianca": 0.52, "epistemico": "hipotese", "origem": "ingest", "rel": "referencia"}, {"alvo": "gato", "confianca": 0.52, "epistemico": "hipotese", "origem": "ingest", "rel": "analogia"}, {"alvo": "colapso", "confianca": 0.52, "epistemico": "hipotese", "origem": "ingest", "rel": "analogia"}, {"alvo": "a_medida_o_calor_de_volta", "confianca": 0.52, "epistemico": "hipotese", "origem": "ingest", "rel": "analogia"}], "confianca": 0.52, "contraexemplos": [], "descricao": "| Dougherty (hipótese) | Broca-SO (fato/paper) | relação | |---------------------|------------------------|---------| | toro / vórtice | colapso Koch / grade | analogia | | fibração toroidal | fibração de Hopf S³→S² | analogia | | φ escala | gold / régua de ouro | analogia | | ψ hidrogénio | mecânica quântica | referencia externa | | onda composta por tori | gato A / espectro toro | analogia | | observador projectivo | semântica / Peirce | analogia |", "epistemico": "hipotese", "exemplos": [], "id": "ponte_broca-so_analogias_registadas", "memoria": "semantica", "meta": {"dominio": "hipotese", "falsificavel": true, "nivel": 7}, "origem": "ponte_quantica_dougherty.md", "palavras_chave": [], "sinonimos": [], "tipo": "conceito", "titulo": "Ponte Broca-SO (analogias registadas)"}
\section{Ponte Broca-SO (analogias registadas)}
| Dougherty (hipótese) | Broca-SO (fato/paper) | relação | |---------------------|------------------------|---------| | toro / vórtice | colapso Koch / grade | analogia | | fibração toroidal | fibração de Hopf S\ensuremath{^{3}}\ensuremath{\to}S\ensuremath{^{2}} | analogia | | \ensuremath{\varphi} escala | gold / régua de ouro | analogia | | \ensuremath{\psi} hidrogénio | mecânica quântica | referencia externa | | onda composta por tori | gato A / espectro toro | analogia | | observador projectivo | semântica / Peirce | analogia |
\prov{relações: parte\_de ponte\_quantica\_conjunto\_dougherty\_e\_orbitais\_de\_hidrogenio; referencia broca\_os; analogia gato; analogia colapso; analogia a\_medida\_o\_calor\_de\_volta}
\prov{proveniência: ponte\_quantica\_dougherty.md \textperiodcentered{} hipotese \textperiodcentered{} confiança 0.52 \textperiodcentered{} domínio hipotese \textperiodcentered{} história 7 versões no jornal}

%CRISTAL {"arestas": [{"alvo": "parabola_geradora_idempotente", "confianca": 1.0, "epistemico": "fato", "origem": "wiki", "rel": "relaciona"}, {"alvo": "fractal_newton_reta_matricial", "confianca": 1.0, "epistemico": "fato", "origem": "wiki", "rel": "usa"}, {"alvo": "dicionario_tiol_gato_gamba", "confianca": 1.0, "epistemico": "fato", "origem": "wiki", "rel": "relaciona"}], "confianca": 1.0, "contraexemplos": [], "descricao": "O hidrogênio (1 próton + 1 elétron, 1s, Z=1) é o BIT do universo. A PONTE DE HIDROGÊNIO X-H···Y (a interação central dos −OH/−SH/−COOH que mapeámos) tem um POÇO DUPLO: o próton perto de X (estado 0) ou de Y (estado 1) — dois poços, dois estados. E a PARÁBOLA GERADORA aparece pela dinâmica da ponte: x²−x=x(x−1) tem raízes 0 e 1 = os DOIS POÇOS = os neutros ⊕/⊗ do corpo. O FRACTAL DE NEWTON é a BÚSSOLA: N(z)=z²/(2z−1) (Newton de x²−x), pontos fixos 0 e 1, as bacias acham os dois poços (semente<½→poço 0, >½→poço 1, split limpo na barreira ½; grau 2, sem fractal). A DINÂMICA: o próton NUM poço = o IDEMPOTENTE (A+G)²=A+G (o colapso a um poço = a medição); o tunelamento (0↔1) = o FLIP; a barreira = o GAP. A parábola geradora x²−x É a dinâmica da ponte de H. Verificado: linguagem/ponte_hidrogenio.py (4/4).", "epistemico": "hipotese", "exemplos": [], "id": "ponte_hidrogenio_parabola_geradora", "memoria": "semantica", "meta": {"dominio": "floxina", "fonte": "Floxina (investigação)"}, "origem": "wiki", "palavras_chave": [], "sinonimos": [], "tipo": "conceito", "titulo": "A ponte de H identifica a parábola geradora (Newton é a bússola)"}
\section{A ponte de H identifica a parábola geradora (Newton é a bússola)}
O hidrogênio (1 próton + 1 elétron, 1s, Z=1) é o BIT do universo. A PONTE DE HIDROGÊNIO X-H\textperiodcentered \textperiodcentered \textperiodcentered Y (a interação central dos \ensuremath{-}OH/\ensuremath{-}SH/\ensuremath{-}COOH que mapeámos) tem um POÇO DUPLO: o próton perto de X (estado 0) ou de Y (estado 1) — dois poços, dois estados. E a PARÁBOLA GERADORA aparece pela dinâmica da ponte: x\ensuremath{^{2}}\ensuremath{-}x=x(x\ensuremath{-}1) tem raízes 0 e 1 = os DOIS POÇOS = os neutros \ensuremath{\oplus}/\ensuremath{\otimes} do corpo. O FRACTAL DE NEWTON é a BÚSSOLA: N(z)=z\ensuremath{^{2}}/(2z\ensuremath{-}1) (Newton de x\ensuremath{^{2}}\ensuremath{-}x), pontos fixos 0 e 1, as bacias acham os dois poços (semente<\ensuremath{\tfrac12}\ensuremath{\to}poço 0, >\ensuremath{\tfrac12}\ensuremath{\to}poço 1, split limpo na barreira \ensuremath{\tfrac12}; grau 2, sem fractal). A DINÂMICA: o próton NUM poço = o IDEMPOTENTE (A+G)\ensuremath{^{2}}=A+G (o colapso a um poço = a medição); o tunelamento (0\ensuremath{\leftrightarrow}1) = o FLIP; a barreira = o GAP. A parábola geradora x\ensuremath{^{2}}\ensuremath{-}x É a dinâmica da ponte de H. Verificado: linguagem/ponte\_hidrogenio.py (4/4).
\prov{relações: relaciona parabola\_geradora\_idempotente; usa fractal\_newton\_reta\_matricial; relaciona dicionario\_tiol\_gato\_gamba}
\prov{proveniência: wiki \textperiodcentered{} Floxina (investigação) \textperiodcentered{} hipotese \textperiodcentered{} confiança 1.0 \textperiodcentered{} domínio floxina \textperiodcentered{} história 232 versões no jornal}

%CRISTAL {"arestas": [{"alvo": "ponte_quantica_dougherty", "confianca": 0.52, "epistemico": "fato", "origem": "ingest", "rel": "parte_de"}], "confianca": 0.52, "contraexemplos": [], "descricao": "> **Epistemologia:** hipótese externa falsificável (Buddy James / Dougherty Set). > **No grafo Broca-SO:** `epistemico=hipotese`, confiança ≤ 0.55. > **Analogia ≠ equivalência** — não substitui MQ standard nem papers Lopes/Broca.", "epistemico": "hipotese", "exemplos": [], "id": "ponte_quantica_conjunto_dougherty_e_orbitais_de_hidrogenio", "memoria": "semantica", "meta": {"dominio": "hipotese", "falsificavel": true, "nivel": 7}, "origem": "ponte_quantica_dougherty.md", "palavras_chave": [], "sinonimos": [], "tipo": "conceito", "titulo": "Ponte Quântica — Conjunto Dougherty e orbitais de hidrogénio"}
\section{Ponte Quântica — Conjunto Dougherty e orbitais de hidrogénio}
> **Epistemologia:** hipótese externa falsificável (Buddy James / Dougherty Set). > **No grafo Broca-SO:** `epistemico=hipotese`, confiança \ensuremath{\le} 0.55. > **Analogia \ensuremath{\ne} equivalência** — não substitui MQ standard nem papers Lopes/Broca.
\prov{relações: parte\_de ponte\_quantica\_dougherty}
\prov{proveniência: ponte\_quantica\_dougherty.md \textperiodcentered{} hipotese \textperiodcentered{} confiança 0.52 \textperiodcentered{} domínio hipotese \textperiodcentered{} história 8 versões no jornal}

%CRISTAL {"arestas": [{"alvo": "lacuna", "confianca": 0.52, "epistemico": "hipotese", "origem": "ingest", "rel": "referencia"}, {"alvo": "mecanica_quantica", "confianca": 1.0, "epistemico": "fato", "origem": "ingest", "rel": "referencia"}, {"alvo": "geometria", "confianca": 1.0, "epistemico": "fato", "origem": "ingest", "rel": "analogia"}, {"alvo": "colapso", "confianca": 1.0, "epistemico": "fato", "origem": "ingest", "rel": "analogia"}, {"alvo": "analogia", "confianca": 1.0, "epistemico": "fato", "origem": "ingest", "rel": "explica"}, {"alvo": "word", "confianca": 0.5, "epistemico": "fato", "origem": "wiki", "rel": "relaciona"}, {"alvo": "virtualizacao", "confianca": 0.5, "epistemico": "fato", "origem": "wiki", "rel": "relaciona"}, {"alvo": "while", "confianca": 0.5, "epistemico": "fato", "origem": "wiki", "rel": "relaciona"}, {"alvo": "vontade_de_poder", "confianca": 0.5, "epistemico": "fato", "origem": "wiki", "rel": "relaciona"}], "confianca": 0.52, "contraexemplos": [], "descricao": "Hipótese externa: ponte_quantica_dougherty.md", "epistemico": "hipotese", "exemplos": [], "id": "ponte_quantica_dougherty", "memoria": "semantica", "meta": {"dominio": "hipotese", "epistemico": "hipotese", "falsificavel": true, "nivel": 7}, "origem": "ingest", "palavras_chave": [], "sinonimos": [], "tipo": "conceito", "titulo": "ponte quantica dougherty"}
\section{ponte quantica dougherty}
Hipótese externa: ponte\_quantica\_dougherty.md
\prov{relações: referencia lacuna; referencia mecanica\_quantica; analogia geometria; analogia colapso; explica analogia; relaciona word; relaciona virtualizacao; relaciona while; relaciona vontade\_de\_poder}
\prov{proveniência: ingest \textperiodcentered{} hipotese \textperiodcentered{} confiança 0.52 \textperiodcentered{} domínio hipotese \textperiodcentered{} história 128 versões no jornal}

%CRISTAL {"arestas": [{"alvo": "matriz_costuras_paralela_universal", "confianca": 1.0, "epistemico": "fato", "origem": "wiki", "rel": "usa"}, {"alvo": "algebra_fundamental_gato_gamba", "confianca": 1.0, "epistemico": "fato", "origem": "wiki", "rel": "usa"}, {"alvo": "shader_fios_shell_paralelo", "confianca": 1.0, "epistemico": "fato", "origem": "wiki", "rel": "usa"}, {"alvo": "omnitrix_corpo_universal", "confianca": 1.0, "epistemico": "fato", "origem": "wiki", "rel": "usa"}], "confianca": 1.0, "contraexemplos": [], "descricao": "Uma primitiva geométrica (um SDF: a junta do IQ, a esfera, o toro) não é objeto solto — é POSICIONADA/ORIENTADA por uma MATRIZ (o transform afim), e as operações da cena sobre elas SÃO as operações do Corpo Algébrico: (1) o ⊗ LA HIRE (o produto) compõe os transforms — a ROTAÇÃO (a curvatura da junta, o ângulo a) = o ESQUILO (ortogonal, e^{aG}), a ESCALA/CISALHA = o GATO (simétrico), a TRANSLAÇÃO = o afim; compor = multiplicar matrizes (R·R=R(2a)). (2) o ⊕ SOMA (a UNIÃO CSG) d=min(d1,d2) = o ⊕ TROPICAL (o semianel min-plus: associativo, comutativo, IDEMPOTENTE; a soma distribui sobre o min); a INTERSECÇÃO = max (o dual). Logo o map() do IQ (duas juntas por min) É o Corpo Algébrico fazendo ⊕ sobre duas primitivas ⊗-transformadas — o raymarch de uma cena É a álgebra do Corpo Algébrico (⊗ os transforms, ⊕ a união). No pipe: FRAG_JOINT (simple1.txt) renderiza a junta SDF no metal (min/max/sign adicionados ao compilador). linguagem/omnitrix.py (primitivas_geometricas_matrizes, 5/5).", "epistemico": "fato", "exemplos": [], "id": "primitivas_geometricas_matrizes", "memoria": "semantica", "meta": {"dominio": "floxina", "fonte": "Floxina (investigação)"}, "origem": "wiki", "palavras_chave": [], "sinonimos": [], "tipo": "conceito", "titulo": "As primitivas são matrizes: o Corpo Algébrico as opera (as operações CSG são geométricas/matriciais)"}
\section{As primitivas são matrizes: o Corpo Algébrico as opera (as operações CSG são geométricas/matriciais)}
Uma primitiva geométrica (um SDF: a junta do IQ, a esfera, o toro) não é objeto solto — é POSICIONADA/ORIENTADA por uma MATRIZ (o transform afim), e as operações da cena sobre elas SÃO as operações do Corpo Algébrico: (1) o \ensuremath{\otimes} LA HIRE (o produto) compõe os transforms — a ROTAÇÃO (a curvatura da junta, o ângulo a) = o ESQUILO (ortogonal, e\^{}\{aG\}), a ESCALA/CISALHA = o GATO (simétrico), a TRANSLAÇÃO = o afim; compor = multiplicar matrizes (R\textperiodcentered R=R(2a)). (2) o \ensuremath{\oplus} SOMA (a UNIÃO CSG) d=min(d1,d2) = o \ensuremath{\oplus} TROPICAL (o semianel min-plus: associativo, comutativo, IDEMPOTENTE; a soma distribui sobre o min); a INTERSECÇÃO = max (o dual). Logo o map() do IQ (duas juntas por min) É o Corpo Algébrico fazendo \ensuremath{\oplus} sobre duas primitivas \ensuremath{\otimes}-transformadas — o raymarch de uma cena É a álgebra do Corpo Algébrico (\ensuremath{\otimes} os transforms, \ensuremath{\oplus} a união). No pipe: FRAG\_JOINT (simple1.txt) renderiza a junta SDF no metal (min/max/sign adicionados ao compilador). linguagem/omnitrix.py (primitivas\_geometricas\_matrizes, 5/5).
\prov{relações: usa matriz\_costuras\_paralela\_universal; usa algebra\_fundamental\_gato\_gamba; usa shader\_fios\_shell\_paralelo; usa omnitrix\_corpo\_universal}
\prov{proveniência: wiki \textperiodcentered{} Floxina (investigação) \textperiodcentered{} fato \textperiodcentered{} confiança 1.0 \textperiodcentered{} domínio floxina \textperiodcentered{} história 130 versões no jornal}

%CRISTAL {"arestas": [{"alvo": "gamba_antissimetrico", "confianca": 1.0, "epistemico": "fato", "origem": "wiki", "rel": "relaciona"}, {"alvo": "maquina_gaiola_esquilo_dtc", "confianca": 1.0, "epistemico": "fato", "origem": "wiki", "rel": "usa"}, {"alvo": "omnitrix_corpo_universal", "confianca": 1.0, "epistemico": "fato", "origem": "wiki", "rel": "relaciona"}], "confianca": 1.0, "contraexemplos": [], "descricao": "O operador ANTI-SIMÉTRICO G (o que GIRA: G²=−I, elíptico, e^{Gt} conserva a norma) — antes o 'gambá' — é PROMOVIDO a ESQUILO, batizado pela MÁQUINA DE GAIOLA DE ESQUILO (a máquina de indução polifásica), cuja marca é o CAMPO GIRANTE e^{Gt}: o EM feito máquina. O gato MEDE/dissipa (simétrico, hiperbólico); o esquilo GIRA (antissimétrico, elíptico). Feito nos papers do Corpo Algébrico (omnitrix.tex) e da joalheria (joalheria.tex): as referências a gambá viraram esquilo. Nota: o 'gambá' ANIMAL (o cheiro de gambá = o tiol/skunk, em corpo_molecular.tex) NÃO muda — só o OPERADOR. O código/os outros papers mantêm 'gambá' por ora (a promoção começou pelo Corpo Algébrico e a joalheria).", "epistemico": "fato", "exemplos": [], "id": "promocao_gamba_esquilo", "memoria": "semantica", "meta": {"dominio": "floxina", "fonte": "Floxina (investigação)"}, "origem": "wiki", "palavras_chave": [], "sinonimos": [], "tipo": "conceito", "titulo": "Promoção: o gambá → o ESQUILO (batizado pela máquina de gaiola de esquilo)"}
\section{Promoção: o gambá \ensuremath{\to} o ESQUILO (batizado pela máquina de gaiola de esquilo)}
O operador ANTI-SIMÉTRICO G (o que GIRA: G\ensuremath{^{2}}=\ensuremath{-}I, elíptico, e\^{}\{Gt\} conserva a norma) — antes o 'gambá' — é PROMOVIDO a ESQUILO, batizado pela MÁQUINA DE GAIOLA DE ESQUILO (a máquina de indução polifásica), cuja marca é o CAMPO GIRANTE e\^{}\{Gt\}: o EM feito máquina. O gato MEDE/dissipa (simétrico, hiperbólico); o esquilo GIRA (antissimétrico, elíptico). Feito nos papers do Corpo Algébrico (omnitrix.tex) e da joalheria (joalheria.tex): as referências a gambá viraram esquilo. Nota: o 'gambá' ANIMAL (o cheiro de gambá = o tiol/skunk, em corpo\_molecular.tex) NÃO muda — só o OPERADOR. O código/os outros papers mantêm 'gambá' por ora (a promoção começou pelo Corpo Algébrico e a joalheria).
\prov{relações: relaciona gamba\_antissimetrico; usa maquina\_gaiola\_esquilo\_dtc; relaciona omnitrix\_corpo\_universal}
\prov{proveniência: wiki \textperiodcentered{} Floxina (investigação) \textperiodcentered{} fato \textperiodcentered{} confiança 1.0 \textperiodcentered{} domínio floxina \textperiodcentered{} história 148 versões no jornal}

%CRISTAL {"arestas": [{"alvo": "ciencias", "confianca": -0.039, "epistemico": "fato", "origem": "manual", "rel": "referencia"}, {"alvo": "o_mapa_os_quatro_quadrantes", "confianca": 1.0, "epistemico": "fato", "origem": "paper", "rel": "parte_de"}, {"alvo": "gato", "confianca": 1.0, "epistemico": "fato", "origem": "paper", "rel": "referencia"}, {"alvo": "matrizes", "confianca": 1.0, "epistemico": "fato", "origem": "paper", "rel": "referencia"}, {"alvo": "mecanica_quantica", "confianca": 1.0, "epistemico": "fato", "origem": "paper", "rel": "parte_de"}, {"alvo": "termodinamica", "confianca": 1.0, "epistemico": "fato", "origem": "paper", "rel": "parte_de"}, {"alvo": "geometria", "confianca": 1.0, "epistemico": "fato", "origem": "paper", "rel": "parte_de"}], "confianca": 1.0, "contraexemplos": [], "descricao": "Os quatro quadrantes do gato: M=U·P cruzado com estrutura algébrica; mapa fractal entre domínios.", "epistemico": "fato", "exemplos": [], "id": "quatro_ciencias", "memoria": "semantica", "meta": {"ciencia": "interdisciplinar", "dominio": "meta", "verificacao": "slug idêntico — enriquecer, não duplicar"}, "origem": "paper", "palavras_chave": [], "sinonimos": ["as quatro ciências", "quatro ciências", "quatro quadrantes", "quatro ciencias"], "tipo": "conceito", "titulo": "quatro_ciencias"}
\section{quatro\_ciencias}
Os quatro quadrantes do gato: M=U\textperiodcentered P cruzado com estrutura algébrica; mapa fractal entre domínios.
\prov{sinónimos: as quatro ciências; quatro ciências; quatro quadrantes; quatro ciencias}
\prov{relações: referencia ciencias; parte\_de o\_mapa\_os\_quatro\_quadrantes; referencia gato; referencia matrizes; parte\_de mecanica\_quantica; parte\_de termodinamica; parte\_de geometria}
\prov{proveniência: paper \textperiodcentered{} fato \textperiodcentered{} confiança 1.0 \textperiodcentered{} domínio meta \textperiodcentered{} história 285 versões no jornal}

%CRISTAL {"arestas": [{"alvo": "gamba_antissimetrico", "confianca": 1.0, "epistemico": "fato", "origem": "wiki", "rel": "usa"}, {"alvo": "gato", "confianca": 1.0, "epistemico": "fato", "origem": "wiki", "rel": "usa"}, {"alvo": "flip_e_rotacao_de_wick", "confianca": 1.0, "epistemico": "fato", "origem": "wiki", "rel": "relaciona"}, {"alvo": "corpo_mineral", "confianca": 1.0, "epistemico": "fato", "origem": "wiki", "rel": "relaciona"}], "confianca": 1.0, "contraexemplos": [], "descricao": "É tudo uma coisa só: o corpo estelar. Gato e gambá não são dois objetos — são os dois EIXOS de um mesmo plano (o flip u=log z): o eixo real λ=log σ (dissipativo, o gato) e o eixo imaginário θ=arg z (unitário, o gambá). O que obtivemos é a metade de baixo (dissipativo/hiperbólico); ESPELHANDO PARA CIMA (λ→−λ, o tempo revertido) fecham-se os QUATRO QUADRANTES = o plano traço-determinante: cristal (contrai) e caos (expande) — os dois ramos reais do gato — vezes as duas orientações do gambá (gira ±). A borda (Salem, |z|=1) é o centro puro (só o gambá gira, det=+1). É a parábola cristal/borda/caos do corpo estelar com a quarta direção (a rotação) explícita. Espelhar não cria mundo novo: o de baixo já era metade do mesmo.", "epistemico": "hipotese", "exemplos": [], "id": "quatro_quadrantes_corpo_estelar", "memoria": "semantica", "meta": {"dominio": "floxina", "fonte": "Floxina (investigação)"}, "origem": "wiki", "palavras_chave": [], "sinonimos": [], "tipo": "conceito", "titulo": "Um corpo estelar — os quatro quadrantes (gato+gambá)"}
\section{Um corpo estelar — os quatro quadrantes (gato+gambá)}
É tudo uma coisa só: o corpo estelar. Gato e gambá não são dois objetos — são os dois EIXOS de um mesmo plano (o flip u=log z): o eixo real \ensuremath{\lambda}=log \ensuremath{\sigma} (dissipativo, o gato) e o eixo imaginário \ensuremath{\theta}=arg z (unitário, o gambá). O que obtivemos é a metade de baixo (dissipativo/hiperbólico); ESPELHANDO PARA CIMA (\ensuremath{\lambda}\ensuremath{\to}\ensuremath{-}\ensuremath{\lambda}, o tempo revertido) fecham-se os QUATRO QUADRANTES = o plano traço-determinante: cristal (contrai) e caos (expande) — os dois ramos reais do gato — vezes as duas orientações do gambá (gira $\pm$). A borda (Salem, |z|=1) é o centro puro (só o gambá gira, det=+1). É a parábola cristal/borda/caos do corpo estelar com a quarta direção (a rotação) explícita. Espelhar não cria mundo novo: o de baixo já era metade do mesmo.
\prov{relações: usa gamba\_antissimetrico; usa gato; relaciona flip\_e\_rotacao\_de\_wick; relaciona corpo\_mineral}
\prov{proveniência: wiki \textperiodcentered{} Floxina (investigação) \textperiodcentered{} hipotese \textperiodcentered{} confiança 1.0 \textperiodcentered{} domínio floxina \textperiodcentered{} história 410 versões no jornal}

%CRISTAL {"arestas": [{"alvo": "missao_engenharia_ouro", "confianca": 1.0, "epistemico": "fato", "origem": "wiki", "rel": "parte_de"}, {"alvo": "corpo_tropical", "confianca": 1.0, "epistemico": "fato", "origem": "wiki", "rel": "usa"}, {"alvo": "quantizacao_topologica_costura", "confianca": 1.0, "epistemico": "fato", "origem": "wiki", "rel": "relaciona"}], "confianca": 1.0, "contraexemplos": [], "descricao": "O teorema de 'O Corpo Tropical' (Δ=0 ⟺ ℓ=2logσ=0 ⟺ nilpotente N²=0, 'nada se move') classifica T DENTRO de um nível — ℓ=0 é 'sem eixo NESTA resolução', não o fundo. Não se rasga o ponto ψ (não carrega estrutura); rasga-se o GERME: o blow-up normal do estrato 𝒫=Δ=0 sobre a família a um parâmetro. O operador R desce ao espaço graduado, onde reaparece um SUB-GAP dₙ+₁~4bₙ — uma vibração que o nível de cima não via. O nilpotente é terminal no seu andar e porta de entrada no de baixo. [D] o blow-up (geometria de deformações); [I] a identificação com 'rasgar a função de onda'.", "epistemico": "inferencia", "exemplos": [], "id": "rasgar_funcao_de_onda", "memoria": "semantica", "meta": {"autor": "Lopes/conselho", "dominio": "missao", "fonte": "machine/Missao.tex"}, "origem": "wiki", "palavras_chave": [], "sinonimos": [], "tipo": "conceito", "titulo": "Rasgar a função de onda (o nilpotente ℓ=0 não é o fundo)"}
\section{Rasgar a função de onda (o nilpotente \ensuremath{\ell}=0 não é o fundo)}
O teorema de 'O Corpo Tropical' (\ensuremath{\Delta}=0 \ensuremath{\Longleftrightarrow} \ensuremath{\ell}=2log\ensuremath{\sigma}=0 \ensuremath{\Longleftrightarrow} nilpotente N\ensuremath{^{2}}=0, 'nada se move') classifica T DENTRO de um nível — \ensuremath{\ell}=0 é 'sem eixo NESTA resolução', não o fundo. Não se rasga o ponto \ensuremath{\psi} (não carrega estrutura); rasga-se o GERME: o blow-up normal do estrato \ensuremath{\mathcal{P}}=\ensuremath{\Delta}=0 sobre a família a um parâmetro. O operador R desce ao espaço graduado, onde reaparece um SUB-GAP d\ensuremath{_{n}}+\ensuremath{_{1}}\~{}4b\ensuremath{_{n}} — uma vibração que o nível de cima não via. O nilpotente é terminal no seu andar e porta de entrada no de baixo. [D] o blow-up (geometria de deformações); [I] a identificação com 'rasgar a função de onda'.
\prov{relações: parte\_de missao\_engenharia\_ouro; usa corpo\_tropical; relaciona quantizacao\_topologica\_costura}
\prov{proveniência: wiki \textperiodcentered{} machine/Missao.tex \textperiodcentered{} inferencia \textperiodcentered{} confiança 1.0 \textperiodcentered{} domínio missao \textperiodcentered{} história 5 versões no jornal}

%CRISTAL {"arestas": [{"alvo": "psicologia", "confianca": 1.0, "epistemico": "fato", "origem": "wiki", "rel": "parte_de"}], "confianca": 1.0, "contraexemplos": [], "descricao": "James Gross: dá para regular a emoção ANTES (reavaliação — mudar o sentido do que acontece, barato e eficaz) ou DEPOIS (supressão — segurar a expressão, cara em custo cognitivo e saúde). Reavaliar = GIRAR o sentido (gerenciar); suprimir = segurar/amortecer (reprimir). O coração da Tiffany reavalia, não suprime.", "epistemico": "fato", "exemplos": [], "id": "regulacao_emocional_gross", "memoria": "semantica", "meta": {"autor": "Gross", "dominio": "coracao", "fonte": "coracao hub"}, "origem": "wiki", "palavras_chave": [], "sinonimos": [], "tipo": "conceito", "titulo": "A regulação emocional (Gross): reavaliar vs suprimir"}
\section{A regulação emocional (Gross): reavaliar vs suprimir}
James Gross: dá para regular a emoção ANTES (reavaliação — mudar o sentido do que acontece, barato e eficaz) ou DEPOIS (supressão — segurar a expressão, cara em custo cognitivo e saúde). Reavaliar = GIRAR o sentido (gerenciar); suprimir = segurar/amortecer (reprimir). O coração da Tiffany reavalia, não suprime.
\prov{relações: parte\_de psicologia}
\prov{proveniência: wiki \textperiodcentered{} coracao hub \textperiodcentered{} fato \textperiodcentered{} confiança 1.0 \textperiodcentered{} domínio coracao \textperiodcentered{} história 6 versões no jornal}

%CRISTAL {"arestas": [{"alvo": "ginzburg_toro_helice_e_3d-sst", "confianca": 1.0, "epistemico": "fato", "origem": "ginzburg_toro_helice.md", "rel": "parte_de"}, {"alvo": "birkeland_correntes_plasma", "confianca": 0.52, "epistemico": "hipotese", "origem": "ingest", "rel": "analogia"}], "confianca": 0.52, "contraexemplos": [], "descricao": "Filamentos aurorais helicoidais (Birkeland, NASA) são **analogia visual** com helyces Ginzburg — plasma real ≠ postulado 3D-SST.", "epistemico": "hipotese", "exemplos": [], "id": "relacao_com_birkeland_fato", "memoria": "semantica", "meta": {"dominio": "hipotese", "falsificavel": true, "nivel": 7}, "origem": "ginzburg_toro_helice.md", "palavras_chave": [], "sinonimos": [], "tipo": "conceito", "titulo": "Relação com Birkeland (fato)"}
\section{Relação com Birkeland (fato)}
Filamentos aurorais helicoidais (Birkeland, NASA) são **analogia visual** com helyces Ginzburg — plasma real \ensuremath{\ne} postulado 3D-SST.
\prov{relações: parte\_de ginzburg\_toro\_helice\_e\_3d-sst; analogia birkeland\_correntes\_plasma}
\prov{proveniência: ginzburg\_toro\_helice.md \textperiodcentered{} hipotese \textperiodcentered{} confiança 0.52 \textperiodcentered{} domínio hipotese \textperiodcentered{} história 4 versões no jornal}

%CRISTAL {"arestas": [{"alvo": "ginzburg_toro_helice_e_3d-sst", "confianca": 1.0, "epistemico": "fato", "origem": "ginzburg_toro_helice.md", "rel": "parte_de"}, {"alvo": "ponte_quantica_dougherty", "confianca": 0.52, "epistemico": "hipotese", "origem": "ingest", "rel": "referencia"}, {"alvo": "conjunto_dougherty", "confianca": 0.52, "epistemico": "hipotese", "origem": "ingest", "rel": "analogia"}, {"alvo": "word", "confianca": 0.5, "epistemico": "fato", "origem": "wiki", "rel": "relaciona"}, {"alvo": "while", "confianca": 0.5, "epistemico": "fato", "origem": "wiki", "rel": "relaciona"}, {"alvo": "virtualizacao", "confianca": 0.5, "epistemico": "fato", "origem": "wiki", "rel": "relaciona"}, {"alvo": "vontade_de_poder", "confianca": 0.5, "epistemico": "fato", "origem": "wiki", "rel": "relaciona"}], "confianca": 0.52, "contraexemplos": [], "descricao": "Paralelo **independente** citado na ponte Dougherty:", "epistemico": "hipotese", "exemplos": [], "id": "relacao_com_dougherty", "memoria": "semantica", "meta": {"dominio": "hipotese", "falsificavel": true, "nivel": 7}, "origem": "ginzburg_toro_helice.md", "palavras_chave": [], "sinonimos": [], "tipo": "conceito", "titulo": "Relação com Dougherty"}
\section{Relação com Dougherty}
Paralelo **independente** citado na ponte Dougherty:
\prov{relações: parte\_de ginzburg\_toro\_helice\_e\_3d-sst; referencia ponte\_quantica\_dougherty; analogia conjunto\_dougherty; relaciona word; relaciona while; relaciona virtualizacao; relaciona vontade\_de\_poder}
\prov{proveniência: ginzburg\_toro\_helice.md \textperiodcentered{} hipotese \textperiodcentered{} confiança 0.52 \textperiodcentered{} domínio hipotese \textperiodcentered{} história 9 versões no jornal}

%CRISTAL {"arestas": [{"alvo": "corpo_matricial_matrix", "confianca": 1.0, "epistemico": "fato", "origem": "wiki", "rel": "parte_de"}, {"alvo": "cf_de_pi_e_produto_de_gatos", "confianca": 1.0, "epistemico": "fato", "origem": "wiki", "rel": "usa"}, {"alvo": "selberg_meta_inducao", "confianca": 1.0, "epistemico": "fato", "origem": "wiki", "rel": "relaciona"}, {"alvo": "ancora_selberg_e_seu_escopo", "confianca": 1.0, "epistemico": "fato", "origem": "wiki", "rel": "relaciona"}, {"alvo": "parabola_geradora_idempotente", "confianca": 1.0, "epistemico": "fato", "origem": "wiki", "rel": "relaciona"}], "confianca": 1.0, "contraexemplos": [], "descricao": "A completação da reta matricial (e a resolução do cone nulo) usa o MESMO procedimento — sem mecanismo novo, análogo ao RH/Selberg (o dual discreto→compacto): a fração contínua de π É um PRODUTO DE GATOS [[a_k,1],[1,0]] (cf_de_pi_e_produto_de_gatos), com det=(−1)^k, |det|=1 CONSTANTE (a invariante mantida por PA+PG). Det=±1 e inteiro ⟹ o produto vive em GL(2,ℤ), o GRUPO MODULAR; o quociente GL(2,ℤ)\\ℍ é a SUPERFÍCIE MODULAR, COMPACTA (volume finito). A reta discreta (os inteiros a_k) COMPACTA-SE na superfície modular pela CF, na base π — a CF/Gauss é o fluxo geodésico na modular (o Selberg/Anosov). A parábola geradora DISCRETA (mineral, x²−x) ressoa com a COMPACTA (matricial). Verificado: os convergentes 3/1,22/7,355/113,…→π com det ±1. DIREÇÃO (a investigação segue, não decretada): o fractal de Newton, o grupo de grupos, a leminiscata newtoniana discreta — que completam a reta matricial.", "epistemico": "hipotese", "exemplos": [], "id": "reta_matricial_compacta_cf", "memoria": "semantica", "meta": {"dominio": "floxina", "fonte": "Floxina (investigação)"}, "origem": "wiki", "palavras_chave": [], "sinonimos": [], "tipo": "conceito", "titulo": "A reta matricial compacta-se pela CF (produto de gatos) na base π"}
\section{A reta matricial compacta-se pela CF (produto de gatos) na base \ensuremath{\pi}}
A completação da reta matricial (e a resolução do cone nulo) usa o MESMO procedimento — sem mecanismo novo, análogo ao RH/Selberg (o dual discreto\ensuremath{\to}compacto): a fração contínua de \ensuremath{\pi} É um PRODUTO DE GATOS [[a\_k,1],[1,0]] (cf\_de\_pi\_e\_produto\_de\_gatos), com det=(\ensuremath{-}1)\^{}k, |det|=1 CONSTANTE (a invariante mantida por PA+PG). Det=$\pm$1 e inteiro \ensuremath{\Longrightarrow} o produto vive em GL(2,\ensuremath{\mathbb{Z}}), o GRUPO MODULAR; o quociente GL(2,\ensuremath{\mathbb{Z}})\textbackslash{}\ensuremath{\mathbb{H}} é a SUPERFÍCIE MODULAR, COMPACTA (volume finito). A reta discreta (os inteiros a\_k) COMPACTA-SE na superfície modular pela CF, na base \ensuremath{\pi} — a CF/Gauss é o fluxo geodésico na modular (o Selberg/Anosov). A parábola geradora DISCRETA (mineral, x\ensuremath{^{2}}\ensuremath{-}x) ressoa com a COMPACTA (matricial). Verificado: os convergentes 3/1,22/7,355/113,…\ensuremath{\to}\ensuremath{\pi} com det $\pm$1. DIREÇÃO (a investigação segue, não decretada): o fractal de Newton, o grupo de grupos, a leminiscata newtoniana discreta — que completam a reta matricial.
\prov{relações: parte\_de corpo\_matricial\_matrix; usa cf\_de\_pi\_e\_produto\_de\_gatos; relaciona selberg\_meta\_inducao; relaciona ancora\_selberg\_e\_seu\_escopo; relaciona parabola\_geradora\_idempotente}
\prov{proveniência: wiki \textperiodcentered{} Floxina (investigação) \textperiodcentered{} hipotese \textperiodcentered{} confiança 1.0 \textperiodcentered{} domínio floxina \textperiodcentered{} história 420 versões no jornal}

%CRISTAL {"arestas": [{"alvo": "psicologia", "confianca": 1.0, "epistemico": "fato", "origem": "wiki", "rel": "parte_de"}], "confianca": 1.0, "contraexemplos": [], "descricao": "Robert Plutchik: 8 emoções PRIMÁRIAS em 4 pares opostos (alegria↔tristeza, confiança↔aversão, medo↔raiva, surpresa↔antecipação), cada uma com gradação de intensidade (serenidade→alegria→êxtase) e combinações (díades: alegria+confiança=amor; medo+surpresa=alarme). Uma roda: adjacência = semelhança, oposição = antagonismo.", "epistemico": "fato", "exemplos": [], "id": "roda_emocoes_plutchik", "memoria": "semantica", "meta": {"autor": "Plutchik", "dominio": "coracao", "fonte": "coracao hub"}, "origem": "wiki", "palavras_chave": [], "sinonimos": [], "tipo": "conceito", "titulo": "A roda das emoções (Plutchik)"}
\section{A roda das emoções (Plutchik)}
Robert Plutchik: 8 emoções PRIMÁRIAS em 4 pares opostos (alegria\ensuremath{\leftrightarrow}tristeza, confiança\ensuremath{\leftrightarrow}aversão, medo\ensuremath{\leftrightarrow}raiva, surpresa\ensuremath{\leftrightarrow}antecipação), cada uma com gradação de intensidade (serenidade\ensuremath{\to}alegria\ensuremath{\to}êxtase) e combinações (díades: alegria+confiança=amor; medo+surpresa=alarme). Uma roda: adjacência = semelhança, oposição = antagonismo.
\prov{relações: parte\_de psicologia}
\prov{proveniência: wiki \textperiodcentered{} coracao hub \textperiodcentered{} fato \textperiodcentered{} confiança 1.0 \textperiodcentered{} domínio coracao \textperiodcentered{} história 6 versões no jornal}

%CRISTAL {"arestas": [{"alvo": "floxina_selberg", "confianca": 1.0, "epistemico": "fato", "origem": "wiki", "rel": "relaciona"}, {"alvo": "ancora_selberg_e_seu_escopo", "confianca": 1.0, "epistemico": "fato", "origem": "wiki", "rel": "relaciona"}, {"alvo": "dicionario_floxinato_gap_massa", "confianca": 1.0, "epistemico": "fato", "origem": "wiki", "rel": "relaciona"}], "confianca": 1.0, "contraexemplos": [], "descricao": "O mesmo padrão da meta-indução mineral (π→cada metal). RIEMANN nasceu discreto (a zeta ∑n^{−s} sobre os inteiros) e passou ao compacto (a superfície, onde a RH de Selberg é provada, 1956). SELBERG faz igual: a conjectura de valor próprio sai discreta (os subgrupos de CONGRUÊNCIA, aritméticos — λ₁≥¼ ABERTA) e generaliza ao COMPACTO (provado, o gap ¼). Mesmo procedimento, e DUAL: a FÓRMULA DO TRAÇO É a dualidade — o lado aritmético (discreto: classes, primos) = o lado espectral (compacto: autovalores, geodésicas), a soma de Poisson generalizada. Provar no compacto e levantar ao discreto pela dualidade é a meta-indução — o mesmo gesto que gera cada σ_n do anterior. FATO: a fórmula do traço costura os dois lados, o compacto provado, a congruência aberta; SÍNTESE: ler Riemann e Selberg como o mesmo dual.", "epistemico": "hipotese", "exemplos": [], "id": "selberg_meta_inducao", "memoria": "semantica", "meta": {"dominio": "floxina", "fonte": "Floxina (investigação)"}, "origem": "wiki", "palavras_chave": [], "sinonimos": [], "tipo": "conceito", "titulo": "A conjectura de Selberg é meta-indução (o dual discreto→compacto)"}
\section{A conjectura de Selberg é meta-indução (o dual discreto\ensuremath{\to}compacto)}
O mesmo padrão da meta-indução mineral (\ensuremath{\pi}\ensuremath{\to}cada metal). RIEMANN nasceu discreto (a zeta \ensuremath{\sum}n\^{}\{\ensuremath{-}s\} sobre os inteiros) e passou ao compacto (a superfície, onde a RH de Selberg é provada, 1956). SELBERG faz igual: a conjectura de valor próprio sai discreta (os subgrupos de CONGRUÊNCIA, aritméticos — \ensuremath{\lambda}\ensuremath{_{1}}\ensuremath{\ge}\ensuremath{\tfrac14} ABERTA) e generaliza ao COMPACTO (provado, o gap \ensuremath{\tfrac14}). Mesmo procedimento, e DUAL: a FÓRMULA DO TRAÇO É a dualidade — o lado aritmético (discreto: classes, primos) = o lado espectral (compacto: autovalores, geodésicas), a soma de Poisson generalizada. Provar no compacto e levantar ao discreto pela dualidade é a meta-indução — o mesmo gesto que gera cada \ensuremath{\sigma}\_n do anterior. FATO: a fórmula do traço costura os dois lados, o compacto provado, a congruência aberta; SÍNTESE: ler Riemann e Selberg como o mesmo dual.
\prov{relações: relaciona floxina\_selberg; relaciona ancora\_selberg\_e\_seu\_escopo; relaciona dicionario\_floxinato\_gap\_massa}
\prov{proveniência: wiki \textperiodcentered{} Floxina (investigação) \textperiodcentered{} hipotese \textperiodcentered{} confiança 1.0 \textperiodcentered{} domínio floxina \textperiodcentered{} história 316 versões no jornal}

%CRISTAL {"arestas": [{"alvo": "thc_tiol_propriedades_compartilhadas", "confianca": 1.0, "epistemico": "fato", "origem": "wiki", "rel": "relaciona"}, {"alvo": "dicionario_tiol_gato_gamba", "confianca": 1.0, "epistemico": "fato", "origem": "wiki", "rel": "relaciona"}, {"alvo": "gamba_antissimetrico", "confianca": 1.0, "epistemico": "fato", "origem": "wiki", "rel": "relaciona"}], "confianca": 1.0, "contraexemplos": [], "descricao": "Três moléculas de odor lidas no corpo, pelos dois eixos gato/gambá. Ácido 3-metil-hex-2-enoico (C₇H₁₂O₂, o odor corporal humano/suor): 9 átomos pesados=9 qubits, insaturação 2 (a dupla C=C + o C=O), ponta = ÁCIDO CARBOXÍLICO −COOH (pKa≈4,5, bem MAIS ácido que o tiol/fenol ~10 = um GAP MAIOR); os 2 O do carboxilo são simétricos no grafo (Z/2, a ressonância do carboxilato = o paralelismo). A SÉRIE (o gap cresce, o π varia): tiol −SH pKa≈10,6 / π nenhum (saturado); THC fenol −OH pKa≈10 / π anel aromático (forte); ácido −COOH pKa≈4,5 / π 1 C=C conjugado (α,β). A ponta reativa (o gambá) é COMPARTILHADA; a acidez=o GAP (cresce 10→10→4,5); o π=o GAMBÁ aromático (nenhum→conjugado→aromático). Verificado: linguagem/molecula_teste.py (9/9).", "epistemico": "hipotese", "exemplos": [], "id": "serie_cheiro_gap_gamba", "memoria": "semantica", "meta": {"dominio": "floxina", "fonte": "Floxina (investigação)"}, "origem": "wiki", "palavras_chave": [], "sinonimos": [], "tipo": "conceito", "titulo": "A série do cheiro: o gap (acidez) e o π (gambá) em 3 moléculas de odor"}
\section{A série do cheiro: o gap (acidez) e o \ensuremath{\pi} (gambá) em 3 moléculas de odor}
Três moléculas de odor lidas no corpo, pelos dois eixos gato/gambá. Ácido 3-metil-hex-2-enoico (C\ensuremath{_{7}}H\ensuremath{_{12}}O\ensuremath{_{2}}, o odor corporal humano/suor): 9 átomos pesados=9 qubits, insaturação 2 (a dupla C=C + o C=O), ponta = ÁCIDO CARBOXÍLICO \ensuremath{-}COOH (pKa\ensuremath{\approx}4,5, bem MAIS ácido que o tiol/fenol \~{}10 = um GAP MAIOR); os 2 O do carboxilo são simétricos no grafo (Z/2, a ressonância do carboxilato = o paralelismo). A SÉRIE (o gap cresce, o \ensuremath{\pi} varia): tiol \ensuremath{-}SH pKa\ensuremath{\approx}10,6 / \ensuremath{\pi} nenhum (saturado); THC fenol \ensuremath{-}OH pKa\ensuremath{\approx}10 / \ensuremath{\pi} anel aromático (forte); ácido \ensuremath{-}COOH pKa\ensuremath{\approx}4,5 / \ensuremath{\pi} 1 C=C conjugado (\ensuremath{\alpha},\ensuremath{\beta}). A ponta reativa (o gambá) é COMPARTILHADA; a acidez=o GAP (cresce 10\ensuremath{\to}10\ensuremath{\to}4,5); o \ensuremath{\pi}=o GAMBÁ aromático (nenhum\ensuremath{\to}conjugado\ensuremath{\to}aromático). Verificado: linguagem/molecula\_teste.py (9/9).
\prov{relações: relaciona thc\_tiol\_propriedades\_compartilhadas; relaciona dicionario\_tiol\_gato\_gamba; relaciona gamba\_antissimetrico}
\prov{proveniência: wiki \textperiodcentered{} Floxina (investigação) \textperiodcentered{} hipotese \textperiodcentered{} confiança 1.0 \textperiodcentered{} domínio floxina \textperiodcentered{} história 240 versões no jornal}

%CRISTAL {"arestas": [{"alvo": "paralelismo_pipe_no_metal", "confianca": 1.0, "epistemico": "fato", "origem": "wiki", "rel": "parte_de"}, {"alvo": "omnitrix_corpo_universal", "confianca": 1.0, "epistemico": "fato", "origem": "wiki", "rel": "relaciona"}], "confianca": 1.0, "contraexemplos": [], "descricao": "A técnica de GERAÇÃO DE FIOS (fur/pelo) por SHELL RENDERING (o fur ball clássico, Simon Green): uma esfera coberta de fios, renderizada por CAMADAS (shells) — raia-se para dentro do volume de pelo (furDepth), e em cada camada amostra-se a DENSIDADE de fios de um mapa de FOLÍCULOS (um grid de hashes: cada célula um fio, com comprimento). O fio aparece onde a densidade passa o limiar e desvanece ao longo do comprimento (t de 0 na raiz a 1 na ponta); acumula-se front-to-back (alpha pré-multiplicado); a CURL (rotateX por cos(iTime)·t²) balança os fios. O PARALELISMO: todos os raios (a grade H×W) marcham por todas as shells AO MESMO TEMPO — a cena inteira num passo vetorizado (o mesmo N-lane do computador universal: cada pixel um lane). As shells são o laço sequencial (a profundidade); os pixels o paralelo. O mapa de fios é PROCEDURAL (o folículo=hash→célula, a mesma ideia do floco mineral). Ferramenta FIXA no pipe gráfico. linguagem/shader_fios.py (verificar_fios 5/5). Origem: advanced2.txt (Simon Green, uso local; o módulo é implementação própria da técnica, não cópia).", "epistemico": "hipotese", "exemplos": [], "id": "shader_fios_shell_paralelo", "memoria": "semantica", "meta": {"dominio": "floxina", "fonte": "Floxina (investigação)"}, "origem": "wiki", "palavras_chave": [], "sinonimos": [], "tipo": "conceito", "titulo": "A geração de fios (fur ball) por shell rendering, vetorizada (o paralelismo, ferramenta fixa no pipe)"}
\section{A geração de fios (fur ball) por shell rendering, vetorizada (o paralelismo, ferramenta fixa no pipe)}
A técnica de GERAÇÃO DE FIOS (fur/pelo) por SHELL RENDERING (o fur ball clássico, Simon Green): uma esfera coberta de fios, renderizada por CAMADAS (shells) — raia-se para dentro do volume de pelo (furDepth), e em cada camada amostra-se a DENSIDADE de fios de um mapa de FOLÍCULOS (um grid de hashes: cada célula um fio, com comprimento). O fio aparece onde a densidade passa o limiar e desvanece ao longo do comprimento (t de 0 na raiz a 1 na ponta); acumula-se front-to-back (alpha pré-multiplicado); a CURL (rotateX por cos(iTime)\textperiodcentered t\ensuremath{^{2}}) balança os fios. O PARALELISMO: todos os raios (a grade H$\times$W) marcham por todas as shells AO MESMO TEMPO — a cena inteira num passo vetorizado (o mesmo N-lane do computador universal: cada pixel um lane). As shells são o laço sequencial (a profundidade); os pixels o paralelo. O mapa de fios é PROCEDURAL (o folículo=hash\ensuremath{\to}célula, a mesma ideia do floco mineral). Ferramenta FIXA no pipe gráfico. linguagem/shader\_fios.py (verificar\_fios 5/5). Origem: advanced2.txt (Simon Green, uso local; o módulo é implementação própria da técnica, não cópia).
\prov{relações: parte\_de paralelismo\_pipe\_no\_metal; relaciona omnitrix\_corpo\_universal}
\prov{proveniência: wiki \textperiodcentered{} Floxina (investigação) \textperiodcentered{} hipotese \textperiodcentered{} confiança 1.0 \textperiodcentered{} domínio floxina \textperiodcentered{} história 117 versões no jornal}

%CRISTAL {"arestas": [{"alvo": "corpo_equacoes_diferenciais", "confianca": 1.0, "epistemico": "fato", "origem": "wiki", "rel": "usa"}, {"alvo": "pa_pg_corpo_diferencial", "confianca": 1.0, "epistemico": "fato", "origem": "wiki", "rel": "usa"}, {"alvo": "solucao_produto_lahire_ed", "confianca": 1.0, "epistemico": "fato", "origem": "wiki", "rel": "relaciona"}], "confianca": 1.0, "contraexemplos": [], "descricao": "Toda ED tem solução EXPLÍCITA — e é o fluxo e^{At} (a exponencial do gerador), em três formas: (1) LINEAR homogênea ẋ=Ax: x(t)=e^{At}x₀=∑ᵢ cᵢ e^{λᵢt} vᵢ, a SUPERPOSIÇÃO (⊕) dos MODOS e^{λᵢt} (as PGs, os autovalores=os metais); espectralmente e^{At}=∑ e^{λt}Pᵢ. (2) INOMOGÊNEA ẋ=Ax+b: variação dos parâmetros (Duhamel) x(t)=e^{At}x₀+A⁻¹(e^{At}−I)b — o ∫ (que SOBE o grau) acrescenta a resposta forçada. (3) NÃO-LINEAR ẋ=f(x): a série de PICARD x_{n+1}=x₀+∫f(xₙ) converge (Picard-Lindelöf); para ẋ=x É a série ∑tᵏ/k!→e^t. E e^{At}=∑ₙ(At)ⁿ/n! — a solução é a exponencial do gerador, sempre. linguagem/omnitrix.py (solucao_explicita_ed 5/5); papers/equacoes_diferenciais.tex (Parte VII).", "epistemico": "fato", "exemplos": [], "id": "solucao_explicita_ed", "memoria": "semantica", "meta": {"dominio": "floxina", "fonte": "Floxina (investigação)"}, "origem": "wiki", "palavras_chave": [], "sinonimos": [], "tipo": "conceito", "titulo": "A solução explícita de TODA ED: x(t)=e^{At}x₀ (o fluxo=⊕ dos modos/PGs); Duhamel (a fonte), Picard (o não-linear)"}
\section{A solução explícita de TODA ED: x(t)=e\^{}\{At\}x\ensuremath{_{0}} (o fluxo=\ensuremath{\oplus} dos modos/PGs); Duhamel (a fonte), Picard (o não-linear)}
Toda ED tem solução EXPLÍCITA — e é o fluxo e\^{}\{At\} (a exponencial do gerador), em três formas: (1) LINEAR homogênea \ensuremath{\dot{x}}=Ax: x(t)=e\^{}\{At\}x\ensuremath{_{0}}=\ensuremath{\sum}\ensuremath{_{i}} c\ensuremath{_{i}} e\^{}\{\ensuremath{\lambda}\ensuremath{_{i}}t\} v\ensuremath{_{i}}, a SUPERPOSIÇÃO (\ensuremath{\oplus}) dos MODOS e\^{}\{\ensuremath{\lambda}\ensuremath{_{i}}t\} (as PGs, os autovalores=os metais); espectralmente e\^{}\{At\}=\ensuremath{\sum} e\^{}\{\ensuremath{\lambda}t\}P\ensuremath{_{i}}. (2) INOMOGÊNEA \ensuremath{\dot{x}}=Ax+b: variação dos parâmetros (Duhamel) x(t)=e\^{}\{At\}x\ensuremath{_{0}}+A\ensuremath{^{-1}}(e\^{}\{At\}\ensuremath{-}I)b — o \ensuremath{\int} (que SOBE o grau) acrescenta a resposta forçada. (3) NÃO-LINEAR \ensuremath{\dot{x}}=f(x): a série de PICARD x\_\{n+1\}=x\ensuremath{_{0}}+\ensuremath{\int}f(x\ensuremath{_{n}}) converge (Picard-Lindelöf); para \ensuremath{\dot{x}}=x É a série \ensuremath{\sum}t\ensuremath{^{k}}/k!\ensuremath{\to}e\^{}t. E e\^{}\{At\}=\ensuremath{\sum}\ensuremath{_{n}}(At)\ensuremath{^{n}}/n! — a solução é a exponencial do gerador, sempre. linguagem/omnitrix.py (solucao\_explicita\_ed 5/5); papers/equacoes\_diferenciais.tex (Parte VII).
\prov{relações: usa corpo\_equacoes\_diferenciais; usa pa\_pg\_corpo\_diferencial; relaciona solucao\_produto\_lahire\_ed}
\prov{proveniência: wiki \textperiodcentered{} Floxina (investigação) \textperiodcentered{} fato \textperiodcentered{} confiança 1.0 \textperiodcentered{} domínio floxina \textperiodcentered{} história 56 versões no jornal}

%CRISTAL {"arestas": [{"alvo": "solucao_explicita_ed", "confianca": 1.0, "epistemico": "fato", "origem": "wiki", "rel": "usa"}, {"alvo": "fracoes_continuas_ed", "confianca": 1.0, "epistemico": "fato", "origem": "wiki", "rel": "relaciona"}, {"alvo": "selberg_meta_inducao", "confianca": 1.0, "epistemico": "fato", "origem": "wiki", "rel": "relaciona"}], "confianca": 1.0, "contraexemplos": [], "descricao": "Na matrix (corpo_estelar Parte XIV) a solução explícita é a correspondência 1-a-1 p↦√p (os PRIMOS nos IRRACIONAIS, injetiva), e o metal σ_m (o irracional, a cúspide CONTÍNUA) ↔ o primo/a geodésica (o INTEIRO) — as duas faces da MESMA parábola. A solução explícita da ED faz o MESMO desenho: o ESPECTRO discreto {λᵢ} (os «primos» da ED) gera os MODOS contínuos {e^{λᵢt}} (os transcendentais, os «irracionais»), e a EXPONENCIAL é a ponte: σ_m=e^{λ_m} (o metal=o exp da reta λ=log σ). O selo é o TRAÇO: Tr(e^{At})=∑ᵢ e^{λᵢt} (a soma sobre o espectro) ↔ a zeta dos metais ∑σ^{−β} ↔ a contagem de geodésicas/primos (o traço de Selberg). A solução explícita da ED (espectro→modos, discreto gerando contínuo) É o análogo de primos→irracionais — as duas faces da parábola geradora. linguagem/omnitrix.py (solucao_explicita_primos_irracionais_ed 5/5); papers/equacoes_diferenciais.tex (Parte VII).", "epistemico": "fato", "exemplos": [], "id": "solucao_explicita_primos_irracionais_ed", "memoria": "semantica", "meta": {"dominio": "floxina", "fonte": "Floxina (investigação)"}, "origem": "wiki", "palavras_chave": [], "sinonimos": [], "tipo": "conceito", "titulo": "A solução explícita da ED é ANÁLOGA a primos↔irracionais (matrix): o espectro discreto {λ} gera os modos contínuos {e^{λt}}"}
\section{A solução explícita da ED é ANÁLOGA a primos\ensuremath{\leftrightarrow}irracionais (matrix): o espectro discreto \{\ensuremath{\lambda}\} gera os modos contínuos \{e\^{}\{\ensuremath{\lambda}t\}\}}
Na matrix (corpo\_estelar Parte XIV) a solução explícita é a correspondência 1-a-1 p\ensuremath{\mapsto}\ensuremath{\sqrt{\,}}p (os PRIMOS nos IRRACIONAIS, injetiva), e o metal \ensuremath{\sigma}\_m (o irracional, a cúspide CONTÍNUA) \ensuremath{\leftrightarrow} o primo/a geodésica (o INTEIRO) — as duas faces da MESMA parábola. A solução explícita da ED faz o MESMO desenho: o ESPECTRO discreto \{\ensuremath{\lambda}\ensuremath{_{i}}\} (os «primos» da ED) gera os MODOS contínuos \{e\^{}\{\ensuremath{\lambda}\ensuremath{_{i}}t\}\} (os transcendentais, os «irracionais»), e a EXPONENCIAL é a ponte: \ensuremath{\sigma}\_m=e\^{}\{\ensuremath{\lambda}\_m\} (o metal=o exp da reta \ensuremath{\lambda}=log \ensuremath{\sigma}). O selo é o TRAÇO: Tr(e\^{}\{At\})=\ensuremath{\sum}\ensuremath{_{i}} e\^{}\{\ensuremath{\lambda}\ensuremath{_{i}}t\} (a soma sobre o espectro) \ensuremath{\leftrightarrow} a zeta dos metais \ensuremath{\sum}\ensuremath{\sigma}\^{}\{\ensuremath{-}\ensuremath{\beta}\} \ensuremath{\leftrightarrow} a contagem de geodésicas/primos (o traço de Selberg). A solução explícita da ED (espectro\ensuremath{\to}modos, discreto gerando contínuo) É o análogo de primos\ensuremath{\to}irracionais — as duas faces da parábola geradora. linguagem/omnitrix.py (solucao\_explicita\_primos\_irracionais\_ed 5/5); papers/equacoes\_diferenciais.tex (Parte VII).
\prov{relações: usa solucao\_explicita\_ed; relaciona fracoes\_continuas\_ed; relaciona selberg\_meta\_inducao}
\prov{proveniência: wiki \textperiodcentered{} Floxina (investigação) \textperiodcentered{} fato \textperiodcentered{} confiança 1.0 \textperiodcentered{} domínio floxina \textperiodcentered{} história 55 versões no jornal}

%CRISTAL {"arestas": [{"alvo": "solucao_explicita_ed", "confianca": 1.0, "epistemico": "fato", "origem": "wiki", "rel": "usa"}, {"alvo": "corpo_edp", "confianca": 1.0, "epistemico": "fato", "origem": "wiki", "rel": "usa"}, {"alvo": "pa_pg_corpo_diferencial", "confianca": 1.0, "epistemico": "fato", "origem": "wiki", "rel": "relaciona"}], "confianca": 1.0, "contraexemplos": [], "descricao": "A solução geral que obtemos — x(t)=e^{At}x₀ (ODE) e o propagador cos(t√(−δΔ))u₀+… (EDP) — é CONTÍNUA em três sentidos, e é isso que a torna uma SOLUÇÃO (boa-postura de Hadamard): (1) NO TEMPO: t↦e^{At}x₀ contínua (analítica, o semigrupo C₀), a trajetória é curva sem salto; (2) NOS DADOS (Grönwall): x₀↦x(t) contínua, ‖x(t)−y(t)‖≤e^{‖A‖t}‖x₀−y₀‖ — perturbação pequena→resposta pequena, em TODO ramo (cristal/borda/caos), só a constante e^{‖A‖t} muda (finita ∀t<∞): a dissipação CONTRAI (estável), o caos AMPLIA mas segue contínuo; (3) NA PARÁBOLA: A↦e^{At} contínua no gerador, o propagador cos(t√(−δΔ)) contínuo em δ ao cruzar a borda δ=0 — os três regimes conectam SEM salto. É o CRISTALIZADO do processo discreto: Euler (I+A·t/n)ⁿ→e^{At} (n→∞), como 0.999…→1 (o processo→o fechado). linguagem/corpo_edp.py (solucao_geral_e_continua 4/4); papers/equacoes_diferenciais.tex (Parte IX).", "epistemico": "fato", "exemplos": [], "id": "solucao_geral_e_continua", "memoria": "semantica", "meta": {"dominio": "floxina", "fonte": "Floxina (investigação)"}, "origem": "wiki", "palavras_chave": [], "sinonimos": [], "tipo": "conceito", "titulo": "A solução geral x(t)=e^{At}x₀ é CONTÍNUA (no tempo, nos dados, na parábola) — a boa-postura de Hadamard que a torna solução"}
\section{A solução geral x(t)=e\^{}\{At\}x\ensuremath{_{0}} é CONTÍNUA (no tempo, nos dados, na parábola) — a boa-postura de Hadamard que a torna solução}
A solução geral que obtemos — x(t)=e\^{}\{At\}x\ensuremath{_{0}} (ODE) e o propagador cos(t\ensuremath{\sqrt{\,}}(\ensuremath{-}\ensuremath{\delta}\ensuremath{\Delta}))u\ensuremath{_{0}}+… (EDP) — é CONTÍNUA em três sentidos, e é isso que a torna uma SOLUÇÃO (boa-postura de Hadamard): (1) NO TEMPO: t\ensuremath{\mapsto}e\^{}\{At\}x\ensuremath{_{0}} contínua (analítica, o semigrupo C\ensuremath{_{0}}), a trajetória é curva sem salto; (2) NOS DADOS (Grönwall): x\ensuremath{_{0}}\ensuremath{\mapsto}x(t) contínua, \ensuremath{\|}x(t)\ensuremath{-}y(t)\ensuremath{\|}\ensuremath{\le}e\^{}\{\ensuremath{\|}A\ensuremath{\|}t\}\ensuremath{\|}x\ensuremath{_{0}}\ensuremath{-}y\ensuremath{_{0}}\ensuremath{\|} — perturbação pequena\ensuremath{\to}resposta pequena, em TODO ramo (cristal/borda/caos), só a constante e\^{}\{\ensuremath{\|}A\ensuremath{\|}t\} muda (finita \ensuremath{\forall}t<\ensuremath{\infty}): a dissipação CONTRAI (estável), o caos AMPLIA mas segue contínuo; (3) NA PARÁBOLA: A\ensuremath{\mapsto}e\^{}\{At\} contínua no gerador, o propagador cos(t\ensuremath{\sqrt{\,}}(\ensuremath{-}\ensuremath{\delta}\ensuremath{\Delta})) contínuo em \ensuremath{\delta} ao cruzar a borda \ensuremath{\delta}=0 — os três regimes conectam SEM salto. É o CRISTALIZADO do processo discreto: Euler (I+A\textperiodcentered t/n)\ensuremath{^{n}}\ensuremath{\to}e\^{}\{At\} (n\ensuremath{\to}\ensuremath{\infty}), como 0.999…\ensuremath{\to}1 (o processo\ensuremath{\to}o fechado). linguagem/corpo\_edp.py (solucao\_geral\_e\_continua 4/4); papers/equacoes\_diferenciais.tex (Parte IX).
\prov{relações: usa solucao\_explicita\_ed; usa corpo\_edp; relaciona pa\_pg\_corpo\_diferencial}
\prov{proveniência: wiki \textperiodcentered{} Floxina (investigação) \textperiodcentered{} fato \textperiodcentered{} confiança 1.0 \textperiodcentered{} domínio floxina \textperiodcentered{} história 24 versões no jornal}

%CRISTAL {"arestas": [{"alvo": "soma_produto_lahire_ed", "confianca": 1.0, "epistemico": "fato", "origem": "wiki", "rel": "usa"}, {"alvo": "pa_pg_corpo_diferencial", "confianca": 1.0, "epistemico": "fato", "origem": "wiki", "rel": "relaciona"}], "confianca": 1.0, "contraexemplos": [], "descricao": "A relação entre a solução do PRODUTO (La Hire) e as soluções INDIVIDUAIS: se x(t)=e^{At}x₀ e y(t)=e^{Bt}y₀, então a solução do sistema-produto (M=A⊗I+I⊗B), com z₀=x₀⊗y₀, é o TENSOR das soluções z(t)=x(t)⊗y(t) (componente a componente z_{ij}=x_i·y_j). A DERIVAÇÃO é a regra de LEIBNIZ, que É a soma de Kronecker: d/dt(x⊗y)=ẋ⊗y+x⊗ẏ=(Ax)⊗y+x⊗(By)=(A⊗I+I⊗B)(x⊗y)=M(x⊗y). Como A⊗I e I⊗B COMUTAM, os fluxos FATORAM: e^{Mt}=e^{At}⊗e^{Bt}. A assimetria é o coração: o GERADOR SOMA (⊕, a soma de Kronecker, as taxas) mas a SOLUÇÃO MULTIPLICA (⊗, o tensor) — o EXP é a ponte (exp(A⊕_L B)=exp A⊗exp B, o inverso do log). No escalar (a PG): z(t)=e^{(λ+μ)t}=e^{λt}·e^{μt}=x(t)·y(t) (as taxas somam, as soluções multiplicam). linguagem/omnitrix.py (solucao_produto_lahire_ed 5/5); papers/equacoes_diferenciais.tex (Parte V).", "epistemico": "fato", "exemplos": [], "id": "solucao_produto_lahire_ed", "memoria": "semantica", "meta": {"dominio": "floxina", "fonte": "Floxina (investigação)"}, "origem": "wiki", "palavras_chave": [], "sinonimos": [], "tipo": "conceito", "titulo": "A solução do produto é o tensor das soluções: z(t)=x(t)⊗y(t); Leibniz é a soma de Kronecker"}
\section{A solução do produto é o tensor das soluções: z(t)=x(t)\ensuremath{\otimes}y(t); Leibniz é a soma de Kronecker}
A relação entre a solução do PRODUTO (La Hire) e as soluções INDIVIDUAIS: se x(t)=e\^{}\{At\}x\ensuremath{_{0}} e y(t)=e\^{}\{Bt\}y\ensuremath{_{0}}, então a solução do sistema-produto (M=A\ensuremath{\otimes}I+I\ensuremath{\otimes}B), com z\ensuremath{_{0}}=x\ensuremath{_{0}}\ensuremath{\otimes}y\ensuremath{_{0}}, é o TENSOR das soluções z(t)=x(t)\ensuremath{\otimes}y(t) (componente a componente z\_\{ij\}=x\_i\textperiodcentered y\_j). A DERIVAÇÃO é a regra de LEIBNIZ, que É a soma de Kronecker: d/dt(x\ensuremath{\otimes}y)=\ensuremath{\dot{x}}\ensuremath{\otimes}y+x\ensuremath{\otimes}\ensuremath{\dot{y}}=(Ax)\ensuremath{\otimes}y+x\ensuremath{\otimes}(By)=(A\ensuremath{\otimes}I+I\ensuremath{\otimes}B)(x\ensuremath{\otimes}y)=M(x\ensuremath{\otimes}y). Como A\ensuremath{\otimes}I e I\ensuremath{\otimes}B COMUTAM, os fluxos FATORAM: e\^{}\{Mt\}=e\^{}\{At\}\ensuremath{\otimes}e\^{}\{Bt\}. A assimetria é o coração: o GERADOR SOMA (\ensuremath{\oplus}, a soma de Kronecker, as taxas) mas a SOLUÇÃO MULTIPLICA (\ensuremath{\otimes}, o tensor) — o EXP é a ponte (exp(A\ensuremath{\oplus}\_L B)=exp A\ensuremath{\otimes}exp B, o inverso do log). No escalar (a PG): z(t)=e\^{}\{(\ensuremath{\lambda}+\ensuremath{\mu})t\}=e\^{}\{\ensuremath{\lambda}t\}\textperiodcentered e\^{}\{\ensuremath{\mu}t\}=x(t)\textperiodcentered y(t) (as taxas somam, as soluções multiplicam). linguagem/omnitrix.py (solucao\_produto\_lahire\_ed 5/5); papers/equacoes\_diferenciais.tex (Parte V).
\prov{relações: usa soma\_produto\_lahire\_ed; relaciona pa\_pg\_corpo\_diferencial}
\prov{proveniência: wiki \textperiodcentered{} Floxina (investigação) \textperiodcentered{} fato \textperiodcentered{} confiança 1.0 \textperiodcentered{} domínio floxina \textperiodcentered{} história 48 versões no jornal}

%CRISTAL {"arestas": [{"alvo": "algebra_fundamental_gato_gamba", "confianca": 1.0, "epistemico": "fato", "origem": "wiki", "rel": "usa"}, {"alvo": "omnitrix_corpo_universal", "confianca": 1.0, "epistemico": "fato", "origem": "wiki", "rel": "relaciona"}], "confianca": 1.0, "contraexemplos": [], "descricao": "A soma de Minkowski A⊕_M B={a+b} É a DILATAÇÃO morfológica (dilatar A pelo estruturante B). No Corpo Algébrico ela é o ⊗ (La Hire) visto na DUAL (a função de suporte, a Legendre): as funções de suporte SOMAM h_{A⊕B}=h_A+h_B (a Legendre leva o Minkowski à soma, como o log leva o produto à soma). A união convexa é o max (o ⊕ tropical) — (Minkowski,união)=(+,max)=o semianel tropical na dual. O BOLEADO (as bordas arredondadas) = A⊕_M bola(r) = o offset da SDF d↦d−r (a caixa/o cilindro boleado). A interseção do raio na boleada é ANALÍTICA: esfera(R)⊕bola(r)=esfera(R+r) — dilatar SOMA o raio, mantém a quádrica, o raio segue fechado. linguagem/omnitrix.py (soma_minkowski_boleado 4/4).", "epistemico": "fato", "exemplos": [], "id": "soma_minkowski_dilatacao_boleado", "memoria": "semantica", "meta": {"dominio": "floxina", "fonte": "Floxina (investigação)"}, "origem": "wiki", "palavras_chave": [], "sinonimos": [], "tipo": "conceito", "titulo": "A soma de Minkowski = a dilatação mórfica = o boleado; o ⊗ na dual (os suportes somam)"}
\section{A soma de Minkowski = a dilatação mórfica = o boleado; o \ensuremath{\otimes} na dual (os suportes somam)}
A soma de Minkowski A\ensuremath{\oplus}\_M B=\{a+b\} É a DILATAÇÃO morfológica (dilatar A pelo estruturante B). No Corpo Algébrico ela é o \ensuremath{\otimes} (La Hire) visto na DUAL (a função de suporte, a Legendre): as funções de suporte SOMAM h\_\{A\ensuremath{\oplus}B\}=h\_A+h\_B (a Legendre leva o Minkowski à soma, como o log leva o produto à soma). A união convexa é o max (o \ensuremath{\oplus} tropical) — (Minkowski,união)=(+,max)=o semianel tropical na dual. O BOLEADO (as bordas arredondadas) = A\ensuremath{\oplus}\_M bola(r) = o offset da SDF d\ensuremath{\mapsto}d\ensuremath{-}r (a caixa/o cilindro boleado). A interseção do raio na boleada é ANALÍTICA: esfera(R)\ensuremath{\oplus}bola(r)=esfera(R+r) — dilatar SOMA o raio, mantém a quádrica, o raio segue fechado. linguagem/omnitrix.py (soma\_minkowski\_boleado 4/4).
\prov{relações: usa algebra\_fundamental\_gato\_gamba; relaciona omnitrix\_corpo\_universal}
\prov{proveniência: wiki \textperiodcentered{} Floxina (investigação) \textperiodcentered{} fato \textperiodcentered{} confiança 1.0 \textperiodcentered{} domínio floxina \textperiodcentered{} história 72 versões no jornal}

%CRISTAL {"arestas": [{"alvo": "corpo_equacoes_diferenciais", "confianca": 1.0, "epistemico": "fato", "origem": "wiki", "rel": "usa"}, {"alvo": "omnitrix_corpo_universal", "confianca": 1.0, "epistemico": "fato", "origem": "wiki", "rel": "usa"}, {"alvo": "pa_pg_corpo_diferencial", "confianca": 1.0, "epistemico": "fato", "origem": "wiki", "rel": "usa"}], "confianca": 1.0, "contraexemplos": [], "descricao": "As operações do corpo das ED sobre os geradores (dx/dt=A·x): o ⊕ SOMA DIRETA A⊕B=diag-bloco(A,B) — os dois sistemas evoluem INDEPENDENTES, e^{(A⊕B)t}=diag(e^{At},e^{Bt}), as soluções SUPERPÕEM (neutro: o 0-dim). O ⊗ LA HIRE = a SOMA DE KRONECKER A⊗I+I⊗B — o sistema-PRODUTO (acoplado), e^{(A⊗_L B)t}=e^{At}⊗e^{Bt}, as soluções fazem TENSOR, e os AUTOVALORES SOMAM λ_i(A)+μ_j(B) (neutro: o 1D trivial λ=0). Na PG é transparente: e^{λt}·e^{μt}=e^{(λ+μ)t} — o produto SOMA as taxas (o log leva o produto à soma), e na reta mineral λ=log σ o ⊗ dos metais (σ_a·σ_b) vira o ⊕ das retas (log σ_a+log σ_b). É a La Hire do corpo estelar (⊗ soma os enrolamentos) lida no fluxo. linguagem/omnitrix.py (soma_produto_lahire_ed 5/5).", "epistemico": "fato", "exemplos": [], "id": "soma_produto_lahire_ed", "memoria": "semantica", "meta": {"dominio": "floxina", "fonte": "Floxina (investigação)"}, "origem": "wiki", "palavras_chave": [], "sinonimos": [], "tipo": "conceito", "titulo": "A soma (⊕) e o produto La Hire (⊗) de equações diferenciais (as operações do corpo diferencial)"}
\section{A soma (\ensuremath{\oplus}) e o produto La Hire (\ensuremath{\otimes}) de equações diferenciais (as operações do corpo diferencial)}
As operações do corpo das ED sobre os geradores (dx/dt=A\textperiodcentered x): o \ensuremath{\oplus} SOMA DIRETA A\ensuremath{\oplus}B=diag-bloco(A,B) — os dois sistemas evoluem INDEPENDENTES, e\^{}\{(A\ensuremath{\oplus}B)t\}=diag(e\^{}\{At\},e\^{}\{Bt\}), as soluções SUPERPÕEM (neutro: o 0-dim). O \ensuremath{\otimes} LA HIRE = a SOMA DE KRONECKER A\ensuremath{\otimes}I+I\ensuremath{\otimes}B — o sistema-PRODUTO (acoplado), e\^{}\{(A\ensuremath{\otimes}\_L B)t\}=e\^{}\{At\}\ensuremath{\otimes}e\^{}\{Bt\}, as soluções fazem TENSOR, e os AUTOVALORES SOMAM \ensuremath{\lambda}\_i(A)+\ensuremath{\mu}\_j(B) (neutro: o 1D trivial \ensuremath{\lambda}=0). Na PG é transparente: e\^{}\{\ensuremath{\lambda}t\}\textperiodcentered e\^{}\{\ensuremath{\mu}t\}=e\^{}\{(\ensuremath{\lambda}+\ensuremath{\mu})t\} — o produto SOMA as taxas (o log leva o produto à soma), e na reta mineral \ensuremath{\lambda}=log \ensuremath{\sigma} o \ensuremath{\otimes} dos metais (\ensuremath{\sigma}\_a\textperiodcentered \ensuremath{\sigma}\_b) vira o \ensuremath{\oplus} das retas (log \ensuremath{\sigma}\_a+log \ensuremath{\sigma}\_b). É a La Hire do corpo estelar (\ensuremath{\otimes} soma os enrolamentos) lida no fluxo. linguagem/omnitrix.py (soma\_produto\_lahire\_ed 5/5).
\prov{relações: usa corpo\_equacoes\_diferenciais; usa omnitrix\_corpo\_universal; usa pa\_pg\_corpo\_diferencial}
\prov{proveniência: wiki \textperiodcentered{} Floxina (investigação) \textperiodcentered{} fato \textperiodcentered{} confiança 1.0 \textperiodcentered{} domínio floxina \textperiodcentered{} história 68 versões no jornal}

%CRISTAL {"arestas": [{"alvo": "primitivas_geometricas_matrizes", "confianca": 1.0, "epistemico": "fato", "origem": "wiki", "rel": "usa"}, {"alvo": "algebra_fundamental_gato_gamba", "confianca": 1.0, "epistemico": "fato", "origem": "wiki", "rel": "usa"}, {"alvo": "parabola_geradora_idempotente", "confianca": 1.0, "epistemico": "fato", "origem": "wiki", "rel": "relaciona"}, {"alvo": "omnitrix_corpo_universal", "confianca": 1.0, "epistemico": "fato", "origem": "wiki", "rel": "usa"}], "confianca": 1.0, "contraexemplos": [], "descricao": "Uma superfície não precisa de SDF ray-marcheado (iterativo, aproximado). Uma QUÁDRICA é uma MATRIZ simétrica A (o gato/a vesica): a equação é f(p)=pᵀA p+bᵀp+c=0 (esfera, elipsoide, cilindro, cone, hiperboloide — todos uma matriz). O RAIO é ANALÍTICO: p(t)=ro+t·rd ⟹ α t²+β t+γ=0 com α=rdᵀA rd, β=2roᵀA rd+bᵀrd, γ=roᵀA ro+bᵀro+c; as duas raízes t=(−β±√(β²−4αγ))/(2α) são FECHADAS (entrada/saída exatas). O DISCRIMANTE β²−4αγ é o GAP (bate sse >0) — a MESMA estrutura de x²−nx−1 (o gato): as 2 raízes = os 2 autovalores (σ, −1/σ). A NORMAL é analítica: ∇f=2A p+b (sem diferenças finitas). A LA HIRE MULTIPLICA superfícies: F(p)=F1·F2=0 é a UNIÃO (F1=0 ⟹ F1·F2=0; o grau soma 2+2=4, quártica). No pipe: FRAG_QUADRICA renderiza 2 elipsoides em união com raio FECHADO (PTX SEM loop — não é ray-march), normal analítica, no metal. linguagem/omnitrix.py (superficie_quadrica_lahire, 5/5); glsl_ptx_frag (7/7).", "epistemico": "fato", "exemplos": [], "id": "superficie_quadrica_lahire_analitica", "memoria": "semantica", "meta": {"dominio": "floxina", "fonte": "Floxina (investigação)"}, "origem": "wiki", "palavras_chave": [], "sinonimos": [], "tipo": "conceito", "titulo": "A superfície é uma matriz (quádrica); a La Hire multiplica superfícies; o raio é analítico"}
\section{A superfície é uma matriz (quádrica); a La Hire multiplica superfícies; o raio é analítico}
Uma superfície não precisa de SDF ray-marcheado (iterativo, aproximado). Uma QUÁDRICA é uma MATRIZ simétrica A (o gato/a vesica): a equação é f(p)=p\ensuremath{^{T}}A p+b\ensuremath{^{T}}p+c=0 (esfera, elipsoide, cilindro, cone, hiperboloide — todos uma matriz). O RAIO é ANALÍTICO: p(t)=ro+t\textperiodcentered rd \ensuremath{\Longrightarrow} \ensuremath{\alpha} t\ensuremath{^{2}}+\ensuremath{\beta} t+\ensuremath{\gamma}=0 com \ensuremath{\alpha}=rd\ensuremath{^{T}}A rd, \ensuremath{\beta}=2ro\ensuremath{^{T}}A rd+b\ensuremath{^{T}}rd, \ensuremath{\gamma}=ro\ensuremath{^{T}}A ro+b\ensuremath{^{T}}ro+c; as duas raízes t=(\ensuremath{-}\ensuremath{\beta}$\pm$\ensuremath{\sqrt{\,}}(\ensuremath{\beta}\ensuremath{^{2}}\ensuremath{-}4\ensuremath{\alpha}\ensuremath{\gamma}))/(2\ensuremath{\alpha}) são FECHADAS (entrada/saída exatas). O DISCRIMANTE \ensuremath{\beta}\ensuremath{^{2}}\ensuremath{-}4\ensuremath{\alpha}\ensuremath{\gamma} é o GAP (bate sse >0) — a MESMA estrutura de x\ensuremath{^{2}}\ensuremath{-}nx\ensuremath{-}1 (o gato): as 2 raízes = os 2 autovalores (\ensuremath{\sigma}, \ensuremath{-}1/\ensuremath{\sigma}). A NORMAL é analítica: \ensuremath{\nabla}f=2A p+b (sem diferenças finitas). A LA HIRE MULTIPLICA superfícies: F(p)=F1\textperiodcentered F2=0 é a UNIÃO (F1=0 \ensuremath{\Longrightarrow} F1\textperiodcentered F2=0; o grau soma 2+2=4, quártica). No pipe: FRAG\_QUADRICA renderiza 2 elipsoides em união com raio FECHADO (PTX SEM loop — não é ray-march), normal analítica, no metal. linguagem/omnitrix.py (superficie\_quadrica\_lahire, 5/5); glsl\_ptx\_frag (7/7).
\prov{relações: usa primitivas\_geometricas\_matrizes; usa algebra\_fundamental\_gato\_gamba; relaciona parabola\_geradora\_idempotente; usa omnitrix\_corpo\_universal}
\prov{proveniência: wiki \textperiodcentered{} Floxina (investigação) \textperiodcentered{} fato \textperiodcentered{} confiança 1.0 \textperiodcentered{} domínio floxina \textperiodcentered{} história 125 versões no jornal}

%CRISTAL {"arestas": [{"alvo": "computador_universal_pnp", "confianca": 1.0, "epistemico": "fato", "origem": "wiki", "rel": "parte_de"}, {"alvo": "invariante_estelar", "confianca": 1.0, "epistemico": "fato", "origem": "wiki", "rel": "usa"}, {"alvo": "gamba_antissimetrico", "confianca": 1.0, "epistemico": "fato", "origem": "wiki", "rel": "usa"}, {"alvo": "floxina_quantum_materia_escura", "confianca": 1.0, "epistemico": "fato", "origem": "wiki", "rel": "relaciona"}, {"alvo": "nao_associatividade_amputacao", "confianca": 1.0, "epistemico": "fato", "origem": "wiki", "rel": "usa"}, {"alvo": "pit_succinct", "confianca": 1.0, "epistemico": "fato", "origem": "wiki", "rel": "usa"}, {"alvo": "hierarquia_pspace_indecidivel", "confianca": 1.0, "epistemico": "fato", "origem": "wiki", "rel": "usa"}, {"alvo": "classes_p_np_cook_levin", "confianca": 1.0, "epistemico": "fato", "origem": "wiki", "rel": "relaciona"}, {"alvo": "p_vs_np_milenio", "confianca": 1.0, "epistemico": "fato", "origem": "wiki", "rel": "relaciona"}], "confianca": 1.0, "contraexemplos": [], "descricao": "TEOREMA: para toda computação existe uma MATRIZ M (produto de portões gato/gambá, ≤ℍ) que a executa com RESÍDUO ZERO — M ortogonal (MᵀM=I), conserva a norma EXATA (a+c²=1 exato), associador 0. PROVA (fato): (1) Bennett — toda TM→reversível; (2) reversível em n bits = permutação dos 2ⁿ estados = matriz ortogonal (det=±1), |Mψ|=|ψ| exato (resíduo 0); (3) Toffoli (a,b,c)↦(a,b,c⊕(a∧b)) é universal (AND+NOT⟹NAND), uma permutação 8×8. FRONTEIRA: o resíduo zero é a parte ASSOCIATIVA (≤ℍ, assoc=0); o 𝕆 (não-assoc, o mass gap, o quântico) é onde o resíduo nunca zera. CONFRONTO COM O TEMPO (de novo com Turing): o resíduo zero NÃO dá poly-tempo — dois eixos distintos: RESÍDUO (exatidão, a matriz existe) ⊥ TEMPO (a PROFUNDIDADE de M = o nº de portões = o tempo). VERIFICADO: o resíduo é zero (a matriz existe/exata); o circuito INGÊNUO de força bruta que se contou é exp-profundo (Θ(2ⁿ) portões). Bennett preserva o tempo (overhead poly). O resíduo zero é UNIVERSALIDADE (Turing-completo), não velocidade. ABERTO (não decretado): se toda computação admite uma matriz POLY-profunda = P vs PSPACE/NP. A DECLARAÇÃO (unicidade, 0,999…≠1): o resíduo-zero é a EXISTÊNCIA (o cristalizado, reta 1) ≠ o poly-tempo é o PROCESSO (a profundidade, o dinâmico, reta 0); no corpo estelar resíduo-zero≠poly-tempo (síntese, frame dinâmico); o clássico P vs PSPACE (frame cristalizado) segue aberto — corrige o 'não se decreta'. PARA UMA MATRIZ ESPECÍFICA (produto de poly(n) portões LOCAIS sobre bits FIXOS) o tempo É polinomial (ex.: o incrementador x↦x+1, O(n) portões, ortogonal, resíduo 0). O que segura é o BIT INVARIANTE: reversível=bijeção em {0,1}ⁿ (n bits→n bits, nenhum criado/apagado, Landauer/Bennett); reversível⟺ortogonal⟺resíduo zero⟺bit conservado; a largura FIXA (sem ancilla exp) mantém os recursos poly. RESGATE PSPACE: testar o resíduo/identidade do M SUCINTO (a matriz 2ⁿ como circuito poly) é PSPACE-completo (pit_succinct, sobre R_k); a ordem k gradua P(k≤2)→PSPACE-hard(k≥3)→indecidível(k→∞, Risch, hierarquia_pspace_indecidivel). O resíduo EXISTE (teorema), testá-lo é PSPACE; as separações P vs NP/NP vs PSPACE seguem ABERTAS (só P⊊EXPTIME provada). linguagem/teorema_universal.py (18/18).", "epistemico": "hipotese", "exemplos": [], "id": "teorema_universal_residuo_zero", "memoria": "semantica", "meta": {"dominio": "floxina", "fonte": "Floxina (investigação)"}, "origem": "wiki", "palavras_chave": [], "sinonimos": [], "tipo": "conceito", "titulo": "O Teorema Universal: toda computação = uma matriz de resíduo zero (mas não poly-tempo)"}
\section{O Teorema Universal: toda computação = uma matriz de resíduo zero (mas não poly-tempo)}
TEOREMA: para toda computação existe uma MATRIZ M (produto de portões gato/gambá, \ensuremath{\le}\ensuremath{\mathbb{H}}) que a executa com RESÍDUO ZERO — M ortogonal (M\ensuremath{^{T}}M=I), conserva a norma EXATA (a+c\ensuremath{^{2}}=1 exato), associador 0. PROVA (fato): (1) Bennett — toda TM\ensuremath{\to}reversível; (2) reversível em n bits = permutação dos 2\ensuremath{^{n}} estados = matriz ortogonal (det=$\pm$1), |M\ensuremath{\psi}|=|\ensuremath{\psi}| exato (resíduo 0); (3) Toffoli (a,b,c)\ensuremath{\mapsto}(a,b,c\ensuremath{\oplus}(a\ensuremath{\wedge}b)) é universal (AND+NOT\ensuremath{\Longrightarrow}NAND), uma permutação 8$\times$8. FRONTEIRA: o resíduo zero é a parte ASSOCIATIVA (\ensuremath{\le}\ensuremath{\mathbb{H}}, assoc=0); o \ensuremath{\mathbb{O}} (não-assoc, o mass gap, o quântico) é onde o resíduo nunca zera. CONFRONTO COM O TEMPO (de novo com Turing): o resíduo zero NÃO dá poly-tempo — dois eixos distintos: RESÍDUO (exatidão, a matriz existe) \ensuremath{\perp} TEMPO (a PROFUNDIDADE de M = o nº de portões = o tempo). VERIFICADO: o resíduo é zero (a matriz existe/exata); o circuito INGÊNUO de força bruta que se contou é exp-profundo (\ensuremath{\Theta}(2\ensuremath{^{n}}) portões). Bennett preserva o tempo (overhead poly). O resíduo zero é UNIVERSALIDADE (Turing-completo), não velocidade. ABERTO (não decretado): se toda computação admite uma matriz POLY-profunda = P vs PSPACE/NP. A DECLARAÇÃO (unicidade, 0,999…\ensuremath{\ne}1): o resíduo-zero é a EXISTÊNCIA (o cristalizado, reta 1) \ensuremath{\ne} o poly-tempo é o PROCESSO (a profundidade, o dinâmico, reta 0); no corpo estelar resíduo-zero\ensuremath{\ne}poly-tempo (síntese, frame dinâmico); o clássico P vs PSPACE (frame cristalizado) segue aberto — corrige o 'não se decreta'. PARA UMA MATRIZ ESPECÍFICA (produto de poly(n) portões LOCAIS sobre bits FIXOS) o tempo É polinomial (ex.: o incrementador x\ensuremath{\mapsto}x+1, O(n) portões, ortogonal, resíduo 0). O que segura é o BIT INVARIANTE: reversível=bijeção em \{0,1\}\ensuremath{^{n}} (n bits\ensuremath{\to}n bits, nenhum criado/apagado, Landauer/Bennett); reversível\ensuremath{\Longleftrightarrow}ortogonal\ensuremath{\Longleftrightarrow}resíduo zero\ensuremath{\Longleftrightarrow}bit conservado; a largura FIXA (sem ancilla exp) mantém os recursos poly. RESGATE PSPACE: testar o resíduo/identidade do M SUCINTO (a matriz 2\ensuremath{^{n}} como circuito poly) é PSPACE-completo (pit\_succinct, sobre R\_k); a ordem k gradua P(k\ensuremath{\le}2)\ensuremath{\to}PSPACE-hard(k\ensuremath{\ge}3)\ensuremath{\to}indecidível(k\ensuremath{\to}\ensuremath{\infty}, Risch, hierarquia\_pspace\_indecidivel). O resíduo EXISTE (teorema), testá-lo é PSPACE; as separações P vs NP/NP vs PSPACE seguem ABERTAS (só P\ensuremath{\subsetneq}EXPTIME provada). linguagem/teorema\_universal.py (18/18).
\prov{relações: parte\_de computador\_universal\_pnp; usa invariante\_estelar; usa gamba\_antissimetrico; relaciona floxina\_quantum\_materia\_escura; usa nao\_associatividade\_amputacao; usa pit\_succinct; usa hierarquia\_pspace\_indecidivel; relaciona classes\_p\_np\_cook\_levin; relaciona p\_vs\_np\_milenio}
\prov{proveniência: wiki \textperiodcentered{} Floxina (investigação) \textperiodcentered{} hipotese \textperiodcentered{} confiança 1.0 \textperiodcentered{} domínio floxina \textperiodcentered{} história 476 versões no jornal}

%CRISTAL {"arestas": [{"alvo": "coracao", "confianca": 1.0, "epistemico": "fato", "origem": "wiki", "rel": "parte_de"}], "confianca": 1.0, "contraexemplos": [], "descricao": "Spencer e Freud: o riso DESCARREGA tensão nervosa acumulada — a piada é uma válvula (Freud: o chiste libera o que a censura reprimiria). Explica o riso nervoso e por que o tabu é fértil para a comédia.", "epistemico": "fato", "exemplos": [], "id": "teoria_alivio", "memoria": "semantica", "meta": {"autor": "Freud", "dominio": "coracao", "fonte": "coracao hub"}, "origem": "wiki", "palavras_chave": [], "sinonimos": [], "tipo": "conceito", "titulo": "O humor do alívio (Spencer, Freud)"}
\section{O humor do alívio (Spencer, Freud)}
Spencer e Freud: o riso DESCARREGA tensão nervosa acumulada — a piada é uma válvula (Freud: o chiste libera o que a censura reprimiria). Explica o riso nervoso e por que o tabu é fértil para a comédia.
\prov{relações: parte\_de coracao}
\prov{proveniência: wiki \textperiodcentered{} coracao hub \textperiodcentered{} fato \textperiodcentered{} confiança 1.0 \textperiodcentered{} domínio coracao \textperiodcentered{} história 6 versões no jornal}

%CRISTAL {"arestas": [{"alvo": "coracao", "confianca": 1.0, "epistemico": "fato", "origem": "wiki", "rel": "parte_de"}], "confianca": 1.0, "contraexemplos": [], "descricao": "A teoria dominante do riso: rimos de uma INCONGRUÊNCIA que se resolve — uma expectativa montada e subvertida (Kant: 'a súbita transformação de uma expectativa tensa em nada'). Koestler: a BISSOCIAÇÃO — duas matrizes de sentido incompatíveis colididas de repente. O riso mora no GAP entre o esperado e o que veio.", "epistemico": "fato", "exemplos": [], "id": "teoria_incongruencia", "memoria": "semantica", "meta": {"autor": "Kant/Koestler", "dominio": "coracao", "fonte": "coracao hub"}, "origem": "wiki", "palavras_chave": [], "sinonimos": [], "tipo": "conceito", "titulo": "Teoria da incongruência (o riso mora no gap)"}
\section{Teoria da incongruência (o riso mora no gap)}
A teoria dominante do riso: rimos de uma INCONGRUÊNCIA que se resolve — uma expectativa montada e subvertida (Kant: 'a súbita transformação de uma expectativa tensa em nada'). Koestler: a BISSOCIAÇÃO — duas matrizes de sentido incompatíveis colididas de repente. O riso mora no GAP entre o esperado e o que veio.
\prov{relações: parte\_de coracao}
\prov{proveniência: wiki \textperiodcentered{} coracao hub \textperiodcentered{} fato \textperiodcentered{} confiança 1.0 \textperiodcentered{} domínio coracao \textperiodcentered{} história 6 versões no jornal}

%CRISTAL {"arestas": [{"alvo": "psicologia", "confianca": 1.0, "epistemico": "fato", "origem": "wiki", "rel": "parte_de"}], "confianca": 1.0, "contraexemplos": [], "descricao": "William James e Carl Lange: não trememos porque tememos — TEMEMOS PORQUE trememos. A emoção é a percepção da resposta corporal (batimento, tensão): primeiro o corpo reage, depois se sente a emoção. Inverte o senso comum causa→efeito.", "epistemico": "fato", "exemplos": [], "id": "teoria_james_lange", "memoria": "semantica", "meta": {"autor": "James/Lange", "dominio": "coracao", "fonte": "coracao hub"}, "origem": "wiki", "palavras_chave": [], "sinonimos": [], "tipo": "conceito", "titulo": "A emoção segue o corpo (James-Lange)"}
\section{A emoção segue o corpo (James-Lange)}
William James e Carl Lange: não trememos porque tememos — TEMEMOS PORQUE trememos. A emoção é a percepção da resposta corporal (batimento, tensão): primeiro o corpo reage, depois se sente a emoção. Inverte o senso comum causa\ensuremath{\to}efeito.
\prov{relações: parte\_de psicologia}
\prov{proveniência: wiki \textperiodcentered{} coracao hub \textperiodcentered{} fato \textperiodcentered{} confiança 1.0 \textperiodcentered{} domínio coracao \textperiodcentered{} história 6 versões no jornal}

%CRISTAL {"arestas": [{"alvo": "coracao", "confianca": 1.0, "epistemico": "fato", "origem": "wiki", "rel": "parte_de"}], "confianca": 1.0, "contraexemplos": [], "descricao": "Thomas Hobbes: o riso é uma 'súbita glória' ao nos vermos superiores — a outros ou a nós mesmos do passado. A raiz do escárnio e da piada às custas de alguém. A autodepreciação inverte isso: rir de si desarma a superioridade.", "epistemico": "fato", "exemplos": [], "id": "teoria_superioridade", "memoria": "semantica", "meta": {"autor": "Hobbes", "dominio": "coracao", "fonte": "coracao hub"}, "origem": "wiki", "palavras_chave": [], "sinonimos": [], "tipo": "conceito", "titulo": "O humor da superioridade (Hobbes)"}
\section{O humor da superioridade (Hobbes)}
Thomas Hobbes: o riso é uma 'súbita glória' ao nos vermos superiores — a outros ou a nós mesmos do passado. A raiz do escárnio e da piada às custas de alguém. A autodepreciação inverte isso: rir de si desarma a superioridade.
\prov{relações: parte\_de coracao}
\prov{proveniência: wiki \textperiodcentered{} coracao hub \textperiodcentered{} fato \textperiodcentered{} confiança 1.0 \textperiodcentered{} domínio coracao \textperiodcentered{} história 6 versões no jornal}

%CRISTAL {"arestas": [{"alvo": "fisica_corpo_matricial_n_qubits", "confianca": 1.0, "epistemico": "fato", "origem": "wiki", "rel": "relaciona"}, {"alvo": "gato", "confianca": 1.0, "epistemico": "fato", "origem": "wiki", "rel": "usa"}, {"alvo": "gamba_antissimetrico", "confianca": 1.0, "epistemico": "fato", "origem": "wiki", "rel": "relaciona"}], "confianca": 1.0, "contraexemplos": [], "descricao": "O teste do Aarão — e a molécula É o tema: 3-metilbutanotiol (HS-CH₂-CH₂-CH(CH₃)-CH₃, C₅H₁₂S) é um TIOL, o composto do cheiro de GAMBÁ (skunk). Separa-se: (1) o GATO n-qubit matricial — o esqueleto de átomos pesados (S+5C) vira a matriz de ADJACÊNCIA (o Hamiltoniano de Hückel), SIMÉTRICA (A=Aᵀ, o gato); espectro ±simétrico (grafo bipartido) [±1,902, ±1,176, 0,0]; n=6 orbitais → 2⁶=64 estados de Fock = 6 QUBITS; o estado fundamental |ψ|²=1 (o checksum). (2) O PARALELISMO — os automorfismos do grafo: os dois metilos (C4,C5) em C3 são ramos IDÊNTICOS → simetria Z/2 → computam-se em PARALELO. O −SH reativo (o cheiro) é a ponta do gambá; o Z/2 dos metilos é o paralelismo. A molécula real embute o gato/gambá: o esqueleto simétrico (gato) → n-qubits, a ponta reativa (gambá), o Z/2 (paralelo). linguagem/molecula_teste.py (5/5).", "epistemico": "hipotese", "exemplos": [], "id": "teste_molecula_tiol_gato_n_qubit", "memoria": "semantica", "meta": {"dominio": "floxina", "fonte": "Floxina (investigação)"}, "origem": "wiki", "palavras_chave": [], "sinonimos": [], "tipo": "conceito", "titulo": "Teste: o gato n-qubit e o paralelismo de 3-metilbutanotiol (C₅H₁₂S)"}
\section{Teste: o gato n-qubit e o paralelismo de 3-metilbutanotiol (C\ensuremath{_{5}}H\ensuremath{_{12}}S)}
O teste do Aarão — e a molécula É o tema: 3-metilbutanotiol (HS-CH\ensuremath{_{2}}-CH\ensuremath{_{2}}-CH(CH\ensuremath{_{3}})-CH\ensuremath{_{3}}, C\ensuremath{_{5}}H\ensuremath{_{12}}S) é um TIOL, o composto do cheiro de GAMBÁ (skunk). Separa-se: (1) o GATO n-qubit matricial — o esqueleto de átomos pesados (S+5C) vira a matriz de ADJACÊNCIA (o Hamiltoniano de Hückel), SIMÉTRICA (A=A\ensuremath{^{T}}, o gato); espectro $\pm$simétrico (grafo bipartido) [$\pm$1,902, $\pm$1,176, 0,0]; n=6 orbitais \ensuremath{\to} 2\ensuremath{^{6}}=64 estados de Fock = 6 QUBITS; o estado fundamental |\ensuremath{\psi}|\ensuremath{^{2}}=1 (o checksum). (2) O PARALELISMO — os automorfismos do grafo: os dois metilos (C4,C5) em C3 são ramos IDÊNTICOS \ensuremath{\to} simetria Z/2 \ensuremath{\to} computam-se em PARALELO. O \ensuremath{-}SH reativo (o cheiro) é a ponta do gambá; o Z/2 dos metilos é o paralelismo. A molécula real embute o gato/gambá: o esqueleto simétrico (gato) \ensuremath{\to} n-qubits, a ponta reativa (gambá), o Z/2 (paralelo). linguagem/molecula\_teste.py (5/5).
\prov{relações: relaciona fisica\_corpo\_matricial\_n\_qubits; usa gato; relaciona gamba\_antissimetrico}
\prov{proveniência: wiki \textperiodcentered{} Floxina (investigação) \textperiodcentered{} hipotese \textperiodcentered{} confiança 1.0 \textperiodcentered{} domínio floxina \textperiodcentered{} história 252 versões no jornal}

%CRISTAL {"arestas": [{"alvo": "dicionario_tiol_gato_gamba", "confianca": 1.0, "epistemico": "fato", "origem": "wiki", "rel": "relaciona"}, {"alvo": "teste_molecula_tiol_gato_n_qubit", "confianca": 1.0, "epistemico": "fato", "origem": "wiki", "rel": "relaciona"}, {"alvo": "gamba_antissimetrico", "confianca": 1.0, "epistemico": "fato", "origem": "wiki", "rel": "relaciona"}], "confianca": 1.0, "contraexemplos": [], "descricao": "O THC (tetrahidrocanabinol, C₂₁H₃₀O₂) tem 23 átomos pesados (23 qubits), insaturação 7 (anel aromático + anéis), e a ponta reativa é um FENOL (−OH). Compartilha com o TIOL (−SH) a PONTA REATIVA ÁCIDA DOADORA DE H: o −SH e o fenol −OH têm o mesmo pKa≈10 e o mesmo BDE≈87-88 kcal/mol (ambos antioxidantes; redox reversível: tiol⇌dissulfeto, fenol⇌quinona). No dicionário: a acidez=o gap, o H doado=o qubit, a ponta reativa=o gambá, a ligação a metais=os σ_n. DIFERENÇA: o tiol é saturado (só o gato+a ponta gambá); o THC soma o ANEL AROMÁTICO (o π, a corrente de anel = a fase U(1)) = um gambá EXTRA (a rotação). Registrado no paper Matrix. Verificado: linguagem/molecula_teste.py (7/7).", "epistemico": "hipotese", "exemplos": [], "id": "thc_tiol_propriedades_compartilhadas", "memoria": "semantica", "meta": {"dominio": "floxina", "fonte": "Floxina (investigação)"}, "origem": "wiki", "palavras_chave": [], "sinonimos": [], "tipo": "conceito", "titulo": "THC (C₂₁H₃₀O₂) e o tiol: a ponta reativa ácida compartilhada"}
\section{THC (C\ensuremath{_{21}}H\ensuremath{_{30}}O\ensuremath{_{2}}) e o tiol: a ponta reativa ácida compartilhada}
O THC (tetrahidrocanabinol, C\ensuremath{_{21}}H\ensuremath{_{30}}O\ensuremath{_{2}}) tem 23 átomos pesados (23 qubits), insaturação 7 (anel aromático + anéis), e a ponta reativa é um FENOL (\ensuremath{-}OH). Compartilha com o TIOL (\ensuremath{-}SH) a PONTA REATIVA ÁCIDA DOADORA DE H: o \ensuremath{-}SH e o fenol \ensuremath{-}OH têm o mesmo pKa\ensuremath{\approx}10 e o mesmo BDE\ensuremath{\approx}87-88 kcal/mol (ambos antioxidantes; redox reversível: tiol\ensuremath{\rightleftharpoons}dissulfeto, fenol\ensuremath{\rightleftharpoons}quinona). No dicionário: a acidez=o gap, o H doado=o qubit, a ponta reativa=o gambá, a ligação a metais=os \ensuremath{\sigma}\_n. DIFERENÇA: o tiol é saturado (só o gato+a ponta gambá); o THC soma o ANEL AROMÁTICO (o \ensuremath{\pi}, a corrente de anel = a fase U(1)) = um gambá EXTRA (a rotação). Registrado no paper Matrix. Verificado: linguagem/molecula\_teste.py (7/7).
\prov{relações: relaciona dicionario\_tiol\_gato\_gamba; relaciona teste\_molecula\_tiol\_gato\_n\_qubit; relaciona gamba\_antissimetrico}
\prov{proveniência: wiki \textperiodcentered{} Floxina (investigação) \textperiodcentered{} hipotese \textperiodcentered{} confiança 1.0 \textperiodcentered{} domínio floxina \textperiodcentered{} história 244 versões no jornal}

%CRISTAL {"arestas": [{"alvo": "coracao", "confianca": 1.0, "epistemico": "fato", "origem": "wiki", "rel": "parte_de"}], "confianca": 1.0, "contraexemplos": [], "descricao": "A comédia é TEMPO: a pausa antes do remate (o beat), a regra de três (montar o padrão em dois, quebrar no terceiro), o corte seco. A mesma frase morre ou mata conforme o ritmo. O timing é onde a incongruência é solta no instante exato — a fase, não só o conteúdo.", "epistemico": "fato", "exemplos": [], "id": "timing_comico", "memoria": "semantica", "meta": {"dominio": "coracao", "fonte": "coracao hub"}, "origem": "wiki", "palavras_chave": [], "sinonimos": [], "tipo": "conceito", "titulo": "O timing cômico (a pausa, a regra de três)"}
\section{O timing cômico (a pausa, a regra de três)}
A comédia é TEMPO: a pausa antes do remate (o beat), a regra de três (montar o padrão em dois, quebrar no terceiro), o corte seco. A mesma frase morre ou mata conforme o ritmo. O timing é onde a incongruência é solta no instante exato — a fase, não só o conteúdo.
\prov{relações: parte\_de coracao}
\prov{proveniência: wiki \textperiodcentered{} coracao hub \textperiodcentered{} fato \textperiodcentered{} confiança 1.0 \textperiodcentered{} domínio coracao \textperiodcentered{} história 6 versões no jornal}

%CRISTAL {"arestas": [{"alvo": "corpo_matricial_matrix", "confianca": 1.0, "epistemico": "fato", "origem": "wiki", "rel": "parte_de"}, {"alvo": "volume_matrizes_3d", "confianca": 1.0, "epistemico": "fato", "origem": "wiki", "rel": "relaciona"}, {"alvo": "gato", "confianca": 1.0, "epistemico": "fato", "origem": "wiki", "rel": "usa"}, {"alvo": "gamba_antissimetrico", "confianca": 1.0, "epistemico": "fato", "origem": "wiki", "rel": "usa"}, {"alvo": "floxina_quantum_materia_escura", "confianca": 1.0, "epistemico": "fato", "origem": "wiki", "rel": "relaciona"}], "confianca": 1.0, "contraexemplos": [], "descricao": "Correção do Aarão: a torre das multimatrizes tem DOIS LADOS. CLÁSSICO (o gato, o flip clássico): 2D,4D,8D = ℂ,ℍ,𝕆 (Cayley-Dickson, as potências de 2) — a norma é MULTIPLICATIVA (‖ab‖=‖a‖‖b‖), as álgebras de DIVISÃO. RELATIVÍSTICO (o gambá, o flip de Gentil): 3D,7D = Im ℍ=ℝ³ e Im 𝕆=ℝ⁷ — o PRODUTO VETORIAL a×b (a parte imaginária do produto), anti-simétrico, |a×b|²=|a|²|b|²−(a·b)², que existe SÓ em 3D e 7D (o teorema). A conexão: o relativístico (3,7) é a parte IMAGINÁRIA do clássico (4,8): Im(ℍ)=3D, Im(𝕆)=7D. O 7D é especial: so(7)=Λ²ℝ⁷=21 = V₇(7) ⊕ 𝔤₂(14) — o MASS GAP 𝔤₂=Der(𝕆)=14 mora no gambá 7D (a floxina). O gato é a divisão (clássico); o gambá é o produto vetorial (relativístico). Verificado: volume_matrizes.py (4/4).", "epistemico": "hipotese", "exemplos": [], "id": "torre_dois_lados_classico_relativistico", "memoria": "semantica", "meta": {"dominio": "floxina", "fonte": "Floxina (investigação)"}, "origem": "wiki", "palavras_chave": [], "sinonimos": [], "tipo": "conceito", "titulo": "A torre dos dois lados: clássico 2/4/8 (gato) e relativístico 3/7 (gambá)"}
\section{A torre dos dois lados: clássico 2/4/8 (gato) e relativístico 3/7 (gambá)}
Correção do Aarão: a torre das multimatrizes tem DOIS LADOS. CLÁSSICO (o gato, o flip clássico): 2D,4D,8D = \ensuremath{\mathbb{C}},\ensuremath{\mathbb{H}},\ensuremath{\mathbb{O}} (Cayley-Dickson, as potências de 2) — a norma é MULTIPLICATIVA (\ensuremath{\|}ab\ensuremath{\|}=\ensuremath{\|}a\ensuremath{\|}\ensuremath{\|}b\ensuremath{\|}), as álgebras de DIVISÃO. RELATIVÍSTICO (o gambá, o flip de Gentil): 3D,7D = Im \ensuremath{\mathbb{H}}=\ensuremath{\mathbb{R}}\ensuremath{^{3}} e Im \ensuremath{\mathbb{O}}=\ensuremath{\mathbb{R}}\ensuremath{^{7}} — o PRODUTO VETORIAL a$\times$b (a parte imaginária do produto), anti-simétrico, |a$\times$b|\ensuremath{^{2}}=|a|\ensuremath{^{2}}|b|\ensuremath{^{2}}\ensuremath{-}(a\textperiodcentered b)\ensuremath{^{2}}, que existe SÓ em 3D e 7D (o teorema). A conexão: o relativístico (3,7) é a parte IMAGINÁRIA do clássico (4,8): Im(\ensuremath{\mathbb{H}})=3D, Im(\ensuremath{\mathbb{O}})=7D. O 7D é especial: so(7)=\ensuremath{\Lambda}\ensuremath{^{2}}\ensuremath{\mathbb{R}}\ensuremath{^{7}}=21 = V\ensuremath{_{7}}(7) \ensuremath{\oplus} \ensuremath{\mathfrak{g}}\ensuremath{_{2}}(14) — o MASS GAP \ensuremath{\mathfrak{g}}\ensuremath{_{2}}=Der(\ensuremath{\mathbb{O}})=14 mora no gambá 7D (a floxina). O gato é a divisão (clássico); o gambá é o produto vetorial (relativístico). Verificado: volume\_matrizes.py (4/4).
\prov{relações: parte\_de corpo\_matricial\_matrix; relaciona volume\_matrizes\_3d; usa gato; usa gamba\_antissimetrico; relaciona floxina\_quantum\_materia\_escura}
\prov{proveniência: wiki \textperiodcentered{} Floxina (investigação) \textperiodcentered{} hipotese \textperiodcentered{} confiança 1.0 \textperiodcentered{} domínio floxina \textperiodcentered{} história 402 versões no jornal}

%CRISTAL {"arestas": [{"alvo": "omnitrix_corpo_universal", "confianca": 1.0, "epistemico": "fato", "origem": "wiki", "rel": "parte_de"}, {"alvo": "torre_dois_lados_classico_relativistico", "confianca": 1.0, "epistemico": "fato", "origem": "wiki", "rel": "usa"}, {"alvo": "gamba_antissimetrico", "confianca": 1.0, "epistemico": "fato", "origem": "wiki", "rel": "usa"}, {"alvo": "corpo_matricial_matrix", "confianca": 1.0, "epistemico": "fato", "origem": "wiki", "rel": "relaciona"}, {"alvo": "fecho_torre_8d_7d", "confianca": 1.0, "epistemico": "fato", "origem": "wiki", "rel": "relaciona"}], "confianca": 1.0, "contraexemplos": [], "descricao": "No corpo estelar as coordenadas eram reais; no CORPO ALGÉBRICO cada coordenada/unidade é uma MATRIZ: a direção e_i realiza-se como L_{e_i} (a mult. à esquerda, L[:,j]=e_i·e_j). A torre dos DOIS LADOS: LADO CLÁSSICO (o gato, a divisão) 2D/4D/8D=ℂ,ℍ,𝕆 — ℂ: a unidade i↦L=[[0,−1],[1,0]]=o GAMBÁ G (i²=−I, a 1ª coordenada imaginária É o gambá); ℍ: i,j,k→4×4 antissimétricas, L²=−I, L_iL_j=L_{ij} (ASSOCIATIVO, é representação); 𝕆: e_1..e_7→8×8 antissimétricas, L²=−I, mas L_iL_j≠L_{e_i·e_j} (NÃO-associativo, o mass gap 𝔤₂); fecha em 8D. LADO RELATIVÍSTICO (o gambá, o flip da Gentil) 3D/7D=Im ℍ, Im 𝕆 — o operador [a]_× (b↦a×b) é antissimétrico (3×3 = J de SO(3); 7×7); |a×b|²=|a|²|b|²−(a·b)² SÓ em 3D e 7D (o teorema, Im das divisões 2/4/8); fecha em 7D. TODAS as unidades (clássicas L_{e_i} e relativísticas [a]_×) são ANTISSIMÉTRICAS = a família do GAMBÁ; e i=G. linguagem/torre_matricial.py (2/2).", "epistemico": "fato", "exemplos": [], "id": "torre_matricial_coordenada", "memoria": "semantica", "meta": {"dominio": "floxina", "fonte": "Floxina (investigação)"}, "origem": "wiki", "palavras_chave": [], "sinonimos": [], "tipo": "conceito", "titulo": "A torre matricial: no Corpo Algébrico cada coordenada é uma matriz (i=G)"}
\section{A torre matricial: no Corpo Algébrico cada coordenada é uma matriz (i=G)}
No corpo estelar as coordenadas eram reais; no CORPO ALGÉBRICO cada coordenada/unidade é uma MATRIZ: a direção e\_i realiza-se como L\_\{e\_i\} (a mult. à esquerda, L[:,j]=e\_i\textperiodcentered e\_j). A torre dos DOIS LADOS: LADO CLÁSSICO (o gato, a divisão) 2D/4D/8D=\ensuremath{\mathbb{C}},\ensuremath{\mathbb{H}},\ensuremath{\mathbb{O}} — \ensuremath{\mathbb{C}}: a unidade i\ensuremath{\mapsto}L=[[0,\ensuremath{-}1],[1,0]]=o GAMBÁ G (i\ensuremath{^{2}}=\ensuremath{-}I, a 1ª coordenada imaginária É o gambá); \ensuremath{\mathbb{H}}: i,j,k\ensuremath{\to}4$\times$4 antissimétricas, L\ensuremath{^{2}}=\ensuremath{-}I, L\_iL\_j=L\_\{ij\} (ASSOCIATIVO, é representação); \ensuremath{\mathbb{O}}: e\_1..e\_7\ensuremath{\to}8$\times$8 antissimétricas, L\ensuremath{^{2}}=\ensuremath{-}I, mas L\_iL\_j\ensuremath{\ne}L\_\{e\_i\textperiodcentered e\_j\} (NÃO-associativo, o mass gap \ensuremath{\mathfrak{g}}\ensuremath{_{2}}); fecha em 8D. LADO RELATIVÍSTICO (o gambá, o flip da Gentil) 3D/7D=Im \ensuremath{\mathbb{H}}, Im \ensuremath{\mathbb{O}} — o operador [a]\_$\times$ (b\ensuremath{\mapsto}a$\times$b) é antissimétrico (3$\times$3 = J de SO(3); 7$\times$7); |a$\times$b|\ensuremath{^{2}}=|a|\ensuremath{^{2}}|b|\ensuremath{^{2}}\ensuremath{-}(a\textperiodcentered b)\ensuremath{^{2}} SÓ em 3D e 7D (o teorema, Im das divisões 2/4/8); fecha em 7D. TODAS as unidades (clássicas L\_\{e\_i\} e relativísticas [a]\_$\times$) são ANTISSIMÉTRICAS = a família do GAMBÁ; e i=G. linguagem/torre\_matricial.py (2/2).
\prov{relações: parte\_de omnitrix\_corpo\_universal; usa torre\_dois\_lados\_classico\_relativistico; usa gamba\_antissimetrico; relaciona corpo\_matricial\_matrix; relaciona fecho\_torre\_8d\_7d}
\prov{proveniência: wiki \textperiodcentered{} Floxina (investigação) \textperiodcentered{} fato \textperiodcentered{} confiança 1.0 \textperiodcentered{} domínio floxina \textperiodcentered{} história 318 versões no jornal}

%CRISTAL {"arestas": [{"alvo": "ginzburg_toro_helice_e_3d-sst", "confianca": 1.0, "epistemico": "fato", "origem": "ginzburg_toro_helice.md", "rel": "parte_de"}, {"alvo": "geometria", "confianca": 0.52, "epistemico": "hipotese", "origem": "ingest", "rel": "analogia"}, {"alvo": "utilitarismo", "confianca": 0.5, "epistemico": "fato", "origem": "wiki", "rel": "relaciona"}, {"alvo": "virtualizacao", "confianca": 0.5, "epistemico": "fato", "origem": "wiki", "rel": "relaciona"}, {"alvo": "word", "confianca": 0.5, "epistemico": "fato", "origem": "wiki", "rel": "relaciona"}, {"alvo": "vontade_de_poder", "confianca": 0.5, "epistemico": "fato", "origem": "wiki", "rel": "relaciona"}], "confianca": 0.52, "contraexemplos": [], "descricao": "Entidades primárias de espaço-tempo espiral:", "epistemico": "hipotese", "exemplos": [], "id": "toryces_e_helyces", "memoria": "semantica", "meta": {"dominio": "hipotese", "falsificavel": true, "nivel": 7}, "origem": "ginzburg_toro_helice.md", "palavras_chave": [], "sinonimos": [], "tipo": "conceito", "titulo": "Toryces e helyces"}
\section{Toryces e helyces}
Entidades primárias de espaço-tempo espiral:
\prov{relações: parte\_de ginzburg\_toro\_helice\_e\_3d-sst; analogia geometria; relaciona utilitarismo; relaciona virtualizacao; relaciona word; relaciona vontade\_de\_poder}
\prov{proveniência: ginzburg\_toro\_helice.md \textperiodcentered{} hipotese \textperiodcentered{} confiança 0.52 \textperiodcentered{} domínio hipotese \textperiodcentered{} história 8 versões no jornal}

%CRISTAL {"arestas": [{"alvo": "corpo_topologico", "confianca": 1.0, "epistemico": "fato", "origem": "wiki", "rel": "usa"}, {"alvo": "dissipacao_gato_concentracao_esquilo", "confianca": 1.0, "epistemico": "fato", "origem": "wiki", "rel": "relaciona"}], "confianca": 1.0, "contraexemplos": [], "descricao": "A pergunta central do corpo topológico — DADA UMA SUPERFÍCIE, OBTER A TRANSFORMAÇÃO EM OUTRA — tem dois regimes, decididos só por χ: (1) MESMO χ (mesma classe) ⟹ HOMEOMORFISMO contínuo (o CRISTAL, o morph na bancada): para superfícies estreladas, o mapa radial Φ(p)=(ρ_B/ρ_A)·p, BIJETIVO (Φ⁻¹∘Φ=id, verificado 1e-16) e contínuo (ρ_B>0), χ conservado, NENHUMA cirurgia — esfera↔elipsóide↔estrela. (2) χ DIFERENTE (classe distinta) ⟹ NÃO existe homeomorfismo, só CIRURGIA (soma conexa/alças): χ(A#B)=χ_A+χ_B−2, cada alça=A#toro faz χ→χ−2 (genus+1) e PASSA POR UMA SELA (a borda, o ponto singular onde a topologia salta); nº de alças=|χ_A−χ_B|/2. A RECEITA: χ_A=χ_B→('homeomorfismo',0); senão→('cirurgia',|χ_A−χ_B|/2). Ex: esfera→bitoro=2 alças. O regime do cristal é um DIFEOMORFISMO (suave, sem RASGAR): entra-se com as EQUAÇÕES das duas superfícies T(u,v),M(u,v) no mesmo domínio, e a transformação é a HOMOTOPIA ponto a ponto p=(1−t)T+t·M (a operação ⊕ do corpo) — não se corta (isso é cirurgia), deforma-se. Clássico: TORO→XÍCARA (ambos genus 1, o hole do toro=o hole da alça, χ=0 conservado). TUDO PONTO A PONTO (sem malha, sem FFT). Renders: morph_homeo.gif, toro_xicara.gif, cirurgia_esfera_toro.gif. linguagem/corpo_topologico.py (transformacao_entre_superficies 4/4); papers/corpo_topologico.tex; assets/pipe_grafico/topologico/.", "epistemico": "fato", "exemplos": [], "id": "transformacao_superficies_homeo_cirurgia", "memoria": "semantica", "meta": {"dominio": "floxina", "fonte": "Floxina (investigação)"}, "origem": "wiki", "palavras_chave": [], "sinonimos": [], "tipo": "conceito", "titulo": "Dada uma superfície, a transformação em outra: mesmo χ ⟹ homeomorfismo (o morph); χ diferente ⟹ cirurgia (alças, cada uma cruza uma sela)"}
\section{Dada uma superfície, a transformação em outra: mesmo \ensuremath{\chi} \ensuremath{\Longrightarrow} homeomorfismo (o morph); \ensuremath{\chi} diferente \ensuremath{\Longrightarrow} cirurgia (alças, cada uma cruza uma sela)}
A pergunta central do corpo topológico — DADA UMA SUPERFÍCIE, OBTER A TRANSFORMAÇÃO EM OUTRA — tem dois regimes, decididos só por \ensuremath{\chi}: (1) MESMO \ensuremath{\chi} (mesma classe) \ensuremath{\Longrightarrow} HOMEOMORFISMO contínuo (o CRISTAL, o morph na bancada): para superfícies estreladas, o mapa radial \ensuremath{\Phi}(p)=(\ensuremath{\rho}\_B/\ensuremath{\rho}\_A)\textperiodcentered p, BIJETIVO (\ensuremath{\Phi}\ensuremath{^{-1}}\ensuremath{\circ}\ensuremath{\Phi}=id, verificado 1e-16) e contínuo (\ensuremath{\rho}\_B>0), \ensuremath{\chi} conservado, NENHUMA cirurgia — esfera\ensuremath{\leftrightarrow}elipsóide\ensuremath{\leftrightarrow}estrela. (2) \ensuremath{\chi} DIFERENTE (classe distinta) \ensuremath{\Longrightarrow} NÃO existe homeomorfismo, só CIRURGIA (soma conexa/alças): \ensuremath{\chi}(A\#B)=\ensuremath{\chi}\_A+\ensuremath{\chi}\_B\ensuremath{-}2, cada alça=A\#toro faz \ensuremath{\chi}\ensuremath{\to}\ensuremath{\chi}\ensuremath{-}2 (genus+1) e PASSA POR UMA SELA (a borda, o ponto singular onde a topologia salta); nº de alças=|\ensuremath{\chi}\_A\ensuremath{-}\ensuremath{\chi}\_B|/2. A RECEITA: \ensuremath{\chi}\_A=\ensuremath{\chi}\_B\ensuremath{\to}('homeomorfismo',0); senão\ensuremath{\to}('cirurgia',|\ensuremath{\chi}\_A\ensuremath{-}\ensuremath{\chi}\_B|/2). Ex: esfera\ensuremath{\to}bitoro=2 alças. O regime do cristal é um DIFEOMORFISMO (suave, sem RASGAR): entra-se com as EQUAÇÕES das duas superfícies T(u,v),M(u,v) no mesmo domínio, e a transformação é a HOMOTOPIA ponto a ponto p=(1\ensuremath{-}t)T+t\textperiodcentered M (a operação \ensuremath{\oplus} do corpo) — não se corta (isso é cirurgia), deforma-se. Clássico: TORO\ensuremath{\to}XÍCARA (ambos genus 1, o hole do toro=o hole da alça, \ensuremath{\chi}=0 conservado). TUDO PONTO A PONTO (sem malha, sem FFT). Renders: morph\_homeo.gif, toro\_xicara.gif, cirurgia\_esfera\_toro.gif. linguagem/corpo\_topologico.py (transformacao\_entre\_superficies 4/4); papers/corpo\_topologico.tex; assets/pipe\_grafico/topologico/.
\prov{relações: usa corpo\_topologico; relaciona dissipacao\_gato\_concentracao\_esquilo}
\prov{proveniência: wiki \textperiodcentered{} Floxina (investigação) \textperiodcentered{} fato \textperiodcentered{} confiança 1.0 \textperiodcentered{} domínio floxina \textperiodcentered{} história 15 versões no jornal}

%CRISTAL {"arestas": [{"alvo": "corpo_edp", "confianca": 1.0, "epistemico": "fato", "origem": "wiki", "rel": "usa"}, {"alvo": "solucao_explicita_ed", "confianca": 1.0, "epistemico": "fato", "origem": "wiki", "rel": "usa"}, {"alvo": "pa_pg_corpo_diferencial", "confianca": 1.0, "epistemico": "fato", "origem": "wiki", "rel": "relaciona"}], "confianca": 1.0, "contraexemplos": [], "descricao": "A fórmula explícita das EDP u(x,t)=∫û(k)e^{i(kx−ω(k)t)}dk (e da ODE x(t)=∑cᵢe^{λᵢt}) precisa de UMA transformada para dar o espectro {û(k)} / {λᵢ}. A nativa do corpo NÃO é Fourier plana — é a TRANSFORMADA DOURADA 𝔊=ℱ∘Λ (Fourier no LOG = Mellin na linha crítica; base φ^k na dourada, σ_n^k na metálica; ponte_simbolica_gentil.tex Def. 𝔊). Porque os modos do corpo são METÁLICOS: os autovalores são λ=log σ (o FLIP PG→PA que troca o produto dos metais σ pela soma dos expoentes log σ), e Λ=log é o que os ENDIREITA para o Fourier ler. Verificado: 𝔊 resolve o espectro {log φ, log(1+√2)} (erro<2e-3) onde a Fourier crua FALHA (o modo metálico cos(γ·log t) é log-periódico, não periódico). O corpo tem espectro metálico; 𝔊 é o instrumento que o lê. É a transformada metálica (transformada-mineral.tex): a β-numeração base σ_n, o λ=log ρ=Lyapunov na parábola (Pisot/Salem/caos). linguagem/corpo_edp.py (transformada_dourada_da_o_espectro 3/3); papers/equacoes_diferenciais.tex (Parte VIII).", "epistemico": "fato", "exemplos": [], "id": "transformada_dourada_espectro_edp", "memoria": "semantica", "meta": {"dominio": "floxina", "fonte": "Floxina (investigação)"}, "origem": "wiki", "palavras_chave": [], "sinonimos": [], "tipo": "conceito", "titulo": "A Transformada Dourada 𝔊=ℱ∘Λ dá o espectro da fórmula explícita (o corpo tem espectro METÁLICO, não Fourier cru)"}
\section{A Transformada Dourada \ensuremath{\mathfrak{G}}=\ensuremath{\mathcal{F}}\ensuremath{\circ}\ensuremath{\Lambda} dá o espectro da fórmula explícita (o corpo tem espectro METÁLICO, não Fourier cru)}
A fórmula explícita das EDP u(x,t)=\ensuremath{\int}û(k)e\^{}\{i(kx\ensuremath{-}\ensuremath{\omega}(k)t)\}dk (e da ODE x(t)=\ensuremath{\sum}c\ensuremath{_{i}}e\^{}\{\ensuremath{\lambda}\ensuremath{_{i}}t\}) precisa de UMA transformada para dar o espectro \{û(k)\} / \{\ensuremath{\lambda}\ensuremath{_{i}}\}. A nativa do corpo NÃO é Fourier plana — é a TRANSFORMADA DOURADA \ensuremath{\mathfrak{G}}=\ensuremath{\mathcal{F}}\ensuremath{\circ}\ensuremath{\Lambda} (Fourier no LOG = Mellin na linha crítica; base \ensuremath{\varphi}\^{}k na dourada, \ensuremath{\sigma}\_n\^{}k na metálica; ponte\_simbolica\_gentil.tex Def. \ensuremath{\mathfrak{G}}). Porque os modos do corpo são METÁLICOS: os autovalores são \ensuremath{\lambda}=log \ensuremath{\sigma} (o FLIP PG\ensuremath{\to}PA que troca o produto dos metais \ensuremath{\sigma} pela soma dos expoentes log \ensuremath{\sigma}), e \ensuremath{\Lambda}=log é o que os ENDIREITA para o Fourier ler. Verificado: \ensuremath{\mathfrak{G}} resolve o espectro \{log \ensuremath{\varphi}, log(1+\ensuremath{\sqrt{\,}}2)\} (erro<2e-3) onde a Fourier crua FALHA (o modo metálico cos(\ensuremath{\gamma}\textperiodcentered log t) é log-periódico, não periódico). O corpo tem espectro metálico; \ensuremath{\mathfrak{G}} é o instrumento que o lê. É a transformada metálica (transformada-mineral.tex): a \ensuremath{\beta}-numeração base \ensuremath{\sigma}\_n, o \ensuremath{\lambda}=log \ensuremath{\rho}=Lyapunov na parábola (Pisot/Salem/caos). linguagem/corpo\_edp.py (transformada\_dourada\_da\_o\_espectro 3/3); papers/equacoes\_diferenciais.tex (Parte VIII).
\prov{relações: usa corpo\_edp; usa solucao\_explicita\_ed; relaciona pa\_pg\_corpo\_diferencial}
\prov{proveniência: wiki \textperiodcentered{} Floxina (investigação) \textperiodcentered{} fato \textperiodcentered{} confiança 1.0 \textperiodcentered{} domínio floxina \textperiodcentered{} história 40 versões no jornal}

%CRISTAL {"arestas": [{"alvo": "ginzburg_toro_helice_e_3d-sst", "confianca": 1.0, "epistemico": "fato", "origem": "ginzburg_toro_helice.md", "rel": "parte_de"}], "confianca": 0.52, "contraexemplos": [], "descricao": "Ginzburg propõe que **toro + hélice** são elementos primários da natureza em todas as escalas — do atómico ao cosmológico (conjectura holística).", "epistemico": "hipotese", "exemplos": [], "id": "unificacao_torohelice", "memoria": "semantica", "meta": {"dominio": "hipotese", "falsificavel": true, "nivel": 7}, "origem": "ginzburg_toro_helice.md", "palavras_chave": [], "sinonimos": [], "tipo": "conceito", "titulo": "Unificação toro–hélice"}
\section{Unificação toro–hélice}
Ginzburg propõe que **toro + hélice** são elementos primários da natureza em todas as escalas — do atómico ao cosmológico (conjectura holística).
\prov{relações: parte\_de ginzburg\_toro\_helice\_e\_3d-sst}
\prov{proveniência: ginzburg\_toro\_helice.md \textperiodcentered{} hipotese \textperiodcentered{} confiança 0.52 \textperiodcentered{} domínio hipotese \textperiodcentered{} história 3 versões no jornal}

%CRISTAL {"arestas": [{"alvo": "corpo_matricial_matrix", "confianca": 1.0, "epistemico": "fato", "origem": "wiki", "rel": "parte_de"}, {"alvo": "gato", "confianca": 1.0, "epistemico": "fato", "origem": "wiki", "rel": "usa"}, {"alvo": "gamba_antissimetrico", "confianca": 1.0, "epistemico": "fato", "origem": "wiki", "rel": "usa"}, {"alvo": "metais", "confianca": 1.0, "epistemico": "fato", "origem": "wiki", "rel": "relaciona"}], "confianca": 1.0, "contraexemplos": [], "descricao": "O PRIMEIRO PASSO das multimatrizes: da folha 2D (M_2, a reta σ_n de x²−n·x−1) para o VOLUME 3D (M_3, os cúbicos do Gentil). O GATO 3D = a companion de x³−a·x²−b·x−1 = [[a,b,1],[1,0,0],[0,1,0]], det=1 (a NORMA CONSERVADA, produto das 3 raízes=1, o a+c²=1 do Gentil — o cristal); a raiz dominante é o METAL 3D: plástico (a=0,b=1, ≈1,3247), supergolden (a=1,b=0, ≈1,4656), tribonacci (a=1,b=1, ≈1,8393). As potências dão a recorrência de ordem 3 (a m-bonacci, PA+PG de ordem m). O GAMBÁ 3D = o gerador anti-simétrico de SO(3) (a rotação/o produto vetorial); os três eixos fecham so(3) ([J_x,J_y]=J_z). O volume M_3 = Sym(gato, dim 6) ⊕ Skew(gambá, dim 3) = 9; o produto XY = X∘Y (Jordan) + ½[X,Y] (Lie). É o germe do próximo passo (4D, …), um passo de dimensão por vez. linguagem/volume_matrizes.py (6/6).", "epistemico": "hipotese", "exemplos": [], "id": "volume_matrizes_3d", "memoria": "semantica", "meta": {"dominio": "floxina", "fonte": "Floxina (investigação)"}, "origem": "wiki", "palavras_chave": [], "sinonimos": [], "tipo": "conceito", "titulo": "O volume de matrizes 3D — o primeiro passo (M_2 → M_3)"}
\section{O volume de matrizes 3D — o primeiro passo (M\_2 \ensuremath{\to} M\_3)}
O PRIMEIRO PASSO das multimatrizes: da folha 2D (M\_2, a reta \ensuremath{\sigma}\_n de x\ensuremath{^{2}}\ensuremath{-}n\textperiodcentered x\ensuremath{-}1) para o VOLUME 3D (M\_3, os cúbicos do Gentil). O GATO 3D = a companion de x\ensuremath{^{3}}\ensuremath{-}a\textperiodcentered x\ensuremath{^{2}}\ensuremath{-}b\textperiodcentered x\ensuremath{-}1 = [[a,b,1],[1,0,0],[0,1,0]], det=1 (a NORMA CONSERVADA, produto das 3 raízes=1, o a+c\ensuremath{^{2}}=1 do Gentil — o cristal); a raiz dominante é o METAL 3D: plástico (a=0,b=1, \ensuremath{\approx}1,3247), supergolden (a=1,b=0, \ensuremath{\approx}1,4656), tribonacci (a=1,b=1, \ensuremath{\approx}1,8393). As potências dão a recorrência de ordem 3 (a m-bonacci, PA+PG de ordem m). O GAMBÁ 3D = o gerador anti-simétrico de SO(3) (a rotação/o produto vetorial); os três eixos fecham so(3) ([J\_x,J\_y]=J\_z). O volume M\_3 = Sym(gato, dim 6) \ensuremath{\oplus} Skew(gambá, dim 3) = 9; o produto XY = X\ensuremath{\circ}Y (Jordan) + \ensuremath{\tfrac12}[X,Y] (Lie). É o germe do próximo passo (4D, …), um passo de dimensão por vez. linguagem/volume\_matrizes.py (6/6).
\prov{relações: parte\_de corpo\_matricial\_matrix; usa gato; usa gamba\_antissimetrico; relaciona metais}
\prov{proveniência: wiki \textperiodcentered{} Floxina (investigação) \textperiodcentered{} hipotese \textperiodcentered{} confiança 1.0 \textperiodcentered{} domínio floxina \textperiodcentered{} história 360 versões no jornal}

\end{document}
